\input{preamble}

% OK, commencer ici.
%
\begin{document}

\title{Morphismes de champs algébriques}


\maketitle

\phantomsection
\label{section-phantom}

\tableofcontents

\section{Introduction}
\label{section-introduction}

\noindent
Dans ce chapitre, nous introduisons certains types de morphismes de champs
algébriques. Une référence pour le cas des champs algébriques quasi-séparés à
diagonale représentable est \cite{LM-B}.

\medskip\noindent
Le but est d'étendre à des morphismes de champs algébriques la définition de
chacun des types de morphismes d'espaces algébriques. Chaque cas présente de
légères différences, et il semble préférable de les traiter séparément.

\medskip\noindent
Pour les morphismes de champs algébriques représentables par des espaces
algébriques, nous avons déjà défini un grand nombre de types de morphismes, voir
Propriétés des champs,
section \ref{stacks-properties-section-properties-morphisms}.
Dans chaque cas correspondant du présent chapitre, nous devons vérifier que la
définition générale est compatible avec celle qui y a été donnée.




\section{Conventions et abus de langage}
\label{section-conventions}

\noindent
Nous continuons d'employer les conventions et les abus de langage introduits
dans Propriétés des champs, section
\ref{stacks-properties-section-conventions}.




\section{Propriétés des diagonales}
\label{section-diagonals}

\noindent
La diagonale d'un champ algébrique est étroitement liée aux faisceaux
$\mathit{Isom}$, voir
Champs algébriques, Lemme \ref{algebraic-lemma-representable-diagonal}.
Par la seconde propriété qui définit un champ algébrique, ces faisceaux
$\mathit{Isom}$ sont toujours des espaces algébriques.

\begin{lemma}
\label{lemma-isom-locally-finite-type}
Soit $\mathcal{X}$ un champ algébrique.
Soit $T$ un schéma et soient $x,y$ des objets de la catégorie fibre de
$\mathcal{X}$ au-dessus de $T$. Alors le morphisme
$\mathit{Isom}_\mathcal{X}(x,y) \to T$ est localement de type fini.
\end{lemma}

\begin{proof}
Par Champs algébriques, Lemme
\ref{algebraic-lemma-stack-presentation}, nous pouvons supposer que
$\mathcal{X}=[U/R]$ pour un groupoïde lisse en espaces algébriques.
Par Descente sur les espaces,
Lemme \ref{spaces-descent-lemma-descending-property-locally-finite-type},
il suffit de vérifier la propriété localement pour la topologie fppf sur $T$.
Nous pouvons donc supposer que $x,y$ proviennent de morphismes
$x',y':T\to U$. Par Groupoïdes en espaces,
Lemme \ref{spaces-groupoids-lemma-quotient-stack-morphisms}, nous avons alors
$\mathit{Isom}_\mathcal{X}(x,y)=T\times_{(y',x'),U\times_S U}R$.
Il suffit donc de prouver que $R\to U\times_S U$ est localement de type fini.
Cela résulte du fait que le composé
$s:R\to U\times_S U\to U$ est lisse (donc localement de type fini, voir
Morphismes d'espaces, Lemmes
\ref{spaces-morphisms-lemma-smooth-locally-finite-presentation} et
\ref{spaces-morphisms-lemma-finite-presentation-finite-type})
et Morphismes d'espaces, Lemme
\ref{spaces-morphisms-lemma-permanence-finite-type}.
\end{proof}

\begin{lemma}
\label{lemma-isom-pseudo-torsor-aut}
Soit $\mathcal{X}$ un champ algébrique.
Soit $T$ un schéma et soient $x,y$ des objets de la catégorie fibre de
$\mathcal{X}$ au-dessus de $T$. Alors
\begin{enumerate}
\item $\mathit{Isom}_\mathcal{X}(y,y)$ est un espace algébrique en groupes
sur $T$, et
\item $\mathit{Isom}_\mathcal{X}(x,y)$ est un pseudo-torseur sous
$\mathit{Isom}_\mathcal{X}(y,y)$ sur $T$.
\end{enumerate}
\end{lemma}

\begin{proof}
Voir Groupoïdes en espaces,
Définitions \ref{spaces-groupoids-definition-group-space} et
\ref{spaces-groupoids-definition-pseudo-torsor}.
Le lemme résulte immédiatement du fait que $\mathcal{X}$ est un champ en
groupoïdes.
\end{proof}

\noindent
Soit $f:\mathcal{X}\to\mathcal{Y}$ un morphisme de champs algébriques.
La {\it diagonale de $f$} est le morphisme
$$
\Delta_f :
\mathcal{X}
\longrightarrow
\mathcal{X} \times_\mathcal{Y} \mathcal{X}
$$
Voici deux propriétés que possède tout morphisme diagonal.

\begin{lemma}
\label{lemma-properties-diagonal}
\begin{slogan}
Les diagonales des morphismes de champs algébriques sont représentables par
des espaces algébriques et localement de type fini.
\end{slogan}
Soit $f:\mathcal{X}\to\mathcal{Y}$ un morphisme de champs algébriques.
Alors
\begin{enumerate}
\item $\Delta_f$ est représentable par des espaces algébriques, et
\item $\Delta_f$ est localement de type fini.
\end{enumerate}
\end{lemma}

\begin{proof}
Soit $T$ un schéma et soit
$a:T\to\mathcal{X}\times_\mathcal{Y}\mathcal{X}$ un morphisme. Par la
définition du produit fibré et le lemme de $2$-Yoneda, le morphisme $a$ est
donné par un triplet $a=(x,x',\alpha)$, où $x,x'$ sont des objets de
$\mathcal{X}$ au-dessus de $T$ et
$\alpha:f(x)\to f(x')$ est un morphisme dans la catégorie fibre de
$\mathcal{Y}$ au-dessus de $T$. Par la définition d'un champ algébrique, les
faisceaux $\mathit{Isom}_\mathcal{X}(x,x')$ et
$\mathit{Isom}_\mathcal{Y}(f(x),f(x'))$ sont des espaces algébriques sur
$T$. Dans ce langage, $\alpha$ définit une section du morphisme
$\mathit{Isom}_\mathcal{Y}(f(x),f(x'))\to T$. Un point à valeurs dans $T'$ de
$\mathcal{X}\times_{\mathcal{X}\times_\mathcal{Y}\mathcal{X},a}T$, où
$T'\to T$ est un schéma au-dessus de $T$, est la même chose qu'un
isomorphisme $x|_{T'}\to x'|_{T'}$ dont l'image par
$f$ est $\alpha|_{T'}$. Nous voyons donc que
\begin{equation}
\label{equation-diagonal}
\vcenter{
\xymatrix{
\mathcal{X} \times_{\mathcal{X} \times_\mathcal{Y} \mathcal{X}, a} T
\ar[d] \ar[r] &
\mathit{Isom}_\mathcal{X}(x, x') \ar[d] \\
T\ar[r]^-\alpha &
\mathit{Isom}_\mathcal{Y}(f(x), f(x'))
}
}
\end{equation}
est un carré cartésien de faisceaux sur $T$. En particulier,
$\mathcal{X} \times_{\mathcal{X} \times_\mathcal{Y} \mathcal{X}, a} T$
est un espace algébrique, ce qui démontre l'assertion (1) du lemme.

\medskip\noindent
Pour démontrer la seconde assertion, nous devons montrer que la flèche
verticale de gauche du Diagramme (\ref{equation-diagonal}) est localement de
type fini. Par le Lemme \ref{lemma-isom-locally-finite-type}, l'espace
algébrique $\mathit{Isom}_\mathcal{X}(x,x')$ est localement de type fini sur
$T$. La flèche verticale de droite du Diagramme
(\ref{equation-diagonal}) est donc localement de type fini, voir
Morphismes d'espaces, Lemme
\ref{spaces-morphisms-lemma-permanence-finite-type}. Nous concluons par
Morphismes d'espaces, Lemme
\ref{spaces-morphisms-lemma-base-change-finite-type}.
\end{proof}

\begin{lemma}
\label{lemma-properties-diagonal-representable}
Soit $f:\mathcal{X}\to\mathcal{Y}$ un morphisme de champs algébriques
représentable par des espaces algébriques. Alors
\begin{enumerate}
\item $\Delta_f$ est représentable (par des schémas),
\item $\Delta_f$ est localement de type fini,
\item $\Delta_f$ est un monomorphisme,
\item $\Delta_f$ est séparé, et
\item $\Delta_f$ est localement quasi-fini.
\end{enumerate}
\end{lemma}

\begin{proof}
Nous avons déjà vu au Lemme \ref{lemma-properties-diagonal} que $\Delta_f$
est représentable par des espaces algébriques. Les assertions (2) -- (5) ont
donc un sens, voir Propriétés des champs,
section \ref{stacks-properties-section-properties-morphisms}. Le Lemme
\ref{lemma-properties-diagonal} garantit aussi (2). Soit
$T\to\mathcal{X}\times_\mathcal{Y}\mathcal{X}$ un morphisme, et considérons
le Diagramme (\ref{equation-diagonal}). Par Champs algébriques, Lemme
\ref{algebraic-lemma-criterion-map-representable-spaces-fibred-in-groupoids},
la flèche verticale de droite est injective comme application de faisceaux,
c'est-à-dire qu'elle est un monomorphisme d'espaces algébriques. Par suite, le
morphisme
$T\times_{\mathcal{X}\times_\mathcal{Y}\mathcal{X}}\mathcal{X}\to T$
est lui aussi un monomorphisme. Ainsi, (3) est satisfaite. Nous savons déjà que
$T\times_{\mathcal{X}\times_\mathcal{Y}\mathcal{X}}\mathcal{X}\to T$ est
localement de type fini. Morphismes d'espaces, Lemme
\ref{spaces-morphisms-lemma-monomorphism-loc-finite-type-loc-quasi-finite},
permet donc de conclure que
$T\times_{\mathcal{X}\times_\mathcal{Y}\mathcal{X}}\mathcal{X}\to T$
est localement quasi-fini et séparé.
Cela démontre (4) et (5). Enfin, Morphismes d'espaces, Proposition
\ref{spaces-morphisms-proposition-locally-quasi-finite-separated-over-scheme},
implique que
$T\times_{\mathcal{X}\times_\mathcal{Y}\mathcal{X}}\mathcal{X}$ est un
schéma, ce qui démontre (1).
\end{proof}

\begin{lemma}
\label{lemma-representable-separated-diagonal-closed}
Soit $f:\mathcal{X}\to\mathcal{Y}$ un morphisme de champs algébriques
représentable par des espaces algébriques. Les assertions suivantes sont
équivalentes :
\begin{enumerate}
\item $f$ est séparé ;
\item $\Delta_f$ est une immersion fermée ;
\item $\Delta_f$ est propre ;
\item $\Delta_f$ est universellement fermé.
\end{enumerate}
\end{lemma}

\begin{proof}
Les énoncés « $f$ est séparé », « $\Delta_f$ est une immersion fermée »,
« $\Delta_f$ est universellement fermé » et « $\Delta_f$ est propre » font
référence aux notions définies dans Propriétés des champs,
section \ref{stacks-properties-section-properties-morphisms}. Choisissons un
schéma $V$ et un morphisme lisse surjectif $V\to\mathcal{Y}$. Posons
$U=\mathcal{X}\times_\mathcal{Y}V$, qui est par hypothèse un espace
algébrique ; le morphisme $U\to\mathcal{X}$ est lisse et surjectif. Par
Catégories, Lemme \ref{categories-lemma-base-change-diagonal}, et Propriétés
des champs, Lemme \ref{stacks-properties-lemma-check-property-covering}, pour
toute propriété $P$ considérée dans ce dernier lemme, $\Delta_f$ possède $P$
si et seulement si $\Delta_{U/V}:U\to U\times_VU$ possède $P$. L'équivalence
de (2), (3) et (4) résulte donc de Morphismes d'espaces, Lemme
\ref{spaces-morphisms-lemma-separated-diagonal-proper}, appliqué à $U\to V$.
De plus, si (1) est satisfaite, $U\to V$ est séparé, donc
$\Delta_{U/V}$ est une immersion fermée : (2) est satisfaite. Supposons enfin
(2). Soient $T$ un schéma et $a:T\to\mathcal{Y}$ un morphisme. Posons
$T'=\mathcal{X}\times_\mathcal{Y}T$. Pour démontrer (1), il faut montrer que
le morphisme d'espaces algébriques $T'\to T$ est séparé. En appliquant de
nouveau Catégories, Lemme \ref{categories-lemma-base-change-diagonal}, nous
voyons que $\Delta_{T'/T}$ est le changement de base de $\Delta_f$.
L'hypothèse (2) implique donc que $\Delta_{T'/T}$ est une immersion fermée,
et ainsi que $T'\to T$ est séparé, comme voulu.
\end{proof}

\begin{lemma}
\label{lemma-representable-quasi-separated-diagonal-quasi-compact}
Soit $f:\mathcal{X}\to\mathcal{Y}$ un morphisme de champs algébriques
représentable par des espaces algébriques. Les assertions suivantes sont
équivalentes :
\begin{enumerate}
\item $f$ est quasi-séparé ;
\item $\Delta_f$ est quasi-compact ;
\item $\Delta_f$ est de type fini.
\end{enumerate}
\end{lemma}

\begin{proof}
Les énoncés « $f$ est quasi-séparé », « $\Delta_f$ est quasi-compact » et
« $\Delta_f$ est de type fini » font référence aux notions définies dans
Propriétés des champs,
section \ref{stacks-properties-section-properties-morphisms}. Remarquons que
(2) et (3) sont équivalentes, puisque $\Delta_f$ est localement de type fini
par le Lemme \ref{lemma-properties-diagonal-representable} (et Champs
algébriques, Lemme
\ref{algebraic-lemma-representable-transformations-property-implication}).
Choisissons un schéma $V$ et un morphisme lisse surjectif
$V\to\mathcal{Y}$. Posons $U=\mathcal{X}\times_\mathcal{Y}V$, qui est par
hypothèse un espace algébrique ; le morphisme $U\to\mathcal{X}$ est lisse et
surjectif. Par Catégories, Lemme
\ref{categories-lemma-base-change-diagonal}, et Propriétés des champs, Lemme
\ref{stacks-properties-lemma-check-property-covering}, $\Delta_f$ est
quasi-compact si et seulement si
$\Delta_{U/V}:U\to U\times_VU$ est quasi-compact. Si (1) est satisfaite,
$U\to V$ est quasi-séparé, donc $\Delta_{U/V}$ est quasi-compact : (2) est
satisfaite. Supposons (2). Soient $T$ un schéma et
$a:T\to\mathcal{Y}$ un morphisme. Posons
$T'=\mathcal{X}\times_\mathcal{Y}T$. Pour démontrer (1), il faut montrer que
le morphisme d'espaces algébriques $T'\to T$ est quasi-séparé. En appliquant
de nouveau Catégories, Lemme \ref{categories-lemma-base-change-diagonal}, nous
voyons que $\Delta_{T'/T}$ est le changement de base de $\Delta_f$.
L'hypothèse (2) implique donc que $\Delta_{T'/T}$ est quasi-compact, et ainsi
que $T'\to T$ est quasi-séparé, comme voulu.
\end{proof}

\begin{lemma}
\label{lemma-representable-locally-separated-diagonal-immersion}
Soit $f:\mathcal{X}\to\mathcal{Y}$ un morphisme de champs algébriques
représentable par des espaces algébriques. Les assertions suivantes sont
équivalentes :
\begin{enumerate}
\item $f$ est localement séparé ;
\item $\Delta_f$ est une immersion.
\end{enumerate}
\end{lemma}

\begin{proof}
Les énoncés « $f$ est localement séparé » et « $\Delta_f$ est une immersion »
font référence aux notions définies dans Propriétés des champs,
section \ref{stacks-properties-section-properties-morphisms}.
Démonstration omise. Indication : raisonner comme dans les démonstrations des
Lemmes \ref{lemma-representable-separated-diagonal-closed} et
\ref{lemma-representable-quasi-separated-diagonal-quasi-compact}.
\end{proof}






\section{Axiomes de séparation}
\label{section-separated}

\noindent
Soit $\mathcal{X}=[U/R]$ une présentation d'un champ algébrique. Les
propriétés de la diagonale de $\mathcal{X}$ sur $S$ sont alors celles du
morphisme $j:R\to U\times_SU$. Par exemple, si
$\mathcal{X}=[S/G]$ pour un groupe lisse $G$ en espaces algébriques sur $S$,
alors $j$ est le morphisme structural $G\to S$. La diagonale elle-même n'est
donc pas automatiquement séparée (contrairement à ce qui se produit pour les
schémas et les espaces algébriques). Dire que $[S/G]$ est quasi-séparé sur
$S$ devrait certainement impliquer que $G\to S$ est quasi-compact ; nous
hésitons toutefois à dire que $[S/G]$ est quasi-séparé sur $S$ sans exiger aussi que
$G\to S$ soit quasi-séparé. Autrement dit, imposer seulement que le morphisme
diagonal soit quasi-compact ne correspond pas vraiment à l'intuition que nous
avons d'un « champ algébrique quasi-séparé » : il faut aussi exiger que la
diagonale elle-même soit quasi-séparée.

\medskip\noindent
Qu'en est-il des « champs algébriques séparés » ? Nous avons vu dans
Morphismes d'espaces,
Lemme \ref{spaces-morphisms-lemma-separated-diagonal-proper}, qu'un espace
algébrique est séparé si et seulement si sa diagonale est propre. C'est aussi
la condition habituellement employée pour définir les champs algébriques
séparés. Dans l'exemple $[S/G]\to S$ ci-dessus, cela signifie que $G\to S$
est un schéma en groupes propre. Ainsi, les champs algébriques de la forme
$[\Spec(k)/E]$ sont propres sur $k$ lorsque $E$ est une courbe elliptique sur
$k$ (insérer ici une référence ultérieure). Dans certaines situations, il peut
être plus naturel de supposer la diagonale finie.

\begin{definition}
\label{definition-separated}
Soit $f:\mathcal{X}\to\mathcal{Y}$ un morphisme de champs algébriques.
\begin{enumerate}
\item Nous disons que $f$ est {\it DM} si $\Delta_f$ est non
ramifié\footnote{Les lettres DM signifient Deligne--Mumford. Si $f$ est DM,
alors, pour tout schéma $T$ et tout morphisme $T\to\mathcal{Y}$, le produit fibré
$\mathcal{X}_T = \mathcal{X} \times_\mathcal{Y} T$
est un champ algébrique sur $T$ dont la diagonale est non ramifiée,
c'est-à-dire que $\mathcal{X}_T$ est DM. Il s'ensuit que $\mathcal{X}_T$ est
un champ de Deligne--Mumford, voir le Théorème \ref{theorem-DM}. Autrement dit,
un morphisme DM est un morphisme dont les « fibres » sont des champs de
Deligne--Mumford. Cela motive, espérons-le, au moins la terminologie.}.
\item Nous disons que $f$ est {\it quasi-DM} si $\Delta_f$ est localement
quasi-fini\footnote{Si $f$ est quasi-DM, alors les « fibres »
$\mathcal{X}_T$ de $\mathcal{X}\to\mathcal{Y}$ sont quasi-DM. Un champ
algébrique $\mathcal{X}$ est quasi-DM si et seulement s'il existe un schéma
$U$ et un morphisme plat surjectif $U\to\mathcal{X}$, de présentation finie
et localement quasi-fini, voir le Théorème \ref{theorem-quasi-DM}. Notons la
ressemblance avec la propriété d'être de Deligne--Mumford, définie par
l'existence d'un recouvrement étale par un schéma.}.
\item Nous disons que $f$ est {\it séparé} si $\Delta_f$ est propre.
\item Nous disons que $f$ est {\it quasi-séparé} si $\Delta_f$ est
quasi-compact et quasi-séparé.
\end{enumerate}
\end{definition}

\noindent
Dans cette définition, nous utilisons le fait que $\Delta_f$ est représentable
par des espaces algébriques ainsi que Propriétés des champs,
section \ref{stacks-properties-section-properties-morphisms}, afin de donner
un sens aux conditions imposées à $\Delta_f$. Notons que ces définitions sont
compatibles avec les notions déjà existantes lorsque $f$ est représentable par
des espaces algébriques, voir les Lemmes
\ref{lemma-representable-quasi-separated-diagonal-quasi-compact} et
\ref{lemma-representable-separated-diagonal-closed}. On peut caractériser ces
conditions d'une manière intéressante en considérant les diagonales d'ordre
supérieur, voir le Lemme \ref{lemma-definition-separated}.

\begin{definition}
\label{definition-absolute-separated}
Soit $\mathcal{X}$ un champ algébrique sur le schéma de base $S$. Notons
$p:\mathcal{X}\to S$ le morphisme structural.
\begin{enumerate}
\item Nous disons que $\mathcal{X}$ est {\it DM sur $S$} si
$p:\mathcal{X}\to S$ est DM.
\item Nous disons que $\mathcal{X}$ est {\it quasi-DM sur $S$} si
$p:\mathcal{X}\to S$ est quasi-DM.
\item Nous disons que $\mathcal{X}$ est {\it séparé sur $S$} si
$p:\mathcal{X}\to S$ est séparé.
\item Nous disons que $\mathcal{X}$ est {\it quasi-séparé sur $S$} si
$p:\mathcal{X}\to S$ est quasi-séparé.
\item Nous disons que $\mathcal{X}$ est {\it DM} si $\mathcal{X}$ est
DM\footnote{Le Théorème \ref{theorem-DM} montre que cela équivaut à dire que
$\mathcal{X}$ est un champ de Deligne--Mumford.} sur $\Spec(\mathbf{Z})$.
\item Nous disons que $\mathcal{X}$ est {\it quasi-DM} si $\mathcal{X}$ est
quasi-DM sur $\Spec(\mathbf{Z})$.
\item Nous disons que $\mathcal{X}$ est {\it séparé} si $\mathcal{X}$ est
séparé sur $\Spec(\mathbf{Z})$.
\item Nous disons que $\mathcal{X}$ est {\it quasi-séparé} si
$\mathcal{X}$ est quasi-séparé sur $\Spec(\mathbf{Z})$.
\end{enumerate}
Dans les quatre dernières définitions, nous considérons $\mathcal{X}$ comme un
champ algébrique sur $\Spec(\mathbf{Z})$ au moyen de Champs algébriques,
Définition \ref{algebraic-definition-viewed-as}.
\end{definition}

\noindent
Dans chaque cas, nous disposons donc d'une notion absolue et d'une notion
relative au schéma de base fixé (que notre abus de notation introduit dans
Propriétés des champs, section \ref{stacks-properties-section-conventions},
conduit généralement à ne pas mentionner). Nous verrons que
(1) $\Leftrightarrow$ (5) et (2) $\Leftrightarrow$ (6) au Lemme
\ref{lemma-separated-implies-morphism-separated}. Nous consacrons quelque
temps à démontrer des résultats usuels sur ces notions.

\begin{lemma}
\label{lemma-trivial-implications}
Soit $f : \mathcal{X} \to \mathcal{Y}$ un morphisme de champs algébriques.
\begin{enumerate}
\item Si $f$ est séparé, alors $f$ est quasi-séparé.
\item Si $f$ est DM, alors $f$ est quasi-DM.
\item Si $f$ est représentable par des espaces algébriques, alors $f$ est DM.
\end{enumerate}
\end{lemma}

\begin{proof}
Pour (1), remarquons qu'un morphisme propre d'espaces algébriques est
quasi-compact et quasi-séparé, voir Morphismes d'espaces, Définition
\ref{spaces-morphisms-definition-proper}. Pour (2), remarquons qu'un
morphisme non ramifié d'espaces algébriques est localement quasi-fini, voir
Morphismes d'espaces, Lemme
\ref{spaces-morphisms-lemma-unramified-quasi-finite}.
Enfin, (3) résulte du Lemme \ref{lemma-properties-diagonal-representable}.
\end{proof}

\begin{lemma}
\label{lemma-base-change-separated}
Tous les axiomes de séparation énumérés dans la Définition
\ref{definition-separated} sont stables par changement de base.
\end{lemma}

\begin{proof}
Soient $f : \mathcal{X} \to \mathcal{Y}$ et
$\mathcal{Y}' \to \mathcal{Y}$ des morphismes de champs algébriques.
Soit $f' : \mathcal{Y}' \times_\mathcal{Y} \mathcal{X} \to \mathcal{Y}'$
le changement de base de $f$ par $\mathcal{Y}' \to \mathcal{Y}$.
Alors $\Delta_{f'}$ est le changement de base de $\Delta_f$ par le morphisme
$\mathcal{X}' \times_{\mathcal{Y}'} \mathcal{X}' \to
\mathcal{X} \times_\mathcal{Y} \mathcal{X}$, voir Catégories, Lemme
\ref{categories-lemma-base-change-diagonal}. D'après les résultats de
Propriétés des champs,
section \ref{stacks-properties-section-properties-morphisms}, chacune des
propriétés de la diagonale employées dans la Définition
\ref{definition-separated} est stable par changement de base. Le lemme en
résulte.
\end{proof}

\begin{lemma}
\label{lemma-check-separated-covering}
Soit $f : \mathcal{X} \to \mathcal{Y}$ un morphisme de champs algébriques.
Soit $W \to \mathcal{Y}$ un morphisme surjectif, plat et localement de
présentation finie, où $W$ est un espace algébrique. Si le changement de base
$W \times_\mathcal{Y} \mathcal{X} \to W$ possède l'une des propriétés de
séparation de la Définition \ref{definition-separated}, alors $f$ la possède
également.
\end{lemma}

\begin{proof}
Notons $g : W \times_\mathcal{Y} \mathcal{X} \to W$ le changement de base.
Alors $\Delta_g$ est le changement de base de $\Delta_f$ par le morphisme
$q : W \times_\mathcal{Y} (\mathcal{X} \times_\mathcal{Y} \mathcal{X})
\to \mathcal{X} \times_\mathcal{Y} \mathcal{X}$. Comme $q$ est le changement
de base de $W \to \mathcal{Y}$, nous voyons que $q$ est représentable par des espaces
algébriques, surjectif, plat et localement de présentation finie. Le résultat
découle donc de Propriétés des champs, Lemme
\ref{stacks-properties-lemma-check-property-weak-covering}.
\end{proof}

\begin{lemma}
\label{lemma-change-of-base-separated}
Soit $S$ un schéma. La propriété d'être quasi-DM sur $S$, quasi-séparé sur
$S$ ou séparé sur $S$ (voir la Définition
\ref{definition-absolute-separated}) est stable par changement du schéma de
base, voir Champs algébriques, Définition
\ref{algebraic-definition-change-of-base}.
\end{lemma}

\begin{proof}
Cela résulte immédiatement du Lemme \ref{lemma-base-change-separated}.
\end{proof}

\begin{lemma}
\label{lemma-fibre-product-after-map}
Soient $f : \mathcal{X} \to \mathcal{Z}$, $g : \mathcal{Y} \to \mathcal{Z}$
et $\mathcal{Z} \to \mathcal{T}$ des morphismes de champs algébriques.
Considérons le morphisme induit
$i : \mathcal{X} \times_\mathcal{Z} \mathcal{Y} \to
\mathcal{X} \times_\mathcal{T} \mathcal{Y}$.
Alors
\begin{enumerate}
\item $i$ est représentable par des espaces algébriques et localement de type
fini ;
\item si $\Delta_{\mathcal{Z}/\mathcal{T}}$ est quasi-séparé, alors $i$ est
quasi-séparé ;
\item si $\Delta_{\mathcal{Z}/\mathcal{T}}$ est séparé, alors $i$ est séparé ;
\item si $\mathcal{Z} \to \mathcal{T}$ est DM, alors $i$ est non ramifié ;
\item si $\mathcal{Z} \to \mathcal{T}$ est quasi-DM, alors $i$ est localement
quasi-fini ;
\item si $\mathcal{Z} \to \mathcal{T}$ est séparé, alors $i$ est propre ;
\item si $\mathcal{Z} \to \mathcal{T}$ est quasi-séparé, alors $i$ est
quasi-compact et quasi-séparé.
\end{enumerate}
\end{lemma}

\begin{proof}
Le diagramme suivant
$$
\xymatrix{
\mathcal{X} \times_\mathcal{Z} \mathcal{Y} \ar[r]_i \ar[d] &
\mathcal{X} \times_\mathcal{T} \mathcal{Y} \ar[d] \\
\mathcal{Z} \ar[r]^-{\Delta_{\mathcal{Z}/\mathcal{T}}} \ar[r] &
\mathcal{Z} \times_\mathcal{T} \mathcal{Z}
}
$$
est un diagramme de $2$-produit fibré, voir Catégories, Lemme
\ref{categories-lemma-fibre-product-after-map}. Par conséquent, $i$ est le
changement de base du morphisme diagonal
$\Delta_{\mathcal{Z}/\mathcal{T}}$. Le lemme résulte donc du Lemme
\ref{lemma-properties-diagonal} et des résultats de Propriétés des champs,
section \ref{stacks-properties-section-properties-morphisms}.
\end{proof}

\begin{lemma}
\label{lemma-semi-diagonal}
Soit $\mathcal{T}$ un champ algébrique. Soit
$g : \mathcal{X} \to \mathcal{Y}$ un morphisme de champs algébriques sur
$\mathcal{T}$. Considérons le graphe
$i : \mathcal{X} \to \mathcal{X} \times_\mathcal{T} \mathcal{Y}$ de $g$. Alors
\begin{enumerate}
\item $i$ est représentable par des espaces algébriques et localement de type
fini ;
\item si $\mathcal{Y} \to \mathcal{T}$ est DM, alors $i$ est non ramifié ;
\item si $\mathcal{Y} \to \mathcal{T}$ est quasi-DM, alors $i$ est localement
quasi-fini ;
\item si $\mathcal{Y} \to \mathcal{T}$ est séparé, alors $i$ est propre ;
\item si $\mathcal{Y} \to \mathcal{T}$ est quasi-séparé, alors $i$ est
quasi-compact et quasi-séparé.
\end{enumerate}
\end{lemma}

\begin{proof}
Il s'agit du cas particulier du Lemme \ref{lemma-fibre-product-after-map}
appliqué au morphisme
$\mathcal{X} = \mathcal{X} \times_\mathcal{Y} \mathcal{Y} \to
\mathcal{X} \times_\mathcal{T} \mathcal{Y}$.
\end{proof}

\begin{lemma}
\label{lemma-section-immersion}
Soit $f : \mathcal{X} \to \mathcal{T}$ un morphisme de champs algébriques.
Soit $s : \mathcal{T} \to \mathcal{X}$ un morphisme tel que
$f \circ s$ soit $2$-isomorphe à $\text{id}_\mathcal{T}$. Alors
\begin{enumerate}
\item $s$ est représentable par des espaces algébriques et localement de type
fini ;
\item si $f$ est DM, alors $s$ est non ramifié ;
\item si $f$ est quasi-DM, alors $s$ est localement quasi-fini ;
\item si $f$ est séparé, alors $s$ est propre ;
\item si $f$ est quasi-séparé, alors $s$ est quasi-compact et quasi-séparé.
\end{enumerate}
\end{lemma}

\begin{proof}
Il s'agit du cas particulier du Lemme \ref{lemma-semi-diagonal} appliqué à
$g = s$ et $\mathcal{Y} = \mathcal{T}$ ; dans ce cas,
$i : \mathcal{T} \to \mathcal{T} \times_\mathcal{T} \mathcal{X}$
est $2$-isomorphe à $s$.
\end{proof}

\begin{lemma}
\label{lemma-composition-separated}
Tous les axiomes de séparation énumérés dans la Définition
\ref{definition-separated} sont stables par composition des morphismes.
\end{lemma}

\begin{proof}
Soient $f : \mathcal{X} \to \mathcal{Y}$ et
$g : \mathcal{Y} \to \mathcal{Z}$ des morphismes de champs algébriques
auxquels s'applique l'axiome considéré. La diagonale
$\Delta_{\mathcal{X}/\mathcal{Z}}$ est le composé
$$
\mathcal{X} \longrightarrow
\mathcal{X} \times_\mathcal{Y} \mathcal{X} \longrightarrow
\mathcal{X} \times_\mathcal{Z} \mathcal{X}.
$$
Notre axiome de séparation est défini en imposant à la diagonale une certaine
propriété $\mathcal{P}$. D'après le Lemme
\ref{lemma-fibre-product-after-map} ci-dessus, la seconde flèche possède elle
aussi cette propriété. Le lemme en résulte, puisque le composé de morphismes
représentables par des espaces algébriques et possédant la propriété
$\mathcal{P}$ possède encore la propriété $\mathcal{P}$, voir la discussion
générale dans Propriétés des champs,
section \ref{stacks-properties-section-properties-morphisms}
et Morphismes d'espaces, Lemmes
\ref{spaces-morphisms-lemma-composition-unramified},
\ref{spaces-morphisms-lemma-composition-quasi-finite},
\ref{spaces-morphisms-lemma-composition-proper},
\ref{spaces-morphisms-lemma-composition-quasi-compact} et
\ref{spaces-morphisms-lemma-composition-separated}.
\end{proof}

\begin{lemma}
\label{lemma-separated-over-separated}
Soit $f : \mathcal{X} \to \mathcal{Y}$ un morphisme de champs algébriques
sur le schéma de base $S$.
\begin{enumerate}
\item Si $\mathcal{Y}$ est DM sur $S$ et $f$ est DM, alors $\mathcal{X}$ est
DM sur $S$.
\item Si $\mathcal{Y}$ est quasi-DM sur $S$ et $f$ est quasi-DM, alors
$\mathcal{X}$ est quasi-DM sur $S$.
\item Si $\mathcal{Y}$ est séparé sur $S$ et $f$ est séparé, alors
$\mathcal{X}$ est séparé sur $S$.
\item Si $\mathcal{Y}$ est quasi-séparé sur $S$ et $f$ est quasi-séparé,
alors $\mathcal{X}$ est quasi-séparé sur $S$.
\item Si $\mathcal{Y}$ est DM et $f$ est DM, alors $\mathcal{X}$ est DM.
\item Si $\mathcal{Y}$ est quasi-DM et $f$ est quasi-DM, alors $\mathcal{X}$
est quasi-DM.
\item Si $\mathcal{Y}$ est séparé et $f$ est séparé, alors $\mathcal{X}$ est
séparé.
\item Si $\mathcal{Y}$ est quasi-séparé et $f$ est quasi-séparé, alors
$\mathcal{X}$ est quasi-séparé.
\end{enumerate}
\end{lemma}

\begin{proof}
Les assertions (1), (2), (3) et (4) résultent immédiatement du Lemme
\ref{lemma-composition-separated} et de la Définition
\ref{definition-absolute-separated}. Pour (5), (6), (7) et (8), considérons
$\mathcal{X}$ et $\mathcal{Y}$ comme des champs algébriques sur
$\Spec(\mathbf{Z})$ et appliquons le Lemme
\ref{lemma-composition-separated}. Détails omis.
\end{proof}

\noindent
Le lemme suivant diffère quelque peu de son analogue pour les espaces
algébriques. Pour les comparer, voir Morphismes d'espaces, Lemme
\ref{spaces-morphisms-lemma-compose-after-separated}.

\begin{lemma}
\label{lemma-compose-after-separated}
Soient $f : \mathcal{X} \to \mathcal{Y}$ et
$g : \mathcal{Y} \to \mathcal{Z}$ des morphismes de champs algébriques.
\begin{enumerate}
\item Si $g \circ f$ est DM, alors $f$ l'est aussi.
\item Si $g \circ f$ est quasi-DM, alors $f$ l'est aussi.
\item Si $g \circ f$ est séparé et $\Delta_g$ est séparé, alors $f$ est
séparé.
\item Si $g \circ f$ est quasi-séparé et si $\Delta_g$ est quasi-séparé,
alors $f$ est quasi-séparé.
\end{enumerate}
\end{lemma}

\begin{proof}
Considérons la factorisation
$$
\mathcal{X} \to
\mathcal{X} \times_\mathcal{Y} \mathcal{X} \to
\mathcal{X} \times_\mathcal{Z} \mathcal{X}
$$
du morphisme diagonal de $g \circ f$. Les deux morphismes sont représentables
par des espaces algébriques, voir les Lemmes
\ref{lemma-properties-diagonal} et \ref{lemma-fibre-product-after-map}. Ainsi,
pour tout schéma $T$ et tout morphisme
$T \to \mathcal{X} \times_\mathcal{Y} \mathcal{X}$, nous obtenons des
morphismes d'espaces algébriques
$$
A = \mathcal{X} \times_{(\mathcal{X} \times_\mathcal{Z} \mathcal{X})} T
\longrightarrow
B = (\mathcal{X} \times_\mathcal{Y} \mathcal{X})
\times_{(\mathcal{X} \times_\mathcal{Z} \mathcal{X})} T
\longrightarrow
T
$$
Si $g \circ f$ est DM (resp.\ quasi-DM), alors le composé $A \to T$ est non
ramifié (resp.\ localement quasi-fini). Par conséquent, $A \to B$ est non
ramifié (resp.\ localement quasi-fini) d'après Morphismes d'espaces, Lemme
\ref{spaces-morphisms-lemma-permanence-unramified}
(resp.\ Morphismes d'espaces, Lemme
\ref{spaces-morphisms-lemma-permanence-quasi-finite}). Considérons maintenant
le diagramme
$$
\xymatrix{
A \times_B T \ar[r] \ar[d] & T \ar[d] \\
A \ar[r] \ar[d] & B \ar[r] \ar[d] & T \ar[d] \\
\mathcal{X} \ar[r] &
\mathcal{X} \times_\mathcal{Y} \mathcal{X} \ar[r] &
\mathcal{X} \times_\mathcal{Z} \mathcal{X}
}
$$
dont tous les carrés sont $2$-cartésiens et où la flèche $T \to B$ provient
de la flèche $T \to \mathcal{X} \times_\mathcal{Y} \mathcal{X}$ de départ.
La flèche $A \times_B T \to T$ est non ramifiée (resp.\ localement
quasi-finie) d'après Morphismes d'espaces, Lemme
\ref{spaces-morphisms-lemma-base-change-unramified}
(resp.\ Morphismes d'espaces, Lemme
\ref{spaces-morphisms-lemma-base-change-quasi-finite}). Comme
$A \times_B T$ est égal à
$\mathcal{X} \times_{\mathcal{X} \times_\mathcal{Y} \mathcal{X}} T$, nous
concluons que $f$ est DM (resp.\ quasi-DM). Cela démontre (1) et (2).

\medskip\noindent
Démonstration de (4). Supposons que $g \circ f$ soit quasi-séparé et que
$\Delta_g$ soit quasi-séparé. Considérons la factorisation
$$
\mathcal{X} \to
\mathcal{X} \times_\mathcal{Y} \mathcal{X} \to
\mathcal{X} \times_\mathcal{Z} \mathcal{X}
$$
du morphisme diagonal de $g \circ f$. Les deux morphismes sont représentables
par des espaces algébriques, et le second est quasi-séparé, voir les Lemmes
\ref{lemma-properties-diagonal} et \ref{lemma-fibre-product-after-map}. Ainsi,
pour tout schéma $T$ et tout morphisme
$T \to \mathcal{X} \times_\mathcal{Y} \mathcal{X}$, nous obtenons des
morphismes d'espaces algébriques
$$
A = \mathcal{X} \times_{(\mathcal{X} \times_\mathcal{Z} \mathcal{X})} T
\longrightarrow
B = (\mathcal{X} \times_\mathcal{Y} \mathcal{X})
\times_{(\mathcal{X} \times_\mathcal{Z} \mathcal{X})} T
\longrightarrow
T
$$
tel que $B \to T$ soit quasi-séparé. Le composé $A \to T$ est quasi-compact
et quasi-séparé, puisque nous avons supposé que $g \circ f$ est quasi-séparé.
Par conséquent, $A \to B$ est quasi-séparé d'après Morphismes d'espaces,
Lemme \ref{spaces-morphisms-lemma-compose-after-separated}. De plus,
$A \to B$ est quasi-compact d'après Morphismes d'espaces, Lemme
\ref{spaces-morphisms-lemma-quasi-compact-permanence}. Ainsi, $f$ est
quasi-séparé.

\medskip\noindent
Démonstration de (3). Supposons que $g \circ f$ soit séparé et que $\Delta_g$
soit séparé. Considérons la factorisation
$$
\mathcal{X} \to
\mathcal{X} \times_\mathcal{Y} \mathcal{X} \to
\mathcal{X} \times_\mathcal{Z} \mathcal{X}
$$
du morphisme diagonal de $g \circ f$. Les deux morphismes sont représentables
par des espaces algébriques, et le second est séparé, voir les Lemmes
\ref{lemma-properties-diagonal} et \ref{lemma-fibre-product-after-map}. Ainsi,
pour tout schéma $T$ et tout morphisme
$T \to \mathcal{X} \times_\mathcal{Y} \mathcal{X}$, nous obtenons des
morphismes d'espaces algébriques
$$
A = \mathcal{X} \times_{(\mathcal{X} \times_\mathcal{Z} \mathcal{X})} T
\longrightarrow
B = (\mathcal{X} \times_\mathcal{Y} \mathcal{X})
\times_{(\mathcal{X} \times_\mathcal{Z} \mathcal{X})} T
\longrightarrow
T
$$
tel que $B \to T$ soit séparé. Le composé $A \to T$ est propre, puisque nous
avons supposé que $g \circ f$ est séparé. Par conséquent, $A \to B$ est propre
d'après Morphismes d'espaces, Lemme
\ref{spaces-morphisms-lemma-universally-closed-permanence}, ce qui signifie
que $f$ est séparé.
\end{proof}

\begin{lemma}
\label{lemma-separated-implies-morphism-separated}
Soit $\mathcal{X}$ un champ algébrique sur le schéma de base $S$.
\begin{enumerate}
\item
$\mathcal{X}$ est DM $\Leftrightarrow$
$\mathcal{X}$ est DM sur $S$.
\item
$\mathcal{X}$ est quasi-DM $\Leftrightarrow$
$\mathcal{X}$ est quasi-DM sur $S$.
\item Si $\mathcal{X}$ est séparé, alors $\mathcal{X}$ est séparé sur $S$.
\item Si $\mathcal{X}$ est quasi-séparé, alors $\mathcal{X}$ est quasi-séparé
sur $S$.
\end{enumerate}
Soit $f : \mathcal{X} \to \mathcal{Y}$ un morphisme de champs algébriques
sur le schéma de base $S$.
\begin{enumerate}
\item[(5)] Si $\mathcal{X}$ est DM sur $S$, alors $f$ est DM.
\item[(6)] Si $\mathcal{X}$ est quasi-DM sur $S$, alors $f$ est quasi-DM.
\item[(7)] Si $\mathcal{X}$ est séparé sur $S$ et si
$\Delta_{\mathcal{Y}/S}$ est séparé, alors $f$ est séparé.
\item[(8)] Si $\mathcal{X}$ est quasi-séparé sur $S$ et si
$\Delta_{\mathcal{Y}/S}$ est quasi-séparé, alors $f$ est quasi-séparé.
\end{enumerate}
\end{lemma}

\begin{proof}
Les assertions (5), (6), (7) et (8) résultent immédiatement du Lemme
\ref{lemma-compose-after-separated} et de Espaces, Définition
\ref{spaces-definition-separated}. Pour démontrer (3) et (4), considérons
$\mathcal{X}$ et $\mathcal{Y}$ comme des champs algébriques sur
$\Spec(\mathbf{Z})$ et appliquons le Lemme
\ref{lemma-compose-after-separated}. De même, pour démontrer (1) et (2),
considérons $\mathcal{X}$ comme un champ algébrique sur
$\Spec(\mathbf{Z})$ et considérons les morphismes
$$
\mathcal{X} \longrightarrow
\mathcal{X} \times_S \mathcal{X} \longrightarrow
\mathcal{X} \times_{\Spec(\mathbf{Z})} \mathcal{X}
$$
Les deux flèches sont représentables par des espaces algébriques. La seconde
est non ramifiée et localement quasi-finie, car elle est le changement de base
de l'immersion $\Delta_{S/\mathbf{Z}}$. Le composé est donc non ramifié
(resp.\ localement quasi-fini) si et seulement si la première flèche est non
ramifiée (resp.\ localement quasi-finie), voir Morphismes d'espaces, Lemmes
\ref{spaces-morphisms-lemma-composition-unramified} et
\ref{spaces-morphisms-lemma-permanence-unramified}
(resp.\ Morphismes d'espaces, Lemmes
\ref{spaces-morphisms-lemma-composition-quasi-finite} et
\ref{spaces-morphisms-lemma-permanence-quasi-finite}).
\end{proof}

\begin{lemma}
\label{lemma-properties-covering-imply-diagonal}
Soit $\mathcal{X}$ un champ algébrique.
Soit $W$ un espace algébrique et soit $f : W \to \mathcal{X}$
un morphisme surjectif, plat et localement de présentation finie.
\begin{enumerate}
\item Si $f$ est non ramifié (c'est-à-dire étale, ou encore si
$\mathcal{X}$ est de Deligne--Mumford), alors $\mathcal{X}$ est DM.
\item Si $f$ est localement quasi-fini, alors $\mathcal{X}$ est quasi-DM.
\end{enumerate}
\end{lemma}

\begin{proof}
Remarquons que, si $f$ est non ramifié, alors il est étale d'après
Morphismes d'espaces, Lemme
\ref{spaces-morphisms-lemma-unramified-flat-lfp-etale}.
Cela explique la parenthèse de (1).
Supposons que $f$ soit non ramifié (resp.\ localement quasi-fini). Il faut
montrer que
$\Delta_\mathcal{X} : \mathcal{X} \to \mathcal{X} \times \mathcal{X}$
est non ramifié (resp.\ localement quasi-fini). Remarquons que
$W \times W \to \mathcal{X} \times \mathcal{X}$ est lui aussi surjectif,
plat et localement de présentation finie. Il suffit donc de montrer que
$$
W \times_{\mathcal{X} \times \mathcal{X}, \Delta_\mathcal{X}} \mathcal{X}
=
W \times_\mathcal{X} W
\longrightarrow
W \times W
$$
est non ramifié (resp.\ localement quasi-fini), voir Propriétés des champs,
Lemme
\ref{stacks-properties-lemma-check-property-covering}.
Par hypothèse, le morphisme $\text{pr}_i : W \times_\mathcal{X} W \to W$
est non ramifié (resp.\ localement quasi-fini). La flèche affichée est donc
non ramifiée (resp.\ localement quasi-finie) d'après Morphismes d'espaces,
Lemme
\ref{spaces-morphisms-lemma-permanence-unramified}
(resp.\ Morphismes d'espaces, Lemme
\ref{spaces-morphisms-lemma-permanence-quasi-finite}).
\end{proof}

\begin{lemma}
\label{lemma-monomorphism-separated}
Un monomorphisme de champs algébriques est séparé et DM.
Il en va de même des immersions de champs algébriques.
\end{lemma}

\begin{proof}
Si $f : \mathcal{X} \to \mathcal{Y}$ est un monomorphisme de champs
algébriques, alors $\Delta_f$ est un isomorphisme, voir Propriétés des
champs, Lemme \ref{stacks-properties-lemma-monomorphism}.
Comme un isomorphisme d'espaces algébriques est propre et non ramifié,
$f$ est séparé et DM. La seconde assertion résulte de la première, puisqu'une
immersion est un monomorphisme, voir Propriétés des champs, Lemme
\ref{stacks-properties-lemma-immersion-monomorphism}.
\end{proof}

\begin{lemma}
\label{lemma-separation-properties-residual-gerbe}
Soit $\mathcal{X}$ un champ algébrique et soit $x \in |\mathcal{X}|$.
Supposons que la gerbe résiduelle $\mathcal{Z}_x$ de $\mathcal{X}$ en $x$
existe. Si $\mathcal{X}$ est DM, resp.\ quasi-DM, resp.\ séparé, resp.\
quasi-séparé, alors $\mathcal{Z}_x$ l'est aussi.
\end{lemma}

\begin{proof}
En effet, $\mathcal{Z}_x \to \mathcal{X}$ est un monomorphisme, donc est
DM et séparé d'après le Lemme \ref{lemma-monomorphism-separated}.
On conclut en appliquant le Lemme \ref{lemma-separated-over-separated}.
\end{proof}


















\section{Champs d'inertie}
\label{section-inertia}

\noindent
Le champ d'inertie (relatif) d'un champ en groupoïdes est défini dans
Champs, section \ref{stacks-section-the-inertia-stack}.
La construction proprement dite, dans le cadre des catégories fibrées, ainsi
que certaines de ses propriétés se trouvent dans Catégories, section
\ref{categories-section-inertia}.

\begin{lemma}
\label{lemma-inertia}
Soit $\mathcal{X}$ un champ algébrique. Alors le champ d'inertie
$\mathcal{I}_\mathcal{X}$ est lui aussi un champ algébrique. Le morphisme
$$
\mathcal{I}_\mathcal{X} \longrightarrow \mathcal{X}
$$
est représentable par des espaces algébriques et localement de type fini.
Plus généralement, soit $f : \mathcal{X} \to \mathcal{Y}$ un morphisme de
champs algébriques. Alors l'inertie relative
$\mathcal{I}_{\mathcal{X}/\mathcal{Y}}$ est un champ algébrique et le
morphisme
$$
\mathcal{I}_{\mathcal{X}/\mathcal{Y}} \longrightarrow \mathcal{X}
$$
est représentable par des espaces algébriques et localement de type fini.
\end{lemma}

\begin{proof}
D'après Catégories, Lemme
\ref{categories-lemma-inertia-fibred-category}, il existe des équivalences
$$
\mathcal{I}_\mathcal{X} \to
\mathcal{X} \times_{\Delta, \mathcal{X} \times_S \mathcal{X}, \Delta}
\mathcal{X}
\quad\text{et}\quad
\mathcal{I}_{\mathcal{X}/\mathcal{Y}} \to
\mathcal{X}
\times_{\Delta, \mathcal{X} \times_\mathcal{Y} \mathcal{X}, \Delta}
\mathcal{X}
$$
qui montrent que les champs d'inertie sont des champs algébriques.
Soit $T \to \mathcal{X}$ le morphisme donné par l'objet $x$ de la catégorie
fibre de $\mathcal{X}$ au-dessus de $T$. Nous obtenons alors un carré de
$2$-produit fibré
$$
\xymatrix{
\mathit{Isom}_\mathcal{X}(x, x) \ar[d] \ar[r] &
\mathcal{I}_\mathcal{X} \ar[d] \\
T \ar[r]^x & \mathcal{X}
}
$$
Cela résulte immédiatement de la définition de $\mathcal{I}_\mathcal{X}$.
Puisque $\mathit{Isom}_\mathcal{X}(x, x)$ est toujours un espace algébrique
localement de type fini sur $T$ (voir le Lemme
\ref{lemma-isom-locally-finite-type}), nous concluons que
$\mathcal{I}_\mathcal{X} \to \mathcal{X}$ est représentable par des espaces
algébriques et localement de type fini. Enfin, pour l'inertie relative, nous
obtenons
$$
\vcenter{
\xymatrix{
\mathit{Isom}_\mathcal{X}(x, x) \ar[d] &
K \ar[l] \ar[d] \ar[r] &
\mathcal{I}_{\mathcal{X}/\mathcal{Y}} \ar[d] \\
\mathit{Isom}_\mathcal{Y}(f(x), f(x)) &
T \ar[l]_-e \ar[r]^x & \mathcal{X}
}
}
$$
où les deux carrés sont des $2$-produits fibrés. Cela résulte de Catégories,
Lemme \ref{categories-lemma-relative-inertia-as-fibre-product}.
La flèche verticale de gauche est un morphisme d'espaces algébriques
localement de type fini sur $T$; elle est donc localement de type fini, voir
Morphismes d'espaces, Lemme
\ref{spaces-morphisms-lemma-permanence-finite-type}.
Ainsi, $K$ est un espace algébrique et $K \to T$ est localement de type fini.
Ceci démontre l'assertion relative à l'inertie relative.
\end{proof}

\begin{remark}
\label{remark-inertia-is-group-in-spaces}
Soit $f : \mathcal{X} \to \mathcal{Y}$ un morphisme de champs algébriques.
Dans Propriétés des champs, Remarque
\ref{stacks-properties-remark-representable-over}, nous avons vu que la
$2$-catégorie des morphismes $\mathcal{Z} \to \mathcal{X}$ représentables
par des espaces algébriques, ayant $\mathcal{X}$ pour but, est une catégorie.
Dans cette catégorie, le champ d'inertie de $\mathcal{X}/\mathcal{Y}$ est un
{\it objet en groupes}. Rappelons qu'un objet de
$\mathcal{I}_{\mathcal{X}/\mathcal{Y}}$
est simplement un couple $(x, \alpha)$, où $x$ est un objet de $\mathcal{X}$
et $\alpha$ un automorphisme de $x$, dans la catégorie fibre de
$\mathcal{X}$ à laquelle appartient $x$, tel que
$f(\alpha) = \text{id}$. La composition
$$
c :
\mathcal{I}_{\mathcal{X}/\mathcal{Y}}
\times_\mathcal{X} \mathcal{I}_{\mathcal{X}/\mathcal{Y}}
\longrightarrow
\mathcal{I}_{\mathcal{X}/\mathcal{Y}}
$$
est donnée sur les objets par la règle
$$
((x, \alpha), (x', \alpha'), \beta) \mapsto
(x, \alpha \circ \beta^{-1} \circ \alpha' \circ \beta)
$$
qui a un sens puisque $\beta : x \to x'$ est un isomorphisme dans la
catégorie fibre, par notre définition des produits fibrés. L'élément neutre
$e : \mathcal{X} \to \mathcal{I}_{\mathcal{X}/\mathcal{Y}}$ est donné par le
foncteur $x \mapsto (x, \text{id}_x)$. Nous omettons la vérification des
axiomes d'un objet en groupes.
\end{remark}

\noindent
Soit $f : \mathcal{X} \to \mathcal{Y}$ un morphisme de champs algébriques
et soit $\mathcal{I}_{\mathcal{X}/\mathcal{Y}}$ son champ d'inertie.
Soit $T$ un schéma et soit $x$ un objet de $\mathcal{X}$ au-dessus de $T$.
Posons $y = f(x)$. Nous avons vu, dans la démonstration du Lemme
\ref{lemma-inertia}, que, pour tout schéma $T$ et tout objet $x$ de
$\mathcal{X}$ au-dessus de $T$, il existe une suite exacte de faisceaux de
groupes
\begin{equation}
\label{equation-exact-sequence-isom}
0 \to
\mathit{Isom}_{\mathcal{X}/\mathcal{Y}}(x, x) \to
\mathit{Isom}_\mathcal{X}(x, x) \to
\mathit{Isom}_\mathcal{Y}(y, y)
\end{equation}
Les structures de groupe sur les deuxième et troisième termes sont celles
définies au Lemme \ref{lemma-isom-pseudo-torsor-aut}, et la suite donne un
sens au premier terme. Il existe en outre un carré cartésien canonique
$$
\xymatrix{
\mathit{Isom}_{\mathcal{X}/\mathcal{Y}}(x, x) \ar[d] \ar[r] &
\mathcal{I}_{\mathcal{X}/\mathcal{Y}} \ar[d] \\
T \ar[r]^x & \mathcal{X}
}
$$
En fait, la structure de groupe sur
$\mathcal{I}_{\mathcal{X}/\mathcal{Y}}$ considérée dans la Remarque
\ref{remark-inertia-is-group-in-spaces} induit la structure de groupe sur
$\mathit{Isom}_{\mathcal{X}/\mathcal{Y}}(x, x)$. Cela nous permet de définir
le faisceau
$\mathit{Isom}_{\mathcal{X}/\mathcal{Y}}(x, x)$
également pour les morphismes d'espaces algébriques vers $\mathcal{X}$.
Nous formalisons ceci dans la définition suivante.

\begin{definition}
\label{definition-isom}
Soit $f : \mathcal{X} \to \mathcal{Y}$ un morphisme de champs algébriques.
Soit $Z$ un espace algébrique.
\begin{enumerate}
\item Soit $x : Z \to \mathcal{X}$ un morphisme. Nous posons
$$
\mathit{Isom}_{\mathcal{X}/\mathcal{Y}}(x, x) =
Z \times_{x, \mathcal{X}} \mathcal{I}_{\mathcal{X}/\mathcal{Y}}
$$
Nous le munissons d'une structure d'espace algébrique en groupes sur $Z$ par
changement de base de la loi de composition considérée dans la Remarque
\ref{remark-inertia-is-group-in-spaces}. Nous appellerons parfois
$\mathit{Isom}_{\mathcal{X}/\mathcal{Y}}(x, x)$ le {\it faisceau relatif
d'automorphismes de $x$}.
\item Soient $x_1, x_2 : Z \to \mathcal{X}$ des morphismes. Posons
$y_i = f \circ x_i$. Soit $\alpha : y_1 \to y_2$ un $2$-morphisme.
Alors $\alpha$ détermine un morphisme
$\Delta^\alpha : Z \to Z \times_{y_1, \mathcal{Y}, y_2} Z$, et nous posons
$$
\mathit{Isom}_{\mathcal{X}/\mathcal{Y}}^\alpha(x_1, x_2) =
(Z \times_{x_1, \mathcal{X}, x_2} Z)
\times_{Z \times_{y_1, \mathcal{Y}, y_2} Z, \Delta^\alpha} Z.
$$
Nous appellerons parfois
$\mathit{Isom}_{\mathcal{X}/\mathcal{Y}}^\alpha(x_1, x_2)$
le {\it faisceau relatif des isomorphismes de $x_1$ vers $x_2$}.
\end{enumerate}
Si $\mathcal{Y} = \Spec(\mathbf{Z})$, ou plus généralement si
$\mathcal{Y}$ est un espace algébrique, nous employons les notations
$\mathit{Isom}_\mathcal{X}(x, x)$ et $\mathit{Isom}_\mathcal{X}(x_1, x_2)$
et nous parlons du {\it faisceau d'automorphismes de $x$} et du
{\it faisceau des isomorphismes de $x_1$ vers $x_2$}.
\end{definition}

\begin{lemma}
\label{lemma-isom-pseudo-torsor-aut-over-space}
Soit $f : \mathcal{X} \to \mathcal{Y}$ un morphisme de champs algébriques.
Soit $Z$ un espace algébrique et soient
$x_i : Z \to \mathcal{X}$, $i = 1, 2$, des morphismes. Alors
\begin{enumerate}
\item $\mathit{Isom}_{\mathcal{X}/\mathcal{Y}}(x_2, x_2)$
est un espace algébrique en groupes sur $Z$,
\item il existe une suite exacte de groupes
$$
0 \to \mathit{Isom}_{\mathcal{X}/\mathcal{Y}}(x_2, x_2)
\to \mathit{Isom}_\mathcal{X}(x_2, x_2)
\to \mathit{Isom}_\mathcal{Y}(f \circ x_2, f \circ x_2)
$$
\item il existe un morphisme d'espaces algébriques
$
\mathit{Isom}_\mathcal{X}(x_1, x_2)
\to \mathit{Isom}_\mathcal{Y}(f \circ x_1, f \circ x_2)
$
telle que, pour tout $2$-morphisme
$\alpha : f \circ x_1 \to f \circ x_2$, nous obtenions un diagramme
cartésien
$$
\xymatrix{
\mathit{Isom}_{\mathcal{X}/\mathcal{Y}}^\alpha(x_1, x_2) \ar[d] \ar[r] &
Z \ar[d]^\alpha \\
\mathit{Isom}_\mathcal{X}(x_1, x_2) \ar[r] &
\mathit{Isom}_\mathcal{Y}(f \circ x_1, f \circ x_2)
}
$$
\item pour tout $2$-morphisme $\alpha : f \circ x_1 \to f \circ x_2$,
l'espace algébrique
$\mathit{Isom}_{\mathcal{X}/\mathcal{Y}}^\alpha(x_1, x_2)$ est un
pseudo-torseur sous
$\mathit{Isom}_{\mathcal{X}/\mathcal{Y}}(x_2, x_2)$ sur $Z$.
\end{enumerate}
\end{lemma}

\begin{proof}
L'assertion (1) résulte de la Définition \ref{definition-isom}.
L'assertion (2) résulte de la suite exacte
(\ref{equation-exact-sequence-isom}) localement pour la topologie étale sur
$Z$. L'assertion (3) se voit en explicitant les définitions. Localement sur
$Z$ pour la topologie étale, l'assertion (4) se réduit à l'assertion (2) du
Lemme \ref{lemma-isom-pseudo-torsor-aut}.
\end{proof}

\begin{lemma}
\label{lemma-cartesian-square-inertia}
Soient $\pi : \mathcal{X} \to \mathcal{Y}$ et
$f : \mathcal{Y}' \to \mathcal{Y}$ des morphismes de champs algébriques.
Posons $\mathcal{X}' = \mathcal{X} \times_\mathcal{Y} \mathcal{Y}'$.
Alors les deux carrés du diagramme
$$
\xymatrix{
\mathcal{I}_{\mathcal{X}'/\mathcal{Y}'} \ar[r]
\ar[d]_{
\text{Catégories, Équation}\ (\ref{categories-equation-functorial})
} &
\mathcal{X}' \ar[r]_{\pi'} \ar[d] & \mathcal{Y}' \ar[d]^f \\
\mathcal{I}_{\mathcal{X}/\mathcal{Y}} \ar[r] &
\mathcal{X} \ar[r]^\pi & \mathcal{Y}
}
$$
sont des carrés de produit fibré.
\end{lemma}

\begin{proof}
Le champ d'inertie $\mathcal{I}_{\mathcal{X}'/\mathcal{Y}'}$ est défini
comme la catégorie des couples $(x', \alpha')$, où $x'$ est un objet de
$\mathcal{X}'$ et $\alpha'$ un automorphisme de $x'$ tel que
$\pi'(\alpha') = \text{id}$, voir Catégories, section
\ref{categories-section-inertia}. Supposons que $x'$ soit au-dessus du
schéma $U$ et s'envoie sur l'objet $x$ de $\mathcal{X}$. Par la construction
du $2$-produit fibré dans Catégories, Lemme
\ref{categories-lemma-2-product-categories-over-C}, nous voyons que
$x' = (U, x, y', \beta)$, où $y'$ est un objet de $\mathcal{Y}'$ au-dessus
de $U$ et $\beta$ un isomorphisme
$\beta : \pi(x) \to f(y')$ dans la catégorie fibre de $\mathcal{Y}$
au-dessus de $U$. Par la construction même du $2$-produit fibré,
l'automorphisme $\alpha'$ est un couple $(\alpha, \gamma)$, où $\alpha$ est
un automorphisme de $x$ au-dessus de $U$ et $\gamma$ un automorphisme de
$y'$ au-dessus de $U$, tels que $\alpha$ et $\gamma$ soient compatibles par
$\beta$. La condition $\pi'(\alpha') = \text{id}$ signifie que
$\gamma = \text{id}$; la compatibilité de $\alpha, \beta, \gamma$ équivaut
alors exactement à la condition $\pi(\alpha) = \text{id}$, c'est-à-dire que
$(x, \alpha)$ est un objet de $\mathcal{I}_{\mathcal{X}/\mathcal{Y}}$.
Nous voyons ainsi que le carré de gauche est un carré de produit fibré
(certains détails sont omis).
\end{proof}

\begin{lemma}
\label{lemma-monomorphism-cartesian-square-inertia}
Soit $f : \mathcal{X} \to \mathcal{Y}$ un monomorphisme de champs
algébriques. Alors le diagramme
$$
\xymatrix{
\mathcal{I}_\mathcal{X} \ar[r] \ar[d] &
\mathcal{X} \ar[d] \\
\mathcal{I}_\mathcal{Y} \ar[r] &
\mathcal{Y}
}
$$
est un carré de produit fibré.
\end{lemma}

\begin{proof}
Cela résulte immédiatement du fait que $f$ est pleinement fidèle (voir
Propriétés des champs, Lemme
\ref{stacks-properties-lemma-monomorphism}) et de la définition de l'inertie
dans Catégories, section \ref{categories-section-inertia}. En effet, un objet
de $\mathcal{I}_\mathcal{X}$ au-dessus d'un schéma $T$ est la même chose
qu'un couple $(x, \alpha)$ constitué d'un objet $x$ de $\mathcal{X}$
au-dessus de $T$ et d'un morphisme $\alpha : x \to x$ dans la catégorie
fibre de $\mathcal{X}$ au-dessus de $T$. Comme $f$ est pleinement fidèle,
se donner $\alpha$ revient à se donner un morphisme
$\beta : f(x) \to f(x)$ dans la catégorie fibre de $\mathcal{Y}$ au-dessus
de $T$. Nous pouvons donc considérer les objets de
$\mathcal{I}_\mathcal{X}$ au-dessus de $T$ comme les triplets
$((y, \beta), x, \gamma)$, où $y$ est un objet de $\mathcal{Y}$ au-dessus
de $T$, où $\beta : y \to y$ appartient à $\mathcal{Y}_T$, et où
$\gamma : y \to f(x)$ est un isomorphisme au-dessus de $T$, c'est-à-dire
comme les objets de $\mathcal{I}_\mathcal{Y} \times_\mathcal{Y}
\mathcal{X}$ au-dessus de $T$.
\end{proof}

\begin{lemma}
\label{lemma-presentation-inertia}
Soit $\mathcal{X}$ un champ algébrique. Soit
$[U/R] \to \mathcal{X}$ une présentation. Soit $G/U$ l'espace algébrique
en groupes stabilisateur associé au groupoïde $(U, R, s, t, c)$. Alors
$$
\xymatrix{
G \ar[d] \ar[r] & U \ar[d] \\
\mathcal{I}_\mathcal{X} \ar[r] & \mathcal{X}
}
$$
est un diagramme de produit fibré.
\end{lemma}

\begin{proof}
C'est immédiat d'après Groupoïdes en espaces, Lemme
\ref{spaces-groupoids-lemma-2-cartesian-inertia}.
\end{proof}









\section{Diagonales d'ordre supérieur}
\label{section-higher-diagonals}

\noindent
Soit $f : \mathcal{X} \to \mathcal{Y}$ un morphisme de champs algébriques.
Dans cette situation, il est naturel de considérer non seulement la diagonale
$$
\Delta_f : \mathcal{X} \to \mathcal{X} \times_\mathcal{Y} \mathcal{X}
$$
mais aussi la diagonale de la diagonale, c'est-à-dire le morphisme
$$
\Delta_{\Delta_f} :
\mathcal{X}
\longrightarrow
\mathcal{X} \times_{(\mathcal{X} \times_\mathcal{Y} \mathcal{X})} \mathcal{X}
$$
C'est pourquoi nous employons parfois la terminologie suivante. Nous appelons
$\Delta_{f, 0} = f$ la {\it diagonale d'ordre zéro},
$\Delta_{f, 1} = \Delta_f$ la {\it première diagonale} et
$\Delta_{f, 2} = \Delta_{\Delta_f}$ la {\it seconde diagonale}.
Remarquons que $\Delta_{f, 1}$ est représentable par des espaces algébriques
et localement de type fini, voir le Lemme \ref{lemma-properties-diagonal}.
Par conséquent, $\Delta_{f, 2}$ est représentable, est un monomorphisme, et
est localement de type fini, séparé et localement quasi-fini, voir le Lemme
\ref{lemma-properties-diagonal-representable}.

\medskip\noindent
Nous pouvons décrire la seconde diagonale à l'aide du champ d'inertie relatif.
En effet, le produit fibré
$\mathcal{X}
\times_{(\mathcal{X} \times_\mathcal{Y} \mathcal{X})} \mathcal{X}$
est équivalent au champ d'inertie relatif
$\mathcal{I}_{\mathcal{X}/\mathcal{Y}}$ d'après Catégories, Lemme
\ref{categories-lemma-inertia-fibred-category}. De plus, par cette
identification, la seconde diagonale devient la {\it section unité}
$$
\Delta_{f, 2} = e : \mathcal{X} \to \mathcal{I}_{\mathcal{X}/\mathcal{Y}}
$$
du champ d'inertie relatif. Par analogie avec ce qui se passe pour les
groupoïdes en espaces algébriques (Groupoïdes en espaces, Lemme
\ref{spaces-groupoids-lemma-diagonal}), nous avons les équivalences suivantes.

\begin{lemma}
\label{lemma-diagonal-diagonal}
Soit $f : \mathcal{X} \to \mathcal{Y}$ un morphisme de champs algébriques.
\begin{enumerate}
\item
Les conditions suivantes sont équivalentes :
\begin{enumerate}
\item $\mathcal{I}_{\mathcal{X}/\mathcal{Y}} \to \mathcal{X}$
est séparé,
\item $\Delta_{f, 1} = \Delta_f :
\mathcal{X} \to \mathcal{X} \times_\mathcal{Y} \mathcal{X}$
est séparé, et
\item $\Delta_{f, 2} = e :
\mathcal{X} \to \mathcal{I}_{\mathcal{X}/\mathcal{Y}}$
est une immersion fermée.
\end{enumerate}
\item
Les conditions suivantes sont équivalentes :
\begin{enumerate}
\item $\mathcal{I}_{\mathcal{X}/\mathcal{Y}} \to \mathcal{X}$
est quasi-séparé,
\item $\Delta_{f, 1} = \Delta_f :
\mathcal{X} \to \mathcal{X} \times_\mathcal{Y} \mathcal{X}$
est quasi-séparé, et
\item $\Delta_{f, 2} = e :
\mathcal{X} \to \mathcal{I}_{\mathcal{X}/\mathcal{Y}}$
est quasi-compact.
\end{enumerate}
\item
Les conditions suivantes sont équivalentes :
\begin{enumerate}
\item $\mathcal{I}_{\mathcal{X}/\mathcal{Y}} \to \mathcal{X}$
est localement séparé,
\item $\Delta_{f, 1} = \Delta_f :
\mathcal{X} \to \mathcal{X} \times_\mathcal{Y} \mathcal{X}$
est localement séparé, et
\item $\Delta_{f, 2} = e :
\mathcal{X} \to \mathcal{I}_{\mathcal{X}/\mathcal{Y}}$
est une immersion.
\end{enumerate}
\item
Les conditions suivantes sont équivalentes :
\begin{enumerate}
\item $\mathcal{I}_{\mathcal{X}/\mathcal{Y}} \to \mathcal{X}$
est non ramifié,
\item $f$ est DM.
\end{enumerate}
\item
Les conditions suivantes sont équivalentes :
\begin{enumerate}
\item $\mathcal{I}_{\mathcal{X}/\mathcal{Y}} \to \mathcal{X}$
est localement quasi-fini,
\item $f$ est quasi-DM.
\end{enumerate}
\end{enumerate}
\end{lemma}

\begin{proof}
Démonstration de (1), (2) et (3).
Choisissons un espace algébrique $U$ et un morphisme lisse surjectif
$U \to \mathcal{X}$. Alors
$G = U \times_\mathcal{X} \mathcal{I}_{\mathcal{X}/\mathcal{Y}}$
est un espace algébrique sur $U$ (Lemme \ref{lemma-inertia}). En fait,
$G$ est un espace algébrique en groupes sur $U$, par la loi de groupe sur
l'inertie relative construite dans la Remarque
\ref{remark-inertia-is-group-in-spaces}. De plus,
$G \to \mathcal{I}_{\mathcal{X}/\mathcal{Y}}$ est surjectif et lisse comme
changement de base de $U \to \mathcal{X}$. Enfin, le changement de base de
$e : \mathcal{X} \to \mathcal{I}_{\mathcal{X}/\mathcal{Y}}$
par $G \to \mathcal{I}_{\mathcal{X}/\mathcal{Y}}$ est la section unité
$U \to G$ de $G/U$. L'équivalence de (a) et (c) résulte donc de Groupoïdes
en espaces, Lemme \ref{spaces-groupoids-lemma-group-scheme-separated}.
Puisque $\Delta_{f, 2}$ est la diagonale de $\Delta_f$, nous avons
(b) $\Leftrightarrow$ (c) par définition.

\medskip\noindent
Démonstration de (4) et (5). Rappelons que (4)(b) signifie que $\Delta_f$
est non ramifié, tandis que (5)(b) signifie que $\Delta_f$ est localement
quasi-fini. Choisissons un schéma $Z$ et un morphisme
$a : Z \to \mathcal{X} \times_\mathcal{Y} \mathcal{X}$.
Alors $a = (x_1, x_2, \alpha)$, où $x_i : Z \to \mathcal{X}$ et où
$\alpha : f \circ x_1  \to f \circ x_2$ est un $2$-morphisme. Rappelons que
$$
\vcenter{
\xymatrix{
\mathit{Isom}_{\mathcal{X}/\mathcal{Y}}^\alpha(x_1, x_2)
\ar[d] \ar[r] &
Z \ar[d] \\
\mathcal{X} \ar[r]^{\Delta_f} &
\mathcal{X} \times_\mathcal{Y} \mathcal{X}
}
}
\quad\text{et}\quad
\vcenter{
\xymatrix{
\mathit{Isom}_{\mathcal{X}/\mathcal{Y}}(x_2, x_2)
\ar[d] \ar[r] &
Z \ar[d]^{x_2} \\
\mathcal{I}_{\mathcal{X}/\mathcal{Y}} \ar[r] &
\mathcal{X}
}
}
$$
sont des carrés cartésiens. D'après le Lemme
\ref{lemma-isom-pseudo-torsor-aut-over-space}, l'espace algébrique
$\mathit{Isom}_{\mathcal{X}/\mathcal{Y}}^\alpha(x_1, x_2)$ est un
pseudo-torseur sous
$\mathit{Isom}_{\mathcal{X}/\mathcal{Y}}(x_2, x_2)$ sur $Z$. Les
équivalences de (4) et (5) résultent donc de Groupoïdes en espaces, Lemme
\ref{spaces-groupoids-lemma-pseudo-torsor-implications}.
\end{proof}

\begin{lemma}
\label{lemma-second-diagonal}
Soit $f : \mathcal{X} \to \mathcal{Y}$ un morphisme de champs algébriques.
Les conditions suivantes sont équivalentes :
\begin{enumerate}
\item le morphisme $f$ est représentable par des espaces algébriques,
\item la seconde diagonale de $f$ est un isomorphisme,
\item le champ en groupes $ \mathcal{I}_{\mathcal{X}/\mathcal{Y}}$
est trivial sur $\mathcal X$, et
\item pour tout schéma $T$ et tout morphisme $x : T \to \mathcal{X}$,
le noyau de $\mathit{Isom}_\mathcal{X}(x, x) \to
\mathit{Isom}_\mathcal{Y}(f(x), f(x))$ est trivial.
\end{enumerate}
\end{lemma}

\begin{proof}
Montrons d'abord l'équivalence de (1) et (2). En effet, $f$ est représentable
par des espaces algébriques si et seulement si $f$ est fidèle, voir Champs
algébriques, Lemme
\ref{algebraic-lemma-characterize-representable-by-algebraic-spaces}.
D'autre part, $f$ est fidèle si et seulement si, pour tout objet $x$ de
$\mathcal{X}$ au-dessus d'un schéma $T$, le foncteur $f$ induit une injection
$\mathit{Isom}_\mathcal{X}(x, x) \to
\mathit{Isom}_\mathcal{Y}(f(x), f(x))$,
ce qui a lieu si et seulement si le noyau $K$ est trivial, ou encore si et
seulement si $e : T \to K$ est un isomorphisme pour tout
$x : T \to \mathcal{X}$. Puisque
$K = T \times_{x, \mathcal{X}} \mathcal{I}_{\mathcal{X}/\mathcal{Y}}$,
comme nous l'avons vu ci-dessus, cela démontre l'équivalence de (1) et (2).
Pour démontrer l'équivalence de (2) et (3), il suffit, d'après ce qui précède,
de remarquer qu'un champ en groupes est trivial si et seulement si sa section
unité est un isomorphisme. Enfin, l'équivalence de (3) et (4) résulte des
définitions : dans la démonstration du Lemme \ref{lemma-inertia}, nous avons
vu que le noyau de (4) correspond au produit fibré
$T \times_{x, \mathcal{X}} \mathcal{I}_{\mathcal{X}/\mathcal{Y}}$ sur $T$.
\end{proof}

\noindent
Ce lemme fournit la hiérarchie suivante pour les morphismes de champs
algébriques.

\begin{lemma}
\label{lemma-hierarchy}
Un morphisme $f : \mathcal{X} \to \mathcal{Y}$ de champs algébriques est
\begin{enumerate}
\item un monomorphisme si et seulement si $\Delta_{f, 1}$ est un
isomorphisme, et
\item représentable par des espaces algébriques si et seulement si
$\Delta_{f, 1}$ est un monomorphisme.
\end{enumerate}
De plus, la seconde diagonale $\Delta_{f, 2}$ est toujours un monomorphisme.
\end{lemma}

\begin{proof}
Rappelons, d'après Propriétés des champs, Lemme
\ref{stacks-properties-lemma-monomorphism}, qu'un morphisme de champs
algébriques est un monomorphisme si et seulement si sa diagonale est un
isomorphisme de champs. Le Lemme \ref{lemma-second-diagonal} peut donc se
reformuler en disant qu'un morphisme est représentable par des espaces
algébriques si sa diagonale est un monomorphisme. En particulier, il montre
que la condition (3) du Lemme
\ref{lemma-properties-diagonal-representable} est en fait une condition
nécessaire et suffisante : un morphisme de champs algébriques est
représentable par des espaces algébriques si et seulement si sa diagonale
est un monomorphisme.
\end{proof}

\begin{lemma}
\label{lemma-first-diagonal-separated-second-diagonal-closed}
Soit $f : \mathcal{X} \to \mathcal{Y}$ un morphisme de champs algébriques.
Alors
\begin{enumerate}
\item $\Delta_{f, 1}$ est séparé $\Leftrightarrow$
$\Delta_{f, 2}$ est une immersion fermée $\Leftrightarrow$
$\Delta_{f, 2}$ est propre $\Leftrightarrow$
$\Delta_{f, 2}$ est universellement fermé,
\item $\Delta_{f, 1}$ est quasi-séparé $\Leftrightarrow$
$\Delta_{f, 2}$ est de type fini $\Leftrightarrow$
$\Delta_{f, 2}$ est quasi-compact, et
\item $\Delta_{f, 1}$ est localement séparé $\Leftrightarrow$
$\Delta_{f, 2}$ est une immersion.
\end{enumerate}
\end{lemma}

\begin{proof}
Cela résulte des Lemmes \ref{lemma-representable-separated-diagonal-closed},
\ref{lemma-representable-quasi-separated-diagonal-quasi-compact} et
\ref{lemma-representable-locally-separated-diagonal-immersion}
appliqués à $\Delta_{f, 1}$.
\end{proof}

\noindent
Le lemme suivant est assez élégant et suggère peut-être une généralisation de
ces conditions aux champs algébriques supérieurs.

\begin{lemma}
\label{lemma-definition-separated}
Soit $f : \mathcal{X} \to \mathcal{Y}$ un morphisme de champs algébriques.
Alors
\begin{enumerate}
\item $f$ est séparé si et seulement si $\Delta_{f, 1}$ et
$\Delta_{f, 2}$ sont universellement fermés, et
\item $f$ est quasi-séparé si et seulement si $\Delta_{f, 1}$ et
$\Delta_{f, 2}$ sont quasi-compacts.
\item $f$ est quasi-DM si et seulement si $\Delta_{f, 1}$ et
$\Delta_{f, 2}$ sont localement quasi-finis.
\item $f$ est DM si et seulement si $\Delta_{f, 1}$ et $\Delta_{f, 2}$
sont non ramifiés.
\end{enumerate}
\end{lemma}

\begin{proof}
Démonstration de (1). Supposons $\Delta_{f, 2}$ et $\Delta_{f, 1}$
universellement fermés. Alors $\Delta_{f, 1}$ est séparé et universellement
fermé d'après le Lemme
\ref{lemma-first-diagonal-separated-second-diagonal-closed}. D'après
Morphismes d'espaces, Lemme
\ref{spaces-morphisms-lemma-universally-closed-quasi-compact}, et Champs
algébriques, Lemme
\ref{algebraic-lemma-representable-transformations-property-implication},
nous voyons que $\Delta_{f, 1}$ est quasi-compact. Il est donc quasi-compact,
séparé, universellement fermé et localement de type fini (d'après le Lemme
\ref{lemma-properties-diagonal}), donc propre. Cela démontre l'implication
``$\Leftarrow$'' de (1). La démonstration de l'implication réciproque est
omise.

\medskip\noindent
Démonstration de (2). Cela résulte immédiatement du Lemme
\ref{lemma-first-diagonal-separated-second-diagonal-closed}.

\medskip\noindent
Démonstration de (3). Cela résulte du fait que $\Delta_{f, 2}$ est toujours
localement quasi-fini, d'après le Lemme
\ref{lemma-properties-diagonal-representable} appliqué à
$\Delta_f = \Delta_{f, 1}$.

\medskip\noindent
Démonstration de (4). Cela résulte du fait que $\Delta_{f, 2}$ est toujours
non ramifié : le Lemme \ref{lemma-properties-diagonal-representable},
appliqué à $\Delta_f = \Delta_{f, 1}$, montre que $\Delta_{f, 2}$ est
localement de type fini et est un monomorphisme. Voir Compléments sur les
morphismes d'espaces, Lemme
\ref{spaces-more-morphisms-lemma-universally-injective-unramified}.
\end{proof}

\begin{lemma}
\label{lemma-separated-implies-isom}
Soit $f : \mathcal{X} \to \mathcal{Y}$ un morphisme séparé
(resp.\ quasi-séparé, resp.\ quasi-DM, resp.\ DM) de champs algébriques.
Alors
\begin{enumerate}
\item pour des espaces algébriques $T_i$, $i = 1, 2$, et des morphismes
$x_i : T_i \to \mathcal{X}$, en posant $y_i = f \circ x_i$, le morphisme
$$
T_1 \times_{x_1, \mathcal{X}, x_2} T_2 \longrightarrow
T_1 \times_{y_1, \mathcal{Y}, y_2} T_2
$$
est propre (resp.\ quasi-compact et quasi-séparé,
resp.\ localement quasi-fini, resp.\ non ramifié),
\item pour un espace algébrique $T$ et des morphismes
$x_i : T \to \mathcal{X}$, $i = 1, 2$, en posant
$y_i = f \circ x_i$, le morphisme
$$
\mathit{Isom}_\mathcal{X}(x_1, x_2) \longrightarrow
\mathit{Isom}_\mathcal{Y}(y_1, y_2)
$$
est propre (resp.\ quasi-compact et quasi-séparé,
resp.\ localement quasi-fini, resp.\ non ramifié).
\end{enumerate}
\end{lemma}

\begin{proof}
Démonstration de (1). Observons que le diagramme
$$
\xymatrix{
T_1 \times_{x_1, \mathcal{X}, x_2} T_2 \ar[d] \ar[r] &
T_1 \times_{y_1, \mathcal{Y}, y_2} T_2 \ar[d] \\
\mathcal{X} \ar[r] & \mathcal{X} \times_\mathcal{Y} \mathcal{X}
}
$$
est cartésien. L'assertion résulte donc du fait que $f$ est séparé
(resp.\ quasi-séparé, resp.\ quasi-DM, resp.\ DM) si et seulement si sa
diagonale est propre (resp.\ quasi-compacte et quasi-séparée,
resp.\ localement quasi-finie, resp.\ non ramifiée).

\medskip\noindent
Démonstration de (2). Cela résulte de l'égalité
$$
\mathit{Isom}_\mathcal{X}(x_1, x_2) =
(T \times_{x_1, \mathcal{X}, x_2} T) \times_{T \times T, \Delta_T} T
$$
et le morphisme de (2) est donc un changement de base du morphisme de (1).
\end{proof}








\section{Morphismes quasi-compacts}
\label{section-quasi-compact}

\noindent
Soit $f$ un morphisme de champs algébriques représentable par des espaces
algébriques. Dans Propriétés des champs, section
\ref{stacks-properties-section-properties-morphisms}
nous avons défini ce que signifie, pour $f$, être quasi-compact.
En voici une autre caractérisation.

\begin{lemma}
\label{lemma-characterize-representable-quasi-compact}
Soit $f : \mathcal{X} \to \mathcal{Y}$ un morphisme de champs algébriques
représentable par des espaces algébriques. Les conditions suivantes sont
équivalentes :
\begin{enumerate}
\item $f$ est quasi-compact (au sens de Propriétés des champs, section
\ref{stacks-properties-section-properties-morphisms}), et
\item pour tout champ algébrique quasi-compact $\mathcal{Z}$ et tout
morphisme $\mathcal{Z} \to \mathcal{Y}$, le champ algébrique
$\mathcal{Z} \times_\mathcal{Y} \mathcal{X}$ est quasi-compact.
\end{enumerate}
\end{lemma}

\begin{proof}
Supposons (1), et soit $\mathcal{Z} \to \mathcal{Y}$ un morphisme de champs
algébriques, avec $\mathcal{Z}$ quasi-compact. D'après Propriétés des champs,
Lemme \ref{stacks-properties-lemma-quasi-compact-stack}, il existe un schéma
quasi-compact $U$ et un morphisme lisse surjectif
$U \to \mathcal{Z}$. Puisque $f$ est représentable par des espaces
algébriques et quasi-compact, la définition montre que
$U \times_\mathcal{Y} \mathcal{X}$ est un espace algébrique et que
$U \times_\mathcal{Y} \mathcal{X} \to U$ est quasi-compact.
Ainsi, $U \times_\mathcal{Y} \mathcal{X}$ est un espace algébrique
quasi-compact. Le morphisme
$U \times_\mathcal{Y} \mathcal{X} \to
\mathcal{Z} \times_\mathcal{Y} \mathcal{X}$
est lisse et surjectif (comme changement de base du morphisme lisse et
surjectif $U \to \mathcal{Z}$). Par conséquent,
$\mathcal{Z} \times_\mathcal{Y} \mathcal{X}$ est quasi-compact par une
nouvelle application de Propriétés des champs, Lemme
\ref{stacks-properties-lemma-quasi-compact-stack}.

\medskip\noindent
Supposons (2). Soit $Z \to \mathcal{Y}$ un morphisme, où $Z$ est un schéma.
Il faut montrer que le morphisme d'espaces algébriques
$p : Z \times_\mathcal{Y} \mathcal{X} \to Z$ est quasi-compact.
Soit $U \subset Z$ un ouvert affine. Alors
$p^{-1}(U) = U \times_\mathcal{Y} \mathcal{X}$,
et l'espace algébrique $U \times_\mathcal{Y} \mathcal{X}$ est quasi-compact
par l'hypothèse (2). Par conséquent, $p$ est quasi-compact, voir Morphismes
d'espaces, Lemme \ref{spaces-morphisms-lemma-quasi-compact-local}.
\end{proof}

\noindent
Ceci motive la définition suivante.

\begin{definition}
\label{definition-quasi-compact}
Soit $f : \mathcal{X} \to \mathcal{Y}$ un morphisme de champs algébriques.
Nous disons que $f$ est {\it quasi-compact} si, pour tout champ algébrique
quasi-compact $\mathcal{Z}$ et tout morphisme
$\mathcal{Z} \to \mathcal{Y}$, le produit fibré
$\mathcal{Z} \times_\mathcal{Y} \mathcal{X}$ est quasi-compact.
\end{definition}

\noindent
D'après le Lemme \ref{lemma-characterize-representable-quasi-compact}
ci-dessus, cette définition coïncide avec la notion déjà existante pour les
morphismes de champs algébriques représentables par des espaces algébriques.
En particulier, cette notion coïncide avec celles déjà définies pour les
morphismes entre champs algébriques et schémas.

\begin{lemma}
\label{lemma-base-change-quasi-compact}
Le changement de base d'un morphisme quasi-compact de champs algébriques par
un morphisme quelconque de champs algébriques est quasi-compact.
\end{lemma}

\begin{proof}
Démonstration omise.
\end{proof}

\begin{lemma}
\label{lemma-composition-quasi-compact}
Le composé de deux morphismes quasi-compacts de champs algébriques est
quasi-compact.
\end{lemma}

\begin{proof}
Démonstration omise.
\end{proof}

\begin{lemma}
\label{lemma-closed-immersion-quasi-compact}
Une immersion fermée de champs algébriques est quasi-compacte.
\end{lemma}

\begin{proof}
Cela résulte du fait que les immersions sont toujours représentables et de
l'énoncé correspondant pour les immersions fermées d'espaces algébriques.
\end{proof}

\begin{lemma}
\label{lemma-surjection-from-quasi-compact}
Soit
$$
\xymatrix{
\mathcal{X} \ar[rr]_f \ar[rd]_p & &
\mathcal{Y} \ar[dl]^q \\
& \mathcal{Z}
}
$$
un diagramme $2$-commutatif de morphismes de champs algébriques.
Si $f$ est surjectif et $p$ quasi-compact, alors $q$ est quasi-compact.
\end{lemma}

\begin{proof}
Soit $\mathcal{T}$ un champ algébrique quasi-compact et soit
$\mathcal{T} \to \mathcal{Z}$ un morphisme. D'après
Propriétés des champs,
Lemme \ref{stacks-properties-lemma-base-change-surjective},
le morphisme
$\mathcal{T} \times_\mathcal{Z} \mathcal{X} \to
\mathcal{T} \times_\mathcal{Z} \mathcal{Y}$
est surjectif et, par hypothèse,
$\mathcal{T} \times_\mathcal{Z} \mathcal{X}$ est quasi-compact. Il s'ensuit
que $\mathcal{T} \times_\mathcal{Z} \mathcal{Y}$ est quasi-compact d'après
Propriétés des champs, Lemme
\ref{stacks-properties-lemma-quasi-compact-stack}.
\end{proof}

\begin{lemma}
\label{lemma-quasi-compact-permanence}
Soient $f : \mathcal{X} \to \mathcal{Y}$ et
$g : \mathcal{Y} \to \mathcal{Z}$ des morphismes de champs algébriques.
Si $g \circ f$ est quasi-compact et si $g$ est quasi-séparé,
alors $f$ est quasi-compact.
\end{lemma}

\begin{proof}
En effet, $f$ est le composé
$(1, f) : \mathcal{X} \to \mathcal{X} \times_\mathcal{Z} \mathcal{Y} \to
\mathcal{Y}$. Le premier morphisme est quasi-compact d'après le
Lemme \ref{lemma-section-immersion}, car c'est une section du morphisme
quasi-séparé
$\mathcal{X} \times_\mathcal{Z} \mathcal{Y} \to \mathcal{X}$
(un changement de base de $g$; voir le
Lemme \ref{lemma-base-change-separated}). Le second morphisme est
quasi-compact puisqu'il est le changement de base de $f$; voir le
Lemme \ref{lemma-base-change-quasi-compact}. Enfin, les composés de
morphismes quasi-compacts sont quasi-compacts; voir le
Lemme \ref{lemma-composition-quasi-compact}.
\end{proof}

\begin{lemma}
\label{lemma-quasi-compact-quasi-separated-permanence}
Soit $f : \mathcal{X} \to \mathcal{Y}$ un morphisme de champs algébriques.
\begin{enumerate}
\item Si $\mathcal{X}$ est quasi-compact et si $\mathcal{Y}$ est
quasi-séparé, alors $f$ est quasi-compact.
\item Si $\mathcal{X}$ est quasi-compact et quasi-séparé et si $\mathcal{Y}$
est quasi-séparé, alors $f$ est quasi-compact et quasi-séparé.
\item Un produit fibré de champs algébriques quasi-compacts et quasi-séparés
est quasi-compact et quasi-séparé.
\end{enumerate}
\end{lemma}

\begin{proof}
L'assertion (1) résulte du
Lemme \ref{lemma-quasi-compact-permanence}. L'assertion (2) résulte de (1)
et du Lemme \ref{lemma-compose-after-separated}. Pour (3), soient
$\mathcal{X} \to \mathcal{Y}$ et $\mathcal{Z} \to \mathcal{Y}$ des
morphismes de champs algébriques quasi-compacts et quasi-séparés. Alors
$\mathcal{X} \times_\mathcal{Y} \mathcal{Z} \to \mathcal{Z}$ est
quasi-compact et quasi-séparé, car il s'obtient par changement de base à
partir de $\mathcal{X} \to \mathcal{Y}$, en vertu de (2) et des Lemmes
\ref{lemma-base-change-quasi-compact} et
\ref{lemma-base-change-separated}. Par conséquent,
$\mathcal{X} \times_\mathcal{Y} \mathcal{Z}$ est quasi-compact et
quasi-séparé en tant que champ algébrique quasi-compact et quasi-séparé sur
$\mathcal{Z}$; voir les Lemmes \ref{lemma-separated-over-separated} et
\ref{lemma-composition-quasi-compact}.
\end{proof}

\begin{lemma}
\label{lemma-reach-points-scheme-theoretic-image}
Soit $f : \mathcal{X} \to \mathcal{Y}$ un morphisme quasi-compact de
champs algébriques. Soit $y \in |\mathcal{Y}|$ un point de l'adhérence de
l'image de $|f|$. Il existe un anneau de valuation $A$, de corps des
fractions $K$, et un diagramme commutatif
$$
\xymatrix{
\Spec(K) \ar[r] \ar[d] & \mathcal{X} \ar[d] \\
\Spec(A) \ar[r] & \mathcal{Y}
}
$$
tels que le point fermé de $\Spec(A)$ s'envoie sur $y$.
\end{lemma}

\begin{proof}
Choisissons un schéma affine $V$, un point $v \in V$ et un morphisme lisse
$V \to \mathcal{Y}$ qui envoie $v$ sur $y$. Considérons le diagramme de
changement de base
$$
\xymatrix{
V \times_\mathcal{Y} \mathcal{X} \ar[r] \ar[d]_g & \mathcal{X} \ar[d]^f \\
V \ar[r] & \mathcal{Y}
}
$$
Rappelons que $|V \times_\mathcal{Y} \mathcal{X}| \to
|V| \times_{|\mathcal{Y}|} |\mathcal{X}|$ est surjectif
(Propriétés des champs, Lemme
\ref{stacks-properties-lemma-points-cartesian}). Comme
$|V| \to |\mathcal{Y}|$ est ouvert (Propriétés des champs, Lemme
\ref{stacks-properties-lemma-topology-points}), nous en déduisons que $v$
appartient à l'adhérence de l'image de $|g|$. Il suffit donc de démontrer le
lemme pour le morphisme quasi-compact $g$
(Lemme \ref{lemma-base-change-quasi-compact}), ce que nous faisons au
paragraphe suivant.

\medskip\noindent
Supposons que $\mathcal{Y} = Y$ soit un schéma affine. Alors $\mathcal{X}$
est quasi-compact, puisque $f$ est quasi-compact
(Définition \ref{definition-quasi-compact}). Choisissons un schéma affine
$W$ et un morphisme lisse surjectif $W \to \mathcal{X}$. L'image de $|f|$
est alors l'image de $W \to Y$. D'après Morphismes, Lemme
\ref{morphisms-lemma-reach-points-scheme-theoretic-image}, nous pouvons
choisir un diagramme
$$
\xymatrix{
\Spec(K) \ar[r] \ar[d] & W \ar[d] \ar[r] & \mathcal{X} \ar[d] \\
\Spec(A) \ar[r] & Y \ar[r] & Y
}
$$
tel que le point fermé de $\Spec(A)$ s'envoie sur $y$.
En composant avec $W \to \mathcal{X}$, nous obtenons le diagramme voulu.
\end{proof}

\begin{lemma}
\label{lemma-check-quasi-compact-covering}
Soit $f : \mathcal{X} \to \mathcal{Y}$ un morphisme de champs algébriques.
Soit $W \to \mathcal{Y}$ un morphisme surjectif, plat et localement de
présentation finie, où $W$ est un espace algébrique. Si le changement de base
$W \times_\mathcal{Y} \mathcal{X} \to W$ est quasi-compact, alors $f$ est
quasi-compact.
\end{lemma}

\begin{proof}
Supposons que $W \times_\mathcal{Y} \mathcal{X} \to W$ soit quasi-compact.
Soit $\mathcal{Z} \to \mathcal{Y}$ un morphisme, où $\mathcal{Z}$ est un
champ algébrique quasi-compact. Choisissons un schéma $U$ et un morphisme
lisse surjectif $U \to W \times_\mathcal{Y} \mathcal{Z}$. Puisque
$U \to \mathcal{Z}$ est plat, surjectif et localement de présentation finie,
et puisque $\mathcal{Z}$ est quasi-compact, il existe un sous-schéma ouvert
quasi-compact $U' \subset U$ tel que $U' \to \mathcal{Z}$ soit surjectif.
Alors
$U' \times_\mathcal{Y} \mathcal{X} =
U' \times_W (W \times_\mathcal{Y} \mathcal{X})$
est quasi-compact par hypothèse et se surjecte sur
$\mathcal{Z} \times_\mathcal{Y} \mathcal{X}$. Par conséquent,
$\mathcal{Z} \times_\mathcal{Y} \mathcal{X}$ est quasi-compact, comme voulu.
\end{proof}











\section{Champs algébriques noethériens}
\label{section-noetherian}

\noindent
Nous avons déjà défini les champs algébriques localement noethériens dans
Propriétés des champs, section
\ref{stacks-properties-section-types-properties}.

\begin{definition}
\label{definition-noetherian}
Soit $\mathcal{X}$ un champ algébrique. Nous disons que $\mathcal{X}$ est
{\it noethérien} si $\mathcal{X}$ est quasi-compact, quasi-séparé et
localement noethérien.
\end{definition}

\noindent
Remarquons qu'un champ algébrique noethérien $\mathcal{X}$ n'est pas seulement
quasi-compact et localement noethérien, mais aussi quasi-séparé. Dans le
langage de la section \ref{section-higher-diagonals}, si l'on note
$p : \mathcal{X} \to \Spec(\mathbf{Z})$ le morphisme structural
``absolu'' (c'est-à-dire le morphisme structural de $\mathcal{X}$ considéré
comme champ algébrique sur $\mathbf{Z}$), alors
$$
\mathcal{X}\text{ noethérien}
\Leftrightarrow
\mathcal{X}\text{ localement noethérien et }
\Delta_{p, 0}, \Delta_{p, 1}, \Delta_{p, 2}
\text{ quasi-compacts}.
$$
Cela impliquera plus loin qu'un champ algébrique de type fini sur un champ
algébrique noethérien n'est pas nécessairement noethérien.

\begin{lemma}
\label{lemma-locally-closed-in-noetherian}
Soit $j : \mathcal{X} \to \mathcal{Y}$ une immersion de champs algébriques.
\begin{enumerate}
\item Si $\mathcal{Y}$ est localement noethérien, alors
$\mathcal{X}$ est localement noethérien et $j$ est quasi-compact.
\item Si $\mathcal{Y}$ est noethérien, alors $\mathcal{X}$ est noethérien.
\end{enumerate}
\end{lemma}

\begin{proof}
Choisissons un schéma $V$ et un morphisme lisse surjectif
$V \to \mathcal{Y}$. Alors $U = \mathcal{X} \times_\mathcal{Y} V$ est un
schéma et $U \to V$ est une immersion; voir Propriétés des champs,
Définition \ref{stacks-properties-definition-immersion}. Rappelons que
$\mathcal{Y}$ est localement noethérien si et seulement si $V$ l'est. Dans ce
cas, $U$ est lui aussi localement noethérien (Morphismes, Lemmes
\ref{morphisms-lemma-immersion-locally-finite-type} et
\ref{morphisms-lemma-finite-type-noetherian}), et $U \to V$ est quasi-compact
(Propriétés, Lemme \ref{properties-lemma-immersion-into-noetherian}). Cela
montre que $j$ est quasi-compact
(Lemme \ref{lemma-check-quasi-compact-covering}) et que $\mathcal{X}$ est
localement noethérien. Enfin, si $\mathcal{Y}$ est noethérien, ce qui précède
montre que $\mathcal{X}$ est quasi-compact et localement noethérien. Pour
achever la démonstration, observons que $j$ est séparé; par conséquent,
$\mathcal{X}$ est quasi-séparé puisque $\mathcal{Y}$ l'est, d'après le
Lemme \ref{lemma-separated-over-separated}.
\end{proof}

\begin{lemma}
\label{lemma-Noetherian-topology}
Soit $\mathcal{X}$ un champ algébrique.
\begin{enumerate}
\item Si $\mathcal{X}$ est localement noethérien, alors $|\mathcal{X}|$ est
un espace topologique localement noethérien.
\item Si $\mathcal{X}$ est quasi-compact et localement noethérien, alors
$|\mathcal{X}|$ est un espace topologique noethérien.
\end{enumerate}
\end{lemma}

\begin{proof}
Supposons $\mathcal{X}$ localement noethérien. Choisissons un schéma $U$ et un
morphisme lisse surjectif $U \to \mathcal{X}$. Comme $\mathcal{X}$ est
localement noethérien, $U$ l'est aussi. D'après Propriétés, Lemme
\ref{properties-lemma-Noetherian-topology}, cela signifie que $|U|$ est un
espace topologique localement noethérien. Puisque
$|U| \to |\mathcal{X}|$ est ouvert et surjectif, nous en déduisons que
$|\mathcal{X}|$ est localement noethérien d'après Topologie, Lemme
\ref{topology-lemma-image-Noetherian}. Cela démontre (1). Si $\mathcal{X}$ est
quasi-compact et localement noethérien, alors $|\mathcal{X}|$ est
quasi-compact et localement noethérien. Par conséquent, $|\mathcal{X}|$ est
noethérien d'après Topologie, Lemme
\ref{topology-lemma-quasi-compact-locally-Noetherian-Noetherian}.
\end{proof}

\begin{lemma}
\label{lemma-locally-Noetherian-quasi-sober}
Soit $\mathcal{X}$ un champ algébrique localement noethérien.
Alors $|\mathcal{X}|$ est quasi-sobre (Topologie, Définition
\ref{topology-definition-generic-point}).
\end{lemma}

\begin{proof}
Il faut démontrer que toute partie fermée irréductible
$T \subset |\mathcal{X}|$ possède un point générique. Choisissons un schéma
affine $U$ et un morphisme lisse $f : U \to \mathcal{X}$ tels que
$f^{-1}(T) \subset |U|$ soit non vide. Puisque $U$ est noethérien, la partie
fermée $f^{-1}(T)$ n'a qu'un nombre fini de composantes irréductibles
(Topologie, Lemme \ref{topology-lemma-Noetherian}). Écrivons
$f^{-1}(T) = Z_1 \cup \ldots \cup Z_n$ pour sa décomposition en composantes
irréductibles. Comme $f$ est ouvert, l'image de
$f|_{f^{-1}(T)} : f^{-1}(T) \to T$ contient une partie ouverte non vide de
$T$. Puisque $T$ est irréductible, cela signifie que $f(f^{-1}(T))$ est
dense. L'irréductibilité de $T$ entraîne alors que $f(Z_i)$ est dense pour un
certain $i$. Si $\xi_i \in Z_i$ est le point générique, $f(\xi_i)$ est donc
un point générique de $T$.
\end{proof}





\section{Morphismes affines}
\label{section-affine}

\noindent
Les morphismes affines de champs algébriques se définissent comme suit.

\begin{definition}
\label{definition-affine}
Un morphisme de champs algébriques est dit {\it affine} s'il est représentable
et affine au sens de Propriétés des champs, section
\ref{stacks-properties-section-properties-morphisms}.
\end{definition}

\noindent
Il nous est un peu plus commode de considérer un morphisme affine de champs
algébriques comme un morphisme de champs algébriques représentable par des
espaces algébriques et affine au sens de Propriétés des champs, section
\ref{stacks-properties-section-properties-morphisms}.
(Rappelons que, par défaut, ``représentable'' signifie dans le Champs Project
représentable par des schémas.) Puisque cela équivaut manifestement à la notion
qui vient d'être définie, nous emploierons cette caractérisation sans autre
mention. Démontrons quelques lemmes simples sur cette notion.

\begin{lemma}
\label{lemma-base-change-affine}
Soit $\mathcal{X} \to \mathcal{Y}$ un morphisme de champs algébriques.
Soit $\mathcal{Z} \to \mathcal{Y}$ un morphisme affine de champs algébriques.
Alors $\mathcal{Z} \times_\mathcal{Y} \mathcal{X} \to \mathcal{X}$ est un
morphisme affine de champs algébriques.
\end{lemma}

\begin{proof}
Cela résulte de la discussion de Propriétés des champs, section
\ref{stacks-properties-section-properties-morphisms}.
\end{proof}

\begin{lemma}
\label{lemma-composition-affine}
Les composés de morphismes affines de champs algébriques sont affines.
\end{lemma}

\begin{proof}
Cela résulte de la discussion de Propriétés des champs, section
\ref{stacks-properties-section-properties-morphisms}, et de Morphismes
d'espaces, Lemme \ref{spaces-morphisms-lemma-composition-affine}.
\end{proof}

\begin{lemma}
\label{lemma-affine-permanence}
Soit
$$
\xymatrix{
\mathcal{X} \ar[rr]_f \ar[rd]_a & & \mathcal{Y} \ar[dl]^b \\
& \mathcal{Z}
}
$$
un diagramme commutatif de morphismes de champs algébriques.
Si $a$ et $\Delta_b$ sont affines, alors $f$ est affine.
\end{lemma}

\begin{proof}
Le changement de base
$\text{pr}_2 : \mathcal{X} \times_\mathcal{Z} \mathcal{Y} \to \mathcal{Y}$
de $a$ est affine d'après le Lemme \ref{lemma-base-change-affine}. Le morphisme
$(1, f) : \mathcal{X} \to \mathcal{X} \times_\mathcal{Z} \mathcal{Y}$
est le changement de base de
$\Delta_b : \mathcal{Y} \to \mathcal{Y} \times_\mathcal{Z} \mathcal{Y}$
par le morphisme $\mathcal{X} \times_\mathcal{Z} \mathcal{Y} \to
\mathcal{Y} \times_\mathcal{Z} \mathcal{Y}$ (voir Catégories, section
\ref{categories-section-2-fibre-products}). Il est donc affine d'après le
Lemme \ref{lemma-base-change-affine}. Le composé
$f = \text{pr}_2 \circ (1, f)$ de morphismes affines est affine d'après le
Lemme \ref{lemma-composition-affine}, ce qui achève la démonstration.
\end{proof}





\section{Morphismes intégraux et finis}
\label{section-integral}

\noindent
Les morphismes intégraux et finis de champs algébriques se définissent comme
suit.

\begin{definition}
\label{definition-integral}
Soit $f : \mathcal{X} \to \mathcal{Y}$ un morphisme de champs algébriques.
\begin{enumerate}
\item Nous disons que $f$ est {\it intégral} si $f$ est représentable et
intégral au sens de Propriétés des champs, section
\ref{stacks-properties-section-properties-morphisms}.
\item Nous disons que $f$ est {\it fini} si $f$ est représentable et fini au
sens de Propriétés des champs, section
\ref{stacks-properties-section-properties-morphisms}.
\end{enumerate}
\end{definition}

\noindent
Il nous est un peu plus commode de considérer un morphisme intégral, resp.\
fini, de champs algébriques comme un morphisme de champs algébriques
représentable par des espaces algébriques et intégral, resp.\ fini, au sens de
Propriétés des champs, section
\ref{stacks-properties-section-properties-morphisms}.
(Rappelons que, par défaut, ``représentable'' signifie dans le Champs Project
représentable par des schémas.) Puisque cela équivaut manifestement à la notion
qui vient d'être définie, nous emploierons cette caractérisation sans autre
mention. Démontrons quelques lemmes simples sur cette notion.

\begin{lemma}
\label{lemma-base-change-integral}
Soit $\mathcal{X} \to \mathcal{Y}$ un morphisme de champs algébriques.
Soit $\mathcal{Z} \to \mathcal{Y}$ un morphisme intégral (ou fini) de champs
algébriques. Alors
$\mathcal{Z} \times_\mathcal{Y} \mathcal{X} \to \mathcal{X}$
est un morphisme intégral (ou fini) de champs algébriques.
\end{lemma}

\begin{proof}
Cela résulte de la discussion de Propriétés des champs, section
\ref{stacks-properties-section-properties-morphisms}.
\end{proof}

\begin{lemma}
\label{lemma-composition-integral}
Les composés de morphismes intégraux, resp.\ finis, de champs algébriques sont
intégraux, resp.\ finis.
\end{lemma}

\begin{proof}
Cela résulte de la discussion de Propriétés des champs, section
\ref{stacks-properties-section-properties-morphisms}, et de Morphismes
d'espaces, Lemme \ref{spaces-morphisms-lemma-composition-integral}.
\end{proof}




\section{Morphismes ouverts}
\label{section-open}

\noindent
Soit $f$ un morphisme de champs algébriques qui est représentable par
des espaces algébriques. Dans Propriétés des champs, section
\ref{stacks-properties-section-properties-morphisms}
nous avons défini ce que signifie pour $f$ d'être universellement ouvert.
En voici une autre caractérisation.

\begin{lemma}
\label{lemma-characterize-representable-universally-open}
Soit $f : \mathcal{X} \to \mathcal{Y}$ un morphisme de
champs algébriques qui est représentable par des espaces algébriques.
Les assertions suivantes sont équivalentes :
\begin{enumerate}
\item $f$ est universellement ouvert (au sens de Propriétés des champs,
section \ref{stacks-properties-section-properties-morphisms}) ;
\item pour tout morphisme de champs algébriques
$\mathcal{Z} \to \mathcal{Y}$, l'application d'espaces topologiques
$|\mathcal{Z} \times_\mathcal{Y} \mathcal{X}| \to |\mathcal{Z}|$ est ouverte.
\end{enumerate}
\end{lemma}

\begin{proof}
Supposons (1), et soit $\mathcal{Z} \to \mathcal{Y}$ comme en (2).
Choisissons un schéma $V$ et un morphisme lisse surjectif
$V \to \mathcal{Z}$. Par hypothèse, le morphisme d'espaces algébriques
$V \times_\mathcal{Y} \mathcal{X} \to V$ est universellement ouvert ;
en particulier, l'application
$|V \times_\mathcal{Y} \mathcal{X}| \to |V|$ est ouverte. D'après
Propriétés des champs, section \ref{stacks-properties-section-points},
dans le diagramme commutatif
$$
\xymatrix{
|V \times_\mathcal{Y} \mathcal{X}| \ar[r] \ar[d] &
|\mathcal{Z} \times_\mathcal{Y} \mathcal{X}| \ar[d] \\
|V| \ar[r] & |\mathcal{Z}|
}
$$
les flèches horizontales sont ouvertes et surjectives ; de plus,
$$
|V \times_\mathcal{Y} \mathcal{X}| \longrightarrow
|V| \times_{|\mathcal{Z}|} |\mathcal{Z} \times_\mathcal{Y} \mathcal{X}|
$$
est surjective. Puisque la flèche verticale de gauche est ouverte,
il s'ensuit que celle de droite est ouverte. Cela démontre (2).
L'implication (2) $\Rightarrow$ (1) résulte des définitions.
\end{proof}

\noindent
Nous pouvons donc adopter la définition naturelle suivante.

\begin{definition}
\label{definition-open}
Soit $f : \mathcal{X} \to \mathcal{Y}$ un morphisme de champs algébriques.
\begin{enumerate}
\item Nous disons que $f$ est {\it ouvert} si l'application d'espaces
topologiques $|\mathcal{X}| \to |\mathcal{Y}|$ est ouverte.
\item Nous disons que $f$ est {\it universellement ouvert} si, pour tout
morphisme de champs algébriques $\mathcal{Z} \to \mathcal{Y}$,
l'application d'espaces topologiques
$$
|\mathcal{Z} \times_\mathcal{Y} \mathcal{X}| \to |\mathcal{Z}|
$$
est ouverte, c'est-à-dire si le changement de base
$\mathcal{Z} \times_\mathcal{Y} \mathcal{X} \to \mathcal{Z}$ est ouvert.
\end{enumerate}
\end{definition}

\begin{lemma}
\label{lemma-base-change-universally-open}
Le changement de base d'un morphisme universellement ouvert de champs
algébriques par un morphisme quelconque de champs algébriques est
universellement ouvert.
\end{lemma}

\begin{proof}
C'est immédiat par définition.
\end{proof}

\begin{lemma}
\label{lemma-composition-universally-open}
Le composé de deux morphismes (universellement) ouverts de champs
algébriques est (universellement) ouvert.
\end{lemma}

\begin{proof}
Démonstration omise.
\end{proof}







\section{Morphismes submersifs}
\label{section-submersive}

\noindent
Soit $f$ un morphisme de champs algébriques qui est représentable par
des espaces algébriques. Dans Propriétés des champs, section
\ref{stacks-properties-section-properties-morphisms}
nous avons défini ce que signifie pour $f$ d'être universellement submersif.
En voici une autre caractérisation.

\begin{lemma}
\label{lemma-characterize-representable-universally-submersive}
Soit $f : \mathcal{X} \to \mathcal{Y}$ un morphisme de
champs algébriques qui est représentable par des espaces algébriques.
Les assertions suivantes sont équivalentes :
\begin{enumerate}
\item $f$ est universellement submersif (au sens de Propriétés des champs,
section \ref{stacks-properties-section-properties-morphisms}) ;
\item pour tout morphisme de champs algébriques
$\mathcal{Z} \to \mathcal{Y}$, l'application d'espaces topologiques
$|\mathcal{Z} \times_\mathcal{Y} \mathcal{X}| \to |\mathcal{Z}|$
est submersive.
\end{enumerate}
\end{lemma}

\begin{proof}
Supposons (1), et soit $\mathcal{Z} \to \mathcal{Y}$ comme en (2).
Choisissons un schéma $V$ et un morphisme lisse surjectif
$V \to \mathcal{Z}$. Par hypothèse, le morphisme d'espaces algébriques
$V \times_\mathcal{Y} \mathcal{X} \to V$ est universellement submersif ;
en particulier, l'application
$|V \times_\mathcal{Y} \mathcal{X}| \to |V|$ est submersive. D'après
Propriétés des champs, section \ref{stacks-properties-section-points},
dans le diagramme commutatif
$$
\xymatrix{
|V \times_\mathcal{Y} \mathcal{X}| \ar[r] \ar[d] &
|\mathcal{Z} \times_\mathcal{Y} \mathcal{X}| \ar[d] \\
|V| \ar[r] & |\mathcal{Z}|
}
$$
les flèches horizontales sont ouvertes et surjectives ; de plus,
$$
|V \times_\mathcal{Y} \mathcal{X}| \longrightarrow
|V| \times_{|\mathcal{Z}|} |\mathcal{Z} \times_\mathcal{Y} \mathcal{X}|
$$
est surjective. Puisque la flèche verticale de gauche est submersive,
il s'ensuit que celle de droite est submersive. Cela démontre (2).
L'implication (2) $\Rightarrow$ (1) résulte des définitions.
\end{proof}

\noindent
Nous pouvons donc adopter la définition naturelle suivante.

\begin{definition}
\label{definition-submersive}
Soit $f : \mathcal{X} \to \mathcal{Y}$ un morphisme de champs algébriques.
\begin{enumerate}
\item Nous disons que $f$ est {\it submersif}\footnote{Cette notion est
très différente de celle d'une submersion entre variétés différentielles.}
si l'application continue $|\mathcal{X}| \to |\mathcal{Y}|$ est submersive ;
voir Topologie, Définition \ref{topology-definition-submersive}.
\item Nous disons que $f$ est {\it universellement submersif} si, pour tout
morphisme de champs algébriques $\mathcal{Y}' \to \mathcal{Y}$,
le changement de base
$\mathcal{Y}' \times_\mathcal{Y} \mathcal{X} \to \mathcal{Y}'$
est submersif.
\end{enumerate}
\end{definition}

\noindent
Notons qu'un morphisme submersif est en particulier surjectif.

\begin{lemma}
\label{lemma-base-change-universally-submersive}
Le changement de base d'un morphisme universellement submersif de champs
algébriques par un morphisme quelconque de champs algébriques est
universellement submersif.
\end{lemma}

\begin{proof}
C'est immédiat par définition.
\end{proof}

\begin{lemma}
\label{lemma-composition-universally-submersive}
Le composé de deux morphismes (universellement) submersifs de champs
algébriques est (universellement) submersif.
\end{lemma}

\begin{proof}
Démonstration omise.
\end{proof}












\section{Morphismes universellement fermés}
\label{section-universally-closed}

\noindent
Soit $f$ un morphisme de champs algébriques qui est représentable par
des espaces algébriques. Dans Propriétés des champs, section
\ref{stacks-properties-section-properties-morphisms}
nous avons défini ce que signifie pour $f$ d'être universellement fermé.
En voici une autre caractérisation.

\begin{lemma}
\label{lemma-characterize-representable-universally-closed}
Soit $f : \mathcal{X} \to \mathcal{Y}$ un morphisme de
champs algébriques qui est représentable par des espaces algébriques.
Les assertions suivantes sont équivalentes :
\begin{enumerate}
\item $f$ est universellement fermé (au sens de Propriétés des champs,
section \ref{stacks-properties-section-properties-morphisms}) ;
\item pour tout morphisme de champs algébriques
$\mathcal{Z} \to \mathcal{Y}$, l'application d'espaces topologiques
$|\mathcal{Z} \times_\mathcal{Y} \mathcal{X}| \to |\mathcal{Z}|$ est fermée.
\end{enumerate}
\end{lemma}

\begin{proof}
Supposons (1), et soit $\mathcal{Z} \to \mathcal{Y}$ comme en (2).
Choisissons un schéma $V$ et un morphisme lisse surjectif
$V \to \mathcal{Z}$. Par hypothèse, le morphisme d'espaces algébriques
$V \times_\mathcal{Y} \mathcal{X} \to V$ est universellement fermé ;
en particulier, l'application
$|V \times_\mathcal{Y} \mathcal{X}| \to |V|$ est fermée. D'après
Propriétés des champs, section \ref{stacks-properties-section-points},
dans le diagramme commutatif
$$
\xymatrix{
|V \times_\mathcal{Y} \mathcal{X}| \ar[r] \ar[d] &
|\mathcal{Z} \times_\mathcal{Y} \mathcal{X}| \ar[d] \\
|V| \ar[r] & |\mathcal{Z}|
}
$$
les flèches horizontales sont ouvertes et surjectives ; de plus,
$$
|V \times_\mathcal{Y} \mathcal{X}| \longrightarrow
|V| \times_{|\mathcal{Z}|} |\mathcal{Z} \times_\mathcal{Y} \mathcal{X}|
$$
est surjective. Puisque la flèche verticale de gauche est fermée,
il s'ensuit que celle de droite est fermée. Cela démontre (2).
L'implication (2) $\Rightarrow$ (1) résulte des définitions.
\end{proof}

\noindent
Nous pouvons donc adopter la définition naturelle suivante.

\begin{definition}
\label{definition-closed}
Soit $f : \mathcal{X} \to \mathcal{Y}$ un morphisme de champs algébriques.
\begin{enumerate}
\item Nous disons que $f$ est {\it fermé} si l'application d'espaces
topologiques $|\mathcal{X}| \to |\mathcal{Y}|$ est fermée.
\item Nous disons que $f$ est {\it universellement fermé} si, pour tout
morphisme de champs algébriques $\mathcal{Z} \to \mathcal{Y}$,
l'application d'espaces topologiques
$$
|\mathcal{Z} \times_\mathcal{Y} \mathcal{X}| \to |\mathcal{Z}|
$$
est fermée, c'est-à-dire si le changement de base
$\mathcal{Z} \times_\mathcal{Y} \mathcal{X} \to \mathcal{Z}$ est fermé.
\end{enumerate}
\end{definition}

\begin{lemma}
\label{lemma-base-change-universally-closed}
Le changement de base d'un morphisme universellement fermé de champs
algébriques par un morphisme quelconque de champs algébriques est
universellement fermé.
\end{lemma}

\begin{proof}
C'est immédiat par définition.
\end{proof}

\begin{lemma}
\label{lemma-composition-universally-closed}
Le composé de deux morphismes (universellement) fermés de champs
algébriques est (universellement) fermé.
\end{lemma}

\begin{proof}
Démonstration omise.
\end{proof}

\begin{lemma}
\label{lemma-universally-closed-local}
Soit $f : \mathcal{X} \to \mathcal{Y}$ un morphisme de champs algébriques.
Les assertions suivantes sont équivalentes :
\begin{enumerate}
\item $f$ est universellement fermé ;
\item pour tout schéma $Z$ et tout morphisme $Z \to \mathcal{Y}$,
la projection $|Z \times_\mathcal{Y} \mathcal{X}| \to |Z|$
est fermée ;
\item pour tout schéma affine $Z$ et tout morphisme $Z \to \mathcal{Y}$,
la projection $|Z \times_\mathcal{Y} \mathcal{X}| \to |Z|$
est fermée ;
\item il existe un espace algébrique $V$ et un morphisme lisse surjectif
$V \to \mathcal{Y}$ tels que $V \times_\mathcal{Y} \mathcal{X} \to V$
soit un morphisme universellement fermé de champs algébriques.
\end{enumerate}
\end{lemma}

\begin{proof}
Nous omettons de démontrer que (1) entraîne (2) et que (2) entraîne (3).

\medskip\noindent
Supposons (3). Choisissons un morphisme lisse surjectif
$V \to \mathcal{Y}$. Nous allons montrer que
$V \times_\mathcal{Y} \mathcal{X} \to V$ est un morphisme
universellement fermé de champs algébriques. Soit
$\mathcal{Z} \to V$ un morphisme d'un champ algébrique vers $V$.
Soit $W \to \mathcal{Z}$ un morphisme lisse surjectif, où
$W = \coprod W_i$ est une réunion disjointe de schémas affines.
Nous avons alors le diagramme commutatif suivant :
$$
\xymatrix{
\coprod_i |W_i \times_\mathcal{Y} \mathcal{X}| \ar@{=}[r] \ar[d] &
|W \times_\mathcal{Y} \mathcal{X}| \ar[r] \ar[d] &
|\mathcal{Z} \times_\mathcal{Y} \mathcal{X}| \ar[d] \ar@{=}[r] &
|\mathcal{Z} \times_V (V \times_\mathcal{Y} \mathcal{X})| \ar[ld] \\
\coprod |W_i| \ar@{=}[r] &
|W| \ar[r] &
|\mathcal{Z}|
}
$$
Il faut montrer que la flèche sud-est est fermée. Les flèches
horizontales du milieu sont surjectives et ouvertes
(Propriétés des champs, Lemme
\ref{stacks-properties-lemma-topology-points}).
En vertu de (3) et du fait que les $W_i$ sont affines, les flèches
verticales de gauche sont fermées. Il s'ensuit que la flèche verticale
de droite est fermée.

\medskip\noindent
Supposons (4). Nous allons montrer que $f$ est universellement fermé.
Soit $\mathcal{Z} \to \mathcal{Y}$ un morphisme de champs algébriques.
Considérons le diagramme
$$
\xymatrix{
|(V \times_\mathcal{Y} \mathcal{Z})
\times_V (V \times_\mathcal{Y} \mathcal{X})| \ar@{=}[r] \ar[rd] &
|V \times_\mathcal{Y} \mathcal{X}| \ar[r] \ar[d] &
|\mathcal{Z} \times_\mathcal{Y} \mathcal{X}| \ar[d] \\
 &
|V \times_\mathcal{Y} \mathcal{Z}| \ar[r] &
|\mathcal{Z}|
}
$$
La flèche sud-ouest est fermée par hypothèse. Les flèches horizontales
sont surjectives et ouvertes, car les morphismes de champs algébriques
correspondants sont lisses et surjectifs (voir la référence ci-dessus).
Il s'ensuit que la flèche verticale de droite est fermée.
\end{proof}










\section{Morphismes universellement injectifs}
\label{section-universally-injective}

\noindent
Soit $f$ un morphisme de champs algébriques qui est représentable par
des espaces algébriques. Dans Propriétés des champs, section
\ref{stacks-properties-section-properties-morphisms}
nous avons défini ce que signifie pour $f$ d'être universellement injectif
(ou radiciel, dans la terminologie classique des morphismes représentables).
En voici une autre caractérisation.

\begin{lemma}
\label{lemma-characterize-representable-universally-injective}
Soit $f : \mathcal{X} \to \mathcal{Y}$ un morphisme de
champs algébriques qui est représentable par des espaces algébriques.
Les assertions suivantes sont équivalentes :
\begin{enumerate}
\item $f$ est universellement injectif (au sens de Propriétés des champs,
section \ref{stacks-properties-section-properties-morphisms}) ;
\item pour tout morphisme de champs algébriques
$\mathcal{Z} \to \mathcal{Y}$, l'application
$|\mathcal{Z} \times_\mathcal{Y} \mathcal{X}| \to |\mathcal{Z}|$
est injective.
\end{enumerate}
\end{lemma}

\begin{proof}
Supposons (1), et soit $\mathcal{Z} \to \mathcal{Y}$ comme en (2).
Choisissons un schéma $V$ et un morphisme lisse surjectif
$V \to \mathcal{Z}$. Par hypothèse, le morphisme d'espaces algébriques
$V \times_\mathcal{Y} \mathcal{X} \to V$ est universellement injectif ;
en particulier, l'application
$|V \times_\mathcal{Y} \mathcal{X}| \to |V|$ est injective. D'après
Propriétés des champs, section \ref{stacks-properties-section-points},
dans le diagramme commutatif
$$
\xymatrix{
|V \times_\mathcal{Y} \mathcal{X}| \ar[r] \ar[d] &
|\mathcal{Z} \times_\mathcal{Y} \mathcal{X}| \ar[d] \\
|V| \ar[r] & |\mathcal{Z}|
}
$$
les flèches horizontales sont ouvertes et surjectives ; de plus,
$$
|V \times_\mathcal{Y} \mathcal{X}| \longrightarrow
|V| \times_{|\mathcal{Z}|} |\mathcal{Z} \times_\mathcal{Y} \mathcal{X}|
$$
est surjective. Puisque la flèche verticale de gauche est injective,
il s'ensuit que celle de droite est injective. Cela démontre (2).
L'implication (2) $\Rightarrow$ (1) résulte des définitions.
\end{proof}

\noindent
Nous pouvons donc adopter la définition naturelle suivante.

\begin{definition}
\label{definition-universally-injective}
Soit $f : \mathcal{X} \to \mathcal{Y}$ un morphisme de champs algébriques.
Nous disons que $f$ est {\it universellement injectif} si, pour tout
morphisme de champs algébriques $\mathcal{Z} \to \mathcal{Y}$,
l'application
$$
|\mathcal{Z} \times_\mathcal{Y} \mathcal{X}| \to |\mathcal{Z}|
$$
est injective.
\end{definition}

\begin{lemma}
\label{lemma-base-change-universally-injective}
Le changement de base d'un morphisme universellement injectif de champs
algébriques par un morphisme quelconque de champs algébriques est
universellement injectif.
\end{lemma}

\begin{proof}
C'est immédiat par définition.
\end{proof}

\begin{lemma}
\label{lemma-composition-universally-injective}
Le composé de deux morphismes universellement injectifs de champs
algébriques est universellement injectif.
\end{lemma}

\begin{proof}
Démonstration omise.
\end{proof}

\begin{lemma}
\label{lemma-universally-injective}
Soit $f : \mathcal{X} \to \mathcal{Y}$ un morphisme de champs algébriques.
Les assertions suivantes sont équivalentes :
\begin{enumerate}
\item $f$ est universellement injectif ;
\item $\Delta : \mathcal{X} \to \mathcal{X} \times_\mathcal{Y} \mathcal{X}$
est surjective ;
\item pour tout corps algébriquement clos, tous
$x_1, x_2 : \Spec(k) \to \mathcal{X}$ et toute $2$-flèche
$\beta : f \circ x_1 \to f \circ x_2$, il existe une
$2$-flèche $\alpha : x_1 \to x_2$ telle que
$\beta = \text{id}_f \star \alpha$.
\end{enumerate}
\end{lemma}

\begin{proof}
(1) $\Rightarrow$ (2). Si $f$ est universellement injectif, la première
projection
$|\mathcal{X} \times_\mathcal{Y} \mathcal{X}| \to |\mathcal{X}|$
est injective, ce qui implique que $|\Delta|$ est surjective.

\medskip\noindent
(2) $\Rightarrow$ (1). Supposons $\Delta$ surjective. Tout changement
de base de $\Delta$ est alors surjectif (voir Propriétés des champs, section
\ref{stacks-properties-section-surjective}).
Comme la diagonale d'un changement de base de $f$ est un changement
de base de $\Delta$, il suffit de montrer que
$|\mathcal{X}| \to |\mathcal{Y}|$ est injective.
Dans le cas contraire, Propriétés des champs, Lemme
\ref{stacks-properties-lemma-points-cartesian}
montre que la première projection
$|\mathcal{X} \times_\mathcal{Y} \mathcal{X}| \to |\mathcal{X}|$
n'est pas injective. Cela signifie évidemment que $|\Delta|$ n'est
pas surjective.

\medskip\noindent
(3) $\Rightarrow$ (2). Soit
$t \in |\mathcal{X} \times_\mathcal{Y} \mathcal{X}|$.
Nous pouvons représenter $t$ par un morphisme
$t : \Spec(k) \to \mathcal{X} \times_\mathcal{Y} \mathcal{X}$
où $k$ est un corps algébriquement clos.
D'après notre construction des $2$-produits fibrés, nous pouvons
représenter $t$ par $(x_1, x_2, \beta)$, où
$x_1, x_2 : \Spec(k) \to \mathcal{X}$ et
$\beta : f \circ x_1 \to f \circ x_2$ est un $2$-morphisme.
L'assertion (3) implique alors qu'il existe un $2$-morphisme
$\alpha : x_1 \to x_2$ dont l'image est $\beta$.
Cela signifie exactement que $\Delta(x_1) = (x_1, x_1, \text{id})$
est isomorphe à $t$. L'assertion (2) est donc satisfaite.

\medskip\noindent
(2) $\Rightarrow$ (3). Soient
$x_1, x_2 : \Spec(k) \to \mathcal{X}$ des morphismes, où $k$ est un
corps algébriquement clos. Soit
$\beta : f \circ x_1 \to f \circ x_2$ un $2$-morphisme.
Comme au paragraphe précédent, nous obtenons un morphisme
$t = (x_1, x_2, \beta) :
\Spec(k) \to \mathcal{X} \times_\mathcal{Y} \mathcal{X}$.
D'après le Lemme \ref{lemma-properties-diagonal},
$$
T = \mathcal{X}
\times_{\Delta, \mathcal{X} \times_\mathcal{Y} \mathcal{X}, t} \Spec(k)
$$
est un espace algébrique localement de type fini sur $\Spec(k)$.
La condition (2) implique que $T$ est non vide. Comme $k$ est
algébriquement clos, il existe alors un point à valeurs dans $k$
appartenant à $T$.
En explicitant les définitions, cela signifie qu'il existe un morphisme
$\alpha : x_1 \to x_2$ dans $\Mor(\Spec(k), \mathcal{X})$
tel que $\beta = \text{id}_f \star \alpha$.
\end{proof}

\begin{lemma}
\label{lemma-universally-injective-point}
Soit $f : \mathcal{X} \to \mathcal{Y}$ un morphisme universellement
injectif de champs algébriques. Soit
$y : \Spec(k) \to \mathcal{Y}$ un morphisme, où $k$ est un corps
algébriquement clos. Si $y$ appartient à l'image de
$|\mathcal{X}| \to |\mathcal{Y}|$, il existe un morphisme
$x : \Spec(k) \to \mathcal{X}$ tel que $y = f \circ x$.
\end{lemma}

\begin{proof}
Remarquons d'abord que ce lemme n'est pas trivial : l'hypothèse selon
laquelle $y$ appartient à l'image de $|f|$ signifie seulement que l'on
peut relever $y$ en un morphisme vers $\mathcal{X}$ après avoir
éventuellement remplacé $k$ par un corps d'extension. Pour démontrer le
lemme, nous pouvons effectuer le changement de base de $f$ par $y$.
Nous pouvons donc supposer que $\mathcal{X}$ est un champ algébrique
non vide muni d'un morphisme universellement injectif
$\mathcal{X} \to \Spec(k)$, et nous voulons trouver, à valeurs dans $k$,
un point de $\mathcal{X}$. Nous pouvons remplacer $\mathcal{X}$ par sa réduction.
Choisissons un corps $k'$ et un morphisme plat surjectif et localement
de type fini $\Spec(k') \to \mathcal{X}$ ; voir Propriétés des champs,
Lemme \ref{stacks-properties-lemma-unique-point}.
Comme $\mathcal{X} \to \Spec(k)$ est universellement injectif,
$$
\Spec(k') \times_\mathcal{X} \Spec(k') \to \Spec(k' \otimes_k k')
$$
est surjectif, car c'est le changement de base du morphisme surjectif
$\Delta : \mathcal{X} \to \mathcal{X} \times_{\Spec(k)} \mathcal{X}$
(Lemme \ref{lemma-universally-injective}).
Comme $k$ est algébriquement clos, $k' \otimes_k k'$ est un domaine
(Algèbre, Lemme
\ref{algebra-lemma-geometrically-integral-any-integral-base-change}).
Soit $\xi \in \Spec(k') \times_\mathcal{X} \Spec(k')$
un point dont l'image est le point générique de $\Spec(k' \otimes_k k')$.
Soit $U$ la composante connexe de
$\Spec(k') \times_\mathcal{X} \Spec(k')$ contenant $\xi$, munie de
sa structure réduite induite de sous-schéma fermé. Les deux projections
$U \to \Spec(k')$ sont alors localement de type fini, puisque c'était
le cas des projections
$\Spec(k') \times_\mathcal{X} \Spec(k') \to \Spec(k')$
qui sont des changements de base du morphisme
$\Spec(k') \to \mathcal{X}$. En appliquant Variétés, Proposition
\ref{varieties-proposition-unique-base-field}, nous voyons que les
clôtures intégrales des deux images de $k'$ dans
$\Gamma(U, \mathcal{O}_U)$ sont égales. Dans $\kappa(\xi)$, cela signifie
que tout élément de la forme $\lambda \otimes 1$ est algébriquement
dépendant du sous-corps
$$
1 \otimes k' \subset 
(\text{corps des fractions de }k' \otimes_k k') \subset
\kappa(\xi).
$$
Comme $k$ est algébriquement clos, cela n'est possible que si
$k' = k$, ce qui achève la démonstration.
\end{proof}

\begin{lemma}
\label{lemma-universally-injective-local}
Soit $f : \mathcal{X} \to \mathcal{Y}$ un morphisme de champs algébriques.
Les assertions suivantes sont équivalentes :
\begin{enumerate}
\item $f$ est universellement injectif ;
\item pour tout schéma affine $Z$ et tout morphisme
$Z \to \mathcal{Y}$, le morphisme
$Z \times_\mathcal{Y} \mathcal{X} \to Z$
est universellement injectif ;
\item compléter ici.
\end{enumerate}
\end{lemma}

\begin{proof}
L'implication (1) $\Rightarrow$ (2) est immédiate.
Supposons (2). Nous allons montrer que
$\Delta_f : \mathcal{X} \to \mathcal{X} \times_\mathcal{Y} \mathcal{X}$
est surjective, ce qui implique (1) d'après le
Lemme \ref{lemma-universally-injective}.
Considérons un schéma affine $V$ et un morphisme lisse
$V \to \mathcal{Y}$. Puisque
$g : V \times_\mathcal{Y} \mathcal{X} \to V$
est universellement injectif d'après (2), sa diagonale
$\Delta_g$ est surjective.
Or $\Delta_g$ est le changement de base de $\Delta_f$
par le morphisme lisse $V \to \mathcal{Y}$.
Comme la famille de ces morphismes $V \to \mathcal{Y}$ est
conjointement surjective, nous concluons que $\Delta_f$ est surjective.
\end{proof}

\begin{lemma}
\label{lemma-check-universally-injective-covering}
Soit $f : \mathcal{X} \to \mathcal{Y}$ un morphisme de champs algébriques.
Soit $W \to \mathcal{Y}$ un morphisme plat surjectif et localement
de présentation finie, où $W$ est un espace algébrique. Si le changement
de base $W \times_\mathcal{Y} \mathcal{X} \to W$ est universellement
injectif, alors $f$ est universellement injectif.
\end{lemma}

\begin{proof}
Observons que la diagonale $\Delta_g$ du morphisme
$g : W \times_\mathcal{Y} \mathcal{X} \to W$
est le changement de base de $\Delta_f$ par $W \to \mathcal{Y}$.
Ainsi, si $\Delta_g$ est surjective, alors $\Delta_f$ l'est aussi
d'après Propriétés des champs, Lemme
\ref{stacks-properties-lemma-check-property-covering}.
Le lemme résulte donc de la caractérisation (2) du
Lemme \ref{lemma-universally-injective}.
\end{proof}








\section{Homéomorphismes universels}
\label{section-universal-homeomorphisms}

\noindent
Soit $f$ un morphisme de champs algébriques qui est représentable par
des espaces algébriques. Dans Propriétés des champs, section
\ref{stacks-properties-section-properties-morphisms}
nous avons défini ce que signifie pour $f$ d'être un homéomorphisme universel.
En voici une autre caractérisation.

\begin{lemma}
\label{lemma-characterize-representable-universal-homeomorphism}
Soit $f : \mathcal{X} \to \mathcal{Y}$ un morphisme de
champs algébriques qui est représentable par des espaces algébriques.
Les assertions suivantes sont équivalentes :
\begin{enumerate}
\item $f$ est un homéomorphisme universel (Propriétés des champs,
section \ref{stacks-properties-section-properties-morphisms}) ;
\item pour tout morphisme de champs algébriques
$\mathcal{Z} \to \mathcal{Y}$, l'application d'espaces topologiques
$|\mathcal{Z} \times_\mathcal{Y} \mathcal{X}| \to |\mathcal{Z}|$ est
un homéomorphisme.
\end{enumerate}
\end{lemma}

\begin{proof}
Supposons (1), et soit $\mathcal{Z} \to \mathcal{Y}$ comme en (2).
Choisissons un schéma $V$ et un morphisme lisse surjectif
$V \to \mathcal{Z}$. Par hypothèse, le morphisme d'espaces algébriques
$V \times_\mathcal{Y} \mathcal{X} \to V$ est un homéomorphisme universel ;
en particulier, l'application
$|V \times_\mathcal{Y} \mathcal{X}| \to |V|$ est un homéomorphisme.
D'après Propriétés des champs, section
\ref{stacks-properties-section-points}, dans le diagramme commutatif
$$
\xymatrix{
|V \times_\mathcal{Y} \mathcal{X}| \ar[r] \ar[d] &
|\mathcal{Z} \times_\mathcal{Y} \mathcal{X}| \ar[d] \\
|V| \ar[r] & |\mathcal{Z}|
}
$$
les flèches horizontales sont ouvertes et surjectives ; de plus,
$$
|V \times_\mathcal{Y} \mathcal{X}| \longrightarrow
|V| \times_{|\mathcal{Z}|} |\mathcal{Z} \times_\mathcal{Y} \mathcal{X}|
$$
est surjective. Puisque la flèche verticale de gauche est un
homéomorphisme, il s'ensuit que celle de droite est un homéomorphisme.
Cela démontre (2). L'implication (2) $\Rightarrow$ (1) résulte
des définitions.
\end{proof}

\noindent
Nous pouvons donc adopter la définition naturelle suivante.

\begin{definition}
\label{definition-universal-homeomorphism}
Soit $f : \mathcal{X} \to \mathcal{Y}$ un morphisme de champs algébriques.
Nous disons que $f$ est un {\it homéomorphisme universel} si, pour tout
morphisme de champs algébriques $\mathcal{Z} \to \mathcal{Y}$,
l'application d'espaces topologiques
$$
|\mathcal{Z} \times_\mathcal{Y} \mathcal{X}| \to |\mathcal{Z}|
$$
est un homéomorphisme.
\end{definition}

\begin{lemma}
\label{lemma-base-change-universal-homeomorphism}
Le changement de base d'un homéomorphisme universel de champs algébriques
par un morphisme quelconque de champs algébriques est un homéomorphisme
universel.
\end{lemma}

\begin{proof}
C'est immédiat par définition.
\end{proof}

\begin{lemma}
\label{lemma-composition-universal-homeomorphism}
Le composé de deux homéomorphismes universels de champs algébriques
est un homéomorphisme universel.
\end{lemma}

\begin{proof}
Démonstration omise.
\end{proof}

\begin{lemma}
\label{lemma-check-universal-homeomorphism-covering}
Soit $f : \mathcal{X} \to \mathcal{Y}$ un morphisme de champs algébriques.
Soit $W \to \mathcal{Y}$ un morphisme plat surjectif et localement
de présentation finie, où $W$ est un espace algébrique. Si le changement
de base $W \times_\mathcal{Y} \mathcal{X} \to W$ est un homéomorphisme
universel, alors $f$ est un homéomorphisme universel.
\end{lemma}

\begin{proof}
Supposons que $g : W \times_\mathcal{Y} \mathcal{X} \to W$ soit un
homéomorphisme universel. Alors $g$ est universellement injectif ;
par conséquent, $f$ est universellement injectif d'après le
Lemme \ref{lemma-check-universally-injective-covering}.
D'autre part, soit $\mathcal{Z} \to \mathcal{Y}$ un morphisme,
où $\mathcal{Z}$ est un champ algébrique. Choisissons un schéma $U$
et un morphisme lisse surjectif
$U \to W \times_\mathcal{Y} \mathcal{Z}$. Considérons le diagramme
$$
\xymatrix{
W \times_\mathcal{Y} \mathcal{X} \ar[d]^g &
U \times_\mathcal{Y} \mathcal{X} \ar[d] \ar[l] \ar[r] &
\mathcal{Z} \times_\mathcal{Y} \mathcal{X} \ar[d] \\
W &
U \ar[l] \ar[r] &
\mathcal{Z}
}
$$
Par l'hypothèse sur $g$, la flèche verticale médiane induit un
homéomorphisme sur les espaces topologiques. Les morphismes
$U \to \mathcal{Z}$ et
$U \times_\mathcal{Y} \mathcal{X} \to
\mathcal{Z} \times_\mathcal{Y} \mathcal{X}$
sont plats, surjectifs et localement de présentation finie ; ils
induisent donc des applications ouvertes sur les espaces topologiques.
Nous concluons que
$|\mathcal{Z} \times_\mathcal{Y} \mathcal{X}| \to |\mathcal{Z}|$
est ouverte. La surjectivité est facile à démontrer ; nous omettons
la démonstration.
\end{proof}










\section{Propriétés de morphismes locales pour la topologie lisse
sur la source et la cible}
\label{section-local-source-target}

\noindent
Étant donnée une propriété des morphismes d'espaces algébriques qui est
{\it locale pour la topologie lisse sur la source et la cible}, voir
Descente sur les espaces,
Définition \ref{spaces-descent-definition-local-source-target},
nous pouvons nous en servir pour définir une propriété correspondante
des morphismes de champs algébriques, en imposant l'une ou l'autre des
conditions équivalentes du lemme suivant.

\begin{lemma}
\label{lemma-local-source-target}
Soit $\mathcal{P}$ une propriété des morphismes d'espaces algébriques
qui est locale pour la topologie lisse sur la source et la cible.
Soit $f : \mathcal{X} \to \mathcal{Y}$ un morphisme de champs algébriques.
Considérons des diagrammes commutatifs
$$
\xymatrix{
U \ar[d]_a \ar[r]_h & V \ar[d]^b \\
\mathcal{X} \ar[r]^f & \mathcal{Y}
}
$$
où $U$ et $V$ sont des espaces algébriques et les flèches verticales
sont lisses. Les assertions suivantes sont équivalentes :
\begin{enumerate}
\item pour tout diagramme comme ci-dessus tel que, de plus,
$U \to \mathcal{X} \times_\mathcal{Y} V$ soit lisse,
le morphisme $h$ possède la propriété $\mathcal{P}$ ;
\item pour un certain diagramme comme ci-dessus où
$a : U \to \mathcal{X}$ est surjectif, le morphisme $h$
possède la propriété $\mathcal{P}$.
\end{enumerate}
Si $\mathcal{X}$ et $\mathcal{Y}$ sont représentables par des espaces
algébriques, ces assertions équivalent encore à ce que $f$, considéré
comme morphisme d'espaces algébriques, possède la propriété $\mathcal{P}$.
Si, de plus, $\mathcal{P}$ est stable par tout changement de base et
locale pour la topologie fppf sur la base, alors, pour les morphismes $f$
représentables par des espaces algébriques, elles équivalent encore à ce
que $f$ possède la propriété $\mathcal{P}$ au sens de Propriétés des champs,
section \ref{stacks-properties-section-properties-morphisms}.
\end{lemma}

\begin{proof}
Démontrons l'implication (1) $\Rightarrow$ (2). Choisissons un espace
algébrique $V$ et un morphisme lisse surjectif $V \to \mathcal{Y}$.
Choisissons un espace algébrique $U$ et un morphisme lisse surjectif
$U \to \mathcal{X} \times_\mathcal{Y} V$. Le morphisme
$U \to \mathcal{X}$ est lui aussi lisse et surjectif, car il est le
composé du changement de base
$\mathcal{X} \times_\mathcal{Y} V \to \mathcal{X}$ et du morphisme choisi
$U \to \mathcal{X} \times_\mathcal{Y} V$. Nous obtenons donc un diagramme
comme en (1). Si (1) est satisfaite, $h : U \to V$ possède la propriété
$\mathcal{P}$ ; puisque $U \to \mathcal{X}$ est surjectif, (2) est satisfaite.

\medskip\noindent
Réciproquement, supposons (2) satisfaite et soient $U, V, a, b, h$
comme en (2). Soit ensuite $U', V', a', b', h'$ un diagramme quelconque
comme en (1). Considérons
$$
\xymatrix{
U \ar[d] \ar[r]_h & V \ar[d] \\
\mathcal{X} \ar[r]^f & \mathcal{Y}
}
\quad\quad
\xymatrix{
U' \ar[d] \ar[r]_{h'} & V' \ar[d] \\
\mathcal{X} \ar[r]^f & \mathcal{Y}
}
$$
Pour montrer que (2) implique (1), il faut prouver que $h'$ possède
$\mathcal{P}$. Considérons à cette fin le diagramme commutatif
$$
\xymatrix{
U \ar[dd]^h &
U \times_\mathcal{X} U' \ar[d] \ar[l] \ar@/^6ex/[dd]^{(h, h')} \ar[r] &
U' \ar[dd]^{h'} \\
& U \times_\mathcal{Y} V' \ar[lu] \ar[d] & \\
V &
V \times_\mathcal{Y} V' \ar[l] \ar[r] &
V'
}
$$
d'espaces algébriques. Les flèches horizontales sont lisses, car ce sont
des changements de base des morphismes lisses
$V \to \mathcal{Y}$, $V' \to \mathcal{Y}$, $U \to \mathcal{X}$ et
$U' \to \mathcal{X}$. D'autre part,
$$
\xymatrix{
U \times_\mathcal{X} U' \ar[d] \ar[r] & U' \ar[d] \\
U \times_\mathcal{Y} V' \ar[r] & \mathcal{X} \times_\mathcal{Y} V'
}
$$
est cartésien ; la flèche verticale de gauche est donc lisse puisque
$U', V', a', b', h'$ sont comme en (1).
Comme $\mathcal{P}$ est locale pour la topologie lisse sur la cible
d'après Descente sur les espaces, Lemme
\ref{spaces-descent-lemma-local-source-target-implies}, partie (2),
le changement de base
$U \times_\mathcal{Y} V' \to V \times_\mathcal{Y} V'$
possède $\mathcal{P}$. Comme $\mathcal{P}$ est locale pour la topologie
lisse sur la source d'après Descente sur les espaces, Lemme
\ref{spaces-descent-lemma-local-source-target-implies}, partie (1),
nous pouvons précomposer par le morphisme lisse
$U \times_\mathcal{X} U' \to U \times_\mathcal{Y} V'$ et conclure que
$(h, h')$ possède $\mathcal{P}$.
Comme $V \times_\mathcal{Y} V' \to V'$ est lisse, Descente sur les espaces,
Lemme \ref{spaces-descent-lemma-local-source-target-implies}, partie (3),
montre que $U \times_\mathcal{X} U' \to V'$ possède $\mathcal{P}$.
Enfin, puisque $U \times_\mathcal{X} U' \to U'$ est lisse et surjectif
et que $\mathcal{P}$ est locale pour la topologie lisse sur la source
(même lemme), nous concluons que $h'$ possède $\mathcal{P}$. Cela achève
la démonstration de l'équivalence de (1) et (2).

\medskip\noindent
Si $\mathcal{X}$ et $\mathcal{Y}$ sont représentables, Descente sur les
espaces, Lemme \ref{spaces-descent-lemma-local-source-target-characterize},
s'applique et montre que (1) et (2) équivalent à ce que $f$ possède
$\mathcal{P}$.

\medskip\noindent
Enfin, supposons $f$ représentable, soient $U, V, a, b, h$ comme dans
la partie (2) du lemme, et supposons $\mathcal{P}$ stable par tout
changement de base. Il faut montrer que, pour tout schéma $Z$ et tout
morphisme $Z \to \mathcal{Y}$, le changement de base
$Z \times_\mathcal{Y} \mathcal{X} \to Z$
possède la propriété $\mathcal{P}$. Considérons le diagramme
$$
\xymatrix{
Z \times_\mathcal{Y} U \ar[d] \ar[r] &
Z \times_\mathcal{Y} V \ar[d] \\
Z \times_\mathcal{Y} \mathcal{X} \ar[r] &
Z
}
$$
La flèche horizontale supérieure est un changement de base de $h$ et
possède donc la propriété $\mathcal{P}$. La flèche verticale de gauche
est lisse et surjective, et celle de droite est lisse. Descente sur les
espaces, Lemme \ref{spaces-descent-lemma-local-source-target-characterize},
s'applique donc et montre que
$Z \times_\mathcal{Y} \mathcal{X} \to Z$ possède la propriété
$\mathcal{P}$.
\end{proof}

\begin{definition}
\label{definition-P}
Soit $\mathcal{P}$ une propriété des morphismes d'espaces algébriques
qui est locale pour la topologie lisse sur la source et la cible.
Nous disons qu'un morphisme $f : \mathcal{X} \to \mathcal{Y}$ de champs
algébriques {\it possède la propriété $\mathcal{P}$} si les conditions
équivalentes du Lemme \ref{lemma-local-source-target} sont satisfaites.
\end{definition}

\begin{remark}
\label{remark-composition}
Soit $\mathcal{P}$ une propriété des morphismes d'espaces algébriques
qui est locale pour la topologie lisse sur la source et la cible et
stable par composition. La propriété des morphismes de champs algébriques
définie dans la Définition \ref{definition-P} est alors stable par
composition. En effet, soient $f : \mathcal{X} \to \mathcal{Y}$ et
$g : \mathcal{Y} \to \mathcal{Z}$ des morphismes de champs algébriques
possédant la propriété $\mathcal{P}$. Choisissons un espace algébrique
$W$ et un morphisme lisse surjectif $W \to \mathcal{Z}$. Choisissons
un espace algébrique $V$ et un morphisme lisse surjectif
$V \to \mathcal{Y} \times_\mathcal{Z} W$. Enfin, choisissons un espace
algébrique $U$ et un morphisme lisse surjectif
$U \to \mathcal{X} \times_\mathcal{Y} V$. Par définition, les morphismes
$V \to W$ et $U \to V$ possèdent la propriété $\mathcal{P}$.
Ainsi, $U \to W$ possède la propriété $\mathcal{P}$, puisque nous avons
supposé $\mathcal{P}$ stable par composition. La définition montre alors
de nouveau que
$g \circ f : \mathcal{X} \to \mathcal{Z}$ possède la propriété
$\mathcal{P}$.
\end{remark}

\begin{remark}
\label{remark-base-change}
Soit $\mathcal{P}$ une propriété des morphismes d'espaces algébriques
qui est locale pour la topologie lisse sur la source et la cible et
stable par changement de base. La propriété des morphismes de champs
algébriques définie dans la Définition \ref{definition-P} est alors stable
par changement de base. En effet, soient
$f : \mathcal{X} \to \mathcal{Y}$ et
$g : \mathcal{Y}' \to \mathcal{Y}$ des morphismes de champs algébriques,
et supposons que $f$ possède la propriété $\mathcal{P}$. Choisissons
un espace algébrique $V$ et un morphisme lisse surjectif
$V \to \mathcal{Y}$. Choisissons un espace algébrique $U$ et un morphisme
lisse surjectif $U \to \mathcal{X} \times_\mathcal{Y} V$. Enfin,
choisissons un espace algébrique $V'$ et un morphisme lisse surjectif
$V' \to \mathcal{Y}' \times_\mathcal{Y} V$. Par définition, le morphisme
$U \to V$ possède la propriété $\mathcal{P}$. Ainsi,
$V' \times_V U \to V'$ possède la propriété $\mathcal{P}$, puisque nous
avons supposé $\mathcal{P}$ stable par changement de base. En considérant
le diagramme
$$
\xymatrix{
V' \times_V U \ar[r] \ar[d] &
\mathcal{Y}' \times_\mathcal{Y} \mathcal{X} \ar[r] \ar[d] &
\mathcal{X} \ar[d] \\
V' \ar[r] & \mathcal{Y}' \ar[r] & \mathcal{Y}
}
$$
nous voyons que la flèche horizontale supérieure gauche est lisse et
surjective ; la définition montre donc que la projection
$\mathcal{Y}' \times_\mathcal{Y} \mathcal{X} \to \mathcal{Y}'$
possède la propriété $\mathcal{P}$.
\end{remark}

\begin{remark}
\label{remark-implication}
Soient $\mathcal{P}, \mathcal{P}'$ des propriétés des morphismes d'espaces
algébriques qui sont locales pour la topologie lisse sur la source et
la cible. Supposons que $\mathcal{P} \Rightarrow \mathcal{P}'$ pour les
morphismes d'espaces algébriques. Alors
$\mathcal{P} \Rightarrow \mathcal{P}'$ également pour les propriétés des
morphismes de champs algébriques définies dans la
Définition \ref{definition-P} à l'aide de $\mathcal{P}$ et
$\mathcal{P}'$. Cela résulte immédiatement de la définition.
\end{remark}









\section{Morphismes de type fini}
\label{section-finite-type}

\noindent
La propriété « localement de type fini » des morphismes d'espaces
algébriques est locale pour la topologie lisse sur la source et la cible ;
voir Descente sur les espaces, Remarque
\ref{spaces-descent-remark-list-local-source-target}.
Elle est également stable par changement de base et locale pour la
topologie fpqc sur la cible ; voir Morphismes d'espaces, Lemme
\ref{spaces-morphisms-lemma-base-change-finite-type}, et Descente sur les
espaces, Lemme
\ref{spaces-descent-lemma-descending-property-locally-finite-type}.
Par conséquent, d'après le Lemme \ref{lemma-local-source-target}
ci-dessus, nous pouvons définir comme suit ce que signifie pour un morphisme
de champs algébriques d'être localement de type fini. Cette définition
coïncide, lorsque le morphisme est représentable par des espaces algébriques,
avec la notion déjà définie dans Propriétés des champs, section
\ref{stacks-properties-section-properties-morphisms}.

\begin{definition}
\label{definition-locally-finite-type}
Soit $f : \mathcal{X} \to \mathcal{Y}$ un morphisme de champs algébriques.
\begin{enumerate}
\item Nous disons que $f$ est {\it localement de type fini} si les
conditions équivalentes du Lemme \ref{lemma-local-source-target}
sont satisfaites avec
$\mathcal{P} = \text{localement de type fini}$.
\item Nous disons que $f$ est {\it de type fini} s'il est localement
de type fini et quasi-compact.
\end{enumerate}
\end{definition}

\begin{lemma}
\label{lemma-composition-finite-type}
Le composé de morphismes de type fini est de type fini.
Il en va de même pour les morphismes localement de type fini.
\end{lemma}

\begin{proof}
Le résultat découle de la Remarque \ref{remark-composition} et du
Lemme \ref{spaces-morphisms-lemma-composition-finite-type} de Morphismes d'espaces.
\end{proof}

\begin{lemma}
\label{lemma-base-change-finite-type}
Un changement de base d'un morphisme de type fini est de type fini.
Il en va de même pour les morphismes localement de type fini.
\end{lemma}

\begin{proof}
Le résultat découle de la Remarque \ref{remark-base-change} et du
Lemme \ref{spaces-morphisms-lemma-base-change-finite-type} de Morphismes d'espaces.
\end{proof}

\begin{lemma}
\label{lemma-immersion-locally-finite-type}
Une immersion est localement de type fini.
\end{lemma}

\begin{proof}
Le résultat découle de la Remarque \ref{remark-implication} et du
Lemme \ref{spaces-morphisms-lemma-immersion-locally-finite-type} de Morphismes d'espaces.
\end{proof}

\begin{lemma}
\label{lemma-locally-finite-type-locally-noetherian}
Soit $f : \mathcal{X} \to \mathcal{Y}$ un morphisme de champs algébriques.
Si $f$ est localement de type fini et si $\mathcal{Y}$ est localement
noethérien, alors $\mathcal{X}$ est localement noethérien.
\end{lemma}

\begin{proof}
Soit
$$
\xymatrix{
U \ar[d] \ar[r] & V \ar[d] \\
\mathcal{X} \ar[r] & \mathcal{Y}
}
$$
un diagramme commutatif où $U$ et $V$ sont des schémas,
$V \to \mathcal{Y}$ est lisse et surjectif, et
$U \to V \times_\mathcal{Y} \mathcal{X}$ est lisse et surjectif.
Alors $U \to V$ est localement de type fini. Si $\mathcal{Y}$ est
localement noethérien, $V$ l'est aussi. D'après Morphismes, Lemme
\ref{morphisms-lemma-finite-type-noetherian}, $U$ est localement
noethérien, ce qui signifie que $\mathcal{X}$ est localement noethérien.
\end{proof}

\noindent
Les deux lemmes suivants seront améliorés plus loin (après l'étude des
morphismes de champs algébriques localement de présentation finie).

\begin{lemma}
\label{lemma-check-finite-type-covering}
Soit $f : \mathcal{X} \to \mathcal{Y}$ un morphisme de champs algébriques.
Soit $W \to \mathcal{Y}$ un morphisme plat surjectif et localement de
présentation finie, où $W$ est un espace algébrique. Si le changement
de base $W \times_\mathcal{Y} \mathcal{X} \to W$ est localement de
type fini, alors $f$ est localement de type fini.
\end{lemma}

\begin{proof}
Choisissons un espace algébrique $V$ et un morphisme lisse surjectif
$V \to \mathcal{Y}$. Choisissons un espace algébrique $U$ et un morphisme
lisse surjectif $U \to V \times_\mathcal{Y} \mathcal{X}$.
Il faut montrer que $U \to V$ est localement de type fini.
Effectuons maintenant partout le changement de base par
$W \to \mathcal{Y}$ : posons
$U' = W \times_\mathcal{Y} U$,
$V' = W \times_\mathcal{Y} V$,
$\mathcal{X}' = W \times_\mathcal{Y} \mathcal{X}$,
et $\mathcal{Y}' = W \times_\mathcal{Y} \mathcal{Y} = W$.
Par changement de base,
$U' \to V' \times_{\mathcal{Y}'} \mathcal{X}'$ est encore lisse.
La définition des morphismes de champs algébriques localement de type
fini et l'hypothèse selon laquelle
$\mathcal{X}' \to \mathcal{Y}'$ est localement de type fini montrent donc
que $U' \to V'$ est localement de type fini. Puisque $V' \to V$ est
plat, surjectif et localement de présentation finie, comme changement
de base de $W \to \mathcal{Y}$, Descente sur les espaces, Lemme
\ref{spaces-descent-lemma-descending-property-locally-finite-type},
montre que $U \to V$ est localement de type fini, ce qui conclut.
\end{proof}

\begin{lemma}
\label{lemma-check-finite-type-precompose}
Soient $\mathcal{X} \to \mathcal{Y} \to \mathcal{Z}$ des morphismes de
champs algébriques. Supposons $\mathcal{X} \to \mathcal{Z}$ localement
de type fini et $\mathcal{X} \to \mathcal{Y}$ représentable par des
espaces algébriques, plat, surjectif et localement de présentation finie.
Alors $\mathcal{Y} \to \mathcal{Z}$ est localement de type fini.
\end{lemma}

\begin{proof}
Choisissons un espace algébrique $W$ et un morphisme lisse surjectif
$W \to \mathcal{Z}$. Choisissons un espace algébrique $V$ et un morphisme
lisse surjectif $V \to W \times_\mathcal{Z} \mathcal{Y}$. Posons
$U = V \times_\mathcal{Y} \mathcal{X}$, qui est un espace algébrique.
Le morphisme $U \to V$ est plat, surjectif et localement de présentation
finie, tandis que $U \to W$ est localement de type fini.
Le lemme se ramène donc au cas des morphismes d'espaces algébriques.
Ce cas est Descente sur les espaces, Lemme
\ref{spaces-descent-lemma-locally-finite-type-fppf-local-source}.
\end{proof}

\begin{lemma}
\label{lemma-finite-type-permanence}
Soient $f : \mathcal{X} \to \mathcal{Y}$ et
$g : \mathcal{Y} \to \mathcal{Z}$ des morphismes de champs algébriques.
Si $g \circ f : \mathcal{X} \to \mathcal{Z}$ est localement de type fini,
alors $f : \mathcal{X} \to \mathcal{Y}$ est localement de type fini.
\end{lemma}

\begin{proof}
On peut trouver un diagramme
$$
\xymatrix{
U \ar[r] \ar[d] & V \ar[r] \ar[d] & W \ar[d] \\
\mathcal{X} \ar[r] & \mathcal{Y} \ar[r] & \mathcal{Z}
}
$$
où $U$, $V$ et $W$ sont des schémas, la flèche verticale
$W \to \mathcal{Z}$ est lisse et surjective, la flèche
$V \to \mathcal{Y} \times_\mathcal{Z} W$ est lisse et surjective,
et la flèche $U \to \mathcal{X} \times_\mathcal{Y} V$ est lisse et
surjective. Alors $U \to \mathcal{X} \times_\mathcal{Z} W$ est lui aussi
lisse et surjectif (comme composé d'un morphisme lisse surjectif et du
changement de base d'un tel morphisme). Par définition, $U \to W$ est
localement de type fini. D'après Morphismes, Lemme
\ref{morphisms-lemma-permanence-finite-type}, $U \to V$ est donc
localement de type fini, ce qui signifie à son tour, par définition,
que $\mathcal{X} \to \mathcal{Y}$ est localement de type fini.
\end{proof}










\section{Points de type fini}
\label{section-points-finite-type}

\noindent
Soit $\mathcal{X}$ un champ algébrique.
Un point de type fini $x \in |\mathcal{X}|$ est un point qui peut être représenté
par un morphisme $\Spec(k) \to \mathcal{X}$ localement de type fini.
Les points de type fini remplacent avantageusement les points fermés pour les
espaces algébriques et les champs algébriques. Par exemple, il y en a toujours
``assez''.

\begin{lemma}
\label{lemma-point-finite-type}
Soit $\mathcal{X}$ un champ algébrique.
Soit $x \in |\mathcal{X}|$. Les conditions suivantes sont équivalentes :
\begin{enumerate}
\item Il existe un morphisme $\Spec(k) \to \mathcal{X}$
localement de type fini qui représente $x$.
\item Il existe un schéma $U$, un point fermé $u \in U$ et un morphisme
lisse $\varphi : U \to \mathcal{X}$ tels que $\varphi(u) = x$.
\end{enumerate}
\end{lemma}

\begin{proof}
Soient $u \in U$ et $U \to \mathcal{X}$ comme en (2). Alors
$\Spec(\kappa(u)) \to U$ est de type fini, et $U \to \mathcal{X}$ est
représentable et localement de type fini (par Morphismes d'espaces,
Lemmes \ref{spaces-morphisms-lemma-etale-locally-finite-presentation} et
\ref{spaces-morphisms-lemma-finite-presentation-finite-type}).
La condition (1) résulte donc du
Lemme \ref{lemma-composition-finite-type}.

\medskip\noindent
Réciproquement, supposons que $\Spec(k) \to \mathcal{X}$ soit localement de
type fini et représente $x$. Soit $U \to \mathcal{X}$ un morphisme lisse
surjectif, où $U$ est un schéma. Par hypothèse,
$U \times_\mathcal{X} \Spec(k) \to U$ est un morphisme d'espaces algébriques
localement de type fini. Choisissons un point de type fini $v$ de
$U \times_\mathcal{X} \Spec(k)$ (il en existe au moins un, voir Morphismes
d'espaces, Lemme \ref{spaces-morphisms-lemma-identify-finite-type-points}).
Par Morphismes d'espaces,
Lemme \ref{spaces-morphisms-lemma-finite-type-points-morphism},
l'image $u \in U$ de $v$ est un point de type fini de $U$.
Par Morphismes, Lemme \ref{morphisms-lemma-identify-finite-type-points},
quitte à restreindre $U$, nous pouvons donc supposer que $u$ est un point
fermé de $U$ ; autrement dit, (2) est satisfaite.
\end{proof}

\begin{definition}
\label{definition-finite-type-point}
Soit $\mathcal{X}$ un champ algébrique. On dit qu'un point
$x \in |\mathcal{X}|$ est un {\it point de type fini}\footnote{Il s'agit
d'un léger abus de langage, puisqu'il serait peut-être plus exact de dire
``point localement de type fini''.} s'il satisfait aux conditions équivalentes
du Lemme \ref{lemma-point-finite-type}. On note
$\mathcal{X}_{\text{ft-pts}}$ l'ensemble des points de type fini de
$\mathcal{X}$.
\end{definition}

\noindent
On peut décrire comme suit l'ensemble des points de type fini.

\begin{lemma}
\label{lemma-identify-finite-type-points}
Soit $\mathcal{X}$ un champ algébrique. On a
$$
\mathcal{X}_{\text{ft-pts}} =
\bigcup\nolimits_{\varphi : U \to \mathcal{X}\text{ lisse}} |\varphi|(U_0)
$$
où $U_0$ est l'ensemble des points fermés de $U$.
Ici, on peut faire parcourir à $U$ soit tous les schémas lisses sur
$\mathcal{X}$, soit tous les schémas affines lisses sur $\mathcal{X}$.
\end{lemma}

\begin{proof}
C'est une conséquence immédiate du
Lemme \ref{lemma-point-finite-type}.
\end{proof}

\begin{lemma}
\label{lemma-finite-type-points-morphism}
Soit $f : \mathcal{X} \to \mathcal{Y}$ un morphisme de champs algébriques.
Si $f$ est localement de type fini, alors
$f(\mathcal{X}_{\text{ft-pts}}) \subset \mathcal{Y}_{\text{ft-pts}}$.
\end{lemma}

\begin{proof}
Soit $x \in \mathcal{X}_{\text{ft-pts}}$. Représentons $x$ par un morphisme
localement de type fini $x : \Spec(k) \to \mathcal{X}$. Alors
$f \circ x$ est localement de type fini par le
Lemme \ref{lemma-composition-finite-type}.
Par conséquent, $f(x) \in \mathcal{Y}_{\text{ft-pts}}$.
\end{proof}

\begin{lemma}
\label{lemma-finite-type-points-surjective-morphism}
Soit $f : \mathcal{X} \to \mathcal{Y}$ un morphisme de champs algébriques.
Si $f$ est localement de type fini et surjectif, alors
$f(\mathcal{X}_{\text{ft-pts}}) = \mathcal{Y}_{\text{ft-pts}}$.
\end{lemma}

\begin{proof}
On a $f(\mathcal{X}_{\text{ft-pts}}) \subset
\mathcal{Y}_{\text{ft-pts}}$ par le
Lemme \ref{lemma-finite-type-points-morphism}.
Soit $y \in |\mathcal{Y}|$ un point de type fini. Représentons $y$ par un
morphisme $\Spec(k) \to \mathcal{Y}$ localement de type fini.
Puisque $f$ est surjectif, le champ algébrique
$\mathcal{X}_k = \Spec(k) \times_\mathcal{Y} \mathcal{X}$ est non vide ; il
possède donc un point de type fini $x \in |\mathcal{X}_k|$ par le
Lemme \ref{lemma-identify-finite-type-points}.
Or $\mathcal{X}_k \to \mathcal{X}$ est localement de type fini, car il se
déduit de $\Spec(k) \to \mathcal{Y}$ par changement de base
(Lemme \ref{lemma-base-change-finite-type}).
L'image de $x$ dans $\mathcal{X}$ est donc un point de type fini par le
Lemme \ref{lemma-finite-type-points-morphism}, et elle s'envoie sur $y$ par
construction.
\end{proof}

\begin{lemma}
\label{lemma-enough-finite-type-points}
Soit $\mathcal{X}$ un champ algébrique.
Pour toute partie localement fermée $T \subset |\mathcal{X}|$, on a
$$
T \not = \emptyset
\Rightarrow
T \cap \mathcal{X}_{\text{ft-pts}} \not = \emptyset.
$$
En particulier, pour toute partie fermée $T \subset |\mathcal{X}|$,
$T \cap \mathcal{X}_{\text{ft-pts}}$ est dense dans $T$.
\end{lemma}

\begin{proof}
Soit $i : \mathcal{Z} \to \mathcal{X}$ la structure réduite de sous-champ
induite sur $T$, voir Propriétés des champs,
Remarque \ref{stacks-properties-remark-stack-structure-locally-closed-subset}.
Une immersion est localement de type fini, voir le
Lemme \ref{lemma-immersion-locally-finite-type}.
Par le Lemme \ref{lemma-finite-type-points-morphism}, on a donc
$\mathcal{Z}_{\text{ft-pts}} \subset \mathcal{X}_{\text{ft-pts}} \cap T$.
Enfin, tout schéma affine non vide $U$ muni d'un morphisme lisse vers
$\mathcal{Z}$ possède au moins un point fermé ; ainsi $\mathcal{Z}$ possède
au moins un point de type fini par le
Lemme \ref{lemma-identify-finite-type-points}.
Le lemme en résulte.
\end{proof}

\noindent
Voici une autre caractérisation, plus technique, d'un point de type fini d'un
champ algébrique. Elle montre notamment que la gerbe résiduelle de
$\mathcal{X}$ en $x$ existe dès que $x$ est un point de type fini !

\begin{lemma}
\label{lemma-point-finite-type-monomorphism}
Soit $\mathcal{X}$ un champ algébrique.
Soit $x \in |\mathcal{X}|$. Les conditions suivantes sont équivalentes :
\begin{enumerate}
\item $x$ est un point de type fini,
\item il existe un champ algébrique $\mathcal{Z}$ dont l'espace topologique
sous-jacent $|\mathcal{Z}|$ est réduit à un point, ainsi qu'un morphisme
localement de type fini $f : \mathcal{Z} \to \mathcal{X}$ tel que
$\{x\} = |f|(|\mathcal{Z}|)$, et
\item la gerbe résiduelle $\mathcal{Z}_x$ de $\mathcal{X}$ en $x$ existe et
le morphisme d'inclusion $\mathcal{Z}_x \to \mathcal{X}$ est localement de
type fini.
\end{enumerate}
\end{lemma}

\begin{proof}
(Tous les morphismes qui interviennent dans ce paragraphe sont représentables
par des espaces algébriques ; les conventions et les résultats de
Propriétés des champs,
section \ref{stacks-properties-section-properties-morphisms},
sont donc applicables.)
Supposons que $x$ soit un point de type fini. Choisissons un schéma affine
$U$, un point fermé $u \in U$ et un morphisme lisse
$\varphi : U \to \mathcal{X}$ tels que $\varphi(u) = x$, voir le
Lemme \ref{lemma-identify-finite-type-points}.
Comme d'habitude, posons $u = \Spec(\kappa(u))$. Posons
$R = u \times_\mathcal{X} u$ ; on obtient ainsi un groupoïde en espaces
algébriques $(u, R, s, t, c)$, voir Champs algébriques,
Lemme \ref{algebraic-lemma-map-space-into-stack}.
Chacun des morphismes de projection $R \to u$ est le composé
$$
R = u \times_\mathcal{X} u \to
u \times_\mathcal{X} U \to
u \times_\mathcal{X} \mathcal{X} = u
$$
où la première flèche est de type fini (c'est un changement de base de
l'immersion fermée de schémas $u \to U$) et la seconde est lisse (c'est un
changement de base du morphisme lisse $U \to \mathcal{X}$). Par conséquent,
$s, t : R \to u$ sont localement de type fini (comme composés, voir
Morphismes d'espaces,
Lemme \ref{spaces-morphisms-lemma-composition-finite-type}).
Comme $u$ est le spectre d'un corps, les morphismes $s, t$ sont plats et localement de
présentation finie (par Morphismes d'espaces, Lemme
\ref{spaces-morphisms-lemma-noetherian-finite-type-finite-presentation}).
Le champ $\mathcal{Z} = [u/R]$ est donc algébrique par Critères de
représentabilité,
Théorème \ref{criteria-theorem-flat-groupoid-gives-algebraic-stack}.
Par Champs algébriques,
Lemme \ref{algebraic-lemma-map-space-into-stack},
on obtient un morphisme canonique
$$
f : \mathcal{Z} \longrightarrow \mathcal{X}
$$
qui est pleinement fidèle. Ce morphisme est donc représentable par des espaces
algébriques, voir Champs algébriques, Lemme
\ref{algebraic-lemma-characterize-representable-by-algebraic-spaces},
et c'est un monomorphisme, voir Propriétés des champs,
Lemme \ref{stacks-properties-lemma-monomorphism}.
Il s'ensuit que la gerbe résiduelle
$\mathcal{Z}_x \subset \mathcal{X}$ de $\mathcal{X}$ en $x$ existe et que
$f$ induit une équivalence $\mathcal{Z} \to \mathcal{Z}_x$, voir
Propriétés des champs,
Lemme \ref{stacks-properties-lemma-residual-gerbe-unique}.
Par construction, le diagramme
$$
\xymatrix{
u \ar[d] \ar[r] & U \ar[d] \\
\mathcal{Z} \ar[r]^f & \mathcal{X}
}
$$
est commutatif. Par Critères de représentabilité,
Lemme \ref{criteria-lemma-flat-quotient-flat-presentation},
la flèche verticale de gauche est surjective, plate et localement de
présentation finie. Considérons
$$
\xymatrix{
u \times_\mathcal{X} U \ar[d] \ar[r] &
\mathcal{Z} \times_\mathcal{X} U \ar[r] \ar[d] & U \ar[d] \\
u \ar[r] & \mathcal{Z} \ar[r]^f & \mathcal{X}
}
$$
Comme $u \to \mathcal{X}$ est localement de type fini, son changement de base
$u \times_\mathcal{X} U \to U$ est localement de type fini. De plus,
$u \times_\mathcal{X} U \to \mathcal{Z} \times_\mathcal{X} U$ est
surjectif, plat et localement de présentation finie, car il se déduit de
$u \to \mathcal{Z}$ par changement de base. Ainsi,
$\{u \times_\mathcal{X} U \to \mathcal{Z} \times_\mathcal{X} U\}$
est un recouvrement fppf d'espaces algébriques ; on en conclut que
$\mathcal{Z} \times_\mathcal{X} U \to U$ est localement de type fini par
Descente sur les espaces, Lemme
\ref{spaces-descent-lemma-locally-finite-presentation-fppf-local-source}.
Par définition, cela signifie que $f$ est localement de type fini (en effet,
la flèche verticale $\mathcal{Z} \times_\mathcal{X} U \to \mathcal{Z}$ est
lisse, comme changement de base de $U \to \mathcal{X}$, et surjective puisque
$\mathcal{Z}$ n'a qu'un point). Comme
$\mathcal{Z} = \mathcal{Z}_x$, la condition (3) est satisfaite.

\medskip\noindent
Il est clair que (3) implique (2).
Si (2) est satisfaite, alors $x$ est un point de type fini de $\mathcal{X}$
par le Lemme \ref{lemma-finite-type-points-morphism} et le
Lemme \ref{lemma-enough-finite-type-points}, qui montre que
$\mathcal{Z}_{\text{ft-pts}}$ est non vide, autrement dit que l'unique point
de $\mathcal{Z}$ est un point de type fini de $\mathcal{Z}$.
\end{proof}







\section{Groupes d'automorphismes}
\label{section-automorphism-groups}

\noindent
Soit $\mathcal{X}$ un champ algébrique. Soit $x \in |\mathcal{X}|$
correspondant à $x : \Spec(k) \to \mathcal{X}$. Dans cette situation, nous
dirons souvent : ``soit $G_x/k$ l'espace algébrique en groupes des
automorphismes de $x$''. Cela signifie simplement que
$$
G_x = \mathit{Isom}_\mathcal{X}(x, x) =
\Spec(k) \times_\mathcal{X} \mathcal{I}_\mathcal{X}
$$
est l'espace algébrique en groupes des automorphismes de $x$. C'est un espace
algébrique en groupes sur $\Spec(k)$. Si $k'/k$ est une extension de corps,
l'espace algébrique en groupes des automorphismes du morphisme induit
$x' : \Spec(k') \to \mathcal{X}$ se déduit de $G_x$ par changement de base
à $\Spec(k')$.

\begin{lemma}
\label{lemma-automorphism-group-scheme}
Dans la situation ci-dessus, $G_x$ est un schéma si l'une des conditions
suivantes est satisfaite :
\begin{enumerate}
\item $\Delta : \mathcal{X} \to \mathcal{X} \times \mathcal{X}$
est quasi-séparée,
\item $\Delta : \mathcal{X} \to \mathcal{X} \times \mathcal{X}$
est localement séparée,
\item $\mathcal{X}$ est quasi-DM,
\item $\mathcal{I}_\mathcal{X} \to \mathcal{X}$
est quasi-séparé,
\item $\mathcal{I}_\mathcal{X} \to \mathcal{X}$
est localement séparé, ou
\item $\mathcal{I}_\mathcal{X} \to \mathcal{X}$
est localement quasi-fini.
\end{enumerate}
\end{lemma}

\begin{proof}
Le Lemme \ref{lemma-diagonal-diagonal} montre que
(1) $\Rightarrow$ (4), (2) $\Rightarrow$ (5) et (3) $\Rightarrow$ (6).
Dans le cas (4), $G_x$ est un espace algébrique quasi-séparé ; dans le cas
(5), $G_x$ est un espace algébrique localement séparé.
Dans les deux cas, $G_x$ est un espace algébrique décent
(Espaces décents,
section \ref{decent-spaces-section-reasonable-decent} et
Lemme \ref{decent-spaces-lemma-locally-separated-decent}).
Ensuite, $G_x$ est séparé par Compléments sur les groupoïdes en espaces,
Lemme \ref{spaces-more-groupoids-lemma-group-scheme-over-field-separated} ;
on en conclut que $G_x$ est un schéma par Compléments sur les groupoïdes en
espaces, Proposition
\ref{spaces-more-groupoids-proposition-group-space-scheme-over-field}.
Dans le cas (6), $G_x \to \Spec(k)$ est localement quasi-fini ; $G_x$ est donc
un schéma par Espaces sur des corps,
Lemme \ref{spaces-over-fields-lemma-locally-quasi-finite-over-field}.
\end{proof}

\begin{lemma}
\label{lemma-property-automorphism-groups}
Soit $\mathcal{X}$ un champ algébrique et soit $x \in |\mathcal{X}|$ un point.
Soit $P$ une propriété des espaces algébriques sur des corps qui soit
invariante par extension du corps de base ; par exemple,
$P(X/k) = X \to \Spec(k)\text{ est fini}$.
Les conditions suivantes sont équivalentes :
\begin{enumerate}
\item pour un certain morphisme $x : \Spec(k) \to \mathcal{X}$ représentant
la classe de $x$, l'espace algébrique en groupes d'automorphismes $G_x/k$
vérifie $P$, et
\item pour tout morphisme $x : \Spec(k) \to \mathcal{X}$ représentant
la classe de $x$, l'espace algébrique en groupes d'automorphismes $G_x/k$
vérifie $P$.
\end{enumerate}
\end{lemma}

\begin{proof}
Démonstration omise.
\end{proof}

\begin{remark}
\label{remark-property-automorphism-groups}
Soit $P$ une propriété des espaces algébriques sur des corps qui soit
invariante par extension du corps de base. Étant donnés un champ algébrique
$\mathcal{X}$ et $x \in |\mathcal{X}|$, on dit que le groupe d'automorphismes
de $\mathcal{X}$ en $x$ vérifie $P$ si les conditions équivalentes du
Lemme \ref{lemma-property-automorphism-groups} sont satisfaites.
Par exemple, on dit que {\it le groupe d'automorphismes de $\mathcal{X}$
en $x$ est fini} si $G_x \to \Spec(k)$ est fini pour tout représentant
$x : \Spec(k) \to \mathcal{X}$ de $x$.
On procède de même pour les propriétés lisse, propre, etc.
(Il s'agit manifestement d'un abus de langage, mais nous pensons qu'il
n'entraînera ni confusion ni imprécision.)
\end{remark}

\begin{lemma}
\label{lemma-iso-automorphism-groups}
Soit $f : \mathcal{X} \to \mathcal{Y}$ un morphisme de champs algébriques.
Soit $x \in |\mathcal{X}|$ un point. Les conditions suivantes sont équivalentes :
\begin{enumerate}
\item pour un certain morphisme $x : \Spec(k) \to \mathcal{X}$ représentant
la classe de $x$, si l'on pose $y = f \circ x$, l'application
$G_x \to G_y$ entre les espaces algébriques en groupes d'automorphismes est
un isomorphisme, et
\item pour tout morphisme $x : \Spec(k) \to \mathcal{X}$ représentant
la classe de $x$, si l'on pose $y = f \circ x$, l'application
$G_x \to G_y$ entre les espaces algébriques en groupes d'automorphismes est
un isomorphisme.
\end{enumerate}
\end{lemma}

\begin{proof}
Cela résulte du fait qu'être un isomorphisme est une propriété locale pour la
topologie fpqc sur la cible, voir Descente sur les espaces, Lemme
\ref{spaces-descent-lemma-descending-property-isomorphism}.
En effet, supposons que $k'/k$ soit une extension de corps, notons
$x' : \Spec(k') \to \mathcal{X}$ le composé et posons $y' = f \circ x'$.
Alors le morphisme $G_{x'} \to G_{y'}$ se déduit de $G_x \to G_y$ par le
changement de base $\Spec(k') \to \Spec(k)$.
Par conséquent, $G_x \to G_y$ est un isomorphisme si et seulement si
$G_{x'} \to G_{y'}$ en est un. La propriété vaut donc pour toute la classe
d'équivalence dès qu'elle vaut pour un représentant.
\end{proof}

\begin{remark}
\label{remark-identify-automorphism-groups}
Soit $f : \mathcal{X} \to \mathcal{Y}$ un morphisme de champs algébriques,
et soit $x \in |\mathcal{X}|$ un point. Pour exprimer que $f$ et $x$
satisfont aux conditions équivalentes du
Lemme \ref{lemma-iso-automorphism-groups}, la littérature dit que
{\it $f$ préserve les stabilisateurs en $x$} ou que
{\it $f$ reflète les points fixes en $x$}.
Nous préférons dire que {\it $f$ induit un isomorphisme entre les groupes
d'automorphismes en $x$ et en $f(x)$}.
\end{remark}





\section{Présentations et propriétés des champs algébriques}
\label{section-presentations-properties}

\noindent
Soit $(U, R, s, t, c)$ un groupoïde en espaces algébriques.
Si $s, t : R \to U$ sont plats et localement de présentation finie,
alors le champ quotient $[U/R]$ est un champ algébrique, voir
Critères de représentabilité, Théorème
\ref{criteria-theorem-flat-groupoid-gives-algebraic-stack}.
Dans cette section, nous étudions les conséquences des propriétés de
$(U, R, s, t, c)$ pour le champ algébrique $[U/R]$.

\begin{lemma}
\label{lemma-properties-diagonal-from-presentation}
Soit $(U, R, s, t, c)$ un groupoïde en espaces algébriques tel que
$s, t : R \to U$ soient plats et localement de présentation finie.
Considérons le champ algébrique $\mathcal{X} = [U/R]$ (voir ci-dessus).
\begin{enumerate}
\item Si $R \to U \times U$ est séparé, alors
$\Delta_\mathcal{X}$ est séparée.
\item Si $U$ et $R$ sont séparés, alors $\Delta_\mathcal{X}$ est séparée.
\item Si $R \to U \times U$ est localement quasi-fini, alors $\mathcal{X}$
est quasi-DM.
\item Si $s, t : R \to U$ sont localement quasi-finis, alors
$\mathcal{X}$ est quasi-DM.
\item Si $R \to U \times U$ est propre, alors $\mathcal{X}$ est séparé.
\item Si $s, t : R \to U$ sont propres et si $U$ est séparé, alors
$\mathcal{X}$ est séparé.
\item Ajouter d'autres cas ici.
\end{enumerate}
\end{lemma}

\begin{proof}
Remarquons que le morphisme $U \to \mathcal{X}$ est surjectif, plat et
localement de présentation finie par Critères de représentabilité, Lemme
\ref{criteria-lemma-flat-quotient-flat-presentation}.
Il en est donc de même de
$U \times U \to \mathcal{X} \times \mathcal{X}$. On a le diagramme cartésien
$$
\xymatrix{
R  = U \times_\mathcal{X} U \ar[r] \ar[d] & U \times U \ar[d] \\
\mathcal{X} \ar[r] & \mathcal{X} \times \mathcal{X}
}
$$
(voir Groupoïdes en espaces, Lemme
\ref{spaces-groupoids-lemma-quotient-stack-2-cartesian}).
Ainsi, $\Delta_\mathcal{X}$ possède l'une des propriétés énumérées dans
Propriétés des champs, section
\ref{stacks-properties-section-properties-morphisms}
si et seulement si le morphisme $R \to U \times U$ la possède, voir
Propriétés des champs, Lemme
\ref{stacks-properties-lemma-check-property-covering}.
Cela démontre (1), (3) et (5).
La condition de (2) implique que $R \to U \times U$ est séparé ; (2) résulte
donc de (1). La condition de (4) implique celle de (3) ; (4) résulte donc de
(3). La condition de (6) implique celle de (5) par Morphismes d'espaces, Lemme
\ref{spaces-morphisms-lemma-universally-closed-permanence} ; (6) résulte donc
de (5).
\end{proof}

\begin{lemma}
\label{lemma-points-presentation}
Soit $(U, R, s, t, c)$ un groupoïde en espaces algébriques tel que
$s, t : R \to U$ soient plats et localement de présentation finie.
Considérons le champ algébrique $\mathcal{X} = [U/R]$ (voir ci-dessus).
Alors l'image de $|R| \to |U| \times |U|$ est une relation d'équivalence et
$|\mathcal{X}|$ est le quotient de $|U|$ par cette relation d'équivalence.
\end{lemma}

\begin{proof}
Le morphisme induit $p : U \to \mathcal{X}$ est surjectif, plat et localement
de présentation finie, voir Critères de représentabilité, Lemme
\ref{criteria-lemma-flat-quotient-flat-presentation}.
Ainsi, $|U| \to |\mathcal{X}|$ est surjectif par Propriétés des champs, Lemme
\ref{stacks-properties-lemma-characterize-surjective}.
Notons que $R = U \times_\mathcal{X} U$, voir Groupoïdes en espaces,
Lemme \ref{spaces-groupoids-lemma-quotient-stack-2-cartesian}.
Par conséquent, Propriétés des champs,
Lemme \ref{stacks-properties-lemma-points-cartesian}, montre que l'application
$$
|R| \longrightarrow |U| \times_{|\mathcal{X}|} |U|
$$
est surjective. L'image de $|R| \to |U| \times |U|$ est donc exactement
l'ensemble des couples $(u_1, u_2) \in |U| \times |U|$ tels que $u_1$ et
$u_2$ aient la même image dans $|\mathcal{X}|$.
En combinant ces deux assertions, on obtient le résultat du lemme.
\end{proof}






\section{Présentations particulières des champs algébriques}
\label{section-presentations}

\noindent
Dans cette section, nous démontrons deux théorèmes importants.
Le premier caractérise les champs quasi-DM $\mathcal{X}$ comme les champs de
la forme $\mathcal{X} = [U/R]$ pour lesquels $s, t : R \to U$ sont localement
quasi-finis (ainsi que plats et localement de présentation finie).
Le second affirme que les champs algébriques DM sont de Deligne--Mumford.

\medskip\noindent
Le lemme suivant donne un critère pour qu'une ``tranche'' d'une présentation
reste plate sur le champ algébrique.

\begin{lemma}
\label{lemma-slice}
Soit $\mathcal{X}$ un champ algébrique.
Considérons un diagramme cartésien
$$
\xymatrix{
U \ar[d] & F \ar[l]^p \ar[d] \\
\mathcal{X} & \Spec(k) \ar[l]
}
$$
où $U$ est un espace algébrique, $k$ un corps, et où $U \to \mathcal{X}$
est plat et localement de présentation finie. Soient
$f_1, \ldots, f_r \in \Gamma(U, \mathcal{O}_U)$
et $z \in |F|$ tels que les images de $f_1, \ldots, f_r$ forment une suite
régulière dans l'anneau local $\mathcal{O}_{F, \overline{z}}$.
Alors, après avoir remplacé $U$ par un sous-espace ouvert contenant $p(z)$,
le morphisme
$$
V(f_1, \ldots, f_r) \longrightarrow \mathcal{X}
$$
est plat et localement de présentation finie.
\end{lemma}

\begin{proof}
Choisissons un schéma $W$ et un morphisme lisse surjectif
$W \to \mathcal{X}$. Choisissons une extension de corps $k'/k$ et un
morphisme $w : \Spec(k') \to W$ tels que
$\Spec(k') \to W \to \mathcal{X}$ soit $2$-isomorphe à
$\Spec(k') \to \Spec(k) \to \mathcal{X}$. Cela est possible puisque
$W \to \mathcal{X}$ est surjectif. Considérons le diagramme commutatif
$$
\xymatrix{
U \ar[d] &
U \times_\mathcal{X} W \ar[l]^-{\text{pr}_0} \ar[d] &
F' \ar[l]^-{p'} \ar[d] \\
\mathcal{X} &
W \ar[l] &
\Spec(k') \ar[l]
}
$$
dont les deux carrés sont cartésiens. D'après le choix de $w$, on a
$F' = F \times_{\Spec(k)} \Spec(k')$. Ainsi $F' \to F$ est surjectif, et
l'on peut choisir un point $z' \in |F'|$ qui s'envoie sur $z$.
Comme $F' \to F$ est plat, l'homomorphisme
$\mathcal{O}_{F, \overline{z}} \to \mathcal{O}_{F', \overline{z}'}$ est
plat, voir Morphismes d'espaces,
Lemme \ref{spaces-morphisms-lemma-flat-at-point-etale-local-rings}.
Les images de $f_1, \ldots, f_r$ forment donc une suite régulière dans
$\mathcal{O}_{F', \overline{z}'}$, voir Algèbre,
Lemme \ref{algebra-lemma-flat-increases-depth}.
Notons que $U \times_\mathcal{X} W \to W$ est un morphisme d'espaces
algébriques plat et localement de présentation finie. Par Compléments sur les
morphismes d'espaces, Lemme \ref{spaces-more-morphisms-lemma-slice}, il existe
donc un sous-espace ouvert $U'$ de $U \times_\mathcal{X} W$, contenant
$p'(z')$, tel que l'intersection
$U' \cap (V(f_1, \ldots, f_r) \times_\mathcal{X} W)$ soit plate et localement
de présentation finie sur $W$. Notons que $\text{pr}_0(U')$ est un
sous-espace ouvert de $U$ contenant $p(z)$, car $\text{pr}_0$ est lisse, donc
ouvert. Or
$U' \cap (V(f_1, \ldots, f_r) \times_\mathcal{X} W) \to \mathcal{X}$
est plat et localement de présentation finie comme composé
$$
U' \cap (V(f_1, \ldots, f_r) \times_\mathcal{X} W) \to W \to \mathcal{X}.
$$
Par conséquent, Propriétés des champs,
Lemme \ref{stacks-properties-lemma-check-property-after-precomposing},
montre que $\text{pr}_0(U') \cap V(f_1, \ldots, f_r) \to \mathcal{X}$ est
plat et localement de présentation finie, comme souhaité.
\end{proof}

\begin{lemma}
\label{lemma-quasi-finite-at-point}
Soit $\mathcal{X}$ un champ algébrique. Considérons un diagramme cartésien
$$
\xymatrix{
U \ar[d] & F \ar[l]^p \ar[d] \\
\mathcal{X} & \Spec(k) \ar[l]
}
$$
où $U$ est un espace algébrique, $k$ un corps, et où $U \to \mathcal{X}$ est
localement de type fini. Soit $z \in |F|$ tel que $\dim_z(F) = 0$.
Alors, après avoir remplacé $U$ par un sous-espace ouvert contenant $p(z)$,
le morphisme
$$
U \longrightarrow \mathcal{X}
$$
est localement quasi-fini.
\end{lemma}

\begin{proof}
Comme $f : U \to \mathcal{X}$ est localement de type fini, il existe un ouvert
maximal $W(f) \subset U$ tel que la restriction
$f|_{W(f)} : W(f) \to \mathcal{X}$ soit localement quasi-finie, voir
Propriétés des champs, Remarque
\ref{stacks-properties-remark-local-source-apply}
(\ref{stacks-properties-item-loc-quasi-finite}).
Il suffit donc de montrer que $p(z)$ appartient à $W(f)$.
De plus, la remarque citée ci-dessus montre aussi que la formation de $W(f)$
commute à tout changement de base par un morphisme représentable par des
espaces algébriques. Il suffit donc de montrer que le morphisme
$F \to \Spec(k)$ est localement quasi-fini en $z$. Cela résulte immédiatement
de Morphismes d'espaces,
Lemme \ref{spaces-morphisms-lemma-locally-quasi-finite-rel-dimension-0}.
\end{proof}

\noindent
Un champ quasi-DM admet un ``recouvrement'' localement quasi-fini par un schéma.

\begin{theorem}
\label{theorem-quasi-DM}
Soit $\mathcal{X}$ un champ algébrique. Les conditions suivantes sont
équivalentes :
\begin{enumerate}
\item $\mathcal{X}$ est quasi-DM, et
\item il existe un schéma $W$ et un morphisme $W \to \mathcal{X}$ surjectif,
plat, localement de présentation finie et localement quasi-fini.
\end{enumerate}
\end{theorem}

\begin{proof}
L'implication (2) $\Rightarrow$ (1) est le contenu du
Lemme \ref{lemma-properties-covering-imply-diagonal}.
Supposons (1).
Soit $x \in |\mathcal{X}|$ un point de type fini. Nous allons construire un
schéma sur $\mathcal{X}$ qui ``convient'' dans un voisinage de $x$. À la fin
de la démonstration, nous prendrons la réunion disjointe de tous ces schémas.

\medskip\noindent
Soient $U$ un schéma affine, $U \to \mathcal{X}$ un morphisme lisse et
$u \in U$ un point fermé qui s'envoie sur $x$, voir le
Lemme \ref{lemma-point-finite-type}.
Notons, comme d'habitude, $u = \Spec(\kappa(u))$. Considérons le diagramme
commutatif suivant
$$
\xymatrix{
u \ar[d] & R \ar[l] \ar[d] \\
U \ar[d] & F \ar[d] \ar[l]^p \\
\mathcal{X} & u \ar[l]
}
$$
dont les deux carrés sont des carrés de produit fibré ; en particulier,
$R = u \times_\mathcal{X} u$. Dans la démonstration du
Lemme \ref{lemma-point-finite-type-monomorphism}, nous avons vu que
$(u, R, s, t, c)$ est un groupoïde en espaces algébriques dont les morphismes
$s, t$ sont localement de type fini. Soit $G \to u$ l'espace algébrique en groupes
stabilisateur (voir Groupoïdes en espaces, Définition
\ref{spaces-groupoids-definition-stabilizer-groupoid}).
Notons que
$$
G = R \times_{(u \times u)} u =
(u \times_\mathcal{X} u) \times_{(u \times u)} u =
\mathcal{X} \times_{\mathcal{X} \times \mathcal{X}} u.
$$
Comme $\mathcal{X}$ est quasi-DM, $G$ est localement quasi-fini sur $u$. Par
Compléments sur les groupoïdes en espaces, Lemme
\ref{spaces-more-groupoids-lemma-groupoid-on-field-dimension-equal-stabilizer}
on a $\dim(R) = 0$.

\medskip\noindent
Soit $e : u \to R$ l'identité du groupoïde. Le composé
$u \to R \to u$ est le morphisme identité de $u$.
Notons que $R \subset F$ est un sous-espace fermé, puisque $u \subset U$ est
un sous-schéma fermé. On peut donc aussi voir $e$ comme un point de $F$ ;
notons-le $z$. Considérons les homomorphismes d'anneaux locaux étales
$$
\mathcal{O}_{U, u}
\xrightarrow{p^\sharp}
\mathcal{O}_{F, \overline{z}}
\longrightarrow
\mathcal{O}_{R, \overline{e}}
$$
L'anneau $\mathcal{O}_{R, \overline{e}}$ est de dimension $0$ d'après le
premier paragraphe. D'autre part, le noyau de la seconde flèche est
$p^\sharp(\mathfrak m_u)\mathcal{O}_{F, \overline{z}}$, puisque $R$ est défini
dans $F$ par $\mathfrak m_u$. Ainsi,
$$
\mathfrak m_{\overline{z}} =
\sqrt{p^\sharp(\mathfrak m_u)\mathcal{O}_{F, \overline{z}}}
$$
D'autre part, puisque le morphisme $U \to \mathcal{X}$ est lisse,
$F \to u$ est un morphisme lisse d'espaces algébriques.
Cela signifie que $F$ est un espace algébrique régulier
(Espaces sur des corps, Lemme \ref{spaces-over-fields-lemma-smooth-regular}).
Par conséquent, $\mathcal{O}_{F, \overline{z}}$ est un anneau local régulier
(Propriétés des espaces, Lemme \ref{spaces-properties-lemma-regular}).
Rappelons qu'un anneau local régulier est de Cohen--Macaulay
(Algèbre, Lemme \ref{algebra-lemma-regular-ring-CM}).
Posons $d = \dim(\mathcal{O}_{F, \overline{z}})$. Par Algèbre,
Lemme \ref{algebra-lemma-find-sequence-image-regular}, on peut trouver
$f_1, \ldots, f_d \in \mathcal{O}_{U, u}$ dont les images
$p^\sharp(f_1), \ldots, p^\sharp(f_d)$ forment une suite régulière dans
$\mathcal{O}_{F, \overline{z}}$. Par le Lemme \ref{lemma-slice}, quitte à
restreindre $U$, nous pouvons supposer que
$Z = V(f_1, \ldots, f_d) \to \mathcal{X}$ est plat et localement de
présentation finie. Par construction,
$F_Z = Z \times_\mathcal{X} u$ est un sous-espace fermé de
$F = U \times_\mathcal{X} u$, le point $z$ appartient à ce sous-espace fermé,
et
$$
\dim(\mathcal{O}_{F_Z, \overline{z}}) = 0.
$$
Par Morphismes d'espaces,
Lemme \ref{spaces-morphisms-lemma-dimension-fibre-at-a-point},
il s'ensuit que $\dim_z(F_Z) = 0$, car le degré de transcendance de $z$
relativement à $u$ est nul. Le Lemme \ref{lemma-quasi-finite-at-point} montre
donc que, quitte à restreindre encore $U$, le morphisme
$Z \to \mathcal{X}$ est localement quasi-fini.

\medskip\noindent
Nous concluons que, pour tout point de type fini $x$ de $\mathcal{X}$, il
existe un morphisme $f_x : Z_x \to \mathcal{X}$ localement quasi-fini, plat
et localement de présentation finie, tel que $x$ appartienne à l'image de
$|f_x|$. Posons $W = \coprod_x Z_x$ et $f = \coprod f_x$. Alors $f$ est plat,
localement de présentation finie et localement quasi-fini. En particulier,
l'image de $|f|$ est ouverte, voir Propriétés des champs,
Lemme \ref{stacks-properties-lemma-topology-points}.
Par construction, l'image contient tous les points de type fini de
$\mathcal{X}$ ; le morphisme $f$ est donc surjectif par le
Lemme \ref{lemma-enough-finite-type-points} (et Propriétés des champs,
Lemme \ref{stacks-properties-lemma-characterize-surjective}).
\end{proof}

\begin{lemma}
\label{lemma-DM-residual-gerbe}
Soit $\mathcal{Z}$ un champ algébrique DM, localement noethérien et réduit,
tel que $|\mathcal{Z}|$ soit réduit à un point. Alors il existe un corps $k$
et un morphisme étale surjectif $\Spec(k) \to \mathcal{Z}$.
\end{lemma}

\begin{proof}
Par Propriétés des champs,
Lemme \ref{stacks-properties-lemma-unique-point-better},
il existe un corps $k$ et un morphisme $\Spec(k) \to \mathcal{Z}$ surjectif,
plat et localement de présentation finie. Posons $U = \Spec(k)$ et
$R = U \times_\mathcal{Z} U$ ; on obtient ainsi un groupoïde en espaces
algébriques $(U, R, s, t, c)$, voir Champs algébriques, Lemme
\ref{algebraic-lemma-criterion-map-representable-spaces-fibred-in-groupoids}.
Notons que, par Champs algébriques,
Remarque \ref{algebraic-remark-flat-fp-presentation}, on a une équivalence
$$
f_{can} : [U/R] \longrightarrow \mathcal{Z}
$$
Les projections $s, t : R \to U$ sont localement de présentation finie.
Comme $\mathcal{Z}$ est DM, l'espace algébrique en groupes stabilisateur
$$
G = U \times_{U \times U} R = U \times_{U \times U} (U \times_\mathcal{Z} U) =
U \times_{\mathcal{Z} \times \mathcal{Z}, \Delta_\mathcal{Z}} \mathcal{Z}
$$
est non ramifié sur $U$. En particulier, $\dim(G) = 0$ et, par Compléments
sur les groupoïdes en espaces, Lemme
\ref{spaces-more-groupoids-lemma-groupoid-on-field-dimension-equal-stabilizer}
on a $\dim(R) = 0$. Il s'ensuit que $R$ est un schéma, voir Espaces sur des
corps, Lemme
\ref{spaces-over-fields-lemma-locally-finite-type-dim-zero}.
Par Variétés, Lemme \ref{varieties-lemma-algebraic-scheme-dim-0},
$R$ (ainsi que $G$) est réunion disjointe de spectres d'anneaux locaux
artiniens, finis sur $k$ aussi bien par $s$ que par $t$. Soit
$P = \Spec(A) \subset R$ le sous-schéma ouvert et fermé dont le point
sous-jacent est l'identité $e$ du schéma en groupoïdes $(U, R, s, t, c)$.
Comme $s \circ e = t \circ e = \text{id}_{\Spec(k)}$, l'anneau $A$ est local
artinien et son corps résiduel s'identifie à $k$ aussi bien par
$s^\sharp : k \to A$ que par $t^\sharp : k \to A$.
Notons que $s, t : \Spec(A) \to \Spec(k)$ sont finis (par le lemme cité
ci-dessus). Puisque $G \to \Spec(k)$ est non ramifié, on a
$$
G \cap P = P \times_{U \times U} U = \Spec(A \otimes_{k \otimes k} k)
$$
qui est non ramifié sur $k$. D'autre part,
$A \otimes_{k \otimes k} k$ est local, comme quotient de $A$, et admet une
surjection sur $k$. On en conclut que $A \otimes_{k \otimes k} k = k$. Il s'ensuit que
$P \to U \times U$ est universellement injectif (car $P$ n'a qu'un point,
de corps résiduel $k$), non ramifié (d'après le calcul ci-dessus de la fibre
sur l'unique point image) et de type fini (parce que $s, t$ le sont) ; c'est
donc un monomorphisme (voir Morphismes étales, Lemme
\ref{etale-lemma-universally-injective-unramified}).
Ainsi, $s|_P, t|_P : P \to U$ définissent une relation d'équivalence finie
plate. La Proposition \ref{groupoids-proposition-finite-flat-equivalence} de
Groupoïdes montre alors que $U/P$ existe et est un schéma $\overline{U}$.
De plus, $U \to \overline{U}$ est fini localement libre et
$P = U \times_{\overline{U}} U$.
En fait, $\overline{U} = \Spec(k_0)$, où $k_0 \subset k$ est l'anneau des
fonctions invariantes par $R$. Comme $k$ est un corps, la définition donnée
dans Groupoïdes, Équation (\ref{groupoids-equation-invariants}), montre que
$k_0$ est un corps.

\medskip\noindent
Nous affirmons que
\begin{equation}
\label{equation-etale-covering}
\Spec(k_0) = \overline{U} = U/P \to [U/R] = \mathcal{Z}
\end{equation}
est le morphisme étale surjectif cherché. Il résulte de Propriétés des champs,
Lemme \ref{stacks-properties-lemma-flat-cover-by-field}, que ce morphisme est
surjectif. Il suffit donc de montrer que (\ref{equation-etale-covering}) est
étale\footnote{Nous invitons le lecteur à trouver sa propre démonstration de
ce fait. L'argument a en réalité beaucoup en commun avec l'argument final de
la démonstration d'Amorçage, Théorème
\ref{bootstrap-theorem-final-bootstrap} ; il devrait donc probablement être
isolé quelque part dans un lemme qui lui soit propre.}.
Au lieu de démontrer directement l'étalité, appliquons d'abord Amorçage,
Lemme \ref{bootstrap-lemma-divide-subgroupoid}. Il existe alors un schéma en
groupoïdes
$(\overline{U}, \overline{R}, \overline{s}, \overline{t}, \overline{c})$
tel que $(U, R, s, t, c)$ soit la restriction de
$(\overline{U}, \overline{R}, \overline{s}, \overline{t}, \overline{c})$
par le morphisme quotient $U \to \overline{U}$.
(Nous avons vérifié ci-dessus toutes les hypothèses du lemme, hormis que
$j : R \to U \times U$ est séparé et localement quasi-fini ; cela résulte du
fait que $R$ est un schéma séparé localement quasi-fini sur $k$.) Comme
$U \to \overline{U}$ est fini localement libre,
$[U/R] \to [\overline{U}/\overline{R}]$ est une équivalence, voir Groupoïdes
en espaces,
Lemme \ref{spaces-groupoids-lemma-quotient-stack-restrict-equivalence}.

\medskip\noindent
Notons que $s, t$ sont les changements de base des morphismes
$\overline{s}, \overline{t}$ par $U \to \overline{U}$.
Comme $\{U \to \overline{U}\}$ est un recouvrement fppf, on en conclut que
$\overline{s}, \overline{t}$ sont plats, localement de présentation finie et
localement quasi-finis, voir Descente, Lemmes
\ref{descent-lemma-descending-property-flat},
\ref{descent-lemma-descending-property-locally-finite-presentation} et
\ref{descent-lemma-descending-property-quasi-finite}.
Considérons le diagramme commutatif
$$
\xymatrix{
U \times_{\overline{U}} U \ar@{=}[r] \ar[rd] & P \ar[r] \ar[d] & R \ar[d] \\
& \overline{U} \ar[r]^{\overline{e}} & \overline{R}
}
$$
Il est général, pour les restrictions, que les quatre sommets extérieurs
forment un diagramme cartésien. L'égalité montre que le carré intérieur est
cartésien. Comme $P$ est ouvert dans $R$, on en conclut que $\overline{e}$
est une immersion ouverte par Descente,
Lemme \ref{descent-lemma-descending-property-open-immersion}.

\medskip\noindent
Bien entendu, si $\overline{e}$ est une immersion ouverte et si
$\overline{s}, \overline{t}$ sont plats et localement de présentation finie,
alors les morphismes $\overline{t}, \overline{s}$ sont étales.
On peut le voir, par exemple, en appliquant Compléments sur les groupoïdes,
Lemme \ref{more-groupoids-lemma-sheaf-differentials}, qui montre que
$\Omega_{\overline{R}/\overline{U}} = 0$ implique que
$\overline{s}, \overline{t} : \overline{R} \to \overline{U}$ sont non
ramifiés (voir Morphismes,
Lemme \ref{morphisms-lemma-unramified-omega-zero}), ce qui implique à son tour
que $\overline{s}, \overline{t}$ sont étales (voir Morphismes,
Lemme \ref{morphisms-lemma-flat-unramified-etale}).
Ainsi, $\mathcal{Z} = [\overline{U}/\overline{R}]$ est une présentation étale
du champ algébrique $\mathcal{Z}$, et l'on conclut que
$\overline{U} \to \mathcal{Z}$ est étale par Propriétés des champs, Lemme
\ref{stacks-properties-lemma-check-property-covering}.
\end{proof}

\begin{lemma}
\label{lemma-etale-at-point}
Soit $\mathcal{X}$ un champ algébrique. Considérons un diagramme cartésien
$$
\xymatrix{
U \ar[d] & F \ar[l]^p \ar[d] \\
\mathcal{X} & \Spec(k) \ar[l]
}
$$
où $U$ est un espace algébrique, $k$ un corps, et où $U \to \mathcal{X}$ est
plat et localement de présentation finie. Soit $z \in |F|$ tel que
$F \to \Spec(k)$ soit non ramifié en $z$. Alors, après avoir remplacé $U$ par
un sous-espace ouvert contenant $p(z)$, le morphisme
$$
U \longrightarrow \mathcal{X}
$$
est étale.
\end{lemma}

\begin{proof}
Comme $f : U \to \mathcal{X}$ est plat et localement de présentation finie,
il existe un ouvert maximal $W(f) \subset U$ tel que la restriction
$f|_{W(f)} : W(f) \to \mathcal{X}$ soit étale, voir Propriétés des champs,
Remarque
\ref{stacks-properties-remark-local-source-apply}
(\ref{stacks-properties-item-etale}).
Il suffit donc de montrer que $p(z)$ appartient à $W(f)$.
De plus, la remarque citée ci-dessus montre que la formation de $W(f)$ commute
à tout changement de base par un morphisme représentable par des espaces
algébriques. Il suffit donc de montrer que $F \to \Spec(k)$ est étale en $z$.
Ce morphisme est plat et localement de présentation finie, comme changement
de base de $U \to \mathcal{X}$, et, puisque $F \to \Spec(k)$ est non ramifié
en $z$ par hypothèse,
le résultat découle donc de Morphismes d'espaces,
Lemme \ref{spaces-morphisms-lemma-unramified-flat-lfp-etale}.
\end{proof}

\noindent
Un champ DM est un champ de Deligne--Mumford.

\begin{theorem}
\label{theorem-DM}
Soit $\mathcal{X}$ un champ algébrique. Les conditions suivantes sont
équivalentes :
\begin{enumerate}
\item $\mathcal{X}$ est DM,
\item $\mathcal{X}$ est de Deligne--Mumford, et
\item il existe un schéma $W$ et un morphisme étale surjectif
$W \to \mathcal{X}$.
\end{enumerate}
\end{theorem}

\begin{proof}
Rappelons que (3) est la définition de (2), voir Champs algébriques,
Définition \ref{algebraic-definition-deligne-mumford}.
L'implication (3) $\Rightarrow$ (1) est le contenu du
Lemme \ref{lemma-properties-covering-imply-diagonal}.
Supposons (1). Soit $x \in |\mathcal{X}|$ un point de type fini.
Nous allons construire un schéma sur $\mathcal{X}$ qui ``convient'' dans un
voisinage de $x$. À la fin de la démonstration, nous prendrons la réunion
disjointe de tous ces schémas.

\medskip\noindent
Par le Lemme \ref{lemma-point-finite-type-monomorphism}, la gerbe résiduelle
$\mathcal{Z}_x$ de $\mathcal{X}$ en $x$ existe et
$\mathcal{Z}_x \to \mathcal{X}$ est localement de type fini. Par le
Lemme \ref{lemma-separation-properties-residual-gerbe}, le champ algébrique
$\mathcal{Z}_x$ est DM. Par le Lemme \ref{lemma-DM-residual-gerbe}, il existe
un corps $k$ et un morphisme étale surjectif
$z : \Spec(k) \to \mathcal{Z}_x$.
En particulier, le composé $x : \Spec(k) \to \mathcal{X}$ est localement de
type fini (par Morphismes d'espaces, Lemmes
\ref{spaces-morphisms-lemma-composition-finite-type} et
\ref{spaces-morphisms-lemma-etale-locally-finite-type}).

\medskip\noindent
Choisissons un schéma $U$ et un morphisme lisse $U \to \mathcal{X}$ tels que
$x$ appartienne à l'image de $|U| \to |\mathcal{X}|$.
Considérons le carré de produit fibré suivant
$$
\xymatrix{
U \ar[d] & F  \ar[l] \ar[d] \\
\mathcal{X} & \Spec(k) \ar[l]_-x
}
$$
autrement dit $F = U \times_{\mathcal{X}, x} \Spec(k)$. Par Propriétés des
champs, Lemme \ref{stacks-properties-lemma-points-cartesian}, l'espace $F$
est non vide. Comme $\mathcal{Z}_x \to \mathcal{X}$ est un monomorphisme,
on a
$$
\Spec(k) \times_{z, \mathcal{Z}_x, z} \Spec(k)
=
\Spec(k) \times_{x, \mathcal{X}, x} \Spec(k)
$$
dont les projections vers $\Spec(k)$ sont étales par construction de $z$.
Comme
$$
F \times_U F =
(\Spec(k) \times_\mathcal{X} \Spec(k))
\times_{\Spec(k)} F
$$
les projections $F \times_U F \to F$ sont elles aussi étales.
Il s'ensuit que $\Delta_{F/U} : F \to F \times_U F$ est étale (voir
Morphismes d'espaces, Lemme \ref{spaces-morphisms-lemma-etale-permanence}).
Par Morphismes d'espaces, Lemme
\ref{spaces-morphisms-lemma-etale-universally-injective-open}
cela implique que $\Delta_{F/U}$ est une immersion ouverte ; enfin, par
Morphismes d'espaces, Lemme
\ref{spaces-morphisms-lemma-diagonal-unramified-morphism}
il s'ensuit que $F \to U$ est non ramifié.

\medskip\noindent
Choisissons un schéma affine non vide $V$ et un morphisme étale $V \to F$.
(On pourrait éviter ce choix en travaillant directement avec $F$, mais il
semble ainsi plus facile d'expliquer ce qui se passe.) Représentons la
situation par le diagramme
$$
\xymatrix{
U \ar[d] & F  \ar[l] \ar[d] & V \ar[l] \ar[ld] \\
\mathcal{X} & \Spec(k) \ar[l]_-x
}
$$
Alors $V \to \Spec(k)$ est un morphisme lisse de schémas et $V \to U$ un
morphisme non ramifié de schémas (voir Morphismes d'espaces, Lemmes
\ref{spaces-morphisms-lemma-composition-smooth} et
\ref{spaces-morphisms-lemma-composition-unramified}).
Choisissons un point fermé $v \in V$ tel que l'extension
$k \subset \kappa(v)$ soit finie séparable, voir Variétés,
Lemme \ref{varieties-lemma-smooth-separable-closed-points-dense}.
Soit $u \in U$ le point image. L'anneau local $\mathcal{O}_{V, v}$ est
régulier (voir Variétés, Lemme \ref{varieties-lemma-smooth-regular}), et
l'homomorphisme d'anneaux locaux
$$
\varphi : \mathcal{O}_{U, u} \longrightarrow \mathcal{O}_{V, v}
$$
induit par le morphisme $V \to U$ vérifie
$\varphi(\mathfrak m_u)\mathcal{O}_{V, v} = \mathfrak m_v$,
voir Morphismes, Lemme \ref{morphisms-lemma-unramified-at-point}.
On peut donc trouver $f_1, \ldots, f_d \in \mathcal{O}_{U, u}$ tels que les
images $\varphi(f_1), \ldots, \varphi(f_d)$ forment une base de
$\mathfrak m_v/\mathfrak m_v^2$ sur $\kappa(v)$.
Comme $\mathcal{O}_{V, v}$ est un anneau local régulier, il s'ensuit que
$\varphi(f_1), \ldots, \varphi(f_d)$ forment une suite régulière dans
$\mathcal{O}_{V, v}$ (voir Algèbre,
Lemme \ref{algebra-lemma-regular-ring-CM}).
Après avoir remplacé $U$ par un voisinage ouvert de $u$, on peut supposer que
$f_1, \ldots, f_d \in \Gamma(U, \mathcal{O}_U)$. En remplaçant encore $U$
par un voisinage ouvert éventuellement plus petit de $u$, on peut supposer que
$V(f_1, \ldots, f_d) \to \mathcal{X}$ est plat et localement de présentation
finie, voir le Lemme \ref{lemma-slice}.
Par construction,
$$
V(f_1, \ldots, f_d) \times_\mathcal{X} \Spec(k)
\longleftarrow
V(f_1, \ldots, f_d) \times_U V
$$
est étale, et $V(f_1, \ldots, f_d) \times_U V$ est le sous-schéma fermé
$T \subset V$ défini par $f_1|_V, \ldots, f_d|_V$.
Par construction, $v \in T$ et
$$
\mathcal{O}_{T, v} =
\mathcal{O}_{V, v}/(\varphi(f_1), \ldots, \varphi(f_d)) = \kappa(v)
$$
qui est une extension finie séparable de $k$. Il s'ensuit que
$T \to \Spec(k)$ est non ramifié en $v$, voir Morphismes,
Lemme \ref{morphisms-lemma-unramified-at-point}.
Par définition d'un morphisme non ramifié d'espaces algébriques, cela signifie
que $V(f_1, \ldots, f_d) \times_\mathcal{X} \Spec(k) \to \Spec(k)$ est non
ramifié au point image de $v$ dans
$V(f_1, \ldots, f_d) \times_\mathcal{X} \Spec(k)$.
En appliquant le Lemme \ref{lemma-etale-at-point}, on voit que, quitte à
restreindre encore $U$ à un voisinage ouvert de $u$, le morphisme
$V(f_1, \ldots, f_d) \to \mathcal{X}$ est étale.

\medskip\noindent
Nous concluons que, pour tout point de type fini $x$ de $\mathcal{X}$, il
existe un morphisme étale $f_x : W_x \to \mathcal{X}$ tel que $x$ appartienne
à l'image de $|f_x|$. Posons $W = \coprod_x W_x$ et $f = \coprod f_x$.
Alors $f$ est étale. En particulier, l'image de $|f|$ est ouverte, voir
Propriétés des champs, Lemme \ref{stacks-properties-lemma-topology-points}.
Par construction, l'image contient tous les points de type fini de
$\mathcal{X}$ ; le morphisme $f$ est donc surjectif par le
Lemme \ref{lemma-enough-finite-type-points} (et Propriétés des champs,
Lemme \ref{stacks-properties-lemma-characterize-surjective}).
\end{proof}

\noindent
Voici un corollaire utile qui montre que les ``fibres'' d'un morphisme DM
de champs algébriques sont de Deligne--Mumford.

\begin{lemma}
\label{lemma-DM}
Soit $f : \mathcal{X} \to \mathcal{Y}$ un morphisme DM de champs
algébriques. Alors
\begin{enumerate}
\item pour tout champ algébrique DM $\mathcal{Z}$ et tout morphisme
$\mathcal{Z} \to \mathcal{Y}$, il existe un schéma et un morphisme
étale surjectif
$U \to \mathcal{X} \times_\mathcal{Y} \mathcal{Z}$.
\item pour tout espace algébrique $Z$ et tout morphisme
$Z \to \mathcal{Y}$, il existe un schéma et un morphisme
étale surjectif
$U \to \mathcal{X} \times_\mathcal{Y} Z$.
\end{enumerate}
\end{lemma}

\begin{proof}
Preuve de (1). Comme $f$ est DM, son changement de base
$\mathcal{X} \times_\mathcal{Y} \mathcal{Z} \to \mathcal{Z}$ est DM
par le Lemme \ref{lemma-base-change-separated}.
Puisque $\mathcal{Z}$ est DM, il en résulte que
$\mathcal{X} \times_\mathcal{Y} \mathcal{Z}$ est DM
par le Lemme \ref{lemma-separated-over-separated}. Il existe donc un
schéma $U$ et un morphisme étale surjectif
$U \to \mathcal{X} \times_\mathcal{Y} \mathcal{Z}$, voir le
Théorème \ref{theorem-DM}.
La partie (2) est un cas particulier de (1), puisqu'un espace algébrique
(considéré comme un champ algébrique) est DM par le
Lemme \ref{lemma-trivial-implications}.
\end{proof}






\section{Le lieu de Deligne--Mumford}
\label{section-DM}

\noindent
Tout champ algébrique admet un plus grand sous-champ ouvert qui soit un champ
de Deligne--Mumford ; cela est plus ou moins clair, mais nous en donnons tout
de même la démonstration ci-dessous. Bien entendu, ce sous-champ peut être
vide, par exemple si $X = [\Spec(\mathbf{Z})/\mathbf{G}_{m, \mathbf{Z}}]$.
Nous caractériserons ci-dessous les points du lieu DM.

\begin{lemma}
\label{lemma-open-DM-locus}
Soit $\mathcal{X}$ un champ algébrique. Il existe des sous-champs ouverts
$$
\mathcal{X}'' \subset \mathcal{X}' \subset \mathcal{X}
$$
tels que $\mathcal{X}''$ soit DM et $\mathcal{X}'$ quasi-DM, et ce sont les
plus grands sous-champs ouverts possédant respectivement ces propriétés.
\end{lemma}

\begin{proof}
Tout ce que nous affirmons ici est que, si
$\mathcal{U} \subset \mathcal{X}$ et $\mathcal{V} \subset \mathcal{X}$
sont des sous-champs ouverts qui sont DM, alors le sous-champ ouvert
$\mathcal{W} \subset \mathcal{X}$ défini par
$|\mathcal{W}| = |\mathcal{U}| \cup |\mathcal{V}|$ est lui aussi DM.
(Il en va de même pour quasi-DM.) Même s'il y a là un léger artifice,
utilisons le Théorème \ref{theorem-DM} pour le démontrer.
D'après ce théorème, on peut choisir des schémas $U$ et $V$ et des morphismes
étales surjectifs $U \to \mathcal{U}$ et $V \to \mathcal{V}$.
Alors, bien entendu, $U \amalg V \to \mathcal{W}$ est étale et surjectif.
Le cas quasi-DM se démontre exactement de la même manière en utilisant le
Théorème \ref{theorem-quasi-DM}.
\end{proof}

\begin{lemma}
\label{lemma-points-DM-locus}
Soit $\mathcal{X}$ un champ algébrique. Soit $x \in |\mathcal{X}|$ le point
correspondant à $x : \Spec(k) \to \mathcal{X}$. Soit $G_x/k$ l'espace
algébrique en groupes des automorphismes de $x$. Alors
\begin{enumerate}
\item $x$ appartient au lieu DM de $\mathcal{X}$ si et seulement si
$G_x \to \Spec(k)$ est non ramifié, et
\item $x$ appartient au lieu quasi-DM de $\mathcal{X}$ si et seulement si
$G_x \to \Spec(k)$ est localement quasi-fini.
\end{enumerate}
\end{lemma}

\begin{proof}
Preuve de (2). Choisissons un schéma $U$ et un morphisme lisse surjectif
$U \to \mathcal{X}$. Considérons le produit fibré
$$
\xymatrix{
G \ar[r] \ar[d] & \mathcal{I}_\mathcal{X} \ar[d] \\
U \ar[r] & \mathcal{X}
}
$$
Rappelons que $G$ est l'espace algébrique en groupes des automorphismes de
$U \to \mathcal{X}$. Par Groupoïdes en espaces, Lemme
\ref{spaces-groupoids-lemma-open-over-which-unramified-or-lqf}
il existe un plus grand sous-schéma ouvert $U' \subset U$ tel que
$G_{U'} \to U'$ soit localement quasi-fini.
En outre, la formation de $U'$ commute à tout changement de base.
En particulier, les deux images réciproques de $U'$ dans
$R = U \times_\mathcal{X} U$ sont le même sous-espace ouvert de $R$
(les deux morphismes $R \to \mathcal{X}$ sont en effet isomorphes, et leurs
espaces algébriques en groupes d'automorphismes sont donc isomorphes).
Par conséquent, $U'$ est l'image réciproque d'un sous-champ ouvert
$\mathcal{X}' \subset \mathcal{X}$ par Propriétés des champs, Lemme
\ref{stacks-properties-lemma-substacks-presentation}
et l'on a un diagramme cartésien
$$
\xymatrix{
G_{U'} \ar[r] \ar[d] & \mathcal{I}_{\mathcal{X}'} \ar[d] \\
U' \ar[r] & \mathcal{X}'
}
$$
Ainsi, le morphisme $\mathcal{I}_{\mathcal{X}'} \to \mathcal{X}'$ est
localement quasi-fini, et le Lemme \ref{lemma-diagonal-diagonal}, partie (5),
montre que $\mathcal{X}'$ est quasi-DM. D'autre part, si
$\mathcal{W} \subset \mathcal{X}$ est un sous-champ ouvert quasi-DM,
l'image réciproque $W \subset U$ de $\mathcal{W}$ est contenue dans $U'$
par construction de $U'$, puisque
$\mathcal{I}_\mathcal{W} =
\mathcal{W} \times_\mathcal{X} \mathcal{I}_\mathcal{X}$
est localement quasi-fini sur $\mathcal{W}$.
Ainsi, $\mathcal{X}'$ est le lieu quasi-DM.
Enfin, choisissons une extension de corps $K/k$ et un diagramme
$2$-commutatif
$$
\xymatrix{
\Spec(K) \ar[r] \ar[d] & \Spec(k) \ar[d]^x \\
U \ar[r] & \mathcal{X}
}
$$
On obtient alors un isomorphisme
$G_x \times_{\Spec(k)} \Spec(K) \cong G \times_U \Spec(K)$
entre espaces algébriques en groupes sur $K$. Par conséquent, $G_x$ est
localement quasi-fini sur $k$ si et seulement si l'image de
$\Spec(K) \to U$ est contenue dans $U'$ (utiliser la commutation de la
formation de $U'$ aux changements de base et Groupoïdes en espaces,
Lemme
\ref{spaces-groupoids-lemma-open-over-which-unramified-or-lqf}
appliqué à $\Spec(K) \to \Spec(k)$ et à $G_x$).
Ceci achève la preuve de (2). Celle de (1) est exactement la même.
\end{proof}









\section{Morphismes localement quasi-finis}
\label{section-locally-quasi-finite}

\noindent
La propriété ``localement quasi-fini'' des morphismes d'espaces algébriques
n'est pas locale pour la topologie lisse sur la source et la cible ; nous ne
pouvons donc pas utiliser les résultats de la
section \ref{section-local-source-target} pour définir les morphismes
localement quasi-finis de champs algébriques. Nous savons déjà ce que signifie
être localement quasi-fini pour un morphisme de champs algébriques
représentable par des espaces algébriques, voir Propriétés des champs, section
\ref{stacks-properties-section-properties-morphisms}.
Pour trouver une condition convenant aux morphismes généraux, faisons
l'observation suivante.

\begin{lemma}
\label{lemma-representable-by-spaces-quasi-finite}
Soit $f : \mathcal{X} \to \mathcal{Y}$ un morphisme de champs algébriques.
Supposons que $f$ soit représentable par des espaces algébriques.
Les conditions suivantes sont équivalentes :
\begin{enumerate}
\item $f$ est localement quasi-fini (au sens de Propriétés des champs,
section \ref{stacks-properties-section-properties-morphisms}), et
\item $f$ est localement de type fini et, pour tout morphisme
$\Spec(k) \to \mathcal{Y}$, où $k$ est un corps, l'espace
$|\Spec(k) \times_\mathcal{Y} \mathcal{X}|$ est discret.
\end{enumerate}
\end{lemma}

\begin{proof}
Supposons (1). Le morphisme d'espaces algébriques
$\mathcal{X}_k \to \Spec(k)$ est alors localement quasi-fini, car il est un
changement de base de $f$. Ainsi, $|\mathcal{X}_k|$ est discret par
Morphismes d'espaces, Lemme
\ref{spaces-morphisms-lemma-locally-quasi-finite}.
Réciproquement, supposons (2). Choisissons un morphisme lisse surjectif
$V \to \mathcal{Y}$, où $V$ est un schéma. Il suffit de montrer que le
morphisme d'espaces algébriques $V \times_\mathcal{Y} \mathcal{X} \to V$
est localement quasi-fini, voir Propriétés des champs, Lemme
\ref{stacks-properties-lemma-check-property-covering}.
Par hypothèse, le morphisme $V \times_\mathcal{Y} \mathcal{X} \to V$ est
localement de type fini. Pour tout morphisme $\Spec(k) \to V$, où $k$ est
un corps,
$$
\Spec(k) \times_V (V \times_\mathcal{Y} \mathcal{X}) =
\Spec(k) \times_\mathcal{Y} \mathcal{X}
$$
a par hypothèse un espace de points discret. On en conclut que
$V \times_\mathcal{Y} \mathcal{X} \to V$ est localement quasi-fini par
Morphismes d'espaces, Lemme
\ref{spaces-morphisms-lemma-locally-quasi-finite}.
\end{proof}

\noindent
Un morphisme de champs algébriques représentable par des espaces algébriques
est quasi-DM, voir le Lemme \ref{lemma-trivial-implications}.
Avec le lemme précédent, cela montre que la définition suivante est compatible
avec la notion déjà existante dans le cas des morphismes représentables par
des espaces algébriques.

\begin{definition}
\label{definition-locally-quasi-finite}
Soit $f : \mathcal{X} \to \mathcal{Y}$ un morphisme de champs algébriques.
On dit que $f$ est {\it localement quasi-fini} si $f$ est quasi-DM, localement
de type fini et si, pour tout morphisme $\Spec(k) \to \mathcal{Y}$, où
$k$ est un corps, l'espace $|\mathcal{X}_k|$ est discret.
\end{definition}

\noindent
La condition que $f$ soit quasi-DM est naturelle. Par exemple, soit $k$ un
corps et considérons le morphisme
$\pi : [\Spec(k)/\mathbf{G}_m] \to \Spec(k)$
dont les fibres sont réduites à un point et qui est localement de type fini.
Comme nous le verrons plus loin, ce morphisme est lisse de dimension relative
$-1$, tandis que nous voulons que les morphismes localement quasi-finis aient
la dimension relative $0$. Notons aussi que la section
$\Spec(k) \to [\Spec(k)/\mathbf{G}_m]$ n'a pas de fibres discrètes et n'est
donc pas localement quasi-finie. Nous souhaitons en outre la propriété de
permanence suivante : si $f : \mathcal{X} \to \mathcal{X}'$ est un morphisme
de champs algébriques localement quasi-fini sur le champ algébrique
$\mathcal{Y}$, alors $f$ est localement quasi-fini (un énoncé légèrement plus
fort est en fait vrai, voir le
Lemme \ref{lemma-locally-quasi-finite-permanence}).

\medskip\noindent
La définition précédente est également justifiée par le
Lemme \ref{lemma-characterize-locally-quasi-finite} ci-dessous, qui caractérise
la propriété d'être localement quasi-fini par l'existence de
``présentations'' ou de ``recouvrements'' convenables de
$\mathcal{X}$ et de $\mathcal{Y}$.

\begin{lemma}
\label{lemma-base-change-locally-quasi-finite}
Tout changement de base d'un morphisme localement quasi-fini est localement
quasi-fini.
\end{lemma}

\begin{proof}
Nous l'avons vu pour les morphismes quasi-DM au
Lemme \ref{lemma-base-change-separated} et pour les morphismes localement de
type fini au Lemme \ref{lemma-base-change-finite-type}.
La condition sur les fibres se conserve manifestement par changement de base.
\end{proof}

\begin{lemma}
\label{lemma-locally-quasi-finite-over-field}
Soit $\mathcal{X} \to \Spec(k)$ un morphisme localement quasi-fini, où
$\mathcal{X}$ est un champ algébrique et $k$ un corps. Soit
$f : V \to \mathcal{X}$ un morphisme localement quasi-fini, où $V$ est un
schéma. Alors $V \to \Spec(k)$ est localement quasi-fini.
\end{lemma}

\begin{proof}
Le Lemme \ref{lemma-composition-finite-type} montre que
$V \to \Spec(k)$ est localement de type fini.
Supposons, afin d'obtenir une contradiction, que $V \to \Spec(k)$ ne soit
pas localement quasi-fini. Il existe alors une spécialisation non triviale
$v \leadsto v'$ de points de $V$, voir Morphismes, Lemme
\ref{morphisms-lemma-quasi-finite-at-point-characterize}.
En particulier,
$\text{trdeg}_k(\kappa(v)) > \text{trdeg}_k(\kappa(v'))$, voir Morphismes,
Lemme \ref{morphisms-lemma-dimension-fibre-specialization}.
Comme $|\mathcal{X}|$ est discret, on a $|f|(v) = |f|(v')$.
Considérons $R = V \times_\mathcal{X} V$. Alors $R$ est un espace algébrique
et les projections $s, t : R \to V$ sont localement quasi-finies comme
changements de base de $V \to \mathcal{X}$ (ce morphisme est représentable
par des espaces algébriques, si bien que cela résulte de la discussion de
Propriétés des champs, section
\ref{stacks-properties-section-properties-morphisms}).
Par Propriétés des champs,
Lemme \ref{stacks-properties-lemma-points-cartesian}, il existe
$r \in |R|$ tel que $s(r) = v$ et $t(r) = v'$.
Par Morphismes d'espaces,
Lemme \ref{spaces-morphisms-lemma-compare-tr-deg}, le degré de transcendance
de $v/k$ est égal à celui de $r/k$, lui-même égal à celui de $v'/k$.
Cette contradiction démontre le lemme.
\end{proof}

\begin{lemma}
\label{lemma-composition-locally-quasi-finite}
Un composé de morphismes localement quasi-finis est localement quasi-fini.
\end{lemma}

\begin{proof}
Nous l'avons vu pour les morphismes quasi-DM au
Lemme \ref{lemma-composition-separated} et pour les morphismes localement de
type fini au Lemme \ref{lemma-composition-finite-type}.
Supposons que $\mathcal{X} \to \mathcal{Y}$ et
$\mathcal{Y} \to \mathcal{Z}$ soient localement quasi-finis. Soit $k$ un
corps et soit $\Spec(k) \to \mathcal{Z}$ un morphisme.
Il suffit de montrer que $|\mathcal{X}_k|$ est discret. Par le
Lemme \ref{lemma-base-change-locally-quasi-finite}, les morphismes
$\mathcal{X}_k \to \mathcal{Y}_k$ et
$\mathcal{Y}_k \to \Spec(k)$ sont localement quasi-finis.
En particulier, $\mathcal{Y}_k$ est un champ algébrique quasi-DM, voir le
Lemme \ref{lemma-separated-implies-morphism-separated}.
Par le Théorème \ref{theorem-quasi-DM}, il existe un schéma $V$ et un
morphisme surjectif, plat, localement de présentation finie et localement
quasi-fini $V \to \mathcal{Y}_k$. Le
Lemme \ref{lemma-locally-quasi-finite-over-field} montre que $V$ est
localement quasi-fini sur $k$ ; en particulier, $|V|$ est discret. Le morphisme
$V \times_{\mathcal{Y}_k} \mathcal{X}_k \to \mathcal{X}_k$ est surjectif,
plat et localement de présentation finie ; l'application
$|V \times_{\mathcal{Y}_k} \mathcal{X}_k| \to |\mathcal{X}_k|$ est donc
surjective et ouverte. Il suffit ainsi de montrer que
$|V \times_{\mathcal{Y}_k} \mathcal{X}_k|$ est discret.
Notons que $V$ est réunion disjointe de spectres de $k$-algèbres locales
artiniennes $A_i$, de corps résiduels $k_i$, voir Variétés,
Lemme \ref{varieties-lemma-algebraic-scheme-dim-0}.
Il suffit donc de montrer que chacun des espaces
$$
|\Spec(A_i) \times_{\mathcal{Y}_k} \mathcal{X}_k| =
|\Spec(k_i) \times_{\mathcal{Y}_k} \mathcal{X}_k| =
|\Spec(k_i) \times_\mathcal{Y} \mathcal{X}|
$$
est discret, ce qui résulte de l'hypothèse que
$\mathcal{X} \to \mathcal{Y}$ est localement quasi-fini.
\end{proof}

\noindent
Avant de caractériser les morphismes localement quasi-finis au moyen de
recouvrements, faisons-le pour les morphismes quasi-DM.

\begin{lemma}
\label{lemma-characterize-quasi-DM}
Soit $f : \mathcal{X} \to \mathcal{Y}$ un morphisme de champs algébriques.
Les conditions suivantes sont équivalentes :
\begin{enumerate}
\item $f$ est quasi-DM,
\item pour tout morphisme $V \to \mathcal{Y}$, où $V$ est un espace
algébrique, il existe un morphisme surjectif, plat, localement de présentation
finie et localement quasi-fini
$U \to \mathcal{X} \times_\mathcal{Y} V$, où $U$ est un espace algébrique,
et
\item il existe des espaces algébriques $U$, $V$, un morphisme
$V \to \mathcal{Y}$ surjectif, plat et localement de présentation finie,
et un morphisme $U \to \mathcal{X} \times_\mathcal{Y} V$ surjectif, plat,
localement de présentation finie et localement quasi-fini.
\end{enumerate}
\end{lemma}

\begin{proof}
L'implication (2) $\Rightarrow$ (3) est immédiate.

\medskip\noindent
Supposons (1) et soit $V \to \mathcal{Y}$ comme en (2). Alors
$\mathcal{X} \times_\mathcal{Y} V \to V$ est quasi-DM, voir le
Lemme \ref{lemma-base-change-separated}.
Par le Lemme \ref{lemma-trivial-implications}, l'espace algébrique $V$ est DM,
donc quasi-DM. Ainsi, $\mathcal{X} \times_\mathcal{Y} V$ est quasi-DM par le
Lemme \ref{lemma-separated-over-separated}.
On peut donc appliquer le Théorème \ref{theorem-quasi-DM} pour obtenir le
morphisme $U \to \mathcal{X} \times_\mathcal{Y} V$ de (2).

\medskip\noindent
Supposons (3). Soient $V \to \mathcal{Y}$ et
$U \to \mathcal{X} \times_\mathcal{Y} V$ comme en (3).
Pour montrer que $f$ est quasi-DM, il suffit de montrer que
$\mathcal{X} \times_\mathcal{Y} V \to V$ est quasi-DM, voir le
Lemme \ref{lemma-check-separated-covering}.
Le Lemme \ref{lemma-properties-covering-imply-diagonal} montre que
$\mathcal{X} \times_\mathcal{Y} V$ est quasi-DM.
Par conséquent, $\mathcal{X} \times_\mathcal{Y} V \to V$ est quasi-DM par
le Lemme \ref{lemma-separated-implies-morphism-separated}, et (1) est
vérifiée. Ceci achève la démonstration du lemme.
\end{proof}

\begin{lemma}
\label{lemma-characterize-locally-quasi-finite}
Soit $f : \mathcal{X} \to \mathcal{Y}$ un morphisme de champs algébriques.
Les conditions suivantes sont équivalentes :
\begin{enumerate}
\item $f$ est localement quasi-fini,
\item $f$ est quasi-DM et, pour tout morphisme $V \to \mathcal{Y}$, où
$V$ est un espace algébrique, et tout morphisme localement quasi-fini
$U \to \mathcal{X} \times_\mathcal{Y} V$, où $U$ est un espace algébrique,
le morphisme $U \to V$ est localement quasi-fini,
\item pour tout morphisme $V \to \mathcal{Y}$ issu d'un espace algébrique
$V$, il existe un morphisme surjectif, plat, localement de présentation finie
et localement quasi-fini $U \to \mathcal{X} \times_\mathcal{Y} V$, où
$U$ est un espace algébrique, tel que $U \to V$ soit localement quasi-fini,
\item il existe des espaces algébriques $U$, $V$, un morphisme surjectif,
plat et localement de présentation finie $V \to \mathcal{Y}$, et un morphisme
$U \to \mathcal{X} \times_\mathcal{Y} V$ surjectif, plat, localement de
présentation finie et localement quasi-fini, tels que $U \to V$ soit
localement quasi-fini.
\end{enumerate}
\end{lemma}

\begin{proof}
Supposons (1). Alors $f$ est quasi-DM par hypothèse. Soient
$V \to \mathcal{Y}$ et $U \to \mathcal{X} \times_\mathcal{Y} V$
comme en (2). Par le Lemme \ref{lemma-composition-locally-quasi-finite}, le
composé $U \to \mathcal{X} \times_\mathcal{Y} V \to V$ est localement
quasi-fini. Ainsi, (1) implique (2).

\medskip\noindent
Supposons (2). Soit $V \to \mathcal{Y}$ comme en (3). Par le
Lemme \ref{lemma-characterize-quasi-DM}, on peut trouver un espace algébrique
$U$ et un morphisme surjectif, plat, localement de présentation finie et
localement quasi-fini $U \to \mathcal{X} \times_\mathcal{Y} V$.
Par (2), le composé $U \to V$ est localement quasi-fini. Ainsi, (2)
implique (3).

\medskip\noindent
Il est immédiat que (3) implique (4).

\medskip\noindent
Supposons (4). Nous allons établir (1), ce qui achèvera la démonstration.
Le Lemme \ref{lemma-characterize-quasi-DM} montre que $f$ est quasi-DM.
Pour montrer que $f$ est localement de type fini, il suffit de montrer que
$g : \mathcal{X} \times_\mathcal{Y} V \to V$ est localement de type fini,
voir le Lemme \ref{lemma-check-finite-type-covering}.
Il suffit ensuite de vérifier que le précomposé de $g$ avec
$h : U \to \mathcal{X} \times_\mathcal{Y} V$ est localement de type fini,
voir le Lemme \ref{lemma-check-finite-type-precompose}.
Cela résulte de l'hypothèse que $g \circ h : U \to V$ est localement
quasi-fini ; ainsi, $f$ est localement de type fini.
Enfin, soit $k$ un corps et soit $\Spec(k) \to \mathcal{Y}$ un morphisme.
Alors $V \times_\mathcal{Y} \Spec(k)$ est un espace algébrique non vide,
localement de présentation finie sur $k$. Il existe donc une extension finie
$k'/k$ et un morphisme $\Spec(k') \to V$ tels que le diagramme
$$
\xymatrix{
\Spec(k') \ar[r] \ar[d] & V \ar[d] \\
\Spec(k) \ar[r] & \mathcal{Y}
}
$$
soit commutatif (détails omis). Alors
$\mathcal{X}_{k'} \to \mathcal{X}_k$ est représentable (par des schémas),
surjectif et fini localement libre. En particulier,
$|\mathcal{X}_{k'}| \to |\mathcal{X}_k|$ est surjectif et ouvert.
Il suffit donc de montrer que $|\mathcal{X}_{k'}|$ est discret. Puisque
$$
U \times_V \Spec(k') =
U \times_{\mathcal{X} \times_\mathcal{Y} V} \mathcal{X}_{k'}
$$
le morphisme $U \times_V \Spec(k') \to \mathcal{X}_{k'}$ est surjectif,
plat et localement de présentation finie (comme changement de base de
$U \to \mathcal{X} \times_\mathcal{Y} V$). L'application
$|U \times_V \Spec(k')| \to |\mathcal{X}_{k'}|$ est donc surjective et
ouverte. Il suffit ainsi de montrer que $|U \times_V \Spec(k')|$ est discret.
Cela résulte du fait que $U \to V$ est localement quasi-fini (soit par notre
définition ci-dessus, soit par la définition originelle pour les morphismes
d'espaces algébriques, au moyen de Morphismes d'espaces,
Lemme \ref{spaces-morphisms-lemma-locally-quasi-finite}).
\end{proof}

\begin{lemma}
\label{lemma-locally-quasi-finite-permanence}
Soient $\mathcal{X} \to \mathcal{Y} \to \mathcal{Z}$ des morphismes de
champs algébriques. Supposons que $\mathcal{X} \to \mathcal{Z}$ soit
localement quasi-fini et que $\mathcal{Y} \to \mathcal{Z}$ soit quasi-DM.
Alors $\mathcal{X} \to \mathcal{Y}$ est localement quasi-fini.
\end{lemma}

\begin{proof}
Écrivons $\mathcal{X} \to \mathcal{Y}$ comme le composé
$$
\mathcal{X} \longrightarrow
\mathcal{X} \times_\mathcal{Z} \mathcal{Y} \longrightarrow
\mathcal{Y}
$$
La seconde flèche est localement quasi-finie comme changement de base de
$\mathcal{X} \to \mathcal{Z}$, voir le
Lemme \ref{lemma-base-change-locally-quasi-finite}.
La première est localement quasi-finie par le
Lemme \ref{lemma-semi-diagonal}, puisque
$\mathcal{Y} \to \mathcal{Z}$ est quasi-DM.
Ainsi, $\mathcal{X} \to \mathcal{Y}$ est localement quasi-fini par le
Lemme \ref{lemma-composition-locally-quasi-finite}.
\end{proof}

















\section{Morphismes quasi-finis}
\label{section-quasi-finite}

\noindent
Nous avons défini les morphismes ``localement quasi-finis'' de champs
algébriques à la section \ref{section-locally-quasi-finite}, et les morphismes
``quasi-compacts'' de champs algébriques à la
section \ref{section-quasi-compact}. Comme un morphisme d'espaces algébriques
est, par définition, quasi-fini si et seulement s'il est à la fois localement
quasi-fini et quasi-compact (Morphismes d'espaces, Définition
\ref{spaces-morphisms-definition-locally-quasi-finite}),
nous pouvons définir comme suit ce que signifie être quasi-fini pour un
morphisme de champs algébriques. Cette définition coïncide avec la notion déjà
définie dans Propriétés des champs,
section \ref{stacks-properties-section-properties-morphisms}, lorsque le
morphisme est représentable par des espaces algébriques.

\begin{definition}
\label{definition-quasi-finite}
\begin{reference}
\cite{rydh_approx}
\end{reference}
Soit $f : \mathcal{X} \to \mathcal{Y}$ un morphisme de champs algébriques.
On dit que $f$ est {\it quasi-fini} si $f$ est localement quasi-fini
(Définition \ref{definition-locally-quasi-finite})
et quasi-compact (Définition \ref{definition-quasi-compact}).
\end{definition}

\begin{lemma}
\label{lemma-composition-quasi-finite}
Le composé de morphismes quasi-finis est quasi-fini.
\end{lemma}

\begin{proof}
Il suffit de combiner les
Lemmes \ref{lemma-composition-locally-quasi-finite} et
\ref{lemma-composition-quasi-compact}.
\end{proof}

\begin{lemma}
\label{lemma-base-change-quasi-finite}
Tout changement de base d'un morphisme quasi-fini est quasi-fini.
\end{lemma}

\begin{proof}
Il suffit de combiner les
Lemmes \ref{lemma-base-change-locally-quasi-finite} et
\ref{lemma-base-change-quasi-compact}.
\end{proof}

\begin{lemma}
\label{lemma-quasi-finite-permanence}
Soient $f : \mathcal{X} \to \mathcal{Y}$ et
$g : \mathcal{Y} \to \mathcal{Z}$ des morphismes de champs algébriques.
Si $g \circ f$ est quasi-fini et si $g$ est quasi-séparé et quasi-DM,
alors $f$ est quasi-fini.
\end{lemma}

\begin{proof}
Il suffit de combiner les
Lemmes \ref{lemma-locally-quasi-finite-permanence} et
\ref{lemma-quasi-compact-permanence}.
\end{proof}












\section{Morphismes plats}
\label{section-flat}

\noindent
La propriété ``être plat'' des morphismes d'espaces algébriques est locale
pour la topologie lisse sur la source et la cible, voir Descente sur les
espaces, Remarque \ref{spaces-descent-remark-list-local-source-target}.
Elle est aussi stable par changement de base et locale pour la topologie fpqc
sur la cible, voir Morphismes d'espaces,
Lemme \ref{spaces-morphisms-lemma-base-change-flat}, et Descente sur les
espaces, Lemme
\ref{spaces-descent-lemma-descending-property-flat}.
Par le Lemme \ref{lemma-local-source-target} ci-dessus, nous pouvons donc
définir comme suit ce que signifie être plat pour un morphisme de champs
algébriques. Cette définition coïncide avec la notion déjà définie dans
Propriétés des champs,
section \ref{stacks-properties-section-properties-morphisms}, lorsque le
morphisme est représentable par des espaces algébriques.

\begin{definition}
\label{definition-flat}
Soit $f : \mathcal{X} \to \mathcal{Y}$ un morphisme de champs algébriques.
On dit que $f$ est {\it plat} si les conditions équivalentes du
Lemme \ref{lemma-local-source-target} sont satisfaites avec
$\mathcal{P} = \text{plat}$.
\end{definition}

\begin{lemma}
\label{lemma-composition-flat}
Le composé de morphismes plats est plat.
\end{lemma}

\begin{proof}
Il suffit de combiner la Remarque \ref{remark-composition} avec Morphismes
d'espaces, Lemme
\ref{spaces-morphisms-lemma-composition-flat}.
\end{proof}

\begin{lemma}
\label{lemma-base-change-flat}
Tout changement de base d'un morphisme plat est plat.
\end{lemma}

\begin{proof}
Il suffit de combiner la Remarque \ref{remark-base-change} avec Morphismes
d'espaces, Lemme
\ref{spaces-morphisms-lemma-base-change-flat}.
\end{proof}

\begin{lemma}
\label{lemma-descent-flat}
Soit $f : \mathcal{X} \to \mathcal{Y}$ un morphisme de champs algébriques.
Soit $\mathcal{Z} \to \mathcal{Y}$ un morphisme plat surjectif de champs
algébriques. Si le changement de base
$\mathcal{Z} \times_\mathcal{Y} \mathcal{X} \to \mathcal{Z}$
est plat, alors $f$ est plat.
\end{lemma}

\begin{proof}
Choisissons un espace algébrique $W$ et un morphisme lisse surjectif
$W \to \mathcal{Z}$. Alors $W \to \mathcal{Z}$ est surjectif et plat
(Morphismes d'espaces, Lemme \ref{spaces-morphisms-lemma-smooth-flat}) ;
par conséquent, $W \to \mathcal{Y}$ est surjectif et plat (par Propriétés des
champs, Lemme \ref{stacks-properties-lemma-composition-surjective}, et le
Lemme \ref{lemma-composition-flat}).
Puisque le changement de base de
$\mathcal{Z} \times_\mathcal{Y} \mathcal{X} \to \mathcal{Z}$
par $W \to \mathcal{Z}$ est un morphisme plat
(Lemme \ref{lemma-base-change-flat}), nous pouvons remplacer
$\mathcal{Z}$ par $W$.

\medskip\noindent
Choisissons un espace algébrique $V$ et un morphisme lisse surjectif
$V \to \mathcal{Y}$. Choisissons un espace algébrique $U$ et un morphisme
lisse surjectif $U \to V \times_\mathcal{Y} \mathcal{X}$.
Il faut montrer que $U \to V$ est plat. Effectuons maintenant partout le
changement de base par $W \to \mathcal{Y}$ : posons
$U' = W \times_\mathcal{Y} U$,
$V' = W \times_\mathcal{Y} V$,
$\mathcal{X}' = W \times_\mathcal{Y} \mathcal{X}$,
et $\mathcal{Y}' = W \times_\mathcal{Y} \mathcal{Y} = W$.
Par changement de base,
$U' \to V' \times_{\mathcal{Y}'} \mathcal{X}'$ est encore lisse.
Notre définition des morphismes plats de champs algébriques et l'hypothèse que
$\mathcal{X}' \to \mathcal{Y}'$ est plat montrent donc que
$U' \to V'$ est plat. Enfin, $V' \to V$ est surjectif comme changement de
base de $W \to \mathcal{Y}$ ; ainsi, $U \to V$ est plat par Morphismes
d'espaces, Lemme
\ref{spaces-morphisms-lemma-base-change-module-flat} (2)
et la démonstration est achevée.
\end{proof}

\begin{lemma}
\label{lemma-flat-permanence}
Soient $\mathcal{X} \to \mathcal{Y} \to \mathcal{Z}$ des morphismes de
champs algébriques. Si $\mathcal{X} \to \mathcal{Z}$ est plat et si
$\mathcal{X} \to \mathcal{Y}$ est surjectif et plat, alors
$\mathcal{Y} \to \mathcal{Z}$ est plat.
\end{lemma}

\begin{proof}
Choisissons un espace algébrique $W$ et un morphisme lisse surjectif
$W \to \mathcal{Z}$. Choisissons un espace algébrique $V$ et un morphisme
lisse surjectif $V \to W \times_\mathcal{Z} \mathcal{Y}$. Choisissons un
espace algébrique $U$ et un morphisme lisse surjectif
$U \to V \times_\mathcal{Y} \mathcal{X}$.
Nous savons que $U \to V$ est plat et que $U \to W$ est plat.
De plus, puisque $\mathcal{X} \to \mathcal{Y}$ est surjectif,
$U \to V$ est surjectif (comme composé de morphismes surjectifs).
Le lemme se ramène donc au cas des morphismes d'espaces algébriques, qui est
Morphismes d'espaces, Lemme
\ref{spaces-morphisms-lemma-flat-permanence}.
\end{proof}

\begin{lemma}
\label{lemma-lift-valuation-ring-through-flat-morphism}
Soit $f : \mathcal{X} \to \mathcal{Y}$ un morphisme plat de champs
algébriques. Soit $\Spec(A) \to \mathcal{Y}$ un morphisme, où $A$ est un
anneau de valuation. Si le point fermé de $\Spec(A)$ s'envoie sur un point de
$|\mathcal{Y}|$ appartenant à l'image de
$|\mathcal{X}| \to |\mathcal{Y}|$, alors il existe un diagramme commutatif
$$
\xymatrix{
\Spec(A') \ar[r] \ar[d] & \mathcal{X} \ar[d] \\
\Spec(A) \ar[r] & \mathcal{Y}
}
$$
où $A \to A'$ est une extension d'anneaux de valuation
(Compléments d'algèbre, Définition
\ref{more-algebra-definition-extension-valuation-rings}).
\end{lemma}

\begin{proof}
Le changement de base $\mathcal{X}_A \to \Spec(A)$ est plat
(Lemme \ref{lemma-base-change-flat}) et le point fermé de
$\Spec(A)$ appartient à l'image de $|\mathcal{X}_A| \to |\Spec(A)|$
(Propriétés des champs,
Lemme \ref{stacks-properties-lemma-points-cartesian}).
Nous pouvons donc supposer que $\mathcal{Y} = \Spec(A)$. Soit
$U \to \mathcal{X}$ un morphisme lisse surjectif, où $U$ est un schéma.
On peut alors appliquer Morphismes d'espaces, Lemme
\ref{spaces-morphisms-lemma-lift-valuation-ring-through-flat-morphism}
au morphisme $U \to \Spec(A)$ pour conclure.
\end{proof}







\section{Plat en un point}
\label{section-flat-at-point}

\noindent
Il nous reste à développer le formalisme général nécessaire pour exprimer
qu'un morphisme de champs algébriques possède une propriété donnée en un
point. Pour l'instant, le lemme suivant suffit.

\begin{lemma}
\label{lemma-flat-at-point}
Soit $f : \mathcal{X} \to \mathcal{Y}$ un morphisme de champs algébriques.
Soit $x \in |\mathcal{X}|$. Considérons des diagrammes commutatifs
$$
\vcenter{
\xymatrix{
U \ar[d]_a \ar[r]_h & V \ar[d]^b \\
\mathcal{X} \ar[r]^f & \mathcal{Y}
}
}
\quad\text{avec des points}
\vcenter{
\xymatrix{
u \in |U| \ar[d] \\
x \in |\mathcal{X}|
}
}
$$
où $U$ et $V$ sont des espaces algébriques, $b$ est plat, et
$(a, h) : U \to \mathcal{X} \times_\mathcal{Y} V$
est plat. Les conditions suivantes sont équivalentes :
\begin{enumerate}
\item $h$ est plat en $u$ pour l'un des diagrammes ci-dessus ;
\item $h$ est plat en $u$ pour tout diagramme comme ci-dessus.
\end{enumerate}
\end{lemma}

\begin{proof}
Supposons donné un second diagramme $U', V', u', a', b', h'$ satisfaisant
aux conditions du lemme. Nous pouvons alors considérer
$$
\xymatrix{
U \ar[d] & U \times_\mathcal{X} U' \ar[l] \ar[d] \ar[r] & U' \ar[d] \\
V & V \times_\mathcal{Y} V' \ar[l] \ar[r] & V'
}
$$
D'après Propriétés des champs, Lemme
\ref{stacks-properties-lemma-points-cartesian}, il existe un point
$u'' \in |U \times_\mathcal{X} U'|$ qui s'envoie sur $u$ et $u'$.
Si $h$ est plat en $u$, alors le changement de base
$U \times_V (V \times_\mathcal{Y} V') \to V \times_\mathcal{Y} V'$
est plat en tout point au-dessus de $u$ ; voir Morphismes d'espaces, Lemme
\ref{spaces-morphisms-lemma-base-change-module-flat}. D'autre part, le
morphisme
$$
U \times_\mathcal{X} U' \to
U \times_\mathcal{X} (\mathcal{X} \times_\mathcal{Y} V') =
U \times_\mathcal{Y} V' =
U \times_V (V \times_\mathcal{Y} V')
$$
est plat, car c'est un changement de base de $(a', h')$ ; voir le
Lemme \ref{lemma-base-change-flat}. En composant et en utilisant Morphismes
d'espaces, Lemme \ref{spaces-morphisms-lemma-composition-module-flat}, nous
concluons que $U \times_\mathcal{X} U' \to V \times_\mathcal{Y} V'$ est plat
en $u''$. En composant ensuite avec le morphisme plat
$V \times_\mathcal{Y} V' \to V'$, nous voyons que
$U \times_\mathcal{X} U' \to V'$ est plat en $u''$. Enfin, puisque
$U \times_\mathcal{X} U' \to U'$ est plat en $u''$ et que $u''$ s'envoie
sur $u'$, nous concluons que $U' \to V'$ est plat en $u'$ d'après Morphismes
d'espaces, Lemme \ref{spaces-morphisms-lemma-flat-permanence}.
\end{proof}

\begin{definition}
\label{definition-flat-at-point}
Soit $f : \mathcal{X} \to \mathcal{Y}$ un morphisme de champs algébriques.
Soit $x \in |\mathcal{X}|$. Nous disons que $f$ est {\it plat en $x$} si les
conditions équivalentes du Lemme \ref{lemma-flat-at-point} sont satisfaites.
\end{definition}





\section{Morphismes de présentation finie}
\label{section-finite-presentation}

\noindent
La propriété ``être localement de présentation finie'' des morphismes
d'espaces algébriques est locale pour la topologie lisse sur la source et la
cible, voir Descente sur les espaces, Remarque
\ref{spaces-descent-remark-list-local-source-target}. Elle est aussi stable
par changement de base et locale pour la topologie fpqc sur la cible, voir
Morphismes d'espaces, Lemme
\ref{spaces-morphisms-lemma-base-change-finite-presentation}, et Descente sur
les espaces, Lemme
\ref{spaces-descent-lemma-descending-property-locally-finite-presentation}.
Par le Lemme \ref{lemma-local-source-target} ci-dessus, nous pouvons donc
définir comme suit ce que signifie être localement de présentation finie pour
un morphisme de champs algébriques. Cette définition coïncide avec la notion
déjà définie dans Propriétés des champs,
section \ref{stacks-properties-section-properties-morphisms}, lorsque le
morphisme est représentable par des espaces algébriques.

\begin{definition}
\label{definition-locally-finite-presentation}
Soit $f : \mathcal{X} \to \mathcal{Y}$ un morphisme de champs algébriques.
\begin{enumerate}
\item On dit que $f$ est {\it localement de présentation finie} si les
conditions équivalentes du Lemme \ref{lemma-local-source-target} sont
satisfaites avec
$\mathcal{P} = \text{localement de présentation finie}$.
\item On dit que $f$ est {\it de présentation finie} s'il est localement de
présentation finie, quasi-compact et quasi-séparé.
\end{enumerate}
\end{definition}

\noindent
Notons qu'un morphisme de présentation finie n'est {\bf pas} simplement un
morphisme quasi-compact qui est localement de présentation finie.

\begin{lemma}
\label{lemma-composition-finite-presentation}
Le composé de morphismes de présentation finie est de présentation finie.
Il en va de même pour les morphismes localement de présentation finie.
\end{lemma}

\begin{proof}
Il suffit de combiner la Remarque \ref{remark-composition} avec Morphismes
d'espaces, Lemme
\ref{spaces-morphisms-lemma-composition-finite-presentation}.
\end{proof}

\begin{lemma}
\label{lemma-base-change-finite-presentation}
Un changement de base d'un morphisme de présentation finie est de présentation
finie. Il en va de même pour les morphismes localement de présentation finie.
\end{lemma}

\begin{proof}
Il suffit de combiner la Remarque \ref{remark-base-change} avec Morphismes
d'espaces, Lemme
\ref{spaces-morphisms-lemma-base-change-finite-presentation}.
\end{proof}

\begin{lemma}
\label{lemma-finite-presentation-finite-type}
Un morphisme localement de présentation finie est localement de type fini.
Un morphisme de présentation finie est de type fini.
\end{lemma}

\begin{proof}
Il suffit de combiner la Remarque \ref{remark-implication} avec Morphismes
d'espaces, Lemme
\ref{spaces-morphisms-lemma-finite-presentation-finite-type}.
\end{proof}

\begin{lemma}
\label{lemma-noetherian-finite-type-finite-presentation}
Soit $f : \mathcal{X} \to \mathcal{Y}$ un morphisme de champs algébriques.
\begin{enumerate}
\item Si $\mathcal{Y}$ est localement noethérien et $f$ localement de type
fini, alors $f$ est localement de présentation finie.
\item Si $\mathcal{Y}$ est localement noethérien et si $f$ est de type fini
et quasi-séparé, alors $f$ est de présentation finie.
\end{enumerate}
\end{lemma}

\begin{proof}
Supposons $f : \mathcal{X} \to \mathcal{Y}$ localement de type fini et
$\mathcal{Y}$ localement noethérien. Il existe donc un diagramme comme dans le
Lemme \ref{lemma-local-source-target}, avec $h$ localement de type fini et la
flèche verticale $a$ surjective. D'après Morphismes d'espaces, Lemme
\ref{spaces-morphisms-lemma-noetherian-finite-type-finite-presentation}
$h$ est localement de présentation finie. Ainsi,
$\mathcal{X} \to \mathcal{Y}$ est localement de présentation finie par
définition. Cela prouve (1). Si $f$ est de type fini et quasi-séparé, il est
aussi quasi-compact et quasi-séparé, et (2) en résulte immédiatement.
\end{proof}

\begin{lemma}
\label{lemma-finite-presentation-permanence}
Soient $f : \mathcal{X} \to \mathcal{Y}$ et
$g : \mathcal{Y} \to \mathcal{Z}$ des morphismes de champs algébriques.
Si $g \circ f$ est localement de présentation finie et si $g$ est localement
de type fini, alors $f$ est localement de présentation finie.
\end{lemma}

\begin{proof}
Choisissons un espace algébrique $W$ et un morphisme lisse surjectif
$W \to \mathcal{Z}$. Choisissons un espace algébrique $V$ et un morphisme
lisse surjectif $V \to \mathcal{Y} \times_\mathcal{Z} W$. Choisissons un
espace algébrique $U$ et un morphisme lisse surjectif
$U \to \mathcal{X} \times_\mathcal{Y} V$. Le lemme résulte de l'application
de Morphismes d'espaces, Lemme
\ref{spaces-morphisms-lemma-finite-presentation-permanence}
aux morphismes $U \to V \to W$.
\end{proof}

\begin{lemma}
\label{lemma-diagonal-morphism-finite-type}
Soit $f : \mathcal{X} \to \mathcal{Y}$ un morphisme de champs algébriques,
de diagonale
$\Delta : \mathcal{X} \to \mathcal{X} \times_\mathcal{Y} \mathcal{X}$.
Si $f$ est localement de type fini, alors $\Delta$ est localement de
présentation finie. Si $f$ est quasi-séparé et localement de type fini, alors
$\Delta$ est de présentation finie.
\end{lemma}

\begin{proof}
Notons que $\Delta$ est un morphisme au-dessus de $\mathcal{X}$ (par la
seconde projection). Si $f$ est localement de type fini, alors
$\mathcal{X}$ est de présentation finie sur $\mathcal{X}$ et
$\text{pr}_2 : \mathcal{X} \times_\mathcal{Y} \mathcal{X} \to \mathcal{X}$
est localement de type fini d'après le
Lemme \ref{lemma-base-change-finite-type}. La première assertion résulte donc
du Lemme \ref{lemma-finite-presentation-permanence}. La seconde résulte de la
première et des définitions (car, par définition, dire que $f$ est
quasi-séparé signifie que $\Delta_f$ est quasi-compact et quasi-séparé).
\end{proof}

\begin{lemma}
\label{lemma-open-immersion-locally-finite-presentation}
Une immersion ouverte est localement de présentation finie.
\end{lemma}

\begin{proof}
Compte tenu de Propriétés des champs, Définition
\ref{stacks-properties-definition-immersion}, cela résulte de Morphismes
d'espaces, Lemme
\ref{spaces-morphisms-lemma-open-immersion-locally-finite-presentation}.
\end{proof}

\begin{lemma}
\label{lemma-check-property-after-fppf-base-change}
Soit $P$ une propriété de morphismes d'espaces algébriques qui est locale pour
la topologie fppf sur la cible et stable par changement de base arbitraire.
Soit $f : \mathcal{X} \to \mathcal{Y}$ un morphisme de champs algébriques
représentable par des espaces algébriques. Soit
$\mathcal{Z} \to \mathcal{Y}$ un morphisme de champs algébriques surjectif,
plat et localement de présentation finie. Posons
$\mathcal{W} = \mathcal{Z} \times_\mathcal{Y} \mathcal{X}$. Alors
$$
(f\text{ possède }P) \Leftrightarrow
(\text{la projection }\mathcal{W} \to \mathcal{Z}\text{ possède }P).
$$
Pour la signification de cet énoncé, voir Propriétés des champs, section
\ref{stacks-properties-section-properties-morphisms}.
\end{lemma}

\begin{proof}
Choisissons un espace algébrique $W$ et un morphisme
$W \to \mathcal{Z}$ surjectif, plat et localement de présentation finie.
D'après Propriétés des champs, Lemme
\ref{stacks-properties-lemma-composition-surjective}
et les Lemmes \ref{lemma-composition-flat} et
\ref{lemma-composition-finite-presentation}, le composé
$W \to \mathcal{Y}$ est encore surjectif, plat et localement de présentation
finie. Posons
$V = W \times_\mathcal{Z} \mathcal{W} = W \times_\mathcal{Y} \mathcal{X}$.
D'après Propriétés des champs, Lemme
\ref{stacks-properties-lemma-check-property-covering}
$f$ possède $P$ si et seulement si $V \to W$ la possède, et
$\mathcal{W} \to \mathcal{Z}$ possède $P$ si et seulement si $V \to W$
la possède. Le lemme en résulte.
\end{proof}

\begin{lemma}
\label{lemma-descent-property}
Soit $\mathcal{P}$ une propriété de morphismes d'espaces algébriques qui est
locale pour la topologie lisse sur la source et la cible, et locale pour la
topologie fppf sur la cible. Soit
$f : \mathcal{X} \to \mathcal{Y}$ un morphisme de champs algébriques. Soit
$\mathcal{Z} \to \mathcal{Y}$ un morphisme de champs algébriques surjectif,
plat et localement de présentation finie. Si le changement de base
$\mathcal{Z} \times_\mathcal{Y} \mathcal{X} \to \mathcal{Z}$
possède $\mathcal{P}$, alors $f$ possède $\mathcal{P}$.
\end{lemma}

\begin{proof}
Supposons que
$\mathcal{Z} \times_\mathcal{Y} \mathcal{X} \to \mathcal{Z}$ possède
$\mathcal{P}$. Choisissons un espace algébrique $W$ et un morphisme lisse
surjectif $W \to \mathcal{Z}$. Observons que
$W \times_\mathcal{Z} \mathcal{Z} \times_\mathcal{Y} \mathcal{X} =
W \times_\mathcal{Y} \mathcal{X}$. Par la définition même de ce que signifie,
pour $\mathcal{Z} \times_\mathcal{Y} \mathcal{X} \to \mathcal{Z}$, posséder
$\mathcal{P}$ (voir la
Définition \ref{definition-P} et le Lemme \ref{lemma-local-source-target}),
le morphisme $W \times_\mathcal{Y} \mathcal{X} \to W$ possède donc
$\mathcal{P}$. D'autre part, $W \to \mathcal{Z}$ est surjectif, plat et
localement de présentation finie (Morphismes d'espaces, Lemmes
\ref{spaces-morphisms-lemma-smooth-flat} et
\ref{spaces-morphisms-lemma-smooth-locally-finite-presentation}) ; par
conséquent, $W \to \mathcal{Y}$ est surjectif, plat et localement de
présentation finie (d'après Propriétés des champs, Lemme
\ref{stacks-properties-lemma-composition-surjective}
et les Lemmes \ref{lemma-composition-flat} et
\ref{lemma-composition-finite-presentation}). Nous pouvons donc remplacer
$\mathcal{Z}$ par $W$.

\medskip\noindent
Choisissons un espace algébrique $V$ et un morphisme lisse surjectif
$V \to \mathcal{Y}$. Choisissons un espace algébrique $U$ et un morphisme
lisse surjectif $U \to V \times_\mathcal{Y} \mathcal{X}$. Il s'agit de
montrer que $U \to V$ possède $\mathcal{P}$. Effectuons partout le changement
de base par $W \to \mathcal{Y}$ et posons
$U' = W \times_\mathcal{Y} U$,
$V' = W \times_\mathcal{Y} V$,
$\mathcal{X}' = W \times_\mathcal{Y} \mathcal{X}$,
et $\mathcal{Y}' = W \times_\mathcal{Y} \mathcal{Y} = W$. Le morphisme
$U' \to V' \times_{\mathcal{Y}'} \mathcal{X}'$ demeure lisse par changement
de base. Le Lemme \ref{lemma-local-source-target}, utilisé dans la définition
de ce que signifie pour $\mathcal{X}' \to \mathcal{Y}' = W$ posséder
$\mathcal{P}$, montre donc que $U' \to V'$ possède
$\mathcal{P}$. Enfin, $V' \to V$ est surjectif, plat et localement de
présentation finie, car c'est un changement de base de
$W \to \mathcal{Y}$ ; il s'ensuit que $U \to V$ possède $\mathcal{P}$,
puisque $\mathcal{P}$ est locale pour la topologie fppf sur la cible.
\end{proof}

\begin{lemma}
\label{lemma-descent-finite-presentation}
Soit $f : \mathcal{X} \to \mathcal{Y}$ un morphisme de champs algébriques.
Soit $\mathcal{Z} \to \mathcal{Y}$ un morphisme de champs algébriques
surjectif, plat et localement de présentation finie. Si le changement de base
$\mathcal{Z} \times_\mathcal{Y} \mathcal{X} \to \mathcal{Z}$
est localement de présentation finie, alors $f$ est localement de présentation
finie.
\end{lemma}

\begin{proof}
La propriété ``être localement de présentation finie'' satisfait aux
conditions du Lemme \ref{lemma-descent-property}. Nous avons vu dans
l'introduction de cette section qu'elle est locale pour la topologie lisse sur
la source et la cible ; sa localité pour la topologie fppf sur la cible est
Descente sur les espaces, Lemme
\ref{spaces-descent-lemma-descending-property-locally-finite-presentation}.
\end{proof}

\begin{lemma}
\label{lemma-flat-finite-presentation-permanence}
Soient $\mathcal{X} \to \mathcal{Y} \to \mathcal{Z}$ des morphismes de
champs algébriques. Si $\mathcal{X} \to \mathcal{Z}$ est localement de
présentation finie et si $\mathcal{X} \to \mathcal{Y}$ est surjectif, plat et
localement de présentation finie, alors $\mathcal{Y} \to \mathcal{Z}$ est
localement de présentation finie.
\end{lemma}

\begin{proof}
Choisissons un espace algébrique $W$ et un morphisme lisse surjectif
$W \to \mathcal{Z}$. Choisissons un espace algébrique $V$ et un morphisme
lisse surjectif $V \to W \times_\mathcal{Z} \mathcal{Y}$. Choisissons un
espace algébrique $U$ et un morphisme lisse surjectif
$U \to V \times_\mathcal{Y} \mathcal{X}$. Nous savons que $U \to V$ est plat
et localement de présentation finie, et que $U \to W$ est localement de
présentation finie. De plus, puisque $\mathcal{X} \to \mathcal{Y}$ est
surjectif, $U \to V$ est surjectif (comme composé de morphismes surjectifs).
Le lemme se ramène donc au cas des morphismes d'espaces algébriques, lequel est
Descente sur les espaces, Lemme
\ref{spaces-descent-lemma-locally-finite-presentation-fppf-local-source}.
\end{proof}

\begin{lemma}
\label{lemma-surjective-flat-locally-finite-presentation}
Soit $f : \mathcal{X} \to \mathcal{Y}$ un morphisme de champs algébriques
surjectif, plat et localement de présentation finie. Pour tout schéma $U$ et
tout objet $y$ de $\mathcal{Y}$ au-dessus de $U$, il existe alors un
recouvrement fppf $\{U_i \to U\}$ et des objets $x_i$ de $\mathcal{X}$
au-dessus de $U_i$ tels que $f(x_i) \cong y|_{U_i}$ dans
$\mathcal{Y}_{U_i}$.
\end{lemma}

\begin{proof}
Nous pouvons considérer $y$ comme un morphisme $U \to \mathcal{Y}$. D'après
Propriétés des champs, Lemme
\ref{stacks-properties-lemma-base-change-surjective}
et les Lemmes \ref{lemma-base-change-finite-presentation} et
\ref{lemma-base-change-flat}, le morphisme
$\mathcal{X} \times_\mathcal{Y} U \to U$ est surjectif, plat et localement de
présentation finie. Soient $V$ un schéma et
$V \to \mathcal{X} \times_\mathcal{Y} U$ un morphisme lisse surjectif.
Le morphisme $V \to \mathcal{X} \times_\mathcal{Y} U$ est aussi surjectif,
plat et localement de présentation finie (voir
Morphismes d'espaces, Lemmes
\ref{spaces-morphisms-lemma-smooth-flat} et
\ref{spaces-morphisms-lemma-smooth-locally-finite-presentation}). Ainsi,
$V \to U$ est lui aussi surjectif, plat et localement de présentation finie ;
voir Propriétés des champs, Lemme
\ref{stacks-properties-lemma-composition-surjective}
et les Lemmes \ref{lemma-composition-finite-presentation} et
\ref{lemma-composition-flat}. Par conséquent, $\{V \to U\}$ est le
recouvrement fppf recherché et $x : V \to \mathcal{X}$ est l'objet recherché.
\end{proof}

\begin{lemma}
\label{lemma-surjective-family-flat-locally-finite-presentation}
Soit $f_j : \mathcal{X}_j \to \mathcal{X}$, $j \in J$, une famille de
morphismes de champs algébriques, plats et localement de présentation finie,
qui sont conjointement surjectifs, c'est-à-dire
$|\mathcal{X}| = \bigcup |f_j|(|\mathcal{X}_j|)$.
Alors, pour tout schéma $U$ et tout objet $x$ de $\mathcal{X}$ au-dessus de
$U$, il existe un recouvrement fppf $\{U_i \to U\}_{i \in I}$, une application
$a : I \to J$ et des objets $x_i$ de $\mathcal{X}_{a(i)}$ au-dessus de $U_i$
tels que $f_{a(i)}(x_i) \cong x|_{U_i}$ dans $\mathcal{X}_{U_i}$.
\end{lemma}

\begin{proof}
Appliquons le Lemme
\ref{lemma-surjective-flat-locally-finite-presentation} au morphisme
$\coprod_{j \in J} \mathcal{X}_j \to \mathcal{X}$. (Il subsiste ici un léger
problème ensembliste, dû à notre formalisme, que nous ignorons.) Enfin, un
morphisme $x_i : U_i \to \coprod_{j \in J} \mathcal{X}_j$ équivaut à la
donnée d'une décomposition en union disjointe
$U_i = \coprod U_{i, j}$ et de morphismes
$U_{i, j} \to \mathcal{X}_j$. Le recouvrement fppf
$\{U_{i, j} \to U\}$ et les morphismes
$U_{i, j} \to \mathcal{X}_j$ conviennent donc.
\end{proof}

\begin{lemma}
\label{lemma-fppf-open}
Soit $f : \mathcal{X} \to \mathcal{Y}$ plat et localement de présentation
finie. Alors $|f| : |\mathcal{X}| \to |\mathcal{Y}|$ est ouverte.
\end{lemma}

\begin{proof}
Choisissons un schéma $V$ et un morphisme lisse surjectif
$V \to \mathcal{Y}$. Choisissons un schéma $U$ et un morphisme lisse
surjectif $U \to V \times_\mathcal{Y} \mathcal{X}$. Par hypothèse, le
morphisme de schémas $U \to V$ est plat et localement de présentation finie.
Ainsi, $U \to V$ est ouvert d'après Morphismes, Lemme
\ref{morphisms-lemma-fppf-open}. Par construction de la topologie sur
$|\mathcal{Y}|$, l'application $|V| \to |\mathcal{Y}|$ est ouverte.
L'application $|U| \to |\mathcal{X}|$ est surjective. Le résultat découle
alors de ces faits par topologie élémentaire.
\end{proof}

\begin{lemma}
\label{lemma-descent-quasi-compact}
Soit $f : \mathcal{X} \to \mathcal{Y}$ un morphisme de champs algébriques.
Soit $\mathcal{Z} \to \mathcal{Y}$ un morphisme de champs algébriques
surjectif, plat et localement de présentation finie. Si le changement de base
$\mathcal{Z} \times_\mathcal{Y} \mathcal{X} \to \mathcal{Z}$
est quasi-compact, alors $f$ est quasi-compact.
\end{lemma}

\begin{proof}
Il faut montrer que, pour tout $\mathcal{Y}' \to \mathcal{Y}$ avec
$\mathcal{Y}'$ quasi-compact, le champ
$\mathcal{Y}' \times_\mathcal{Y} \mathcal{X}$ est quasi-compact. Posons
$\mathcal{Z}' = \mathcal{Z} \times_\mathcal{Y} \mathcal{Y}'$. Alors
$|\mathcal{Z}'| \to |\mathcal{Y}'|$ est ouverte ; voir le
Lemme \ref{lemma-fppf-open}. Il existe donc un sous-champ ouvert quasi-compact
$\mathcal{W} \subset \mathcal{Z}'$ dont l'image contient
$\mathcal{Y}'$. Puisque
$\mathcal{Z} \times_\mathcal{Y} \mathcal{X} \to \mathcal{Z}$
est quasi-compact, nous savons que
$$
\mathcal{W} \times_\mathcal{Z} \mathcal{Z} \times_\mathcal{Y} \mathcal{X} =
\mathcal{W} \times_\mathcal{Y} \mathcal{X}
$$
est quasi-compact. Or l'application
$\mathcal{W} \times_\mathcal{Y} \mathcal{X} \to
\mathcal{Y}' \times_\mathcal{Y} \mathcal{X}$
est surjective, ce qui conclut. Quelques détails sont omis.
\end{proof}

\begin{lemma}
\label{lemma-check-separated-on-ui-cover}
Soient $f : \mathcal{X} \to \mathcal{Y}$ et
$g : \mathcal{Y} \to \mathcal{Z}$ des morphismes composables de champs
algébriques, de composé
$h = g \circ f : \mathcal{X} \to \mathcal{Z}$.
Si $f$ est surjectif, plat, localement de présentation finie et
universellement injectif, et si $h$ est séparé, alors $g$ est séparé.
\end{lemma}

\begin{proof}
Considérons le diagramme
$$
\xymatrix{
\mathcal{X} \ar[r]_\Delta \ar[rd] &
\mathcal{X} \times_\mathcal{Y} \mathcal{X} \ar[r] \ar[d] &
\mathcal{X} \times_\mathcal{Z} \mathcal{X} \ar[d] \\
& \mathcal{Y} \ar[r] & \mathcal{Y} \times_\mathcal{Z} \mathcal{Y}
}
$$
Le carré est cartésien. Il faut montrer que la flèche horizontale inférieure
est propre. Nous savons déjà qu'elle est représentable par des espaces
algébriques et localement de type fini
(Lemme \ref{lemma-properties-diagonal}). Puisque la flèche verticale de droite
est surjective, plate et localement de présentation finie, il suffit de
montrer que la flèche horizontale supérieure de droite est propre
(Lemme \ref{lemma-check-property-after-fppf-base-change}). Comme $h$ est
séparé, le composé des flèches horizontales supérieures est propre.

\medskip\noindent
Puisque $f$ est universellement injectif, $\Delta$ est surjective
(Lemme \ref{lemma-universally-injective}). Comme le composé de $\Delta$ avec
la projection
$\mathcal{X} \times_\mathcal{Y} \mathcal{X} \to \mathcal{X}$
est l'identité, $\Delta$ est universellement fermé. D'après Morphismes
d'espaces, Lemme
\ref{spaces-morphisms-lemma-image-universally-closed-separated}
nous concluons que
$\mathcal{X} \times_\mathcal{Y} \mathcal{X} \to
\mathcal{X} \times_\mathcal{Z} \mathcal{X}$ est séparé, puisque
$\mathcal{X} \to \mathcal{X} \times_\mathcal{Z} \mathcal{X}$ est séparé.
Nous utilisons ici le fait que les implications entre propriétés de
morphismes d'espaces algébriques se transportent aux mêmes implications entre
propriétés de morphismes de champs algébriques représentables par des espaces
algébriques ; ceci est expliqué dans Propriétés des champs, section
\ref{stacks-properties-section-properties-morphisms}.
Enfin, le même principe et Morphismes d'espaces, Lemme
\ref{spaces-morphisms-lemma-image-proper-is-proper}, permettent de conclure que
$\mathcal{X} \times_\mathcal{Y} \mathcal{X} \to
\mathcal{X} \times_\mathcal{Z} \mathcal{X}$ est propre.
\end{proof}














\section{Gerbes}
\label{section-gerbes}

\noindent
Une classe importante de champs algébriques est formée des champs de la forme
$[B/G]$, où $B$ est un espace algébrique et $G$ un espace algébrique en groupes
sur $B$, plat et localement de présentation finie (agissant trivialement sur
$B$) ; voir Critères de représentabilité, Lemme
\ref{criteria-lemma-BG-algebraic}. Un champ algébrique est une gerbe lorsqu'il
est localement de cette forme pour la topologie fppf ; voir le
Lemme \ref{lemma-local-structure-gerbe}. Dans cette section, nous examinons
brièvement cette notion et la notion relative correspondante.

\begin{definition}
\label{definition-gerbe}
Soit $f : \mathcal{X} \to \mathcal{Y}$ un morphisme de champs algébriques.
On dit que $\mathcal{X}$ est une {\it gerbe sur} $\mathcal{Y}$ si
$\mathcal{X}$ est une gerbe sur $\mathcal{Y}$ comme champs en
groupoïdes sur $(\Sch/S)_{fppf}$ ; voir Champs, Définition
\ref{stacks-definition-gerbe-over-stack-in-groupoids}. On dit qu'un champ
algébrique $\mathcal{X}$ est une {\it gerbe} s'il existe un morphisme
$\mathcal{X} \to X$, où $X$ est un espace algébrique, qui fasse de
$\mathcal{X}$ une gerbe sur $X$.
\end{definition}

\noindent
La condition que $\mathcal{X}$ soit une gerbe sur $\mathcal{Y}$ se
définit uniquement à partir de la topologie et de la théorie des catégories
sous-jacentes aux champs algébriques considérés ; mais, comme nous le verrons,
elle a des conséquences géométriques. Elle implique par exemple que
$\mathcal{X} \to \mathcal{Y}$ est surjectif, plat et localement de présentation
finie ; voir le Lemme \ref{lemma-gerbe-fppf}. La notion absolue est plus
délicate à interpréter, car il n'est pas évident à première vue que $X$ soit
bien déterminé. Il l'est néanmoins.

\begin{lemma}
\label{lemma-gerbe-over-iso-classes}
Soit $\mathcal{X}$ un champ algébrique. Si $\mathcal{X}$ est une gerbe, alors
le faisceau associé au préfaisceau
$$
(\Sch/S)_{fppf}^{opp} \to \textit{Ens}, \quad
U \mapsto \Ob(\mathcal{X}_U)/\!\!\cong
$$
est un espace algébrique, et $\mathcal{X}$ est une gerbe sur celui-ci.
\end{lemma}

\begin{proof}
(Dans cette démonstration, l'abus de langage introduit à la
section \ref{section-conventions} se révèle particulièrement utile.)
Choisissons un morphisme $\pi : \mathcal{X} \to X$, où $X$ est un espace
algébrique, qui fasse de $\mathcal{X}$ une gerbe sur $X$. Il suffit de
montrer que $X$ est le faisceau associé au préfaisceau $\mathcal{F}$ affiché
dans le lemme. Il existe manifestement une application
$c : \mathcal{F} \to X$. Nous utiliserons les propriétés (2)(a) et (2)(b) de
Champs, Lemme \ref{stacks-lemma-when-gerbe}, pour montrer que l'application
$c^\# : \mathcal{F}^\# \to X$ est surjective et injective, donc un
isomorphisme ; voir Sites, Lemme \ref{sites-lemma-mono-epi-sheaves}.

Surjectivité : soit $T$ un schéma et $f : T \to X$. D'après la propriété
(2)(a), il existe un recouvrement fppf $\{h_i : T_i \to T\}$ et des morphismes
$x_i : T_i \to \mathcal{X}$ tels que $f \circ h_i$ corresponde à
$\pi \circ x_i$. Ainsi, $f|_{T_i}$ appartient à l'image de $c$.

Injectivité : soient $T$ un schéma et $x, y : T \to \mathcal{X}$ des
morphismes tels que $c \circ x = c \circ y$. D'après (2)(b), il existe un
recouvrement $\{T_i \to T\}$ et des morphismes
$x|_{T_i} \to y|_{T_i}$ dans la catégorie fibre $\mathcal{X}_{T_i}$.
Les restrictions $x|_{T_i}, y|_{T_i}$ sont donc égales dans
$\mathcal{F}(T_i)$. Cela prouve, comme voulu, que $x, y$ définissent la même
section de $\mathcal{F}^\#$ au-dessus de $T$.
\end{proof}

\begin{lemma}
\label{lemma-base-change-gerbe}
Soit
$$
\xymatrix{
\mathcal{X}' \ar[r] \ar[d] & \mathcal{X} \ar[d] \\
\mathcal{Y}' \ar[r] & \mathcal{Y}
}
$$
un produit fibré de champs algébriques. Si $\mathcal{X}$ est une gerbe
sur $\mathcal{Y}$, alors $\mathcal{X}'$ est une gerbe sur
$\mathcal{Y}'$.
\end{lemma}

\begin{proof}
Cela résulte immédiatement des définitions et de Champs, Lemme
\ref{stacks-lemma-base-change-gerbe}.
\end{proof}

\begin{lemma}
\label{lemma-composition-gerbe}
Soient $\mathcal{X} \to \mathcal{Y}$ et
$\mathcal{Y} \to \mathcal{Z}$ des morphismes de champs algébriques. Si
$\mathcal{X}$ est une gerbe sur $\mathcal{Y}$ et $\mathcal{Y}$ une
gerbe sur $\mathcal{Z}$, alors $\mathcal{X}$ est une gerbe sur
de $\mathcal{Z}$.
\end{lemma}

\begin{proof}
Cela résulte immédiatement de Champs, Lemme
\ref{stacks-lemma-composition-gerbe}.
\end{proof}

\begin{lemma}
\label{lemma-gerbe-descent}
Soit
$$
\xymatrix{
\mathcal{X}' \ar[r] \ar[d] & \mathcal{X} \ar[d] \\
\mathcal{Y}' \ar[r] & \mathcal{Y}
}
$$
un produit fibré de champs algébriques. Si
$\mathcal{Y}' \to \mathcal{Y}$ est surjectif, plat et localement de
présentation finie, et si $\mathcal{X}'$ est une gerbe sur
$\mathcal{Y}'$, alors $\mathcal{X}$ est une gerbe sur $\mathcal{Y}$.
\end{lemma}

\begin{proof}
Cela résulte immédiatement du
Lemme \ref{lemma-surjective-flat-locally-finite-presentation} et de Champs,
Lemme \ref{stacks-lemma-gerbe-descent}.
\end{proof}

\begin{lemma}
\label{lemma-gerbe-with-section}
Soit $\pi : \mathcal{X} \to U$ un morphisme d'un champ algébrique vers un
espace algébrique, et soit $x : U \to \mathcal{X}$ une section de $\pi$.
Posons $G = \mathit{Isom}_\mathcal{X}(x, x)$ ; voir la
Définition \ref{definition-isom}. Si $\mathcal{X}$ est une gerbe sur
$U$, alors :
\begin{enumerate}
\item il existe une équivalence canonique de champs en groupoïdes
$$
x_{can} : [U/G] \longrightarrow \mathcal{X}.
$$
où $[U/G]$ est le champ quotient pour l'action triviale de $G$ sur $U$ ;
\item $G \to U$ est plat et localement de présentation finie ;
\item $U \to \mathcal{X}$ est surjectif, plat et localement de présentation
finie.
\end{enumerate}
\end{lemma}

\begin{proof}
Posons $R = U \times_{x, \mathcal{X}, x} U$. Le morphisme
$R \to U \times U$ se factorise par la diagonale
$\Delta_U : U \to U \times U$, puisqu'il se factorise par
$U \times_U U = U$. Ainsi $R = G$, car
\begin{align*}
G & = \mathit{Isom}_\mathcal{X}(x, x) \\
& = U \times_{x, \mathcal{X}} \mathcal{I}_\mathcal{X} \\
& = U \times_{x, \mathcal{X}}
(\mathcal{X}
\times_{\Delta, \mathcal{X} \times_S \mathcal{X}, \Delta}
\mathcal{X}) \\
& = (U \times_{x, \mathcal{X}, x} U) \times_{U \times U, \Delta_U} U \\
& = R \times_{U \times U, \Delta_U} U \\
& = R
\end{align*}
où, pour la quatrième égalité, on utilise Catégories, Lemme
\ref{categories-lemma-diagonal-2}.
Soient $t, s : R \to U$ les projections. La loi de composition
$c : R \times_{s, U, t} R \to R$ construite sur $R$ dans Champs algébriques,
Lemme \ref{algebraic-lemma-map-space-into-stack}, coïncide avec la loi de
groupe sur $G$ (démonstration omise). Ainsi, Champs algébriques, Lemme
\ref{algebraic-lemma-map-space-into-stack}, fournit un $1$-morphisme canonique
pleinement fidèle
$$
x_{can} : [U/G] \longrightarrow \mathcal{X}
$$
de champs en groupoïdes sur $(\Sch/S)_{fppf}$. Pour voir que c'est une
équivalence, il suffit de montrer qu'il est essentiellement surjectif. Il
suffit pour cela de montrer que tout objet de $\mathcal{X}$ au-dessus d'un
schéma $T$ provient, localement pour la topologie fppf, de $x$ par un morphisme
$T \to U$ ; voir Champs, Lemme
\ref{stacks-lemma-characterize-essentially-surjective-when-ff}. Or cela résulte
de la condition selon laquelle $\pi$ fait de $\mathcal{X}$ une gerbe sur
$U$ ; voir la propriété (2)(a) de Champs, Lemme
\ref{stacks-lemma-when-gerbe}.

\medskip\noindent
D'après Critères de représentabilité, Lemme
\ref{criteria-lemma-BG-algebraic}, le morphisme $G \to U$ est plat et
localement de présentation finie. Enfin, $U \to \mathcal{X}$ est surjectif,
plat et localement de présentation finie d'après Critères de représentabilité,
Lemme
\ref{criteria-lemma-flat-quotient-flat-presentation}.
\end{proof}

\begin{lemma}
\label{lemma-local-structure-gerbe}
Soit $\pi : \mathcal{X} \to \mathcal{Y}$ un morphisme de champs algébriques.
Les conditions suivantes sont équivalentes :
\begin{enumerate}
\item $\mathcal{X}$ est une gerbe sur $\mathcal{Y}$ ;
\item il existe un espace algébrique $U$, un espace algébrique en groupes $G$
sur $U$, plat et localement de présentation finie, et un morphisme
$U \to \mathcal{Y}$ surjectif, plat et localement de présentation finie, tels
que $\mathcal{X} \times_\mathcal{Y} U \cong [U/G]$ au-dessus de $U$.
\end{enumerate}
\end{lemma}

\begin{proof}
Supposons (2). D'après le Lemme \ref{lemma-gerbe-descent}, il suffit, pour
prouver (1), de montrer que $[U/G]$ est une gerbe sur $U$. Cela
résulte immédiatement de Groupoïdes en espaces, Lemme
\ref{spaces-groupoids-lemma-group-quotient-gerbe}.

\medskip\noindent
Supposons (1). Tout changement de base de $\pi$ est une gerbe ; voir le
Lemme \ref{lemma-base-change-gerbe}. Choisissons d'abord un schéma $V$ et un
morphisme lisse surjectif $V \to \mathcal{Y}$. Nous pouvons ainsi supposer que
$\pi : \mathcal{X} \to V$ est une gerbe sur un schéma. Il existe alors
un recouvrement fppf $\{V_i \to V\}$ tel que la catégorie fibre
$\mathcal{X}_{V_i}$ soit non vide ; voir Champs, Lemme
\ref{stacks-lemma-when-gerbe} (2)(a). Notons que
$U = \coprod V_i \to V$ est surjectif, plat et localement de présentation
finie. Nous pouvons donc remplacer $V$ par $U$ et supposer que
$\pi : \mathcal{X} \to U$ est une gerbe sur le schéma $U$, et qu'il
existe un objet $x$ de $\mathcal{X}$ au-dessus de $U$. D'après le
Lemme \ref{lemma-gerbe-with-section}, on a alors
$\mathcal{X} = [U/G]$ au-dessus de $U$, pour un certain espace algébrique en
groupes $G$ sur $U$, plat et localement de présentation finie.
\end{proof}

\begin{lemma}
\label{lemma-gerbe-fppf}
Soit $\pi : \mathcal{X} \to \mathcal{Y}$ un morphisme de champs algébriques.
Si $\mathcal{X}$ est une gerbe sur $\mathcal{Y}$, alors $\pi$ est
surjectif, plat et localement de présentation finie.
\end{lemma}

\begin{proof}
D'après Propriétés des champs, Lemme
\ref{stacks-properties-lemma-descent-surjective}
et les Lemmes \ref{lemma-descent-flat} et
\ref{lemma-descent-finite-presentation}, il suffit de démontrer le lemme après
avoir remplacé $\pi$ par son changement de base par un morphisme
$\mathcal{Y}' \to \mathcal{Y}$ surjectif, plat et localement de présentation
finie. D'après le Lemme \ref{lemma-local-structure-gerbe}, nous pouvons
supposer que $\mathcal{Y} = U$ est un espace algébrique et que
$\mathcal{X} = [U/G]$ au-dessus de $U$. Le morphisme $U \to [U/G]$ est alors
surjectif, plat et localement de présentation finie ; voir le
Lemme \ref{lemma-gerbe-with-section}. Il en résulte que $\pi$ est surjectif,
plat et localement de présentation finie d'après Propriétés des champs, Lemme
\ref{stacks-properties-lemma-surjective-permanence}, et les
Lemmes \ref{lemma-flat-permanence} et
\ref{lemma-flat-finite-presentation-permanence}.
\end{proof}

\begin{proposition}
\label{proposition-when-gerbe}
Soit $\mathcal{X}$ un champ algébrique. Les conditions suivantes sont
équivalentes :
\begin{enumerate}
\item $\mathcal{X}$ est une gerbe ;
\item $\mathcal{I}_\mathcal{X} \to \mathcal{X}$ est plat et localement de
présentation finie.
\end{enumerate}
\end{proposition}

\begin{proof}
Supposons (1). Choisissons un morphisme $\mathcal{X} \to X$ vers un espace
algébrique $X$ qui fasse de $\mathcal{X}$ une gerbe sur $X$. Soit
$X' \to X$ un morphisme surjectif, plat et localement de présentation finie,
et posons $\mathcal{X}' = X' \times_X \mathcal{X}$. Le champ $\mathcal{X}'$
est une gerbe sur $X'$ d'après le
Lemme \ref{lemma-base-change-gerbe}. Les deux carrés du diagramme
$$
\xymatrix{
\mathcal{I}_{\mathcal{X}'} \ar[r] \ar[d] &
\mathcal{X}' \ar[r] \ar[d] & X' \ar[d] \\
\mathcal{I}_\mathcal{X} \ar[r] &
\mathcal{X} \ar[r] & X
}
$$
sont cartésiens ; voir le Lemme \ref{lemma-cartesian-square-inertia}. Pour
montrer que $\mathcal{I}_\mathcal{X} \to \mathcal{X}$ est plat et localement
de présentation finie, il suffit donc de le vérifier après un tel changement
de base, d'après les Lemmes \ref{lemma-descent-flat} et
\ref{lemma-descent-finite-presentation}. Le
Lemme \ref{lemma-local-structure-gerbe} permet alors de supposer que
$\mathcal{X} = [U/G]$. D'après le Lemme \ref{lemma-gerbe-with-section}, $G$
est plat et localement de présentation finie sur $U$, et
$x : U \to [U/G]$ est surjectif, plat et localement de présentation finie.
De plus, le changement de base de $\mathcal{I}_\mathcal{X}$ par $x$ est $G$ ;
une nouvelle descente, c'est-à-dire les Lemmes \ref{lemma-descent-flat} et
\ref{lemma-descent-finite-presentation}, montre donc (2).

\medskip\noindent
Réciproquement, supposons (2). Choisissons une présentation lisse
$\mathcal{X} = [U/R]$ ; voir Champs algébriques, section
\ref{algebraic-section-stack-to-presentation}. Notons $G \to U$ l'espace
algébrique en groupes stabilisateur du groupoïde
$(U, R, s, t, c, e, i)$ ; voir Groupoïdes en espaces, Définition
\ref{spaces-groupoids-definition-stabilizer-groupoid}.
D'après le Lemme \ref{lemma-presentation-inertia}, $G \to U$ est plat et
localement de présentation finie, comme changement de base de
$\mathcal{I}_\mathcal{X} \to \mathcal{X}$ ; voir les
Lemmes \ref{lemma-base-change-flat} et
\ref{lemma-base-change-finite-presentation}. Considérons l'action
$$
a : G \times_{U, t} R \to R, \quad (g, r) \mapsto c(g, r)
$$
de $G$ sur $R$. Pour tout schéma $T$, cette action est libre sur les points à
valeurs dans $T$, puisque $R$ est un groupoïde. Ainsi $R' = R/G$ est un espace
algébrique, et le morphisme quotient $\pi : R \to R'$ est surjectif, plat et
localement de présentation finie d'après Amorçage, Lemme
\ref{bootstrap-lemma-quotient-free-action}. Les projections
$s, t : R \to U$ sont $G$-invariantes ; il existe donc des morphismes
$s' , t' : R' \to U$ tels que $s = s' \circ \pi$ et
$t = t' \circ \pi$. Puisque $s, t : R \to U$ sont plats et localement de
présentation finie, $s', t'$ le sont aussi ; voir Morphismes d'espaces, Lemme
\ref{spaces-morphisms-lemma-flat-permanence}, et Descente sur les espaces,
Lemme
\ref{spaces-descent-lemma-locally-finite-presentation-fppf-local-source}.
Considérons le morphisme
$$
j' = (t', s') : R' \longrightarrow U \times U.
$$
Nous affirmons que c'est un monomorphisme. Soient en effet $T$ un schéma et
$a, b : T \to R'$ des morphismes ayant la même image dans $U \times U$. Par
définition du quotient $R' = R/G$, il existe un recouvrement fppf
$\{h_j : T_j \to T\}$ tel que
$a \circ h_j = \pi \circ a_j$ et $b \circ h_j = \pi \circ b_j$, pour des
morphismes $a_j, b_j : T_j \to R$. Puisque $a_j, b_j$ ont la même image dans
$U \times U$, l'élément $g_j = c(a_j, i(b_j))$ est un point à valeurs dans
$T_j$ de $G$ tel que $c(g_j, b_j) = a_j$. Autrement dit, $a_j$ et $b_j$ ont la
même image dans $R'$, ce qui prouve l'affirmation. Comme
$j : R \to U \times U$ est une prérelation d'équivalence (voir Groupoïdes en
espaces, Lemme
\ref{spaces-groupoids-lemma-groupoid-pre-equivalence})
et que $R \to R'$ est surjectif (comme morphisme de faisceaux), le morphisme
$j' : R' \to U \times U$ est une relation d'équivalence. Par conséquent,
Amorçage, Théorème \ref{bootstrap-theorem-final-bootstrap}, montre que
$X = U/R'$ est un espace algébrique. Enfin, nous affirmons que le morphisme
$$
\mathcal{X} = [U/R] \longrightarrow X = U/R'
$$
fait de $\mathcal{X}$ une gerbe sur $X$. Cela résulte de Groupoïdes
en espaces, Lemme \ref{spaces-groupoids-lemma-when-gerbe}, puisque
$R \to R'$ est surjectif, plat et localement de présentation finie (au besoin,
utiliser Amorçage, Lemme
\ref{bootstrap-lemma-surjective-flat-locally-finite-presentation}, pour voir
que cela entraîne l'hypothèse requise).
\end{proof}

\begin{lemma}
\label{lemma-gerbe-isom-fppf}
Soit $f : \mathcal{X} \to \mathcal{Y}$ un morphisme de champs algébriques
qui fait de $\mathcal{X}$ une gerbe sur $\mathcal{Y}$. Alors :
\begin{enumerate}
\item $\mathcal{I}_{\mathcal{X}/\mathcal{Y}} \to \mathcal{X}$
est plat et localement de présentation finie ;
\item $\mathcal{X} \to \mathcal{X} \times_\mathcal{Y} \mathcal{X}$
est surjectif, plat et localement de présentation finie ;
\item étant donnés des espaces algébriques $T_i$, $i = 1, 2$, et des
morphismes $x_i : T_i \to \mathcal{X}$, avec $y_i = f \circ x_i$, le
morphisme
$$
T_1 \times_{x_1, \mathcal{X}, x_2} T_2 \longrightarrow
T_1 \times_{y_1, \mathcal{Y}, y_2} T_2
$$
est surjectif, plat et localement de présentation finie ;
\item étant donnés un espace algébrique $T$ et des morphismes
$x_i : T \to \mathcal{X}$, $i = 1, 2$, avec $y_i = f \circ x_i$, le morphisme
$$
\mathit{Isom}_\mathcal{X}(x_1, x_2) \longrightarrow
\mathit{Isom}_\mathcal{Y}(y_1, y_2)
$$
est surjectif, plat et localement de présentation finie.
\end{enumerate}
\end{lemma}

\begin{proof}
Démonstration de (1). Choisissons un schéma $Y$ et un morphisme lisse
surjectif $Y \to \mathcal{Y}$. Posons
$\mathcal{X}' = \mathcal{X} \times_\mathcal{Y} Y$. Le
Lemme \ref{lemma-cartesian-square-inertia} fournit les carrés cartésiens
$$
\xymatrix{
\mathcal{I}_{\mathcal{X}'} \ar[r] \ar[d] &
\mathcal{X}' \ar[r] \ar[d] & Y \ar[d] \\
\mathcal{I}_{\mathcal{X}/\mathcal{Y}} \ar[r] &
\mathcal{X} \ar[r] & \mathcal{Y}
}
$$
D'après les Lemmes \ref{lemma-descent-flat} et
\ref{lemma-descent-finite-presentation}, il suffit de montrer que
$\mathcal{I}_{\mathcal{X}'} \to \mathcal{X}'$ est plat et localement de
présentation finie. Cela résulte de la
Proposition \ref{proposition-when-gerbe} (car $\mathcal{X}'$ est une gerbe
au-dessus de $Y$ d'après le Lemme \ref{lemma-base-change-gerbe}).

\medskip\noindent
Démonstration de (2). Avec les notations ci-dessus, nous pouvons supposer que
$\mathcal{X}' = [Y/G]$ pour un espace algébrique en groupes $G$ sur $Y$, plat
et localement de présentation finie ; voir le
Lemme \ref{lemma-local-structure-gerbe}. Le changement de base du morphisme
$\Delta : \mathcal{X} \to \mathcal{X} \times_\mathcal{Y} \mathcal{X}$
au-dessus de $\mathcal{Y}$ par $Y \to \mathcal{Y}$ est le morphisme
$\Delta' : \mathcal{X}' \to \mathcal{X}' \times_Y \mathcal{X}'$.
Il suffit donc de montrer que $\Delta'$ est surjectif, plat et localement de
présentation finie (voir les Lemmes \ref{lemma-descent-flat} et
\ref{lemma-descent-finite-presentation}). Autrement dit, il faut montrer que
$$
[Y/G] \longrightarrow [Y/G \times_Y G]
$$
est surjectif, plat et localement de présentation finie. En effet, son
changement de base par le morphisme surjectif, plat et localement de
présentation finie $Y \to [Y/G \times_Y G]$ est le morphisme $G \to Y$.

\medskip\noindent
Démonstration de (3). Observons que le diagramme
$$
\xymatrix{
T_1 \times_{x_1, \mathcal{X}, x_2} T_2 \ar[d] \ar[r] &
T_1 \times_{y_1, \mathcal{Y}, y_2} T_2 \ar[d] \\
\mathcal{X} \ar[r] & \mathcal{X} \times_\mathcal{Y} \mathcal{X}
}
$$
est cartésien. Ainsi (3) résulte de (2).

\medskip\noindent
Démonstration de (4). Cela résulte de l'égalité
$$
\mathit{Isom}_\mathcal{X}(x_1, x_2) =
(T \times_{x_1, \mathcal{X}, x_2} T) \times_{T \times T, \Delta_T} T
$$
car le morphisme de (4) est ainsi un changement de base du morphisme de (3).
\end{proof}

\begin{proposition}
\label{proposition-when-gerbe-over}
Soit $f : \mathcal{X} \to \mathcal{Y}$ un morphisme de champs algébriques.
Les conditions suivantes sont équivalentes :
\begin{enumerate}
\item $\mathcal{X}$ est une gerbe sur $\mathcal{Y}$ ;
\item $f : \mathcal{X} \to \mathcal{Y}$ et
$\Delta : \mathcal{X} \to \mathcal{X} \times_\mathcal{Y} \mathcal{X}$
sont surjectifs, plats et localement de présentation finie.
\end{enumerate}
\end{proposition}

\begin{proof}
L'implication (1) $\Rightarrow$ (2) résulte des
Lemmes \ref{lemma-gerbe-fppf} et \ref{lemma-gerbe-isom-fppf}.

\medskip\noindent
Supposons (2). Il suffit de prouver (1) après avoir changé $f$ de base par un
morphisme $\mathcal{Y}' \to \mathcal{Y}$ surjectif, plat et localement de
présentation finie ; voir le Lemme \ref{lemma-gerbe-descent} (le changement
de base de la diagonale de $f$ est la diagonale du changement de base). Nous
pouvons donc supposer que $\mathcal{Y}$ est un schéma $Y$. Dans ce cas,
$\mathcal{I}_\mathcal{X} \to \mathcal{X}$ est un changement de base de
$\Delta$, et la Proposition \ref{proposition-when-gerbe} montre que
$\mathcal{X}$ est une gerbe. Il reste à montrer que $\mathcal{X}$ est une
gerbe sur $Y$. Soit $\mathcal{X} \to X$ le morphisme du
Lemme \ref{lemma-gerbe-over-iso-classes} qui fait de $\mathcal{X}$ une gerbe
au-dessus de l'espace algébrique $X$ classifiant les classes d'isomorphisme
d'objets de $\mathcal{X}$. Le morphisme $f : \mathcal{X} \to Y$ se factorise
manifestement en $\mathcal{X} \to X \to Y$. Comme $f$ est surjectif, plat et
localement de présentation finie, $X \to Y$ est surjectif comme morphisme de
faisceaux fppf (utiliser par exemple le
Lemme \ref{lemma-surjective-flat-locally-finite-presentation}). D'autre part,
$X \to Y$ est aussi injectif : pour tout schéma $T$ et tous points à valeurs
dans $T$, $x_1, x_2$, de $X$ ayant même image dans $Y$, on peut
d'abord, localement pour la topologie fppf sur $T$, relever $x_1, x_2$ en des
objets $\xi_1, \xi_2$ de $\mathcal{X}$ au-dessus de $T$, puis déduire que
$\xi_1$ et $\xi_2$ sont localement isomorphes pour la topologie fppf, puisque
$\Delta : \mathcal{X} \to \mathcal{X} \times_Y \mathcal{X}$
est surjectif, plat et localement de présentation finie. On a donc
$x_1 = x_2$ par construction de $X$. Ainsi $X = Y$, ce qui achève la
démonstration.
\end{proof}

\noindent
Nous disposons maintenant d'un formalisme suffisant pour montrer que les
gerbes résiduelles, lorsqu'elles existent, sont des gerbes.

\begin{lemma}
\label{lemma-noetherian-singleton-stack-gerbe}
Soit $\mathcal{Z}$ un champ algébrique réduit et localement noethérien tel que
$|\mathcal{Z}|$ soit un singleton. Alors $\mathcal{Z}$ est une gerbe sur
un espace algébrique $Z$ réduit et localement noethérien, dont $|Z|$ est un
singleton.
\end{lemma}

\begin{proof}
D'après Propriétés des champs, Lemme
\ref{stacks-properties-lemma-unique-point-better}, il existe un morphisme
surjectif, plat et localement de présentation finie
$\Spec(k) \to \mathcal{Z}$, où $k$ est un corps. Le morphisme
$\mathcal{I}_\mathcal{Z} \times_\mathcal{Z} \Spec(k) \to \Spec(k)$ est alors
représentable par des espaces algébriques et localement de type fini (comme
changement de base de $\mathcal{I}_\mathcal{Z} \to \mathcal{Z}$ ; voir les
Lemmes \ref{lemma-inertia} et \ref{lemma-base-change-finite-type}). Il est donc
localement de présentation finie ; voir Morphismes d'espaces, Lemme
\ref{spaces-morphisms-lemma-noetherian-finite-type-finite-presentation}.
Il est aussi plat, puisque $k$ est un corps. Les
Lemmes \ref{lemma-descent-flat} et
\ref{lemma-descent-finite-presentation} montrent donc que
$\mathcal{I}_\mathcal{Z} \to \mathcal{Z}$ est plat et localement de
présentation finie. La Proposition \ref{proposition-when-gerbe} implique que
$\mathcal{Z}$ est une gerbe. Soit $\pi : \mathcal{Z} \to Z$ un morphisme vers
un espace algébrique tel que $\mathcal{Z}$ soit une gerbe sur $Z$.
Le morphisme $\pi$ est surjectif, plat et localement de présentation finie
d'après le Lemme \ref{lemma-gerbe-fppf}. Par composition,
$\Spec(k) \to Z$ est donc surjectif, plat et localement de présentation finie ;
voir Propriétés des champs, Lemme
\ref{stacks-properties-lemma-composition-surjective}, et les
Lemmes \ref{lemma-composition-flat} et
\ref{lemma-composition-finite-presentation}. Enfin, Propriétés des champs,
Lemme \ref{stacks-properties-lemma-unique-point-better}, montre que $|Z|$ est
un singleton et que $Z$ est localement noethérien et réduit.
\end{proof}

\begin{lemma}
\label{lemma-gerbe-bijection-points}
Soit $f : \mathcal{X} \to \mathcal{Y}$ un morphisme de champs algébriques.
Si $\mathcal{X}$ est une gerbe sur $\mathcal{Y}$, alors $f$ est un
homéomorphisme universel.
\end{lemma}

\begin{proof}
D'après le Lemme \ref{lemma-base-change-gerbe}, l'hypothèse sur $f$ est stable
par changement de base. Il suffit donc de montrer que l'application
$|\mathcal{X}| \to |\mathcal{Y}|$ est un homéomorphisme d'espaces
topologiques. Soient $k$ un corps et $y$ un objet de $\mathcal{Y}$ au-dessus de
$\Spec(k)$. D'après la propriété (2)(a) de Champs, Lemme
\ref{stacks-lemma-when-gerbe}, il existe un recouvrement fppf
$\{T_i \to \Spec(k)\}$ et des objets $x_i$ de $\mathcal{X}$ au-dessus de
$T_i$, avec $f(x_i) \cong y|_{T_i}$. Choisissons $i$ tel que
$T_i \not = \emptyset$, puis un morphisme $\Spec(K) \to T_i$ pour un certain
corps $K$. Alors $k \subset K$ et $x_i|_K$ est un objet de $\mathcal{X}$
au-dessus de $y|_K$. L'application
$|\mathcal{X}| \to |\mathcal{Y}|$ est donc surjective. L'application
$|\mathcal{X}| \to |\mathcal{Y}|$ est aussi injective. En effet, si
$x, x'$ sont des objets de $\mathcal{X}$ au-dessus de
$\Spec(k)$ dont les images $f(x), f(x')$ deviennent isomorphes (après une
extension) dans $\mathcal{Y}$, la propriété (2)(b) de Champs, Lemme
\ref{stacks-lemma-when-gerbe}, garantit l'existence d'une extension de $k$ sur
laquelle $x$ et $x'$ deviennent isomorphes (détails omis). Ainsi,
$|\mathcal{X}| \to |\mathcal{Y}|$ est continue et bijective ; il suffit de
montrer qu'elle est ouverte. Cela résulte des
Lemmes \ref{lemma-gerbe-fppf} et \ref{lemma-fppf-open}.
\end{proof}

\begin{lemma}
\label{lemma-gerbe-diagonal-quasi-compact}
Soit $f : \mathcal{X} \to \mathcal{Y}$ un morphisme de champs algébriques tel
que $\mathcal{X}$ soit une gerbe sur $\mathcal{Y}$. Si
$\Delta_\mathcal{X}$ est quasi-compact, alors $\Delta_\mathcal{Y}$ l'est aussi.
\end{lemma}

\begin{proof}
Considérons le diagramme
$$
\xymatrix{
\mathcal{X} \ar[r] &
\mathcal{X} \times_\mathcal{Y} \mathcal{X} \ar[r] \ar[d] &
\mathcal{X} \times \mathcal{X} \ar[d] \\
&
\mathcal{Y} \ar[r] &
\mathcal{Y} \times \mathcal{Y}
}
$$
La Proposition \ref{proposition-when-gerbe-over} montre que la flèche
supérieure de gauche est surjective. Puisque le composé des flèches
horizontales supérieures est quasi-compact, le
Lemme \ref{lemma-surjection-from-quasi-compact} montre que la flèche supérieure
de droite est quasi-compact. Le carré est cartésien et la flèche verticale de
droite est surjective, plate et localement de présentation finie. On conclut
donc par le Lemme \ref{lemma-descent-quasi-compact}.
\end{proof}

\noindent
Le lemme suivant affirme que les gerbes résiduelles existent en tout point
d'un champ algébrique qui est une gerbe.

\begin{lemma}
\label{lemma-gerbe-residual-gerbe-exists}
Soit $\mathcal{X}$ un champ algébrique. Si $\mathcal{X}$ est une gerbe, alors,
pour tout $x \in |\mathcal{X}|$, la gerbe résiduelle de $\mathcal{X}$ en $x$
existe.
\end{lemma}

\begin{proof}
Soit $\pi : \mathcal{X} \to X$ un morphisme de $\mathcal{X}$ vers un espace
algébrique $X$ qui fasse de $\mathcal{X}$ une gerbe sur $X$. Soit
$Z_x \to X$ l'espace résiduel de $X$ en $x$ ; voir Espaces décents, Définition
\ref{decent-spaces-definition-residual-space}. Posons
$\mathcal{Z} = \mathcal{X} \times_X Z_x$. D'après le
Lemme \ref{lemma-base-change-gerbe}, le champ algébrique $\mathcal{Z}$ est une
gerbe sur $Z_x$. Ainsi $|\mathcal{Z}| = |Z_x|$ est un singleton
(Lemme \ref{lemma-gerbe-bijection-points}). Comme
$\mathcal{Z} \to Z_x$ est localement de présentation finie, en tant que
changement de base de $\pi$ (voir les Lemmes \ref{lemma-gerbe-fppf} et
\ref{lemma-base-change-finite-presentation}), le champ $\mathcal{Z}$ est
localement noethérien ; voir le
Lemme \ref{lemma-locally-finite-type-locally-noetherian}. La gerbe résiduelle
$\mathcal{Z}_x$ de $\mathcal{X}$ en $x$ existe donc et vaut
$\mathcal{Z}_x = \mathcal{Z}_{red}$, la réduction du champ algébrique
$\mathcal{Z}$. En effet, nous avons vu que $|\mathcal{Z}_{red}|$ est un
singleton qui s'envoie sur $x \in |\mathcal{X}|$ ; ce champ est réduit par
construction et localement noethérien (la réduction d'un champ algébrique
localement noethérien est localement noethérienne ; détails omis).
\end{proof}









\section{Stratification par des gerbes}
\label{section-stratify}

\noindent
Le but de cette section est de montrer que de nombreux champs algébriques
$\mathcal{X}$ possèdent une « stratification » par des sous-champs
localement fermés $\mathcal{X}_i \subset \mathcal{X}$ tels que chaque
$\mathcal{X}_i$ soit une gerbe. Cela montre qu'en un certain sens les gerbes
sont les briques élémentaires dont est construit tout champ algébrique.
Notons que, par stratification, nous entendons ici uniquement que
$$
|\mathcal{X}| = \bigcup\nolimits_i |\mathcal{X}_i|
$$
est une stratification de l'espace topologique associé à $\mathcal{X}$,
et rien de plus (dans cette section). On peut donc, sans inconvénient,
remplacer $\mathcal{X}$ par sa réduction (voir Propriétés des champs,
section \ref{stacks-properties-section-reduced}) pour étudier cette
stratification.

\medskip\noindent
La proposition suivante affirme qu'il existe (presque toujours) un
sous-champ ouvert dense de la réduction de $\mathcal{X}$.

\begin{proposition}
\label{proposition-open-stratum}
Soit $\mathcal{X}$ un champ algébrique réduit tel que
$\mathcal{I}_\mathcal{X} \to \mathcal{X}$ soit quasi-compact.
Il existe alors un sous-champ ouvert dense
$\mathcal{U} \subset \mathcal{X}$ qui est une gerbe.
\end{proposition}

\begin{proof}
D'après la Proposition \ref{proposition-when-gerbe}, il suffit de trouver un
sous-champ ouvert dense $\mathcal{U}$ tel que
$\mathcal{I}_\mathcal{U} \to \mathcal{U}$ soit plat et localement de
présentation finie. Notons que
$\mathcal{I}_\mathcal{U} =
\mathcal{I}_\mathcal{X} \times_\mathcal{X} \mathcal{U}$ ; voir le
Lemme \ref{lemma-cartesian-square-inertia}.

\medskip\noindent
Choisissons une présentation $\mathcal{X} = [U/R]$. Soit $G \to U$
l'espace algébrique en groupes stabilisateur du groupoïde $R$. D'après le
Lemme \ref{lemma-presentation-inertia}, $G \to U$ est le changement de base
de $\mathcal{I}_\mathcal{X} \to \mathcal{X}$ ; il est donc quasi-compact
(par hypothèse) et localement de type fini (d'après le
Lemme \ref{lemma-inertia}). Soit $W \subset U$ le plus grand sous-schéma
ouvert (éventuellement vide) tel que la restriction $G_W \to W$ soit plate
et localement de présentation finie (nous omettons la preuve de l'existence de
$W$ ; indication : utiliser le caractère local des propriétés). D'après
Morphismes d'espaces, Proposition
\ref{spaces-morphisms-proposition-generic-flatness-reduced}
$W \subset U$ est dense. Notons que $W \subset U$ est $R$-invariant d'après
Compléments sur les groupoïdes en espaces, Lemme
\ref{spaces-more-groupoids-lemma-property-G-invariant}.
Ainsi, $W$ correspond à un sous-champ ouvert
$\mathcal{U} \subset \mathcal{X}$ d'après Propriétés des champs, Lemme
\ref{stacks-properties-lemma-substacks-presentation}.
Comme l'application $|U| \to |\mathcal{X}|$ est ouverte et que
$|W| \subset |U|$ est dense, $\mathcal{U}$ est dense dans $\mathcal{X}$.
Enfin, le morphisme $\mathcal{I}_\mathcal{U} \to \mathcal{U}$ est plat et
localement de présentation finie, car son changement de base par le morphisme
lisse surjectif $W \to \mathcal{U}$ est le morphisme $G_W \to W$, qui est
plat et localement de présentation finie par construction. Voir les
Lemmes \ref{lemma-descent-flat} et
\ref{lemma-descent-finite-presentation}.
\end{proof}

\noindent
La proposition précédente implique immédiatement que tout point d'un champ
algébrique à inertie quasi-compacte possède une gerbe résiduelle, comme nous
le montrerons dans le Lemme \ref{lemma-every-point-residual-gerbe}.
Il n'existe cependant pas toujours de stratification finie par des gerbes.
En voici un exemple.

\begin{example}
\label{example-infinite-stratification}
Soit $k$ un corps. Prenons $U = \Spec(k[x_0, x_1, x_2, \ldots])$ et faisons
agir $\mathbf{G}_m$ par
$t(x_0, x_1, x_2, \ldots) = (tx_0, t^p x_1, t^{p^2} x_2, \ldots)$
où $p$ est un nombre premier. Posons
$\mathcal{X} = [U/\mathbf{G}_m]$. C'est un champ algébrique. Il existe une
stratification de $\mathcal{X}$ dont les strates sont les suivantes :
\begin{enumerate}
\item $\mathcal{X}_0$ est le lieu où $x_0$ est non nul,
\item $\mathcal{X}_1$ est le lieu où $x_0$ est nul mais $x_1$ est non nul,
\item $\mathcal{X}_2$ est le lieu où $x_0, x_1$ sont nuls mais $x_2$ est non
nul,
\item et ainsi de suite, et
\item $\mathcal{X}_{\infty}$ est le lieu où tous les $x_i$ sont nuls.
\end{enumerate}
Chaque strate est une gerbe sur un schéma, de groupe $\mu_{p^i}$ pour
$\mathcal{X}_i$ et $\mathbf{G}_m$ pour $\mathcal{X}_{\infty}$.
Les strates sont des sous-champs localement fermés réduits. Il n'existe pas
de stratification plus grossière ayant les mêmes propriétés.
\end{example}

\noindent
Néanmoins, par induction transfinie, la
Proposition \ref{proposition-open-stratum} permet de construire des
stratifications par des gerbes, éventuellement infinies... !

\begin{lemma}
\label{lemma-every-point-in-a-stratum}
Soit $\mathcal{X}$ un champ algébrique tel que
$\mathcal{I}_\mathcal{X} \to \mathcal{X}$ soit quasi-compact.
Il existe alors un ensemble d'indices bien ordonné $I$ et, pour tout
$i \in I$, un sous-champ localement fermé réduit
$\mathcal{U}_i \subset \mathcal{X}$ tels que
\begin{enumerate}
\item chaque $\mathcal{U}_i$ est une gerbe,
\item on a $|\mathcal{X}| = \bigcup_{i \in I} |\mathcal{U}_i|$,
\item $T_i = |\mathcal{X}| \setminus \bigcup_{i' < i} |\mathcal{U}_{i'}|$
est fermé dans $|\mathcal{X}|$ pour tout $i \in I$, et
\item $|\mathcal{U}_i|$ est ouvert dans $T_i$.
\end{enumerate}
On peut de plus s'arranger pour que, soit (a)
$|\mathcal{U}_i| \subset T_i$ soit dense, soit (b) $\mathcal{U}_i$ soit
quasi-compact. Dans le cas (a), si l'on choisit $\mathcal{U}_i$ aussi grand
que possible (voir la démonstration pour les détails), la stratification est
canonique.
\end{lemma}

\begin{proof}
Soit $T \subset |\mathcal{X}|$ une partie fermée non vide. Pour tout $T$ de
cette sorte, nous allons trouver (resp.\ choisir) un sous-champ localement
fermé réduit $\mathcal{U}(T) \subset \mathcal{X}$ tel que
$|\mathcal{U}(T)| \subset T$ soit une partie ouverte dense
(resp.\ une partie ouverte non vide et quasi-compacte) de cette dernière. En effet,
d'après Propriétés des champs, Lemme
\ref{stacks-properties-lemma-reduced-closed-substack}
il existe un unique sous-champ fermé réduit
$\mathcal{X}' \subset \mathcal{X}$ tel que $T = |\mathcal{X}'|$.
Notons que $\mathcal{I}_{\mathcal{X}'} =
\mathcal{I}_\mathcal{X} \times_\mathcal{X} \mathcal{X}'$ d'après le
Lemme \ref{lemma-monomorphism-cartesian-square-inertia}.
Ainsi, $\mathcal{I}_{\mathcal{X}'} \to \mathcal{X}'$ est quasi-compact en
tant que changement de base ; voir le
Lemme \ref{lemma-base-change-quasi-compact}. La
Proposition \ref{proposition-open-stratum} implique donc qu'il existe un
sous-champ ouvert dense maximal $\mathcal{U} \subset \mathcal{X}'$ qui est
une gerbe (voir la démonstration de cette proposition). Dans le cas (a),
posons $\mathcal{U}(T) = \mathcal{U}$ (ce choix est canonique) ; dans le cas
(b), choisissons simplement un ouvert non vide quasi-compact
$\mathcal{U}(T) \subset \mathcal{U}$ ; voir Propriétés des champs, Lemme
\ref{stacks-properties-lemma-space-locally-quasi-compact}
(l'axiome du choix permet d'effectuer ces choix simultanément pour tous les
$T$).

\medskip\noindent
Par récurrence transfinie, construisons pour tout ordinal $\alpha$ une partie
fermée $T_\alpha \subset |\mathcal{X}|$. Pour $\alpha = 0$, posons
$T_0 = |\mathcal{X}|$. Une fois $T_\alpha$ construit, posons
$$
T_{\alpha + 1} = T_\alpha \setminus |\mathcal{U}(T_\alpha)|.
$$
Si $\beta$ est un ordinal limite, posons
$$
T_\beta = \bigcap\nolimits_{\alpha < \beta} T_\alpha.
$$
Nous affirmons que $T_\alpha = \emptyset$ pour tout $\alpha$ assez grand.
Supposons en effet que $T_\alpha \not = \emptyset$ pour tout $\alpha$.
On obtiendrait alors une application injective de la classe des ordinaux dans
l'ensemble des parties de $|\mathcal{X}|$, ce qui est contradictoire.

\medskip\noindent
L'affirmation implique le lemme. Posons en effet
$$
I = \{\alpha \mid \mathcal{U}(T_\alpha) \not = \emptyset \}.
$$
D'après l'affirmation, c'est un ensemble bien ordonné. Pour
$i = \alpha \in I$, posons $\mathcal{U}_i = \mathcal{U}(T_\alpha)$.
Ainsi, $\mathcal{U}_i$ est un sous-champ localement fermé réduit et une
gerbe, ce qui donne (1). Par construction, $T_i = T_\alpha$ si
$i = \alpha \in I$, d'où (3). De même, (4) ainsi que (a) ou (b) résultent
de notre choix de $\mathcal{U}(T)$. Enfin, pour établir (2), soit
$x \in |\mathcal{X}|$. Il existe un plus petit ordinal $\beta$ tel que
$x \not \in T_\beta$ (puisque les ordinaux sont bien ordonnés). Cet ordinal
$\beta$ est nécessairement un successeur d'après la définition de $T_\beta$
pour les ordinaux limites. Ainsi, $\beta = \alpha + 1$ et
$x \in |\mathcal{U}(T_\alpha)|$, ce qui achève la preuve.
\end{proof}

\begin{remark}
\label{remark-order-type}
On peut s'interroger sur le type d'ordre des stratifications canoniques
obtenues par la construction de type (a) du
Lemme \ref{lemma-every-point-in-a-stratum}. Il est naturel de conjecturer que
l'ensemble bien ordonné $I$ est de {\it cardinal} au plus $\aleph_0$.
Nous ignorons si cela est vrai ou faux. Si vous le savez, veuillez écrire à
\href{mailto:stacks.project@gmail.com}{stacks.project@gmail.com}.
\end{remark}







\section{L'espace topologique d'un champ algébrique}
\label{section-topology}

\noindent
Dans cette section, nous appliquons les résultats précédents à l'espace
topologique $|\mathcal{X}|$ associé à un champ algébrique.

\begin{lemma}
\label{lemma-spectral-qc-diagonal-qc}
Soit $\mathcal{X}$ un champ algébrique quasi-compact dont la diagonale
$\Delta$ est quasi-compacte. Alors $|\mathcal{X}|$ est un espace spectral.
\end{lemma}

\begin{proof}
Choisissons un schéma affine $U$ et un morphisme lisse surjectif
$U \to \mathcal{X}$ ; voir Propriétés des champs, Lemme
\ref{stacks-properties-lemma-quasi-compact-stack}. L'application
$|U| \to |\mathcal{X}|$ est alors continue, ouverte et surjective ; voir
Propriétés des champs, Lemme
\ref{stacks-properties-lemma-topology-points}. Les ouverts quasi-compacts de
$|\mathcal{X}|$ forment donc une base de la topologie. Pour deux ouverts
quasi-compacts $W_1, W_2 \subset |\mathcal{X}|$, on peut choisir des ouverts
quasi-compacts $V_1, V_2$ de $U$ dont les images sont respectivement $W_1$ et
$W_2$. Comme $\Delta$ est quasi-compacte,
$$
V_1 \times_\mathcal{X} V_2 =
(V_1 \times V_2) \times_{\mathcal{X} \times \mathcal{X}, \Delta} \mathcal{X}
$$
est quasi-compact. L'image de $|V_1 \times_\mathcal{X} V_2|$ dans
$|\mathcal{X}|$ est $W_1 \cap W_2$ d'après Propriétés des champs, Lemme
\ref{stacks-properties-lemma-points-cartesian}. Ainsi,
$W_1 \cap W_2$ est quasi-compact. Pour conclure, il suffit de montrer que
$|\mathcal{X}|$ est sobre ; voir Topologie, Définition
\ref{topology-definition-spectral-space}.

\medskip\noindent
Soit $T \subset |\mathcal{X}|$ une partie fermée irréductible. Il faut
montrer que $T$ possède un unique point générique. Soit
$\mathcal{Z} \subset \mathcal{X}$ le sous-champ fermé muni de la structure
réduite induite correspondant à $T$ ; voir Propriétés des champs, Définition
\ref{stacks-properties-definition-reduced-induced-stack}.
Comme $\mathcal{Z} \to \mathcal{X}$ est une immersion fermée,
$\Delta_\mathcal{Z}$ est quasi-compacte : montrons d'abord que
$\mathcal{Z} \to \mathcal{X} \times \mathcal{X}$ est quasi-compact comme
composé de $\mathcal{Z} \to \mathcal{X}$ et de $\Delta$, puis écrivons
$\mathcal{Z} \to \mathcal{X} \times \mathcal{X}$ comme le composé de
$\Delta_\mathcal{Z}$ et de
$\mathcal{Z} \times \mathcal{Z} \to \mathcal{X} \times \mathcal{X}$, et
appliquons le Lemme \ref{lemma-quasi-compact-permanence} ainsi que le fait que
$\mathcal{Z} \times \mathcal{Z} \to \mathcal{X} \times \mathcal{X}$ est
séparé. Nous sommes ainsi ramenés au cas traité dans le paragraphe suivant.

\medskip\noindent
Supposons $\mathcal{X}$ quasi-compact, $\Delta$ quasi-compacte,
$\mathcal{X}$ réduit et $|\mathcal{X}|$ irréductible. Il faut montrer que
$|\mathcal{X}|$ possède un unique point générique. Comme
$\mathcal{I}_\mathcal{X} \to \mathcal{X}$ est un changement de base de
$\Delta$, le morphisme
$\mathcal{I}_\mathcal{X} \to \mathcal{X}$ est quasi-compact
(Lemme \ref{lemma-base-change-quasi-compact}). D'après la
Proposition \ref{proposition-open-stratum}, il existe donc un sous-champ
ouvert dense $\mathcal{U} \subset \mathcal{X}$ qui est une gerbe. Autrement
dit, $|\mathcal{U}| \subset |\mathcal{X}|$ est ouvert dense. Nous pouvons
donc supposer que $\mathcal{X}$ est une gerbe. Soit
$\mathcal{X} \to X$ un morphisme faisant de $\mathcal{X}$ une gerbe sur
l'espace algébrique $X$.
D'après le Lemme \ref{lemma-gerbe-bijection-points}, on a alors
$|\mathcal{X}| \cong |X|$. En particulier, $X$ est quasi-compact. D'après le
Lemme \ref{lemma-gerbe-diagonal-quasi-compact}, la diagonale de $X$ est
quasi-compacte, c'est-à-dire que $X$ est un espace algébrique quasi-séparé.
Alors $|X|$ est spectral d'après Propriétés des espaces, Lemme
\ref{spaces-properties-lemma-quasi-compact-quasi-separated-spectral}
et le résultat en découle.
\end{proof}

\begin{lemma}
\label{lemma-spectral-qcqs}
Soit $\mathcal{X}$ un champ algébrique quasi-compact et quasi-séparé.
Alors $|\mathcal{X}|$ est un espace spectral.
\end{lemma}

\begin{proof}
C'est un cas particulier du Lemme \ref{lemma-spectral-qc-diagonal-qc}.
\end{proof}

\begin{lemma}
\label{lemma-sober-qs}
Soit $\mathcal{X}$ un champ algébrique dont la diagonale est quasi-compacte
(par exemple, si $\mathcal{X}$ est quasi-séparé). Il existe alors un
recouvrement ouvert $|\mathcal{X}| = \bigcup U_i$ tel que chaque $U_i$ soit
spectral. En particulier, $|\mathcal{X}|$ est un espace sobre.
\end{lemma}

\begin{proof}
Conséquence immédiate du Lemme \ref{lemma-spectral-qc-diagonal-qc}.
\end{proof}






\section{Existence des gerbes résiduelles}
\label{section-existence-residual-gerbes}

\noindent
La gerbe résiduelle d'un point d'un champ algébrique est définie dans
Propriétés des champs, Définition
\ref{stacks-properties-definition-residual-gerbe}.
Nous avons déjà montré que les gerbes résiduelles existent aux points
de type fini (Lemme \ref{lemma-point-finite-type-monomorphism})
et en tout point d'une gerbe (Lemme \ref{lemma-gerbe-residual-gerbe-exists}).
Dans cette section, nous montrons que les gerbes résiduelles existent sur
de nombreux champs algébriques. Commençons par l'application annoncée de la
Proposition \ref{proposition-open-stratum}.

\begin{lemma}
\label{lemma-every-point-residual-gerbe}
Soit $\mathcal{X}$ un champ algébrique tel que
$\mathcal{I}_\mathcal{X} \to \mathcal{X}$ soit quasi-compact.
Alors la gerbe résiduelle de $\mathcal{X}$ en $x$ existe pour tout
$x \in |\mathcal{X}|$.
\end{lemma}

\begin{proof}
Soit $T = \overline{\{x\}} \subset |\mathcal{X}|$ l'adhérence de $x$.
D'après Propriétés des champs, Lemme
\ref{stacks-properties-lemma-reduced-closed-substack}
il existe un sous-champ fermé réduit $\mathcal{X}' \subset \mathcal{X}$
tel que $T = |\mathcal{X}'|$. Notons que
$\mathcal{I}_{\mathcal{X}'} =
\mathcal{I}_\mathcal{X} \times_\mathcal{X} \mathcal{X}'$ d'après le
Lemme \ref{lemma-monomorphism-cartesian-square-inertia}.
Ainsi $\mathcal{I}_{\mathcal{X}'} \to \mathcal{X}'$ est quasi-compact
comme changement de base, voir le
Lemme \ref{lemma-base-change-quasi-compact}.
La Proposition \ref{proposition-open-stratum} implique donc qu'il existe
un sous-champ ouvert dense
$\mathcal{U} \subset \mathcal{X}'$
qui est une gerbe. Notons que $x \in |\mathcal{U}|$, car
$\{x\} \subset T$ est lui aussi dense. Il existe donc une gerbe résiduelle
$\mathcal{Z}_x \subset \mathcal{U}$ de $\mathcal{U}$ en $x$, d'après le
Lemme \ref{lemma-gerbe-residual-gerbe-exists}.
Il résulte immédiatement des définitions que $\mathcal{Z}_x \to \mathcal{X}$
est une gerbe résiduelle de $\mathcal{X}$ en $x$.
\end{proof}

\noindent
Si le champ est quasi-DM, les gerbes résiduelles existent également.
En particulier, elles existent toujours sur les champs de Deligne--Mumford.

\begin{lemma}
\label{lemma-every-point-residual-gerbe-quasi-DM}
Soit $\mathcal{X}$ un champ algébrique quasi-DM.
Alors la gerbe résiduelle de $\mathcal{X}$ en $x$ existe pour tout
$x \in |\mathcal{X}|$.
\end{lemma}

\begin{proof}
Choisissons un schéma $U$ et un morphisme $U \to \mathcal{X}$ surjectif,
plat, localement de présentation finie et localement quasi-fini, voir le
Théorème \ref{theorem-quasi-DM}.
Posons $R = U \times_\mathcal{X} U$. Les projections $s, t : R \to U$
sont surjectives, plates, localement de présentation finie et localement
quasi-finies, car ce sont des changements de base du morphisme
$U \to \mathcal{X}$. Il existe un morphisme canonique
$[U/R] \to \mathcal{X}$ (voir Champs algébriques, Lemme
\ref{algebraic-lemma-map-space-into-stack}), qui est une équivalence puisque
$U \to \mathcal{X}$ est surjectif, plat et localement de présentation finie,
voir Champs algébriques, Remarque
\ref{algebraic-remark-flat-fp-presentation}.
Nous pouvons donc supposer que $\mathcal{X} = [U/R]$, où
$(U, R, s, t, c)$ est un groupoïde en espaces algébriques tel que
$s, t : R \to U$ soient surjectifs, plats, localement de présentation finie
et localement quasi-finis. Posons
$$
U' = \coprod\nolimits_{u \in U\text{ se projetant sur }x} \Spec(\kappa(u)).
$$
Le morphisme canonique $U' \to U$ est un monomorphisme. Posons
$$
R' = U' \times_\mathcal{X} U' =
R \times_{(U \times U)} (U' \times U')
$$
Comme $U' \to U$ est un monomorphisme, chacune des deux projections
$s', t' : R' \to U'$ se factorise en un monomorphisme suivi d'un morphisme
localement quasi-fini. Puisque $U'$ est une somme disjointe de spectres
de corps, le Lemme
\ref{spaces-over-fields-lemma-mono-towards-locally-quasi-finite-over-field}
d'Espaces sur des corps montre que les morphismes
$s', t' : R' \to U'$ sont localement quasi-finis. Comme $U'$ est une
somme disjointe de spectres de corps, les morphismes $s', t'$ sont aussi
plats. Enfin, comme $s', t'$ sont localement quasi-finis, $s', t'$ sont
localement de type fini et, par suite, $s', t'$ sont localement de
présentation finie (en effet,
$U'$ est une somme disjointe de spectres de corps et, en particulier,
est localement noethérien, si bien que le Lemme
\ref{spaces-morphisms-lemma-noetherian-finite-type-finite-presentation}
de Morphismes d'espaces s'applique). Ainsi $\mathcal{Z} = [U'/R']$ est un
champ algébrique d'après Critères de représentabilité, Théorème
\ref{criteria-theorem-flat-groupoid-gives-algebraic-stack}.
Comme $R'$ est la restriction de $R$ par $U' \to U$, le morphisme
$\mathcal{Z} \to \mathcal{X}$ est un monomorphisme d'après Groupoïdes en
espaces, Lemme
\ref{spaces-groupoids-lemma-quotient-stack-restrict}
et Propriétés des champs, Lemme
\ref{stacks-properties-lemma-monomorphism}.
Puisque $\mathcal{Z} \to \mathcal{X}$ est un monomorphisme, l'application
$|\mathcal{Z}| \to |\mathcal{X}|$ est injective, voir Propriétés des champs,
Lemme
\ref{stacks-properties-lemma-monomorphism-injective-points}.
D'après Propriétés des champs, Lemme
\ref{stacks-properties-lemma-points-cartesian}, l'application
$$
|U'| = |\mathcal{Z} \times_\mathcal{X} U'|
\longrightarrow
|\mathcal{Z}| \times_{|\mathcal{X}|} |U'|
$$
est surjective. Par le choix de $U'$, il en résulte que l'image de
$|\mathcal{Z}| \to |\mathcal{X}|$ est $\{x\}$.
Ainsi $|\mathcal{Z}|$ est réduit à un point.
Enfin, $U'$ est par construction localement noethérien et réduit ; autrement
dit, $\mathcal{Z}$ est réduit et localement noethérien. L'image essentielle
de $\mathcal{Z} \to \mathcal{X}$ est donc la gerbe résiduelle de
$\mathcal{X}$ en $x$, voir Propriétés des champs, Lemme
\ref{stacks-properties-lemma-residual-gerbe-unique}.
\end{proof}

\begin{lemma}
\label{lemma-every-point-residual-gerbe-locally-Noetherian}
Soit $\mathcal{X}$ un champ algébrique localement noethérien.
Alors la gerbe résiduelle de $\mathcal{X}$ en $x$ existe pour tout
$x \in |\mathcal{X}|$.
\end{lemma}

\begin{proof}
Choisissons un schéma affine $U$ et un morphisme lisse
$U \to \mathcal{X}$ tels que $x$ appartienne à l'image de l'application
continue ouverte $|U| \to |\mathcal{X}|$. Remplaçons $\mathcal{X}$ par le
sous-champ ouvert correspondant à l'image de $|U| \to |\mathcal{X}|$, voir
Propriétés des champs, Lemme \ref{stacks-properties-lemma-open-substacks}.
Nous pouvons ainsi supposer que $\mathcal{X} = [U/R]$ pour un groupoïde lisse
$(U, R, s, t, c)$ en espaces algébriques, où $U$ est un schéma affine
noethérien, voir Champs algébriques, section
\ref{algebraic-section-stack-to-presentation}.

\medskip\noindent
Soit $E \subset |U|$ l'image réciproque de
$\{x\} \subset |\mathcal{X}|$. Bien sûr, $E \not = \emptyset$.
Puisque $|U|$ est un espace topologique noethérien, nous pouvons choisir
$u \in E$ tel que $\overline{\{u\}} \cap E = \{u\}$.
Comme d'habitude, nous identifions $u = \Spec(\kappa(u))$ au spectre de son
corps résiduel. Écrivons
$$
F = u \times_{U, t} R = u \times_\mathcal{X} U
\quad\text{et}\quad
R' = (u \times u) \times_{(U \times U), (t, s)} R = u \times_\mathcal{X} u
$$
Munissons en outre $Z = \overline{\{u\}} \subset U$ de la structure réduite
de sous-schéma induite. Notons $p : F \to U$ le morphisme induit par la
seconde projection (on utilise $s : R \to U$ dans la première description
de $F$ comme produit fibré). Alors $E$ est l'image ensembliste de $p$.
Le morphisme $R' \to F$ est un monomorphisme qui se factorise par l'image
réciproque $p^{-1}(Z)$ de $Z$. Cette image réciproque
$p^{-1}(Z) \subset F$ est un sous-schéma fermé ; l'image de la restriction
$p|_{p^{-1}(Z)} : p^{-1}(Z) \to Z$ est, ensemblistement, contenue dans
$\{u\} \subset Z$ par le choix soigneux de $u \in E$ fait plus haut.
Puisque $u = \lim W$, où la limite porte sur les
sous-schémas ouverts affines non vides du schéma réduit irréductible $Z$,
le morphisme $p|_{p^{-1}(Z)} : p^{-1}(Z) \to Z$ se factorise par
$u \to Z$. Cela implique clairement que $R' = p^{-1}(Z)$. En particulier,
le morphisme $t' : R' \to u$ est localement de présentation finie : c'est la
composée de l'immersion fermée $p^{-1}(Z) \to F$ d'espaces algébriques
localement noethériens et du morphisme lisse $\text{pr}_1 : F \to u$ ; voir
Morphismes d'espaces, Lemmes
\ref{spaces-morphisms-lemma-locally-finite-type-locally-noetherian},
\ref{spaces-morphisms-lemma-closed-immersion-finite-presentation} et
\ref{spaces-morphisms-lemma-composition-finite-presentation}.
La restriction $(u, R', s', t', c')$ de $(U, R, s, t, c)$ par
$u \to U$ est donc un groupoïde en espaces algébriques dont $s'$ et $t'$
sont plats et localement de présentation finie. Ainsi
$\mathcal{Z} = [u/R']$ est un champ algébrique d'après Critères de
représentabilité, Théorème
\ref{criteria-theorem-flat-groupoid-gives-algebraic-stack}.
Comme $R'$ est la restriction de $R$ par $u \to U$, le morphisme
$\mathcal{Z} \to \mathcal{X}$ est un monomorphisme d'après Groupoïdes en
espaces, Lemme
\ref{spaces-groupoids-lemma-quotient-stack-restrict}
et Propriétés des champs, Lemme
\ref{stacks-properties-lemma-monomorphism}.
Le champ $\mathcal{Z}$ est alors isomorphe à la gerbe résiduelle d'après les
résultats de Propriétés des champs, section
\ref{stacks-properties-section-residual-gerbe}.
\end{proof}













\section{Structure locale \'etale}
\label{section-etale-local}

\noindent
Dans cette section, nous commen\c{c}ons \`a examiner ce que l'on peut dire
de la structure locale \'etale d'un champ alg\'ebrique.

\begin{lemma}
\label{lemma-quotient-etale}
Soit $Y$ un espace alg\'ebrique.
Soit $(U, R, s, t, c)$ un groupo\"ide en espaces alg\'ebriques sur $Y$.
Supposons que $U \to Y$ soit plat et localement de pr\'esentation finie
et que $R \to U \times_Y U$ soit une immersion ouverte.
Alors $X = [U/R] = U/R$ est un espace alg\'ebrique et $X \to Y$
est \'etale.
\end{lemma}

\begin{proof}
Le champ quotient $[U/R]$ est un champ alg\'ebrique d'apr\`es
Crit\`eres de repr\'esentabilit\'e, th\'eor\`eme
\ref{criteria-theorem-flat-groupoid-gives-algebraic-stack}.
D'autre part, puisque $R \to U \times U$ est un monomorphisme,
ce champ est un espace alg\'ebrique (en vertu de notre abus de langage et d'apr\`es
Champs alg\'ebriques, proposition
\ref{algebraic-proposition-algebraic-stack-no-automorphisms})
et il est bien s\^ur \'egal \`a l'espace alg\'ebrique $U/R$
(d'Amor\c{c}age, th\'eor\`eme \ref{bootstrap-theorem-final-bootstrap}).
Comme $U \to X$ est surjectif, plat et localement de pr\'esentation finie
(Amor\c{c}age, lemme \ref{bootstrap-lemma-covering-quotient}),
nous en d\'eduisons que $X \to Y$ est plat et localement de pr\'esentation finie d'apr\`es
Morphismes d'espaces, lemme \ref{spaces-morphisms-lemma-flat-permanence}
et
Descente sur les espaces, lemme
\ref{spaces-descent-lemma-flat-finitely-presented-permanence}.
Enfin, consid\'erons le diagramme cart\'esien
$$
\xymatrix{
R \ar[d] \ar[r] & U \times_Y U \ar[d] \\
X \ar[r] & X \times_Y X
}
$$
Puisque la fl\`eche verticale de droite est surjective, plate et
localement de pr\'esentation finie (petit d\'etail omis), nous
voyons que $X \to X \times_Y X$ est une immersion ouverte, car la fl\`eche
horizontale sup\'erieure poss\`ede cette propri\'et\'e par hypoth\`ese (utiliser
Propri\'et\'es des champs, lemme
\ref{stacks-properties-lemma-check-property-covering}).
Ainsi $X \to Y$ est non ramifi\'e d'apr\`es
Morphismes d'espaces, lemme
\ref{spaces-morphisms-lemma-diagonal-unramified-morphism}.
Nous concluons \`a l'aide de
Morphismes d'espaces, lemme
\ref{spaces-morphisms-lemma-unramified-flat-lfp-etale}.
\end{proof}

\begin{lemma}
\label{lemma-quasi-splitting-etale}
Soit $S$ un sch\'ema.
Soit $(U, R, s, t, c)$ un groupo\"ide en espaces alg\'ebriques sur $S$.
Supposons que $s, t$ soient plats et localement de pr\'esentation finie.
Soit $P \subset R$ un sous-espace ouvert tel que
$(U, P, s|_P, t|_P, c|_{P \times_{s, U, t} P})$ soit un
groupo\"ide en espaces alg\'ebriques sur $S$. Alors
$$
[U/P] \longrightarrow [U/R]
$$
est un morphisme de champs alg\'ebriques qui
est repr\'esentable par des espaces alg\'ebriques, surjectif et \'etale.
\end{lemma}

\begin{proof}
Puisque $P \subset R$ est ouvert, $s|_P$ et $t|_P$
sont plats et localement de pr\'esentation finie.
Ainsi $[U/R]$ et $[U/P]$ sont des champs alg\'ebriques d'apr\`es
Crit\`eres de repr\'esentabilit\'e, th\'eor\`eme
\ref{criteria-theorem-flat-groupoid-gives-algebraic-stack}.
Pour voir que le morphisme est repr\'esentable par des espaces alg\'ebriques,
il suffit de montrer que $[U/P] \to [U/R]$ est fid\`ele sur
chaque cat\'egorie fibre ; voir
Champs alg\'ebriques, lemme
\ref{algebraic-lemma-characterize-representable-by-algebraic-spaces}.
Cela r\'esulte imm\'ediatement du fait que $P \to R$ est un monomorphisme
et de la description explicite des champs quotients donn\'ee dans
Groupo\"ides en espaces, lemme
\ref{spaces-groupoids-lemma-quotient-stack-objects}.
Ceci \'etabli, nous savons ce que signifie le fait que
$[U/P] \to [U/R]$ soit surjectif et \'etale, d'apr\`es
Champs alg\'ebriques, d\'efinition
\ref{algebraic-definition-relative-representable-property}.
La surjectivit\'e r\'esulte par exemple de
Crit\`eres de repr\'esentabilit\'e,
lemme \ref{criteria-lemma-representable-by-spaces-surjective}
et de la description des objets des champs quotients
(voir le lemme cit\'e ci-dessus) sur les spectres de corps.
Il reste \`a prouver que notre morphisme est \'etale.

\medskip\noindent
Pour cela, il suffit de montrer que $U \times_{[U/R]} [U/P] \to U$
est \'etale ; voir Propri\'et\'es des champs, lemme
\ref{stacks-properties-lemma-check-property-covering}.
D'apr\`es Groupo\"ides en espaces, lemme
\ref{spaces-groupoids-lemma-cartesian-square-of-morphism},
le produit fibr\'e s'identifie \`a $[R/P \times_{s, U, t} R]$,
le morphisme vers $U$ \'etant induit par $s : R \to U$.
Les applications $s', t' : P \times_{s, U, t} R \to R$ sont d\'efinies par
$s' : (p, r) \mapsto r$ et $t' : (p, r) \mapsto c(p, r)$. Puisque
$P \subset R$ est ouvert, nous en d\'eduisons que
$(t', s') : P \times_{s, U, t} R \to R \times_{s, U, s} R$
est une immersion ouverte.
Nous pouvons donc appliquer le lemme \ref{lemma-quotient-etale}
pour conclure.
\end{proof}

\begin{lemma}
\label{lemma-etale-local-quasi-DM}
Soit $\mathcal{X}$ un champ alg\'ebrique. Supposons que $\mathcal{X}$ soit
quasi-DM \`a diagonale s\'epar\'ee (de mani\`ere \'equivalente,
$\mathcal{I}_\mathcal{X} \to \mathcal{X}$ est localement quasi-fini et
s\'epar\'e). Soit $x \in |\mathcal{X}|$. Il existe alors un
morphisme de champs alg\'ebriques
$$
\mathcal{U} \longrightarrow \mathcal{X}
$$
ayant les propri\'et\'es suivantes :
\begin{enumerate}
\item il existe un point $u \in |\mathcal{U}|$ dont l'image est $x$,
\item $\mathcal{U} \to \mathcal{X}$ est repr\'esentable par des espaces alg\'ebriques et
\'etale,
\item $\mathcal{U} = [U/R]$, o\`u $(U, R, s, t, c)$ est un sch\'ema
en groupo\"ides, avec $U$ et $R$ affines, et $s, t$ finis, plats et
localement de pr\'esentation finie.
\end{enumerate}
\end{lemma}

\begin{proof}
(L'assertion entre parenth\`eses r\'esulte des \'equivalences du
lemme \ref{lemma-diagonal-diagonal}).
Choisissons un sch\'ema affine $U$ et un morphisme $U \to \mathcal{X}$
plat, localement de pr\'esentation finie et localement quasi-fini, tel que $x$
soit l'image d'un point $u \in U$. Cela est possible d'apr\`es
le th\'eor\`eme \ref{theorem-quasi-DM} et l'hypoth\`ese que $\mathcal{X}$
est quasi-DM. Soit $(U, R, s, t, c)$ le groupo\"ide en espaces alg\'ebriques
obtenu en posant $R = U \times_\mathcal{X} U$ ; voir
Champs alg\'ebriques, lemme \ref{algebraic-lemma-map-space-into-stack}.
Soit $\mathcal{X}' \subset \mathcal{X}$ le sous-champ ouvert correspondant
\`a l'image ouverte de $|U| \to |\mathcal{X}|$
(Propri\'et\'es des champs, lemmes
\ref{stacks-properties-lemma-topology-points} et
\ref{stacks-properties-lemma-open-substacks}).
Il est clair que nous pouvons remplacer $\mathcal{X}$ par le sous-champ ouvert $\mathcal{X}'$.
Nous pouvons donc supposer que $U \to \mathcal{X}$ est surjectif ; alors
Champs alg\'ebriques, remarque \ref{algebraic-remark-flat-fp-presentation}
donne $\mathcal{X} = [U/R]$.
Remarquons que $s, t : R \to U$ sont plats, localement de pr\'esentation finie
et localement quasi-finis.
Comme $R = U \times U \times_{(\mathcal{X} \times \mathcal{X})} \mathcal{X}$
et que la diagonale de $\mathcal{X}$ est s\'epar\'ee, nous voyons que
$R$ est s\'epar\'e. Par cons\'equent, $s, t : R \to U$ sont s\'epar\'es. Il s'ensuit
que $R$ est un sch\'ema d'apr\`es
Morphismes d'espaces, proposition
\ref{spaces-morphisms-proposition-locally-quasi-finite-separated-over-scheme},
appliqu\'ee \`a $s : R \to U$.

\medskip\noindent
Nous avons v\'erifi\'e ci-dessus que toutes les hypoth\`eses de
Compl\'ements sur les groupo\"ides en espaces, lemme
\ref{spaces-more-groupoids-lemma-quasi-splitting-affine-scheme}
sont satisfaites pour $(U, R, s, t, c)$ et $u$.
Nous pouvons donc trouver un voisinage \'etale \'el\'ementaire
$(U', u') \to (U, u)$ tel que la restriction $R'$ de $R$ \`a $U'$
soit quasi-scind\'ee au-dessus de $u'$. Remarquons que $R' = U' \times_\mathcal{X} U'$
(petit d\'etail omis ; indication : transitivit\'e des produits fibr\'es).
En rempla\c{c}ant $(U, R, s, t, c)$ par $(U', R', s', t', c')$ et en r\'etr\'ecissant
$\mathcal{X}$ comme ci-dessus, nous pouvons supposer que $(U, R, s, t, c)$ admet
un quasi-scindage au-dessus de $u$ (le point $u$ ne joue d\'esormais plus
aucun r\^ole, comme le montre la note de bas de page de
Compl\'ements sur les groupo\"ides en espaces, d\'efinition
\ref{spaces-more-groupoids-definition-split-at-point}).
Soit $P \subset R$ un quasi-scindage de $R$ au-dessus de $u$.
Le lemme \ref{lemma-quasi-splitting-etale}
montre que
$$
\mathcal{U} = [U/P] \longrightarrow [U/R] = \mathcal{X}
$$
poss\`ede toutes les propri\'et\'es voulues.
\end{proof}

\begin{lemma}
\label{lemma-etale-local-quasi-DM-at-x}
Soit $\mathcal{X}$ un champ alg\'ebrique. Supposons que $\mathcal{X}$ soit
quasi-DM \`a diagonale s\'epar\'ee (de mani\`ere \'equivalente,
$\mathcal{I}_\mathcal{X} \to \mathcal{X}$ est localement quasi-fini et
s\'epar\'e). Soit $x \in |\mathcal{X}|$. Supposons que le
groupe d'automorphismes de $\mathcal{X}$ en $x$ soit fini
(remarque \ref{remark-property-automorphism-groups}).
Il existe alors un morphisme de champs alg\'ebriques
$$
g : \mathcal{U} \longrightarrow \mathcal{X}
$$
ayant les propri\'et\'es suivantes :
\begin{enumerate}
\item il existe un point $u \in |\mathcal{U}|$ dont l'image est $x$, et
$g$ induit un isomorphisme entre les groupes d'automorphismes en $u$ et en $x$
(remarque \ref{remark-identify-automorphism-groups}),
\item $\mathcal{U} \to \mathcal{X}$ est repr\'esentable par des espaces alg\'ebriques et
\'etale,
\item $\mathcal{U} = [U/R]$, o\`u $(U, R, s, t, c)$ est un sch\'ema
en groupo\"ides, avec $U$ et $R$ affines, et $s, t$ finis, plats et
localement de pr\'esentation finie.
\end{enumerate}
\end{lemma}

\begin{proof}
Remarquons que $G_x$ est un sch\'ema en groupes d'apr\`es
le lemme \ref{lemma-automorphism-group-scheme}.
Le d\'ebut de cette preuve est {\bf exactement} le m\^eme que celui de la preuve
du lemme \ref{lemma-etale-local-quasi-DM}.
Nous pouvons donc supposer que $\mathcal{X} = [U/R]$, o\`u $(U, R, s, t, c)$
et $u \in U$, d'image $x$, satisfont \`a toutes les hypoth\`eses de
Compl\'ements sur les groupo\"ides en espaces, lemme
\ref{spaces-more-groupoids-lemma-quasi-splitting-affine-scheme}.
Notre hypoth\`ese sur $G_x$ implique que $G_u$ est fini sur $u$.
Toutes les hypoth\`eses de
Compl\'ements sur les groupo\"ides en espaces, lemme
\ref{spaces-more-groupoids-lemma-splitting-affine-scheme}
sont donc satisfaites.
Nous pouvons par cons\'equent trouver un voisinage \'etale \'el\'ementaire
$(U', u') \to (U, u)$ tel que la restriction $R'$ de $R$ \`a $U'$
soit scind\'ee au-dessus de $u'$. Remarquons que $R' = U' \times_\mathcal{X} U'$
(petit d\'etail omis ; indication : transitivit\'e des produits fibr\'es).
En rempla\c{c}ant $(U, R, s, t, c)$ par $(U', R', s', t', c')$ et en r\'etr\'ecissant
$\mathcal{X}$ comme ci-dessus, nous pouvons supposer que $(U, R, s, t, c)$ admet
un scindage au-dessus de $u$. Soit $P \subset R$ un scindage de $R$ au-dessus de $u$.
Le lemme \ref{lemma-quasi-splitting-etale} montre que
$$
\mathcal{U} = [U/P] \longrightarrow [U/R] = \mathcal{X}
$$
est repr\'esentable par des espaces alg\'ebriques et \'etale. Par construction,
$G_u$ est contenu dans $P$ ; ce morphisme induit donc l'isomorphisme voulu
du groupe d'automorphismes en $u$ sur celui du point image.
\end{proof}

\begin{lemma}
\label{lemma-etale-local-quasi-DM-at-x-inertia}
Soit $\mathcal{X}$ un champ alg\'ebrique. Supposons que $\mathcal{X}$ soit
quasi-DM \`a diagonale s\'epar\'ee (de mani\`ere \'equivalente,
$\mathcal{I}_\mathcal{X} \to \mathcal{X}$ est localement quasi-fini et
s\'epar\'e). Soit $x \in |\mathcal{X}|$. Supposons que $x$ puisse \^etre repr\'esent\'e
par un morphisme quasi-compact $\Spec(k) \to \mathcal{X}$.
Il existe alors un morphisme de champs alg\'ebriques
$$
g : \mathcal{U} \longrightarrow \mathcal{X}
$$
ayant les propri\'et\'es suivantes :
\begin{enumerate}
\item il existe un point $u \in |\mathcal{U}|$ dont l'image est $x$, et
$g$ induit un isomorphisme entre les gerbes r\'esiduelles en $u$ et en $x$,
\item $\mathcal{U} \to \mathcal{X}$ est repr\'esentable par des espaces alg\'ebriques et
\'etale,
\item $\mathcal{U} = [U/R]$, o\`u $(U, R, s, t, c)$ est un sch\'ema
en groupo\"ides, avec $U$ et $R$ affines, et $s, t$ finis, plats et
localement de pr\'esentation finie.
\end{enumerate}
\end{lemma}

\begin{proof}
Le d\'ebut de cette preuve est {\bf exactement} le m\^eme que celui de la preuve
du lemme \ref{lemma-etale-local-quasi-DM}.
Nous pouvons donc supposer que $\mathcal{X} = [U/R]$, o\`u $(U, R, s, t, c)$
et $u \in U$, d'image $x$, satisfont \`a toutes les hypoth\`eses de
Compl\'ements sur les groupo\"ides en espaces, lemme
\ref{spaces-more-groupoids-lemma-quasi-splitting-affine-scheme}.
Remarquons que $u = \Spec(\kappa(u)) \to \mathcal{X}$ est quasi-compact ; voir
Propri\'et\'es des champs, lemme
\ref{stacks-properties-lemma-UR-quasi-compact-above-x}.
Consid\'erons le diagramme cart\'esien
$$
\xymatrix{
F \ar[d] \ar[r] & U \ar[d] \\
u \ar[r]^u & \mathcal{X}
}
$$
Puisque $U$ est un sch\'ema affine et que $F \to U$ est quasi-compact,
$F$ est quasi-compact. Puisque $U \to \mathcal{X}$
est localement quasi-fini, $F \to u$ est
localement quasi-fini. Ainsi $F \to u$ est quasi-fini,
et $F$ est un sch\'ema affine dont l'espace topologique
sous-jacent est fini et discret (Espaces sur des corps, lemme
\ref{spaces-over-fields-lemma-locally-quasi-finite-over-field}).
Remarquons que nous avons un monomorphisme $u \times_\mathcal{X} u \to F$.
En particulier, l'ensemble $\{r \in R : s(r) = u, t(r) = u\}$,
qui est l'image de $|u \times_\mathcal{X} u| \to |R|$,
est fini. Nous en concluons que toutes les hypoth\`eses de
Compl\'ements sur les groupo\"ides en espaces, lemme
\ref{spaces-more-groupoids-lemma-strong-splitting-affine-scheme}
sont satisfaites.

\medskip\noindent
Nous pouvons donc trouver un voisinage \'etale \'el\'ementaire
$(U', u') \to (U, u)$ tel que la restriction $R'$ de $R$ \`a $U'$
soit fortement scind\'ee au-dessus de $u'$. Remarquons que $R' = U' \times_\mathcal{X} U'$
(petit d\'etail omis ; indication : transitivit\'e des produits fibr\'es).
En rempla\c{c}ant $(U, R, s, t, c)$ par $(U', R', s', t', c')$ et en r\'etr\'ecissant
$\mathcal{X}$ comme ci-dessus, nous pouvons supposer que $(U, R, s, t, c)$ admet
un scindage fort au-dessus de $u$. Soit $P \subset R$ un scindage fort
de $R$ au-dessus de $u$. Le lemme \ref{lemma-quasi-splitting-etale} montre que
$$
\mathcal{U} = [U/P] \longrightarrow [U/R] = \mathcal{X}
$$
est repr\'esentable par des espaces alg\'ebriques et \'etale. Puisque $P \subset R$
est ouvert et contient $\{r \in R : s(r) = u, t(r) = u\}$ par construction,
nous voyons que
$u \times_\mathcal{U} u \to u \times_\mathcal{X} u$ est un isomorphisme.
L'assertion sur les gerbes r\'esiduelles r\'esulte alors de
Propri\'et\'es des champs, lemme
\ref{stacks-properties-lemma-residual-gerbe-isomorphic}
(nous remarquons que les gerbes r\'esiduelles en question existent d'apr\`es
le lemme \ref{lemma-every-point-residual-gerbe-quasi-DM}).
\end{proof}



\section{Morphismes lisses}
\label{section-smooth}

\noindent
La propriété ``être lisse'' des morphismes d'espaces
algébriques est locale pour la topologie lisse sur la source et la cible, voir
Descente sur les espaces, Remarque \ref{spaces-descent-remark-list-local-source-target}.
Elle est également stable par changement de base et locale pour la topologie fpqc sur la cible, voir
Morphismes d'espaces,
Lemme \ref{spaces-morphisms-lemma-base-change-smooth}
et
Descente sur les espaces, Lemme
\ref{spaces-descent-lemma-descending-property-smooth}.
Par conséquent, d'après le
Lemme \ref{lemma-local-source-target}
ci-dessus, nous pouvons définir comme suit ce que signifie être lisse pour un
morphisme de champs algébriques, et cette définition coïncide avec la notion
déjà définie dans
Propriétés des champs,
section \ref{stacks-properties-section-properties-morphisms}
lorsque le morphisme est représentable par des espaces algébriques.

\begin{definition}
\label{definition-smooth}
Soit $f : \mathcal{X} \to \mathcal{Y}$ un morphisme de champs algébriques.
Nous disons que $f$ est {\it lisse} si les conditions équivalentes du
Lemme \ref{lemma-local-source-target}
sont satisfaites avec $\mathcal{P} = \text{lisse}$.
\end{definition}

\begin{lemma}
\label{lemma-composition-smooth}
Le composé de morphismes lisses est lisse.
\end{lemma}

\begin{proof}
Le résultat découle de
la Remarque \ref{remark-composition}
et du
Lemme \ref{spaces-morphisms-lemma-composition-smooth}
de Morphismes d'espaces.
\end{proof}

\begin{lemma}
\label{lemma-base-change-smooth}
Un changement de base d'un morphisme lisse est lisse.
\end{lemma}

\begin{proof}
Le résultat découle de
la Remarque \ref{remark-base-change}
et du
Lemme \ref{spaces-morphisms-lemma-base-change-smooth}
de Morphismes d'espaces.
\end{proof}

\begin{lemma}
\label{lemma-descent-smooth}
Soit $f : \mathcal{X} \to \mathcal{Y}$ un morphisme de champs algébriques.
Soit $\mathcal{Z} \to \mathcal{Y}$ un morphisme de champs algébriques surjectif,
plat et localement de présentation finie. Si le changement de base
$\mathcal{Z} \times_\mathcal{Y} \mathcal{X} \to \mathcal{Z}$
est lisse, alors $f$ est lisse.
\end{lemma}

\begin{proof}
La propriété ``être lisse''
satisfait aux conditions du Lemme \ref{lemma-descent-property}.
Nous avons vu dans l'introduction de cette section qu'elle est locale pour la
topologie lisse sur la source et la cible ; le caractère local pour la topologie fppf sur la cible résulte de
Descente sur les espaces, Lemme
\ref{spaces-descent-lemma-descending-property-smooth}.
\end{proof}

\begin{lemma}
\label{lemma-smooth-locally-finite-presentation}
Un morphisme lisse de champs algébriques est localement de présentation finie.
\end{lemma}

\begin{proof}
Démonstration omise.
\end{proof}

\begin{lemma}
\label{lemma-where-smooth}
Soit $f : \mathcal{X} \to \mathcal{Y}$ un morphisme de champs algébriques.
Il existe un plus grand sous-champ ouvert $\mathcal{U} \subset \mathcal{X}$
tel que $f|_\mathcal{U} : \mathcal{U} \to \mathcal{Y}$ soit lisse.
En outre, la formation de cet ouvert commute aux opérations suivantes :
\begin{enumerate}
\item la précomposition par des morphismes lisses,
\item le changement de base par des morphismes plats et localement de
présentation finie,
\item le changement de base par des morphismes plats, pourvu que $f$ soit localement de
présentation finie.
\end{enumerate}
\end{lemma}

\begin{proof}
Choisissons un diagramme commutatif
$$
\xymatrix{
U \ar[d]_a \ar[r]_h & V \ar[d]^b \\
\mathcal{X} \ar[r]^f & \mathcal{Y}
}
$$
où $U$ et $V$ sont des espaces algébriques, les flèches verticales sont lisses
et $a : U \to \mathcal{X}$ est surjectif. Il existe un ouvert maximal
$U' \subset U$ tel que $h_{U'} : U' \to V$ soit lisse, voir
Morphismes d'espaces, Lemme \ref{spaces-morphisms-lemma-where-smooth}.
Soit $\mathcal{U} \subset \mathcal{X}$ le sous-champ ouvert
correspondant à l'image de $|U'| \to |\mathcal{X}|$
(Propriétés des champs, Lemmes \ref{stacks-properties-lemma-topology-points} et
\ref{stacks-properties-lemma-open-substacks}).
D'après l'équivalence du Lemme \ref{lemma-local-source-target},
nous voyons que $f|_\mathcal{U} : \mathcal{U} \to \mathcal{Y}$ est lisse
et que $\mathcal{U}$ est le plus grand sous-champ ouvert possédant cette
propriété.

\medskip\noindent
L'assertion (1) résulte du fait que la propriété d'être lisse
est locale pour la topologie lisse sur la source (cette propriété a même servi à définir les
morphismes lisses de champs algébriques).
Les assertions (2) et (3) résultent du cas des espaces algébriques, voir
Morphismes d'espaces, Lemme \ref{spaces-morphisms-lemma-where-smooth}.
\end{proof}

\begin{lemma}
\label{lemma-smooth-quotient-stack}
Soit $X \to Y$ un morphisme lisse d'espaces algébriques.
Soit $G$ un espace algébrique en groupes sur $Y$, plat
et localement de présentation finie sur $Y$. Soit donnée une action de $G$ sur $X$ au-dessus de $Y$.
Alors le champ quotient $[X/G]$ est lisse sur $Y$.
\end{lemma}

\noindent
Cela reste vrai même si $G$ n'est pas lisse sur $Y$ !

\begin{proof}
Le champ quotient $[X/G]$ est un champ algébrique d'après
Critères de représentabilité, Théorème
\ref{criteria-theorem-flat-groupoid-gives-algebraic-stack}.
La lissité de $[X/G]$ sur $Y$ résulte du fait que la lissité
descend par les recouvrements fppf :
Choisissons un morphisme lisse surjectif $U \to [X/G]$, où $U$ est un schéma.
La lissité de $[X/G]$ sur $Y$ équivaut à celle de $U$ sur $Y$.
Notons que $U \times_{[X/G]} X$ est lisse sur $X$, donc lisse
sur $Y$ (car un composé de morphismes lisses est lisse).
D'autre part, $U \times_{[X/G]} X \to U$ est localement de
présentation finie, plat et surjectif (car c'est
le changement de base de $X \to [X/G]$, qui possède ces propriétés,
par exemple d'après Critères de représentabilité, Lemme
\ref{criteria-lemma-flat-quotient-flat-presentation}).
Nous pouvons donc appliquer Descente sur les espaces,
Lemme \ref{spaces-descent-lemma-smooth-permanence}.
\end{proof}

\begin{lemma}
\label{lemma-gerbe-smooth}
Soit $\pi : \mathcal{X} \to \mathcal{Y}$ un morphisme de champs algébriques.
Si $\mathcal{X}$ est une gerbe sur $\mathcal{Y}$, alors $\pi$ est surjectif
et lisse.
\end{lemma}

\begin{proof}
Nous avons établi la surjectivité au Lemme \ref{lemma-gerbe-fppf}.
D'après le Lemme \ref{lemma-descent-smooth},
il suffit de démontrer le lemme après avoir remplacé $\pi$ par son changement de base
par un morphisme surjectif, plat et localement de présentation finie
$\mathcal{Y}' \to \mathcal{Y}$. D'après le
Lemme \ref{lemma-local-structure-gerbe},
nous pouvons supposer que $\mathcal{Y} = U$ est un espace algébrique et que
$\mathcal{X} = [U/G]$ au-dessus de $U$, avec $G \to U$ plat et
localement de présentation finie.
Le résultat découle alors du Lemme \ref{lemma-smooth-quotient-stack}.
\end{proof}




\section[Morphismes DM : localité étale-lisse]{Types de morphismes dont la propriété est de nature locale pour la topologie \'etale sur la source et pour la topologie lisse sur le but}
\label{section-etale-smooth-local-source-target}

\noindent
Soit une propriété de morphismes d'espaces algébriques qui est
{\it de nature locale pour la topologie \'etale sur la source et pour la topologie lisse sur le but}, voir
Descente sur les espaces,
Définition \ref{spaces-descent-definition-etale-smooth-local-source-target}.
Nous pouvons l'utiliser pour définir une propriété correspondante
des morphismes DM de champs algébriques, en imposant l'une ou l'autre
des conditions équivalentes du lemme ci-dessous.

\begin{lemma}
\label{lemma-etale-smooth-local-source-target}
Soit $\mathcal{P}$ une propriété de morphismes d'espaces algébriques
qui est de nature locale pour la topologie \'etale sur la source et pour la topologie lisse sur le but.
Soit $f : \mathcal{X} \to \mathcal{Y}$ un morphisme DM de champs algébriques.
Considérons des diagrammes commutatifs
$$
\xymatrix{
U \ar[d]_a \ar[r]_h & V \ar[d]^b \\
\mathcal{X} \ar[r]^f & \mathcal{Y}
}
$$
où $U$ et $V$ sont des espaces algébriques, $V \to \mathcal{Y}$ est lisse,
et $U \to \mathcal{X} \times_\mathcal{Y} V$ est \'etale.
Les conditions suivantes sont équivalentes :
\begin{enumerate}
\item pour tout diagramme comme ci-dessus, le morphisme $h$ possède la propriété $\mathcal{P}$ ;
\item pour un diagramme comme ci-dessus dans lequel $a : U \to \mathcal{X}$ est surjectif,
le morphisme $h$ possède la propriété $\mathcal{P}$.
\end{enumerate}
Si $\mathcal{X}$ et $\mathcal{Y}$ sont représentables par des espaces algébriques,
cela équivaut également au fait que $f$ (considéré comme morphisme d'espaces algébriques)
possède la propriété $\mathcal{P}$. Si $\mathcal{P}$ est en outre stable par
tout changement de base et de nature locale pour la topologie fppf sur la base, alors, pour les morphismes $f$
qui sont représentables par des espaces algébriques,
cela équivaut également au fait que $f$ possède la propriété $\mathcal{P}$ au sens
de
Propriétés des champs,
section \ref{stacks-properties-section-properties-morphisms}.
\end{lemma}

\begin{proof}
Montrons l'implication (1) $\Rightarrow$ (2). Choisissons un espace
algébrique $V$ et un morphisme lisse surjectif $V \to \mathcal{Y}$.
Comme $f$ est DM, il existe un schéma $U$ et un morphisme \'etale surjectif
$U \to V \times_\mathcal{Y} \mathcal{X}$, voir
Lemme \ref{lemma-DM}. Nous disposons ainsi des données géométriques requises en (2).
Notons que $U \to \mathcal{X}$ est également lisse et surjectif, puisqu'il est le
composé du changement de base
$\mathcal{X} \times_\mathcal{Y} V \to \mathcal{X}$ et du morphisme choisi
$U \to \mathcal{X} \times_\mathcal{Y} V$. Nous obtenons donc un
diagramme comme en (1). Par conséquent, si (1) est vérifiée, $h : U \to V$ possède la propriété
$\mathcal{P}$, ce qui signifie que (2) est vérifiée puisque $U \to \mathcal{X}$ est surjectif.

\medskip\noindent
Réciproquement, supposons que (2) soit vérifiée et soient $U, V, a, b, h$ comme en (2).
Soient ensuite $U', V', a', b', h'$ les données d'un diagramme quelconque comme en (1).
Considérons
$$
\xymatrix{
U \ar[d] \ar[r]_h & V \ar[d] \\
\mathcal{X} \ar[r]^f & \mathcal{Y}
}
\quad\quad
\xymatrix{
U' \ar[d] \ar[r]_{h'} & V' \ar[d] \\
\mathcal{X} \ar[r]^f & \mathcal{Y}
}
$$
Pour montrer que (2) implique (1), nous devons prouver que $h'$ possède $\mathcal{P}$.
À cette fin, considérons le diagramme commutatif
$$
\xymatrix{
U \ar[d]^h &
U \times_\mathcal{X} U' \ar[l] \ar[d]^{(h, h')} \ar[r] &
U' \ar[d]^{h'} \\
V &
V \times_\mathcal{Y} V' \ar[l] \ar[r] &
V'
}
$$
d'espaces algébriques. Notons que les flèches horizontales sont
lisses comme changements de base des morphismes lisses
$V \to \mathcal{Y}$, $V' \to \mathcal{Y}$, $U \to \mathcal{X}$ et
$U' \to \mathcal{X}$. Notons que les carrés
$$
\xymatrix{
U \ar[d] & U \times_\mathcal{X} U' \ar[l] \ar[d] &
U \times_\mathcal{X} U' \ar[d] \ar[r] & U' \ar[d] \\
V \times_\mathcal{Y} \mathcal{X} &
V \times_\mathcal{Y} U' \ar[l] &
U \times_\mathcal{Y} V' \ar[r] &
\mathcal{X} \times_\mathcal{Y} V'
}
$$
sont cartésiens ; les flèches verticales sont donc \'etales par nos hypothèses sur
$U', V', a', b', h'$ et $U, V, a, b, h$.
Puisque $\mathcal{P}$ est de nature locale pour la topologie lisse sur le but d'après
Descente sur les espaces, Lemme
\ref{spaces-descent-lemma-etale-smooth-local-source-target-implies}, partie (2),
nous voyons que le changement de base
$t : U \times_\mathcal{Y} V' \to V \times_\mathcal{Y} V'$ de $h$
possède $\mathcal{P}$. Puisque $\mathcal{P}$ est de nature locale pour la topologie \'etale sur la source d'après
Descente sur les espaces, Lemme
\ref{spaces-descent-lemma-etale-smooth-local-source-target-implies}, partie (1),
et que $s : U \times_\mathcal{X} U' \to U \times_\mathcal{Y} V'$ est \'etale,
nous voyons que le morphisme $(h, h') = t \circ s$ possède $\mathcal{P}$.
Considérons le diagramme
$$
\xymatrix{
U \times_\mathcal{X} U' \ar[r]_{(h, h')} \ar[d] &
V \times_\mathcal{Y} V' \ar[d] \\
U' \ar[r]^{h'} & V'
}
$$
La flèche verticale gauche est surjective, la flèche verticale droite est lisse et
le morphisme induit
$$
U \times_\mathcal{X} U'
\longrightarrow
U' \times_{V'} (V \times_\mathcal{Y} V') = V \times_\mathcal{Y} U'
$$
est \'etale, comme on l'a vu ci-dessus. Par conséquent, d'après
Descente sur les espaces, Définition
\ref{spaces-descent-definition-etale-smooth-local-source-target}, partie (3),
nous concluons que $h'$ possède $\mathcal{P}$. Ceci achève la démonstration de
l'équivalence de (1) et (2).

\medskip\noindent
Si $\mathcal{X}$ et $\mathcal{Y}$ sont représentables, alors
Descente sur les espaces,
Lemme \ref{spaces-descent-lemma-etale-smooth-local-source-target-characterize},
s'applique et montre que les conditions (1) et (2) équivalent au fait que $f$ possède
$\mathcal{P}$.

\medskip\noindent
Enfin, supposons $f$ représentable, soient $U, V, a, b, h$
comme dans la partie (2) du lemme, et supposons que $\mathcal{P}$ est stable par
changement de base arbitraire. Nous devons montrer que pour tout schéma
$Z$ et tout morphisme $Z \to \mathcal{Y}$, le changement de base
$Z \times_\mathcal{Y} \mathcal{X} \to Z$
possède la propriété $\mathcal{P}$. Considérons le diagramme
$$
\xymatrix{
Z \times_\mathcal{Y} U \ar[d] \ar[r] &
Z \times_\mathcal{Y} V \ar[d] \\
Z \times_\mathcal{Y} \mathcal{X} \ar[r] &
Z
}
$$
Notons que la flèche horizontale supérieure est un changement de base de $h$ et
possède donc la propriété $\mathcal{P}$. La flèche verticale gauche est
surjective, et le morphisme induit
$$
Z \times_\mathcal{Y} U \longrightarrow
(Z \times_\mathcal{Y} \mathcal{X}) \times_Z (Z \times_\mathcal{Y} V)
$$
est \'etale, tandis que la flèche verticale droite est lisse. Ainsi,
Descente sur les espaces,
Lemme \ref{spaces-descent-lemma-etale-smooth-local-source-target-characterize}
s'applique et montre que $Z \times_\mathcal{Y} \mathcal{X} \to Z$
possède la propriété $\mathcal{P}$.
\end{proof}

\begin{definition}
\label{definition-etale-smooth-P}
Soit $\mathcal{P}$ une propriété de morphismes d'espaces algébriques
qui est de nature locale pour la topologie \'etale sur la source et pour la topologie lisse sur le but.
On dit qu'un morphisme DM $f : \mathcal{X} \to \mathcal{Y}$ de champs algébriques
{\it possède la propriété $\mathcal{P}$} si les conditions équivalentes du
Lemme \ref{lemma-etale-smooth-local-source-target}
sont vérifiées.
\end{definition}

\begin{remark}
\label{remark-etale-smooth-composition}
Soit $\mathcal{P}$ une propriété de morphismes d'espaces algébriques
qui est de nature locale pour la topologie \'etale sur la source et pour la topologie lisse sur le but, et
stable par composition. Alors la propriété des morphismes DM de champs algébriques
définie dans la Définition \ref{definition-etale-smooth-P}
est stable par composition. En effet, soient $f : \mathcal{X} \to \mathcal{Y}$
et $g : \mathcal{Y} \to \mathcal{Z}$ des morphismes DM de champs algébriques
possédant la propriété $\mathcal{P}$. D'après le Lemme \ref{lemma-composition-separated},
le composé $g \circ f$ est DM. Choisissons un espace algébrique $W$ et un
morphisme lisse surjectif $W \to \mathcal{Z}$. Choisissons un espace algébrique
$V$ et un morphisme \'etale surjectif $V \to \mathcal{Y} \times_\mathcal{Z} W$
(Lemme \ref{lemma-DM}).
Choisissons un espace algébrique $U$ et un morphisme \'etale surjectif
$U \to \mathcal{X} \times_\mathcal{Y} V$. Alors les morphismes
$V \to W$ et $U \to V$ possèdent la propriété $\mathcal{P}$ par définition.
Par suite, $U \to W$ possède la propriété $\mathcal{P}$, puisque nous avons supposé que
$\mathcal{P}$ est stable par composition. Ainsi, par définition encore,
nous voyons que $g \circ f : \mathcal{X} \to \mathcal{Z}$ possède la
propriété $\mathcal{P}$.
\end{remark}

\begin{remark}
\label{remark-etale-smooth-base-change}
Soit $\mathcal{P}$ une propriété de morphismes d'espaces algébriques
qui est de nature locale pour la topologie \'etale sur la source et pour la topologie lisse sur le but, et
stable par changement de base. Alors la propriété des
morphismes DM de champs algébriques définie dans la
Définition \ref{definition-etale-smooth-P}
est stable par changement de base arbitraire.
En effet, soit $f : \mathcal{X} \to \mathcal{Y}$
un morphisme DM de champs algébriques,
soit $g : \mathcal{Y}' \to \mathcal{Y}$ un morphisme de champs algébriques,
et supposons que $f$ possède la propriété $\mathcal{P}$.
Alors le changement de base
$\mathcal{Y}' \times_\mathcal{Y} \mathcal{X} \to \mathcal{Y}'$
est un morphisme DM d'après le Lemme \ref{lemma-base-change-separated}.
Choisissons un espace algébrique $V$
et un morphisme lisse surjectif $V \to \mathcal{Y}$. Choisissons un espace
algébrique $U$ et un morphisme \'etale surjectif
$U \to \mathcal{X} \times_\mathcal{Y} V$ (Lemme \ref{lemma-DM}).
Enfin, choisissons un espace algébrique
$V'$ et un morphisme lisse surjectif
$V' \to \mathcal{Y}' \times_\mathcal{Y} V$. Alors le morphisme
$U \to V$ possède la propriété $\mathcal{P}$ par définition.
Par suite, $V' \times_V U \to V'$ possède la propriété $\mathcal{P}$, puisque nous avons supposé que
$\mathcal{P}$ est stable par changement de base. Considérons le diagramme
$$
\xymatrix{
V' \times_V U \ar[r] \ar[d] &
\mathcal{Y}' \times_\mathcal{Y} \mathcal{X} \ar[r] \ar[d] &
\mathcal{X} \ar[d] \\
V' \ar[r] & \mathcal{Y}' \ar[r] & \mathcal{Y}
}
$$
nous voyons que la flèche horizontale supérieure de gauche est surjective et que
$$
V' \times_V U \to V' \times_{\mathcal{Y}'}
(\mathcal{Y}' \times_\mathcal{Y} \mathcal{X}) =
V' \times_V (\mathcal{X} \times_\mathcal{Y} V)
$$
est \'etale comme changement de base de $U \to \mathcal{X} \times_\mathcal{Y} V$ ;
d'où, par définition, la projection
$\mathcal{Y}' \times_\mathcal{Y} \mathcal{X} \to \mathcal{Y}'$ possède la
propriété $\mathcal{P}$.
\end{remark}

\begin{remark}
\label{remark-etale-smooth-implication}
Soient $\mathcal{P}, \mathcal{P}'$ des propriétés de morphismes d'espaces algébriques
qui sont de nature locale pour la topologie \'etale sur la source et pour la topologie lisse sur le but.
Supposons que l'on ait $\mathcal{P} \Rightarrow \mathcal{P}'$ pour les morphismes
d'espaces algébriques. Alors on a également $\mathcal{P} \Rightarrow \mathcal{P}'$
pour les propriétés des morphismes de champs algébriques définies dans la
Définition \ref{definition-etale-smooth-P}
à partir de $\mathcal{P}$ et de $\mathcal{P}'$. Cela résulte immédiatement de la définition.
\end{remark}



\section{Morphismes \'etales}
\label{section-etale}

\noindent
Un morphisme \'etale de champs alg\'ebriques ne saurait \^etre d\'efini comme un
morphisme lisse de dimension relative $0$. En effet, le morphisme
$$
[\mathbf{A}^1_k/\mathbf{G}_{m, k}] \longrightarrow \Spec(k)
$$
est lisse de dimension relative $0$, quelle que soit l'action choisie
du sch\'ema en groupes $\mathbf{G}_{m, k}$ sur $\mathbf{A}^1_k$.
Cela ne correspond pas \`a l'id\'ee usuelle selon laquelle les morphismes \'etales
devraient identifier les espaces tangents. L'exemple ci-dessus n'est pas quasi-fini,
mais le morphisme
$$
\mathcal{X} = [\Spec(k)/\mu_{p, k}] \longrightarrow \Spec(k)
$$
est lisse et quasi-fini (section \ref{section-locally-quasi-finite}).
Toutefois, si la caract\'eristique de $k$ est $p > 0$, on constate
que le morphisme repr\'esentable $\Spec(k) \to \mathcal{X}$
n'est pas \'etale, puisque son changement de base
$\mu_{p, k} = \Spec(k) \times_\mathcal{X} \Spec(k) \to \Spec(k)$
est un morphisme d'un sch\'ema non r\'eduit vers le spectre d'un corps.
Ainsi, si l'on d\'efinissait un morphisme \'etale comme un morphisme lisse et localement quasi-fini,
alors l'analogue du lemme de Morphismes d'espaces
\ref{spaces-morphisms-lemma-etale-permanence} serait faux.

\medskip\noindent
Notre d\'emarche consistera plut\^ot \`a partir des exigences suivantes :
``\^etre \'etale'' devra \^etre une propri\'et\'e stable par changement de base et, si
$\mathcal{X} \to X$ est un morphisme \'etale d'un champ alg\'ebrique
vers un sch\'ema, alors $\mathcal{X}$ devra \^etre de Deligne--Mumford.
Autrement dit, nous exigerons qu'un morphisme \'etale soit DM et
nous utiliserons les r\'esultats de la
section \ref{section-etale-smooth-local-source-target} pour d\'efinir
les morphismes \'etales de champs alg\'ebriques.

\medskip\noindent
Dans le lemme \ref{lemma-characterize-etale}, nous caract\'eriserons
les morphismes \'etales de champs alg\'ebriques comme ceux qui sont
(a) localement de pr\'esentation finie,
(b) plats et (c) \`a diagonale \'etale.

\medskip\noindent
La propri\'et\'e ``\^etre \'etale'' des morphismes d'espaces alg\'ebriques
est de nature locale pour la topologie \'etale sur la source et pour la topologie lisse
sur le but; voir Descente sur les espaces, remarque
\ref{spaces-descent-remark-list-etale-smooth-local-source-target}.
Elle est aussi stable par changement de base et de nature locale pour la topologie fpqc sur le but; voir
Morphismes d'espaces,
lemme \ref{spaces-morphisms-lemma-base-change-etale}
et
Descente sur les espaces,
lemme \ref{spaces-descent-lemma-descending-property-etale}.
Par cons\'equent, d'apr\`es
le lemme \ref{lemma-etale-smooth-local-source-target}
ci-dessus, nous pouvons d\'efinir comme suit ce que signifie, pour un morphisme
de champs alg\'ebriques, \^etre \'etale; cette d\'efinition co\"incide avec la notion
d\'ej\`a d\'efinie dans Propri\'et\'es des champs, section
\ref{stacks-properties-section-properties-morphisms},
lorsque le morphisme est repr\'esentable par des espaces alg\'ebriques,
car un tel morphisme est automatiquement DM d'apr\`es
le lemme \ref{lemma-trivial-implications}.

\begin{definition}
\label{definition-etale}
Soit $f : \mathcal{X} \to \mathcal{Y}$ un morphisme de champs alg\'ebriques.
On dit que $f$ est {\it \'etale} si $f$ est DM et si les conditions \'equivalentes du
lemme \ref{lemma-etale-smooth-local-source-target}
sont v\'erifi\'ees avec $\mathcal{P} = \etale$.
\end{definition}

\noindent
Nous utiliserons sans autre mention le fait que cette d\'efinition co\"incide avec la notion
d\'ej\`a existante de morphisme \'etale lorsque $f$ est repr\'esentable
par des espaces alg\'ebriques ou lorsque $\mathcal{X}$ et $\mathcal{Y}$ sont
repr\'esentables par des espaces alg\'ebriques.

\begin{lemma}
\label{lemma-composition-etale}
Le compos\'e de deux morphismes \'etales est \'etale.
\end{lemma}

\begin{proof}
Cela résulte de la remarque \ref{remark-etale-smooth-composition}
et du lemme \ref{spaces-morphisms-lemma-composition-etale} de
Morphismes d'espaces.
\end{proof}

\begin{lemma}
\label{lemma-base-change-etale}
Tout changement de base d'un morphisme \'etale est \'etale.
\end{lemma}

\begin{proof}
Cela résulte de la remarque \ref{remark-etale-smooth-base-change}
et du lemme \ref{spaces-morphisms-lemma-base-change-etale} de
Morphismes d'espaces.
\end{proof}

\begin{lemma}
\label{lemma-open-immersion-etale}
Toute immersion ouverte est \'etale.
\end{lemma}

\begin{proof}
Soit $j : \mathcal{U} \to \mathcal{X}$ une immersion ouverte de
champs alg\'ebriques. Puisque $j$ est repr\'esentable, il est DM d'apr\`es
le lemme \ref{lemma-trivial-implications}. D'autre part,
si $X \to \mathcal{X}$ est un morphisme lisse et surjectif,
o\`u $X$ est un sch\'ema, alors $U = \mathcal{U} \times_\mathcal{X} X$
est un sous-sch\'ema ouvert de $X$. Par cons\'equent, $U \to X$ est \'etale
(Morphismes, lemme \ref{morphisms-lemma-open-immersion-etale})
et la d\'efinition montre que $j$ est \'etale.
\end{proof}

\begin{lemma}
\label{lemma-etale}
Soit $f : \mathcal{X} \to \mathcal{Y}$ un morphisme de champs alg\'ebriques.
Les conditions suivantes sont \'equivalentes :
\begin{enumerate}
\item $f$ est \'etale,
\item $f$ est DM et, pour tout morphisme $V \to \mathcal{Y}$
o\`u $V$ est un espace alg\'ebrique, et tout morphisme \'etale
$U \to V \times_\mathcal{Y} \mathcal{X}$ o\`u $U$ est un espace alg\'ebrique,
le morphisme $U \to V$ est \'etale,
\item il existe un morphisme surjectif, localement de pr\'esentation finie et plat
$W \to \mathcal{Y}$, o\`u $W$ est un espace alg\'ebrique, et un
morphisme \'etale surjectif $T \to W \times_\mathcal{Y} \mathcal{X}$
o\`u $T$ est un espace alg\'ebrique, tel que le morphisme $T \to W$ soit \'etale.
\end{enumerate}
\end{lemma}

\begin{proof}
Supposons (1). Alors $f$ est DM et, comme la propri\'et\'e d'\^etre \'etale est stable
par changement de base, la condition (2) est v\'erifi\'ee.

\medskip\noindent
Supposons (2). Choisissons un sch\'ema $V$ et un morphisme \'etale surjectif
$V \to \mathcal{Y}$. Choisissons un sch\'ema $U$ et un morphisme \'etale surjectif
$U \to V \times_\mathcal{Y} \mathcal{X}$ (lemme \ref{lemma-DM}).
La condition (3) est donc v\'erifi\'ee.

\medskip\noindent
Supposons que $W \to \mathcal{Y}$ et $T \to W \times_\mathcal{Y} \mathcal{X}$
soient comme dans (3). V\'erifions d'abord que $f$ est DM. En effet, il suffit de v\'erifier
que $W \times_\mathcal{Y} \mathcal{X} \to W$ est DM; voir
le lemme \ref{lemma-check-separated-covering}.
D'apr\`es le lemme \ref{lemma-compose-after-separated},
il suffit de v\'erifier que $W \times_\mathcal{Y} \mathcal{X}$ est de Deligne--Mumford.
Cela r\'esulte de l'existence de $T \to W \times_\mathcal{Y} \mathcal{X}$
par le sens facile du th\'eor\`eme \ref{theorem-DM}.

\medskip\noindent
Supposons que $f$ soit DM et que $W \to \mathcal{Y}$ et
$T \to W \times_\mathcal{Y} \mathcal{X}$ soient comme dans (3).
Soit $V$ un espace alg\'ebrique, soit $V \to \mathcal{Y}$ un morphisme lisse surjectif,
soit $U$ un espace alg\'ebrique, et soit
$U \to V \times_\mathcal{Y} \mathcal{X}$ un morphisme \'etale surjectif
(lemme \ref{lemma-DM}). Il faut v\'erifier que $U \to V$ est \'etale.
Il suffit de montrer que $U \times_\mathcal{Y} W \to V \times_\mathcal{Y} W$
est \'etale d'apr\`es Descente sur les espaces, lemme
\ref{spaces-descent-lemma-descending-property-etale}.
Nous pouvons remplacer $\mathcal{X}, \mathcal{Y}, W, T, U, V$ respectivement par
$\mathcal{X} \times_\mathcal{Y} W, W, W, T, U \times_\mathcal{Y} W,
V \times_\mathcal{Y} W$ (un petit d\'etail est omis).
Nous pouvons donc supposer que $Y = \mathcal{Y}$ est un espace alg\'ebrique, qu'il existe
un espace alg\'ebrique $T$ et un morphisme \'etale surjectif
$T \to \mathcal{X}$ tel que $T \to Y$ soit \'etale, et que $U$ et $V$
soient comme ci-dessus. Dans ce cas, on sait que
$$
U \to V\text{ est \'etale}
\Leftrightarrow
\mathcal{X} \to Y\text{ est \'etale}
\Leftrightarrow
T \to Y\text{ est \'etale}
$$
par l'\'equivalence des propri\'et\'es (1) et (2) du
lemme \ref{lemma-etale-smooth-local-source-target}
et la d\'efinition \ref{definition-etale}.
Cela ach\`eve la d\'emonstration.
\end{proof}

\begin{lemma}
\label{lemma-etale-permanence}
Soient $\mathcal{X}, \mathcal{Y}$ des champs alg\'ebriques \'etales sur
un champ alg\'ebrique $\mathcal{Z}$. Tout morphisme
$\mathcal{X} \to \mathcal{Y}$ au-dessus de $\mathcal{Z}$ est \'etale.
\end{lemma}

\begin{proof}
Le morphisme $\mathcal{X} \to \mathcal{Y}$ est DM d'apr\`es
le lemme \ref{lemma-compose-after-separated}.
Soit $W \to \mathcal{Z}$ un morphisme lisse surjectif
dont la source est un espace alg\'ebrique. Soit
$V \to \mathcal{Y} \times_\mathcal{Z} W$ un morphisme \'etale surjectif
dont la source est un espace alg\'ebrique
(lemme \ref{lemma-DM}).
Soit $U \to \mathcal{X} \times_\mathcal{Y} V$ un morphisme \'etale surjectif
dont la source est un espace alg\'ebrique
(lemme \ref{lemma-DM}).
Alors
$$
U \longrightarrow \mathcal{X} \times_\mathcal{Z} W
$$
est \'etale surjectif en tant que compos\'e de
$U \to \mathcal{X} \times_\mathcal{Y} V$
et du changement de base de $V \to \mathcal{Y} \times_\mathcal{Z} W$
par $\mathcal{X} \times_\mathcal{Z} W \to \mathcal{Y} \times_\mathcal{Z} W$.
Il suffit donc de montrer que $U \to V$ est \'etale.
Or $U \to W$ et $V \to W$ sont \'etales, puisque
$\mathcal{X} \to \mathcal{Z}$ et $\mathcal{Y} \to \mathcal{Z}$
le sont; la conclusion r\'esulte donc de la version du lemme pour les
espaces alg\'ebriques, \`a savoir
Morphismes d'espaces, lemme \ref{spaces-morphisms-lemma-etale-permanence}.
\end{proof}








\section{Morphismes non ramifi\'es}
\label{section-unramified}

\noindent
Pour justifier notre choix de la d\'efinition des morphismes
non ramifi\'es, nous renvoyons le lecteur \`a la discussion
de la section sur les morphismes \'etales, section \ref{section-etale}.

\medskip\noindent
Dans le Lemme \ref{lemma-characterize-unramified}, nous caract\'eriserons
les morphismes non ramifi\'es de champs alg\'ebriques comme les morphismes
qui sont localement de type fini et dont la diagonale est \'etale.

\medskip\noindent
La propri\'et\'e ``non ramifi\'e'' des morphismes d'espaces alg\'ebriques est de
nature locale pour la topologie \'etale sur la source et pour la topologie lisse sur le but, voir
Descente sur les espaces, Remarque
\ref{spaces-descent-remark-list-etale-smooth-local-source-target}.
Elle est aussi stable par changement de base et locale sur la cible pour la
topologie fpqc, voir Morphismes d'espaces,
Lemme \ref{spaces-morphisms-lemma-base-change-unramified}
et
Descente sur les espaces,
Lemme \ref{spaces-descent-lemma-descending-property-unramified}.
Par cons\'equent, d'apr\`es le
Lemme \ref{lemma-etale-smooth-local-source-target}
ci-dessus, nous pouvons d\'efinir comme suit ce que signifie, pour un morphisme
de champs alg\'ebriques, d'\^etre non ramifi\'e; cette d\'efinition co\"incide avec
la notion d\'ej\`a d\'efinie dans Propri\'et\'es des champs, section
\ref{stacks-properties-section-properties-morphisms},
lorsque le morphisme est repr\'esentable par des espaces alg\'ebriques,
car un tel morphisme est automatiquement DM d'apr\`es le
Lemme \ref{lemma-trivial-implications}.

\begin{definition}
\label{definition-unramified}
Soit $f : \mathcal{X} \to \mathcal{Y}$ un morphisme de champs alg\'ebriques.
On dit que $f$ est {\it non ramifi\'e} si $f$ est DM et si les conditions
\'equivalentes du Lemme \ref{lemma-etale-smooth-local-source-target}
sont satisfaites pour $\mathcal{P} =$``non ramifi\'e''.
\end{definition}

\noindent
Nous utiliserons sans autre mention le fait que cette d\'efinition co\"incide
avec la notion d\'ej\`a existante de morphisme non ramifi\'e lorsque $f$ est
repr\'esentable par des espaces alg\'ebriques, ou lorsque $\mathcal{X}$ et
$\mathcal{Y}$ sont repr\'esentables par des espaces alg\'ebriques.

\begin{lemma}
\label{lemma-composition-unramified}
Le compos\'e de morphismes non ramifi\'es est non ramifi\'e.
\end{lemma}

\begin{proof}
Cela r\'esulte de la Remarque \ref{remark-etale-smooth-composition}
et de
Morphismes d'espaces, Lemme \ref{spaces-morphisms-lemma-composition-unramified}.
\end{proof}

\begin{lemma}
\label{lemma-base-change-unramified}
Tout morphisme d\'eduit par changement de base d'un morphisme non ramifi\'e est non ramifi\'e.
\end{lemma}

\begin{proof}
Cela r\'esulte de
la Remarque \ref{remark-etale-smooth-base-change}
et de
Morphismes d'espaces, Lemme \ref{spaces-morphisms-lemma-base-change-unramified}.
\end{proof}

\begin{lemma}
\label{lemma-etale-unramified}
Un morphisme \'etale est non ramifi\'e.
\end{lemma}

\begin{proof}
Cela r\'esulte de la Remarque \ref{remark-etale-smooth-implication}
et de Morphismes d'espaces, Lemme
\ref{spaces-morphisms-lemma-etale-unramified}.
\end{proof}

\begin{lemma}
\label{lemma-immersion-unramified}
Une immersion est non ramifi\'ee.
\end{lemma}

\begin{proof}
Soit $j : \mathcal{Z} \to \mathcal{X}$ une immersion de
champs alg\'ebriques. Comme $j$ est repr\'esentable, il est DM d'apr\`es le
Lemme \ref{lemma-trivial-implications}. D'autre part,
si $X \to \mathcal{X}$ est un morphisme lisse et surjectif,
o\`u $X$ est un sch\'ema, alors $Z = \mathcal{Z} \times_\mathcal{X} X$
est un sous-sch\'ema localement ferm\'e de $X$. Par cons\'equent, $Z \to X$
est non ramifi\'e (Morphismes, Lemmes \ref{morphisms-lemma-open-immersion-unramified} et
\ref{morphisms-lemma-closed-immersion-unramified}),
et la d\'efinition montre que $j$ est non ramifi\'e.
\end{proof}

\begin{lemma}
\label{lemma-unramified}
Soit $f : \mathcal{X} \to \mathcal{Y}$ un morphisme de champs alg\'ebriques.
Les conditions suivantes sont \'equivalentes :
\begin{enumerate}
\item $f$ est non ramifi\'e,
\item $f$ est DM et, pour tout morphisme $V \to \mathcal{Y}$
o\`u $V$ est un espace alg\'ebrique, et tout morphisme \'etale
$U \to V \times_\mathcal{Y} \mathcal{X}$ o\`u $U$ est un espace alg\'ebrique,
le morphisme $U \to V$ est non ramifi\'e,
\item il existe un morphisme $W \to \mathcal{Y}$ surjectif,
localement de pr\'esentation finie et plat, o\`u $W$ est un espace alg\'ebrique,
et un morphisme $T \to W \times_\mathcal{Y} \mathcal{X}$ surjectif et \'etale,
o\`u $T$ est un espace alg\'ebrique, tel que le morphisme $T \to W$ soit non ramifi\'e.
\end{enumerate}
\end{lemma}

\begin{proof}
Supposons (1). Alors $f$ est DM et, puisque la propri\'et\'e d'\^etre non
ramifi\'e est stable par changement de base, la condition (2) est satisfaite.

\medskip\noindent
Supposons (2). Choisissons un sch\'ema $V$ et un morphisme \'etale surjectif
$V \to \mathcal{Y}$. Choisissons un sch\'ema $U$ et un morphisme \'etale surjectif
$U \to V \times_\mathcal{Y} \mathcal{X}$ (Lemme \ref{lemma-DM}).
La condition (3) est donc satisfaite.

\medskip\noindent
Supposons que $W \to \mathcal{Y}$ et $T \to W \times_\mathcal{Y} \mathcal{X}$
soient comme en (3). Montrons d'abord que $f$ est DM. Il suffit en effet de
v\'erifier que $W \times_\mathcal{Y} \mathcal{X} \to W$ est DM, voir le
Lemme \ref{lemma-check-separated-covering}.
D'apr\`es le Lemme \ref{lemma-compose-after-separated},
il suffit de v\'erifier que $W \times_\mathcal{Y} \mathcal{X}$ est un champ
de Deligne--Mumford. Cela r\'esulte de l'existence de
$T \to W \times_\mathcal{Y} \mathcal{X}$ par le sens facile du
Th\'eor\`eme \ref{theorem-DM}.

\medskip\noindent
Supposons que $f$ soit DM et que $W \to \mathcal{Y}$ et
$T \to W \times_\mathcal{Y} \mathcal{X}$ soient comme en (3).
Soit $V$ un espace alg\'ebrique, soit $V \to \mathcal{Y}$ un morphisme
lisse surjectif, soit $U$ un espace alg\'ebrique, et soit
$U \to V \times_\mathcal{Y} \mathcal{X}$ un morphisme surjectif et \'etale
(Lemme \ref{lemma-DM}). Nous devons v\'erifier que $U \to V$ est non ramifi\'e.
Il suffit de montrer que $U \times_\mathcal{Y} W \to V \times_\mathcal{Y} W$
est non ramifi\'e d'apr\`es Descente sur les espaces, Lemme
\ref{spaces-descent-lemma-descending-property-unramified}.
Nous pouvons remplacer $\mathcal{X}, \mathcal{Y}, W, T, U, V$ par
$\mathcal{X} \times_\mathcal{Y} W, W, W, T, U \times_\mathcal{Y} W,
V \times_\mathcal{Y} W$ (un d\'etail mineur est omis).
Nous pouvons donc supposer que $Y = \mathcal{Y}$ est un espace alg\'ebrique,
qu'il existe un espace alg\'ebrique $T$ et un morphisme \'etale surjectif
$T \to \mathcal{X}$ tel que $T \to Y$ soit non ramifi\'e, et que $U$ et $V$
soient comme ci-dessus. Dans ce cas, nous savons que
$$
U \to V\text{ est non ramifi\'e}
\Leftrightarrow
\mathcal{X} \to Y\text{ est non ramifi\'e}
\Leftrightarrow
T \to Y\text{ est non ramifi\'e}
$$
par l'\'equivalence des propri\'et\'es (1) et (2) du
Lemme \ref{lemma-etale-smooth-local-source-target}
et par la D\'efinition \ref{definition-unramified}.
Cela ach\`eve la d\'emonstration.
\end{proof}

\begin{lemma}
\label{lemma-unramified-quasi-finite}
Un morphisme non ramifi\'e de champs alg\'ebriques est localement quasi-fini.
\end{lemma}

\begin{proof}
Cela r\'esulte du
Lemme \ref{lemma-unramified} (qui caract\'erise les morphismes non ramifi\'es),
du Lemme \ref{lemma-characterize-locally-quasi-finite} (qui caract\'erise
les morphismes localement quasi-finis), et de
Morphismes d'espaces, Lemme \ref{spaces-morphisms-lemma-unramified-quasi-finite}
(qui est l'\'enonc\'e correspondant pour les espaces alg\'ebriques).
\end{proof}

\begin{lemma}
\label{lemma-permanence-unramified}
Soient $\mathcal{X} \to \mathcal{Y} \to \mathcal{Z}$ des
morphismes de champs alg\'ebriques.
Si $\mathcal{X} \to \mathcal{Z}$ est non ramifi\'e et si
$\mathcal{Y} \to \mathcal{Z}$ est DM, alors
$\mathcal{X} \to \mathcal{Y}$ est non ramifi\'e.
\end{lemma}

\begin{proof}
Supposons que $\mathcal{X} \to \mathcal{Z}$ soit non ramifi\'e. D'apr\`es le
Lemme \ref{lemma-compose-after-separated}, le morphisme
$\mathcal{X} \to \mathcal{Y}$ est DM. Choisissons un diagramme commutatif
$$
\xymatrix{
U \ar[d] \ar[r] & V \ar[d] \ar[r] & W \ar[d] \\
\mathcal{X} \ar[r] & \mathcal{Y} \ar[r] & \mathcal{Z}
}
$$
o\`u $U, V, W$ sont des espaces alg\'ebriques,
o\`u $W \to \mathcal{Z}$ est lisse et surjectif,
$V \to \mathcal{Y} \times_\mathcal{Z} W$ est \'etale et surjectif, et
$U \to \mathcal{X} \times_\mathcal{Y} V$ est \'etale et surjectif
(voir le Lemme \ref{lemma-DM}). Alors
$U \to \mathcal{X} \times_\mathcal{Z} W$ est lui aussi surjectif et \'etale.
Nous savons donc que $U \to W$ est non ramifi\'e, et il nous faut montrer que
$U \to V$ est non ramifi\'e. Cela r\'esulte de
Morphismes d'espaces, Lemme \ref{spaces-morphisms-lemma-permanence-unramified}.
\end{proof}

\begin{lemma}
\label{lemma-characterize-unramified}
Soit $f : \mathcal{X} \to \mathcal{Y}$ un morphisme de champs alg\'ebriques.
Les conditions suivantes sont \'equivalentes :
\begin{enumerate}
\item $f$ est non ramifi\'e, et
\item $f$ est localement de type fini et sa diagonale est \'etale.
\end{enumerate}
\end{lemma}

\begin{proof}
Supposons que $f$ soit non ramifi\'e. Alors $f$ est DM; nous pouvons donc choisir
des espaces alg\'ebriques $U$, $V$, un morphisme lisse surjectif
$V \to \mathcal{Y}$ et un morphisme \'etale surjectif
$U \to \mathcal{X} \times_\mathcal{Y} V$ (Lemme \ref{lemma-DM}).
Comme $f$ est non ramifi\'e, le morphisme induit $U \to V$ est non ramifi\'e.
Ainsi, $U \to V$ est localement de type fini
(Morphismes d'espaces, Lemme
\ref{spaces-morphisms-lemma-unramified-locally-finite-type}),
et nous en concluons que $f$ est localement de type fini. La diagonale
$\Delta : \mathcal{X} \to \mathcal{X} \times_\mathcal{Y} \mathcal{X}$
est un morphisme de champs alg\'ebriques au-dessus de $\mathcal{Y}$.
Le changement de base de $\Delta$ par le morphisme lisse surjectif
$V \to \mathcal{Y}$ est la diagonale du changement de base de
$f$, c'est-\`a-dire de $\mathcal{X}_V = \mathcal{X} \times_\mathcal{Y} V \to V$.
Autrement dit, le diagramme
$$
\xymatrix{
\mathcal{X}_V \ar[r] \ar[d] &  \mathcal{X}_V \times_V \mathcal{X}_V \ar[d] \\
\mathcal{X} \ar[r] & \mathcal{X} \times_\mathcal{Y} \mathcal{X}
}
$$
est cart\'esien. Comme la fl\`eche verticale de droite est lisse et surjective,
il suffit de montrer que la fl\`eche horizontale sup\'erieure est \'etale, d'apr\`es
Propri\'et\'es des champs, Lemme
\ref{stacks-properties-lemma-check-property-weak-covering}.
Consid\'erons le diagramme commutatif
$$
\xymatrix{
U \ar[d] \ar[r] & U \times_V U \ar[d] \\
\mathcal{X}_V \ar[r] & \mathcal{X}_V \times_V \mathcal{X}_V
}
$$
Toutes les fl\`eches sont repr\'esentables par des espaces alg\'ebriques,
les fl\`eches verticales sont \'etales, celle de gauche est surjective, et
la fl\`eche horizontale sup\'erieure est une immersion ouverte d'apr\`es
Morphismes d'espaces, Lemme
\ref{spaces-morphisms-lemma-diagonal-unramified-morphism}.
Cela donne le r\'esultat voulu : tout d'abord,
$U \to \mathcal{X}_V \times_V \mathcal{X}_V$ est \'etale
comme compos\'e de morphismes \'etales; nous pouvons ensuite appliquer
Propri\'et\'es des champs, Lemme
\ref{stacks-properties-lemma-check-property-after-precomposing}
pour voir que $\mathcal{X}_V \to \mathcal{X}_V \times_V \mathcal{X}_V$
est \'etale, puisque la propri\'et\'e d'\^etre \'etale (pour les morphismes
d'espaces alg\'ebriques) est locale sur la source pour la topologie \'etale
(Descente sur les espaces, Lemme \ref{spaces-descent-lemma-etale-etale-local-source}).

\medskip\noindent
Supposons que $f$ soit localement de type fini et que sa diagonale soit \'etale.
Alors $f$ est DM par d\'efinition (puisque les morphismes \'etales d'espaces
alg\'ebriques sont non ramifi\'es). Comme ci-dessus, cela signifie que nous
pouvons choisir des espaces alg\'ebriques $U$, $V$, un morphisme lisse surjectif
$V \to \mathcal{Y}$ et un morphisme \'etale surjectif
$U \to \mathcal{X} \times_\mathcal{Y} V$ (Lemme \ref{lemma-DM}).
Pour achever la d\'emonstration, nous devons montrer que $U \to V$ est non ramifi\'e.
Nous savons d\'ej\`a que $U \to V$ est localement de type fini.
En raisonnant comme ci-dessus, nous obtenons un diagramme commutatif
$$
\xymatrix{
U \ar[d] \ar[r] & U \times_V U \ar[d] \\
\mathcal{X}_V \ar[r] & \mathcal{X}_V \times_V \mathcal{X}_V
}
$$
o\`u toutes les fl\`eches sont repr\'esentables par des espaces alg\'ebriques,
les fl\`eches verticales sont \'etales, et la fl\`eche horizontale inf\'erieure
est \'etale comme changement de base de $\Delta$.
Il s'ensuit que $U \to U \times_V U$ est \'etale,
par exemple d'apr\`es le Lemme \ref{lemma-etale-permanence}\footnote{On peut
aussi le d\'eduire directement, et assez facilement, de
Morphismes d'espaces, Lemme \ref{spaces-morphisms-lemma-etale-permanence}.}.
Ainsi, $U \to U \times_V U$ est un monomorphisme \'etale,
donc une immersion ouverte (Morphismes d'espaces, Lemme
\ref{spaces-morphisms-lemma-etale-universally-injective-open}).
Le morphisme $U \to V$ est alors non ramifi\'e d'apr\`es
Morphismes d'espaces, Lemme
\ref{spaces-morphisms-lemma-diagonal-unramified-morphism}.
\end{proof}

\begin{lemma}
\label{lemma-characterize-etale}
Soit $f : \mathcal{X} \to \mathcal{Y}$ un morphisme de champs alg\'ebriques.
Les conditions suivantes sont \'equivalentes :
\begin{enumerate}
\item $f$ est \'etale, et
\item $f$ est localement de pr\'esentation finie, plat et non ramifi\'e,
\item $f$ est localement de pr\'esentation finie, plat, et sa diagonale
est \'etale.
\end{enumerate}
\end{lemma}

\begin{proof}
L'\'equivalence de (2) et (3) r\'esulte imm\'ediatement du
Lemme \ref{lemma-characterize-unramified}. Ainsi, dans chaque cas,
le morphisme $f$ est DM.
Nous pouvons donc choisir des espaces alg\'ebriques $U$, $V$,
un morphisme lisse surjectif
$V \to \mathcal{Y}$ et un morphisme \'etale surjectif
$U \to \mathcal{X} \times_\mathcal{Y} V$ (Lemme \ref{lemma-DM}).
Pour achever la d\'emonstration, nous devons montrer que
$U \to V$ est \'etale si et seulement s'il est localement de pr\'esentation
finie, plat et non ramifi\'e.
Cela r\'esulte de Morphismes d'espaces, Lemme
\ref{spaces-morphisms-lemma-unramified-flat-lfp-etale}
(et, plus imm\'ediatement, de Morphismes d'espaces, Lemmes
\ref{spaces-morphisms-lemma-etale-unramified},
\ref{spaces-morphisms-lemma-etale-locally-finite-presentation} et
\ref{spaces-morphisms-lemma-etale-flat}).
\end{proof}






\section{Morphismes propres}
\label{section-proper}

\noindent
La notion de morphisme propre joue un rôle important en géométrie algébrique.
Voici la définition d'un morphisme propre de champs algébriques.

\begin{definition}
\label{definition-proper}
Soit $f : \mathcal{X} \to \mathcal{Y}$
un morphisme de champs algébriques.
On dit que $f$ est {\it propre} si $f$ est séparé, de type fini et
universellement fermé.
\end{definition}

\noindent
Cette définition est compatible avec la notion déjà existante de morphisme propre
d'espaces algébriques : un morphisme d'espaces algébriques est propre si et seulement si
il est séparé, de type fini et universellement fermé
(Morphismes d'espaces, Définition \ref{spaces-morphisms-definition-proper}),
et nous avons déjà vérifié la compatibilité de ces notions dans le
Lemme \ref{lemma-representable-separated-diagonal-closed},
la section \ref{section-finite-type} et le
Lemme \ref{lemma-characterize-representable-universally-closed}.
De même, si $f : \mathcal{X} \to \mathcal{Y}$ est un
morphisme de champs algébriques représentable par des espaces algébriques,
nous avons défini ce que signifie pour $f$ d'être propre dans
Propriétés des champs, section
\ref{stacks-properties-section-properties-morphisms}.
Cependant, la discussion de cette section montre que cela
équivaut à exiger que $f$ soit séparé, de type fini et
universellement fermé, et les mêmes références que ci-dessus établissent la
compatibilité.

\begin{lemma}
\label{lemma-base-change-proper}
Tout changement de base d'un morphisme propre est propre.
\end{lemma}

\begin{proof}
Voir les
Lemmes \ref{lemma-base-change-separated},
\ref{lemma-base-change-finite-type} et
\ref{lemma-base-change-universally-closed}.
\end{proof}

\begin{lemma}
\label{lemma-composition-proper}
Le composé de morphismes propres est propre.
\end{lemma}

\begin{proof}
Voir les
Lemmes \ref{lemma-composition-separated},
\ref{lemma-composition-finite-type} et
\ref{lemma-composition-universally-closed}.
\end{proof}

\begin{lemma}
\label{lemma-closed-immersion-proper}
Une immersion fermée de champs algébriques est un morphisme propre de
champs algébriques.
\end{lemma}

\begin{proof}
Une immersion fermée est, par définition, représentable
(Propriétés des champs, Définition
\ref{stacks-properties-definition-immersion}).
L'assertion résulte donc de la discussion de
Propriétés des champs, section
\ref{stacks-properties-section-properties-morphisms}
et du résultat correspondant pour les morphismes d'espaces algébriques ; voir
Morphismes d'espaces, Lemme
\ref{spaces-morphisms-lemma-closed-immersion-proper}.
\end{proof}

\begin{lemma}
\label{lemma-universally-closed-permanence}
Considérons un diagramme commutatif
$$
\xymatrix{
\mathcal{X} \ar[rr] \ar[rd] & &
\mathcal{Y} \ar[ld] \\
& \mathcal{Z} &
}
$$
de champs algébriques.
\begin{enumerate}
\item Si $\mathcal{X} \to \mathcal{Z}$ est universellement fermé et
$\mathcal{Y} \to \mathcal{Z}$ est séparé,
alors le morphisme $\mathcal{X} \to \mathcal{Y}$ est universellement fermé.
En particulier, l'image de $|\mathcal{X}|$ dans $|\mathcal{Y}|$ est fermée.
\item Si $\mathcal{X} \to \mathcal{Z}$ est propre et
$\mathcal{Y} \to \mathcal{Z}$ est séparé, alors
le morphisme $\mathcal{X} \to \mathcal{Y}$ est propre.
\end{enumerate}
\end{lemma}

\begin{proof}
Supposons que $\mathcal{X} \to \mathcal{Z}$ soit universellement fermé et que
$\mathcal{Y} \to \mathcal{Z}$ soit séparé.
Factorisons le morphisme sous la forme
$\mathcal{X} \to \mathcal{X} \times_\mathcal{Z} \mathcal{Y} \to \mathcal{Y}$.
Le premier morphisme est propre (Lemme \ref{lemma-semi-diagonal}),
donc universellement fermé.
La projection $\mathcal{X} \times_\mathcal{Z} \mathcal{Y} \to \mathcal{Y}$
est déduite par changement de base d'un morphisme universellement fermé, donc
elle est universellement fermée ; voir le
Lemme \ref{lemma-base-change-universally-closed}.
Ainsi, $\mathcal{X} \to \mathcal{Y}$ est universellement fermé comme composé
de morphismes universellement fermés ; voir le
Lemme \ref{lemma-composition-universally-closed}.
Cela prouve (1). Pour en déduire (2), combinons (1) avec les
Lemmes \ref{lemma-compose-after-separated},
\ref{lemma-quasi-compact-permanence} et
\ref{lemma-finite-type-permanence}.
\end{proof}

\begin{lemma}
\label{lemma-image-proper-is-proper}
Soit $\mathcal{Z}$ un champ algébrique.
Soit $f : \mathcal{X} \to \mathcal{Y}$
un morphisme de champs algébriques au-dessus de $\mathcal{Z}$.
Si $\mathcal{X}$ est universellement fermé au-dessus de $\mathcal{Z}$
et si $f$ est surjectif, alors $\mathcal{Y}$
est universellement fermé au-dessus de $\mathcal{Z}$.
En particulier, si de plus $\mathcal{Y}$ est
séparé et de type fini au-dessus de $\mathcal{Z}$,
alors $\mathcal{Y}$ est propre au-dessus de $\mathcal{Z}$.
\end{lemma}

\begin{proof}
Supposons que $\mathcal{X}$ soit universellement fermé et que $f$ soit surjectif.
Notons $p : \mathcal{X} \to \mathcal{Z}$ et
$q : \mathcal{Y} \to \mathcal{Z}$ les morphismes structuraux.
Soit $\mathcal{Z}' \to \mathcal{Z}$ un morphisme de champs algébriques.
Le changement de base $f' : \mathcal{X}' \to \mathcal{Y}'$
de $f$ par $\mathcal{Z}' \to \mathcal{Z}$ est surjectif
(Propriétés des champs, Lemme
\ref{stacks-properties-lemma-base-change-surjective}) et le changement
de base $p' : \mathcal{X}' \to \mathcal{Z}'$ de $p$ est fermé.
Notons $q' : \mathcal{Y}' \to \mathcal{Z}'$ le changement de base de $q$.
Si $T \subset |\mathcal{Y}'|$ est fermé, alors
$(f')^{-1}(T) \subset |\mathcal{X}'|$ est fermé, donc
$p'((f')^{-1}(T)) = q'(T)$ est fermé. Ainsi, $q'$ est fermé.
\end{proof}






\section{Image schématique}
\label{section-scheme-theoretic-image}

\noindent
Voici la définition.

\begin{definition}
\label{definition-scheme-theoretic-image}
Soit $f : \mathcal{X} \to \mathcal{Y}$ un morphisme de champs algébriques.
L'{\it image schématique} de $f$ est le plus petit sous-champ fermé
$\mathcal{Z} \subset \mathcal{Y}$ par lequel $f$
se factorise\footnote{Nous verrons au
lemme \ref{lemma-scheme-theoretic-image-existence}
que l'image schématique existe toujours.}.
\end{definition}

\noindent
Nous désignons souvent encore par $f : \mathcal{X} \to \mathcal{Z}$ la factorisation de $f$.
Si le morphisme $f$ n'est pas quasi-compact, alors, en général, la
construction de l'image schématique ne commute pas à la
restriction aux sous-champs ouverts de $\mathcal{Y}$. Toutefois, si $f$ est
quasi-compact, l'image schématique commute au changement de base plat
(lemme \ref{lemma-existence-plus-flat-base-change}).

\begin{lemma}
\label{lemma-cover-upstairs}
Soit $f : \mathcal{X} \to \mathcal{Y}$ un morphisme de champs algébriques.
Soit $g : \mathcal{W} \to \mathcal{X}$ un morphisme de champs algébriques
surjectif, plat et localement de présentation finie.
Alors l'image schématique de $f$ existe si et seulement si celle de
$f \circ g$ existe et, dans ce cas, ces images schématiques
sont les mêmes.
\end{lemma}

\begin{proof}
Supposons que $\mathcal{Z} \subset \mathcal{Y}$
soit un sous-champ fermé et que $f \circ g$ se factorise par $\mathcal{Z}$.
Pour démontrer le lemme, il suffit de montrer
que $f$ se factorise par $\mathcal{Z}$.
Considérons un schéma $T$ et un morphisme $T \to \mathcal{X}$
donné par un objet $x$ de la catégorie fibre de $\mathcal{X}$ au-dessus de $T$.
Nous allons montrer que $f(x)$ appartient en fait à la catégorie fibre de $\mathcal{Z}$
au-dessus de $T$. En effet, la projection $T \times_\mathcal{X} \mathcal{W} \to T$
est un morphisme surjectif, plat et localement de présentation finie.
Il existe donc un recouvrement fppf $\{T_i \to T\}$ tel que
$T_i \to T$ se factorise par $T \times_\mathcal{X} \mathcal{W} \to T$
pour tout $i$. Alors $T_i \to \mathcal{X}$ se factorise par $\mathcal{W}$
et donc $T_i \to \mathcal{Y}$ se factorise par $\mathcal{Z}$.
Ainsi $f(x)|_{T_i}$ est un objet de $\mathcal{Z}$.
Puisque $\mathcal{Z}$ est un sous-champ strictement plein, nous en concluons
que $f(x)$ est un objet de $\mathcal{Z}$, comme voulu.
\end{proof}

\begin{lemma}
\label{lemma-scheme-theoretic-image-existence}
Soit $f : \mathcal{Y} \to \mathcal{X}$ un morphisme de champs algébriques.
Alors l'image schématique de $f$ existe.
\end{lemma}

\begin{proof}
Choisissons un schéma $V$ et un morphisme lisse surjectif $V \to \mathcal{Y}$.
D'après le lemme \ref{lemma-cover-upstairs}, nous pouvons remplacer $\mathcal{Y}$ par $V$.
Il suffit donc de montrer que, si $X \to \mathcal{X}$ est un morphisme d'un
schéma vers un champ algébrique, alors l'image schématique existe.
Choisissons un schéma $U$ et un morphisme lisse surjectif $U \to \mathcal{X}$.
Posons $R = U \times_\mathcal{X} U$.
Nous avons $\mathcal{X} = [U/R]$ d'après
Champs algébriques, lemme \ref{algebraic-lemma-stack-presentation}.
D'après Propriétés des champs, lemme
\ref{stacks-properties-lemma-substacks-presentation},
les sous-champs fermés $\mathcal{Z}$ de $\mathcal{X}$
sont en correspondance $1$ à $1$ avec les sous-schémas fermés invariants
par $R$, $Z \subset U$.
Soit $Z_1 \subset U$ l'image schématique de
$X \times_\mathcal{X} U \to U$.
Remarquons que $X \to \mathcal{X}$ se factorise par $\mathcal{Z}$
si et seulement si $X \times_\mathcal{X} U \to U$ se factorise par
le sous-schéma fermé correspondant, invariant par $R$, à savoir $Z$
(détails omis ; indication : ceci résulte de ce que
$X \times_\mathcal{X} U \to X$ est lisse et surjectif).
Il nous faut donc montrer qu'il existe un plus petit sous-schéma fermé
invariant par $R$, $Z \subset U$, et contenant $Z_1$.

\medskip\noindent
Soit $\mathcal{I}_1 \subset \mathcal{O}_U$ le faisceau quasi-cohérent
d'idéaux correspondant au sous-schéma fermé $Z_1 \subset U$.
Soient $Z_\alpha$, $\alpha \in A$, tous les sous-schémas fermés invariants
par $R$ de $U$ et contenant $Z_1$.
Pour $\alpha \in A$, soit $\mathcal{I}_\alpha \subset \mathcal{O}_U$
le faisceau quasi-cohérent d'idéaux correspondant au sous-schéma fermé
$Z_\alpha \subset U$. L'inclusion $Z_1 \subset Z_\alpha$
signifie que $\mathcal{I}_\alpha \subset \mathcal{I}_1$.
L'invariance par $R$ de $Z_\alpha$ signifie que
$$
s^{-1}\mathcal{I}_\alpha \cdot \mathcal{O}_R =
t^{-1}\mathcal{I}_\alpha \cdot \mathcal{O}_R
$$
comme faisceaux quasi-cohérents d'idéaux sur l'espace algébrique $R$.
Considérons l'image
$$
\mathcal{I} =
\Im\left(
\bigoplus\nolimits_{\alpha \in A} \mathcal{I}_\alpha \to \mathcal{I}_1
\right) =
\Im\left(
\bigoplus\nolimits_{\alpha \in A} \mathcal{I}_\alpha \to \mathcal{O}_U
\right)
$$
Comme les sommes directes de faisceaux quasi-cohérents sont quasi-cohérentes
et que les images de morphismes entre faisceaux quasi-cohérents sont
quasi-cohérentes, nous voyons que $\mathcal{I}$ est quasi-cohérent.
Comme le foncteur image inverse est exact et commute aux sommes directes, nous obtenons
$$
s^{-1}\mathcal{I} \cdot \mathcal{O}_R =
t^{-1}\mathcal{I} \cdot \mathcal{O}_R
$$
Par conséquent, $\mathcal{I}$ définit un sous-schéma fermé invariant par $R$,
$Z \subset U$ qui est contenu dans tout $Z_\alpha$ et contient
$Z_1$, comme voulu.
\end{proof}

\begin{lemma}
\label{lemma-factor-factor}
Soit
$$
\xymatrix{
\mathcal{X}_1 \ar[d] \ar[r]_{f_1} & \mathcal{Y}_1 \ar[d] \\
\mathcal{X}_2 \ar[r]^{f_2} & \mathcal{Y}_2
}
$$
un diagramme commutatif de champs algébriques.
Soient $\mathcal{Z}_i \subset \mathcal{Y}_i$, $i = 1, 2$, les
images schématiques de $f_i$. Alors le morphisme
$\mathcal{Y}_1 \to \mathcal{Y}_2$ induit un morphisme
$\mathcal{Z}_1 \to \mathcal{Z}_2$ et un
diagramme commutatif
$$
\xymatrix{
\mathcal{X}_1 \ar[r] \ar[d] &
\mathcal{Z}_1 \ar[d] \ar[r] &
\mathcal{Y}_1 \ar[d] \\
\mathcal{X}_2 \ar[r] &
\mathcal{Z}_2 \ar[r] &
\mathcal{Y}_2
}
$$
\end{lemma}

\begin{proof}
L'image réciproque schématique de $\mathcal{Z}_2$ dans $\mathcal{Y}_1$
est un sous-champ fermé de $\mathcal{Y}_1$ par lequel
$f_1$ se factorise. Par conséquent, $\mathcal{Z}_1$ y est contenu.
Ceci démontre le lemme.
\end{proof}

\begin{lemma}
\label{lemma-existence-plus-flat-base-change}
Soit $f : \mathcal{X} \to \mathcal{Y}$ un morphisme quasi-compact
de champs algébriques. Alors la formation de l'image schématique
commute au changement de base plat.
\end{lemma}

\begin{proof}
Soit $\mathcal{Y}' \to \mathcal{Y}$ un morphisme plat de champs algébriques.
Choisissons un schéma $V$ et un morphisme lisse surjectif $V \to \mathcal{Y}$.
Choisissons un schéma $V'$ et un
morphisme lisse surjectif $V' \to \mathcal{Y}' \times_\mathcal{Y} V$.
Nous supposerons que $V = \coprod_{i \in I} V_i$ est une somme
disjointe de schémas affines et que
$V' = \coprod_{i \in I} \coprod_{j \in J_i} V_{i, j}$
est une somme disjointe de schémas affines, chaque $V_{i, j}$ s'envoyant dans $V_i$.
Soient
\begin{enumerate}
\item $\mathcal{Z} \subset \mathcal{Y}$ l'image schématique de $f$,
\item $\mathcal{Z}' \subset \mathcal{Y}'$ l'image schématique
du changement de base de $f$ par $\mathcal{Y}' \to \mathcal{Y}$,
\item $Z \subset V$ l'image schématique
du changement de base de $f$ par $V \to \mathcal{Y}$,
\item $Z' \subset V'$ l'image schématique
du changement de base de $f$ par $V' \to \mathcal{Y}$.
\end{enumerate}
Si nous pouvons montrer que
(a) $Z = V \times_\mathcal{Y} \mathcal{Z}$,
(b) $Z' = V' \times_{\mathcal{Y}'} \mathcal{Z}'$, et
(c) $Z' = V' \times_V Z$,
alors le lemme en résulte : l'inclusion
$\mathcal{Z}' \to \mathcal{Z} \times_\mathcal{Y} \mathcal{Y}'$
(lemme \ref{lemma-factor-factor})
est nécessairement un isomorphisme, car son changement de base par le morphisme
lisse surjectif $V' \to \mathcal{Y}'$ en est un.

\medskip\noindent
Démonstration de (a). Posons $R = V \times_\mathcal{Y} V$.
D'après Propriétés des champs, lemme
\ref{stacks-properties-lemma-substacks-presentation},
l'application $\mathcal{Z} \mapsto \mathcal{Z} \times_\mathcal{Y} V$
définit une correspondance bijective entre les sous-champs fermés
de $\mathcal{Y}$ et les sous-espaces fermés invariants par $R$ dans $V$.
De plus, $f : \mathcal{X} \to \mathcal{Y}$ se factorise par $\mathcal{Z}$
si et seulement si le changement de base
$g : \mathcal{X} \times_\mathcal{Y} V \to V$ se factorise par
$\mathcal{Z} \times_\mathcal{Y} V$.
Nous affirmons que l'image schématique $Z \subset V$ de $g$
est invariante par $R$. Cette assertion implique (a) d'après ce qui précède.

\medskip\noindent
Pour chaque $i$, le morphisme $\mathcal{X} \times_\mathcal{Y} V_i \to V_i$
est quasi-compact ; par conséquent, $\mathcal{X} \times_\mathcal{Y} V_i$ est
quasi-compact. Nous pouvons donc choisir un schéma affine $W_i$ et un
morphisme lisse surjectif $W_i \to \mathcal{X} \times_\mathcal{Y} V_i$.
Remarquons que $W = \coprod W_i$ est un schéma muni d'un
morphisme lisse surjectif $W \to \mathcal{X} \times_\mathcal{Y} V$
tel que le composé $W \to V$ défini par $g$ soit quasi-compact.
Soit $Z \to V$ l'image schématique de $W \to V$, voir
Morphismes, section
\ref{morphisms-section-scheme-theoretic-image}, et
Morphismes d'espaces, section
\ref{spaces-morphisms-section-scheme-theoretic-image}.
Il résulte du lemme \ref{lemma-cover-upstairs}
que $Z \subset V$ est l'image schématique de $g$.
Pour montrer que $Z$ est invariant par $R$, nous affirmons que les deux
$$
\text{pr}_0^{-1}(Z), \text{pr}_1^{-1}(Z) \subset R = V \times_\mathcal{Y} V
$$
sont l'image schématique de $\mathcal{X} \times_\mathcal{Y} R \to R$.
En effet, Morphismes d'espaces, lemme
\ref{spaces-morphisms-lemma-flat-base-change-scheme-theoretic-image},
montre d'abord que $\text{pr}_0^{-1}(Z)$ est l'image schématique
du composé
$$
W \times_{V, \text{pr}_0} R = W \times_\mathcal{Y} V \to
\mathcal{X} \times_\mathcal{Y} R \to R
$$
Comme la première flèche est lisse et surjective, nous voyons que
$\text{pr}_0^{-1}(Z)$ est l'image schématique de
$\mathcal{X} \times_\mathcal{Y} R \to R$. Le même argument montre que
$\text{pr}_1^{-1}(Z)$ l'est aussi. Ainsi $Z$ est invariant par $R$.

\medskip\noindent
L'assertion (b) se démontre exactement de la même manière que (a).

\medskip\noindent
Démonstration de (c). Soit $Z_i \subset V_i$ l'image schématique
de $\mathcal{X} \times_\mathcal{Y} V_i \to V_i$ et soit
$Z_{i, j} \subset V_{i, j}$ l'image schématique de
$\mathcal{X} \times_\mathcal{Y} V_{i, j} \to V_{i, j}$.
Il suffit clairement de montrer que l'image réciproque de $Z_i$
dans $V_{i, j}$ est $Z_{i, j}$. Nous avons vu ci-dessus que
$Z_i$ est l'image schématique de $W_i \to V_i$
et, de même, que $Z_{i, j}$ est l'image
schématique de $W_i \times_{V_i} V_{i, j} \to V_{i, j}$.
L'égalité résulte donc du cas des schémas
(Morphismes, lemme
\ref{morphisms-lemma-flat-base-change-scheme-theoretic-image})
et du fait que $V_{i, j} \to V_i$ est plat.
\end{proof}

\begin{lemma}
\label{lemma-topology-scheme-theoretic-image}
Soit $f : \mathcal{X} \to \mathcal{Y}$ un morphisme quasi-compact
de champs algébriques. Soit $\mathcal{Z} \subset \mathcal{Y}$
l'image schématique de $f$. Alors $|\mathcal{Z}|$
est l'adhérence de l'image de $|f|$.
\end{lemma}

\begin{proof}
Soit $z \in |\mathcal{Z}|$ un point.
Choisissons un schéma affine $V$, un point $v \in V$ et un morphisme lisse
$V \to \mathcal{Y}$ envoyant $v$ sur $z$.
Alors $\mathcal{X} \times_\mathcal{Y} V$ est un champ algébrique quasi-compact.
Nous pouvons donc trouver un schéma affine $W$ et un morphisme lisse
surjectif $W \to \mathcal{X} \times_\mathcal{Y} V$.
D'après le lemme \ref{lemma-existence-plus-flat-base-change},
l'image schématique de
$\mathcal{X} \times_\mathcal{Y} V \to V$ est
$Z = \mathcal{Z} \times_\mathcal{Y} V$.
Ainsi, l'image réciproque de $|\mathcal{Z}|$ dans $|V|$ est $|Z|$ d'après
Propriétés des champs, lemme \ref{stacks-properties-lemma-points-cartesian}.
D'après le lemme \ref{lemma-cover-upstairs}, $Z$ est
l'image schématique de $W \to V$.
D'après Morphismes d'espaces, lemme
\ref{spaces-morphisms-lemma-quasi-compact-scheme-theoretic-image},
nous voyons que l'image de $|W| \to |Z|$ est dense.
Par conséquent, l'image de $|\mathcal{X} \times_\mathcal{Y} V| \to |Z|$
est dense. Remarquons que $v \in |Z|$.
Puisque $|V| \to |\mathcal{Y}|$ est ouvert, un argument topologique
montre que $z$ appartient à l'adhérence de l'image de $|f|$, comme voulu.
\end{proof}

\begin{lemma}
\label{lemma-scheme-theoretic-image-of-partial-section}
Soit $f : \mathcal{X} \to \mathcal{Y}$ un morphisme de champs algébriques
représentable par des espaces algébriques et séparé.
Soit $\mathcal{V} \subset \mathcal{Y}$ un sous-champ ouvert tel que
$\mathcal{V} \to \mathcal{Y}$ soit quasi-compact.
Soit $s : \mathcal{V} \to \mathcal{X}$ un morphisme tel que
$f \circ s = \text{id}_\mathcal{V}$.
Soit $\mathcal{Y}'$ l'image schématique de $s$.
Alors $\mathcal{Y}' \to \mathcal{Y}$ est un isomorphisme au-dessus de $\mathcal{V}$.
\end{lemma}

\begin{proof}
D'après le lemme \ref{lemma-quasi-compact-permanence},
le morphisme $s : \mathcal{V} \to \mathcal{X}$ est quasi-compact.
Par conséquent, la construction de l'image schématique $\mathcal{Y}'$
de $s$ commute au changement de base plat d'après
le lemme \ref{lemma-existence-plus-flat-base-change}.
Ainsi, pour démontrer le lemme,
nous pouvons supposer que $\mathcal{Y}$ est représentable par un espace algébrique,
et nous sommes ramenés au cas des espaces algébriques, qui est
Morphismes d'espaces, lemme
\ref{spaces-morphisms-lemma-scheme-theoretic-image-of-partial-section}.
\end{proof}






\section{Crit\`eres valuatifs}
\label{section-valuative}

\noindent
Il faut prendre garde lors de l'\'enonc\'e du crit\`ere valuatif. En effet, dans
la formulation, nous devons parler de diagrammes commutatifs, mais nous
travaillons dans une $2$-cat\'egorie et nous devons nous assurer que les $2$-morphismes
se composent eux aussi correctement !

\begin{definition}
\label{definition-fill-in-diagram}
Soit $f : \mathcal{X} \to \mathcal{Y}$ un morphisme de champs alg\'ebriques.
Consid\'erons le diagramme $2$-commutatif en fl\`eches pleines
\begin{equation}
\label{equation-diagram}
\vcenter{
\xymatrix{
\Spec(K) \ar[r]_-x \ar[d]_j & \mathcal{X} \ar[d]^f \\
\Spec(A) \ar[r]^-y \ar@{..>}[ru] & \mathcal{Y}
}
}
\end{equation}
o\`u $A$ est un anneau de valuation et $K$ son corps des fractions. Soit
$$
\gamma : y \circ j \longrightarrow f \circ x
$$
un $2$-morphisme t\'emoignant de la $2$-commutativit\'e du diagramme.
(Notations comme dans Cat\'egories, sections \ref{categories-section-formal-cat-cat}
et \ref{categories-section-2-categories}.)
\'Etant donn\'es (\ref{equation-diagram}) et $\gamma$,
une {\it fl\`eche pointill\'ee} est un triplet $(a, \alpha, \beta)$ constitu\'e d'un
morphisme $a : \Spec(A) \to \mathcal{X}$ et de $2$-fl\`eches
$\alpha : a \circ j \to x$, $\beta : y \to f \circ a$
telles que
$\gamma = (\text{id}_f \star \alpha) \circ (\beta \star \text{id}_j)$,
autrement dit telles que
$$
\xymatrix{
& f \circ a \circ j \ar[rd]^{\text{id}_f \star \alpha} \\
y \circ j \ar[ru]^{\beta \star \text{id}_j} \ar[rr]^\gamma & &
f \circ x
}
$$
soit commutatif. Un {\it morphisme de fl\`eches pointill\'ees}
$(a, \alpha, \beta) \to (a', \alpha', \beta')$ est une
$2$-fl\`eche $\theta : a \to a'$ telle que
$\alpha = \alpha' \circ (\theta \star \text{id}_j)$ et
$\beta' = (\text{id}_f \star \theta) \circ \beta$.
\end{definition}

\noindent
La d\'efinition pr\'ec\'edente est un cas particulier de Cat\'egories, 
D\'efinition \ref{categories-definition-dotted-arrows}.
La cat\'egorie des fl\`eches pointill\'ees d\'epend en g\'en\'eral de $\gamma$.
Si $\mathcal{Y}$ est repr\'esentable par un espace alg\'ebrique
(ou si les groupes d'automorphismes des objets au-dessus des corps sont triviaux),
il existe bien entendu au plus un $\gamma$, et il n'est pas n\'ecessaire
de v\'erifier la commutativit\'e du triangle. Plus g\'en\'eralement, nous
disposons du Lemme \ref{lemma-cat-dotted-arrows-independent}.
La commutativit\'e du triangle est importante dans la d\'emonstration
de la compatibilit\'e au changement de base, voir la
d\'emonstration du Lemme \ref{lemma-cat-dotted-arrows-base-change}.

\begin{lemma}
\label{lemma-cat-dotted-arrows}
Dans la situation de la D\'efinition \ref{definition-fill-in-diagram},
la cat\'egorie des fl\`eches pointill\'ees est un groupo\"\i de. Si $\Delta_f$
est s\'epar\'e, alors c'est un s\'eto\"\i de.
\end{lemma}

\begin{proof}
Comme les $2$-fl\`eches sont inversibles, il est clair que la cat\'egorie des
fl\`eches pointill\'ees est un groupo\"\i de. Pour une fl\`eche pointill\'ee $(a, \alpha, \beta)$,
un automorphisme de $(a, \alpha, \beta)$ est un $2$-morphisme
$\theta : a \to a$ qui v\'erifie deux conditions. La premi\`ere condition
$\beta = (\text{id}_f \star \theta) \circ \beta$ exprime que
$\theta$ d\'efinit un morphisme
$(a, \theta) : \Spec(A) \to \mathcal{I}_{\mathcal{X}/\mathcal{Y}}$.
La seconde condition
$\alpha = \alpha \circ (\theta \star \text{id}_j)$
implique que la restriction de $(a, \theta)$ \`a $\Spec(K)$
est l'identit\'e. On a le diagramme
$$
\xymatrix{
\mathcal{I}_{\mathcal{X}/\mathcal{Y}} \ar[d] & &
\Spec(K) \ar[d]^j \ar[ll]_{(a \circ j, \text{id})} \\
\mathcal{X} & & \Spec(A) \ar[ll]_a \ar[llu]_{(a, \theta)}
}
$$
Autrement dit, si $G \to \Spec(A)$ est l'espace alg\'ebrique en groupes
obtenu par changement de base de l'inertie relative suivant $a$, alors
$\theta$ d\'efinit un point $\theta \in G(A)$ dont l'image
dans $G(K)$ est l'\'el\'ement neutre. Bien entendu, si la section unit\'e $e : \Spec(A) \to G$
est une immersion ferm\'ee, cela ne peut se produire que si
$\theta$ est l'identit\'e.
Le Lemme \ref{lemma-diagonal-diagonal}
donne le r\'esultat voulu.
\end{proof}

\begin{lemma}
\label{lemma-cat-dotted-arrows-independent}
Dans la D\'efinition \ref{definition-fill-in-diagram}, supposons que
$\mathcal{I}_\mathcal{Y} \to \mathcal{Y}$ soit propre
(par exemple si $\mathcal{Y}$ est s\'epar\'e, ou si $\mathcal{Y}$
est s\'epar\'e au-dessus d'un espace alg\'ebrique). Alors la cat\'egorie des fl\`eches pointill\'ees
est ind\'ependante, \`a une \'equivalence non canonique pr\`es, du choix de $\gamma$,
et l'existence d'une fl\`eche pointill\'ee
(pour un certain choix, donc, de mani\`ere \`equivalente, pour tout $\gamma$)
\'equivaut \`a l'existence d'un diagramme
$$
\xymatrix{
\Spec(K) \ar[r]_-x \ar[d]_j & \mathcal{X} \ar[d]^f \\
\Spec(A) \ar[r]^-y \ar[ru]_a & \mathcal{Y}
}
$$
dont les triangles sont $2$-commutatifs
(sans v\'erifier que les $2$-morphismes se composent correctement).
\end{lemma}

\begin{proof}
Soient $\gamma, \gamma' : y \circ j \longrightarrow f \circ x$
deux $2$-morphismes. Alors $\gamma^{-1} \circ \gamma'$
est un automorphisme de $y$ au-dessus de $\Spec(K)$.
Ainsi, si $\mathit{Isom}_\mathcal{Y}(y, y) \to \Spec(A)$
est propre, le crit\`ere valuatif de propret\'e
(Morphismes d'espaces, Lemme \ref{spaces-morphisms-lemma-characterize-proper})
permet de trouver $\delta : y \to y$ dont la restriction \`a
$\Spec(K)$ est $\gamma^{-1} \circ \gamma'$.
Nous pouvons alors utiliser $\delta$ pour d\'efinir une \`equivalence
entre la cat\'egorie des fl\`eches pointill\'ees pour $\gamma$
et la cat\'egorie des fl\`eches pointill\'ees pour $\gamma'$, en
envoyant $(a, \alpha, \beta)$ sur $(a, \alpha, \beta \circ \delta)$.
La derni\`ere assertion est claire.
\end{proof}

\begin{lemma}
\label{lemma-cat-dotted-arrows-base-change}
Supposons donn\'e un diagramme $2$-commutatif
$$
\xymatrix{
\Spec(K) \ar[r]_-{x'} \ar[d]_j &
\mathcal{X}' \ar[d]^p \ar[r]_q &
\mathcal{X} \ar[d]^f \\
\Spec(A) \ar[r]^-{y'} &
\mathcal{Y}' \ar[r]^g &
\mathcal{Y}
}
$$
dont le carr\'e de droite est $2$-cart\'esien. Choisissons une $2$-fl\`eche
$\gamma' : y' \circ j \to p \circ x'$. Posons
$x = q \circ x'$, $y = g \circ y'$ et soit
$\gamma : y \circ j \to f \circ x$ le compos\'e de
$\gamma'$ et de la $2$-fl\`eche implicite dans la $2$-commutativit\'e
du carr\'e de droite. Alors la cat\'egorie des fl\`eches pointill\'ees
pour le carr\'e de gauche et $\gamma'$ est \`equivalente \`a la cat\'egorie des fl\`eches
pointill\'ees pour le rectangle ext\'erieur et $\gamma$.
\end{lemma}

\begin{proof}
(Nous ne savons pas d\'emontrer l'analogue de ce lemme si, au lieu
de la cat\'egorie des fl\`eches pointill\'ees, nous consid\'erons l'ensemble des classes
d'isomorphie des morphismes donnant deux triangles
$2$-commutatifs comme dans le Lemme \ref{lemma-cat-dotted-arrows-independent};
il se peut fort bien, en fait, que cet analogue soit faux.)
Premi\`ere d\'emonstration : ce lemme est un cas particulier de Cat\'egories, 
Lemme \ref{categories-lemma-cat-dotted-arrows-base-change}.
Seconde d\'emonstration : nous pouvons remplacer
$\mathcal{X}'$ par le $2$-produit fibr\'e
$\mathcal{Y}' \times_\mathcal{Y} \mathcal{X}$
tel qu'il est d\'ecrit dans Cat\'egories, Lemme
\ref{categories-lemma-2-product-categories-over-C}.
L'objet $x'$ devient alors le triplet $(y' \circ j, x, \gamma)$.
On passe alors d'une fl\`eche pointill\'ee $(a, \alpha, \beta)$ pour le
rectangle ext\'erieur \`a une fl\`eche pointill\'ee $(a', \alpha', \beta')$
pour le carr\'e de gauche en prenant $a' = (y', a, \beta)$,
$\alpha' = (\text{id}_{y' \circ j}, \alpha)$ et
$\beta' = \text{id}_{y'}$. D\'etails omis.
\end{proof}

\begin{lemma}
\label{lemma-cat-dotted-arrows-composition}
Supposons donn\'e un diagramme $2$-commutatif
$$
\xymatrix{
\Spec(K) \ar[r]_-x \ar[dd]_j & \mathcal{X} \ar[d]^f \\
& \mathcal{Y} \ar[d]^g \\
\Spec(A) \ar[r]^-z & \mathcal{Z}
}
$$
Choisissons une $2$-fl\`eche $\gamma : z \circ j \to g \circ f \circ x$.
Soit $\mathcal{C}$ la cat\'egorie des fl\`eches pointill\'ees pour
le rectangle ext\'erieur et $\gamma$. Soit $\mathcal{C}'$ la
cat\'egorie des fl\`eches pointill\'ees pour le carr\'e
$$
\xymatrix{
\Spec(K) \ar[r]_-{f \circ x} \ar[d]_j & \mathcal{Y} \ar[d]^g \\
\Spec(A) \ar[r]^-z & \mathcal{Z}
}
$$
et $\gamma$. Alors $\mathcal{C}$ est \`equivalente \`a une cat\'egorie $\mathcal{C}''$ 
qui poss\`ede la propri\'et\'e suivante : il existe 
un foncteur $\mathcal{C}'' \to \mathcal{C}'$
qui fait de $\mathcal{C}''$ une cat\'egorie fibr\'ee en groupo\"\i des sur
$\mathcal{C}'$ et dont les cat\'egories-fibres sont des cat\'egories de fl\`eches pointill\'ees
pour certains carr\'es de la forme
$$
\xymatrix{
\Spec(K) \ar[r]_-x \ar[d]_j & \mathcal{X} \ar[d]^f \\
\Spec(A) \ar[r]^-y & \mathcal{Y}
}
$$
et certains choix de $y \circ j \to f \circ x$.
\end{lemma}

\begin{proof}
Ce lemme est un cas particulier de Cat\'egories, 
Lemme \ref{categories-lemma-cat-dotted-arrows-composition}.
\end{proof}

\begin{definition}
\label{definition-uniqueness}
Soit $f : \mathcal{X} \to \mathcal{Y}$ un morphisme de champs alg\'ebriques.
Nous disons que $f$ v\'erifie la {\it partie d'unicit\'e du crit\`ere valuatif}
si, pour tout diagramme (\ref{equation-diagram}) et tout $\gamma$
comme dans la D\'efinition \ref{definition-fill-in-diagram},
la cat\'egorie des fl\`eches pointill\'ees est soit vide, soit
un s\'eto\"\i de ayant exactement une classe d'isomorphie.
\end{definition}

\begin{lemma}
\label{lemma-base-change-uniqueness}
Le changement de base, par un morphisme quelconque de champs alg\'ebriques,
d'un morphisme de champs alg\'ebriques qui v\'erifie la
partie d'unicit\'e du crit\`ere valuatif est un morphisme de champs alg\'ebriques
qui v\'erifie la partie d'unicit\'e du crit\`ere valuatif.
\end{lemma}

\begin{proof}
Cela r\'esulte du Lemme \ref{lemma-cat-dotted-arrows-base-change}
et de la d\'efinition.
\end{proof}

\begin{lemma}
\label{lemma-composition-uniqueness}
Le compos\'e de morphismes de champs alg\'ebriques qui v\'erifient la
partie d'unicit\'e du crit\`ere valuatif est encore un
morphisme de champs alg\'ebriques qui v\'erifie la
partie d'unicit\'e du crit\`ere valuatif.
\end{lemma}

\begin{proof}
Cela r\'esulte du Lemme \ref{lemma-cat-dotted-arrows-composition}
et de la d\'efinition.
\end{proof}

\begin{lemma}
\label{lemma-uniqueness-representable}
Soit $f : \mathcal{X} \to \mathcal{Y}$ un morphisme de champs alg\'ebriques
qui est repr\'esentable par des espaces alg\'ebriques. Alors les conditions suivantes sont \`equivalentes :
\begin{enumerate}
\item $f$ v\'erifie la partie d'unicit\'e du crit\`ere valuatif ;
\item pour tout sch\'ema $T$ et tout morphisme $T \to \mathcal{Y}$,
le morphisme $\mathcal{X} \times_\mathcal{Y} T \to T$ v\'erifie
la partie d'unicit\'e du crit\`ere valuatif en tant que morphisme
d'espaces alg\'ebriques.
\end{enumerate}
\end{lemma}

\begin{proof}
Cela r\'esulte du Lemme \ref{lemma-cat-dotted-arrows-base-change}
et de la d\'efinition.
\end{proof}

\begin{definition}
\label{definition-existence}
Soit $f : \mathcal{X} \to \mathcal{Y}$ un morphisme de champs alg\'ebriques.
Nous disons que $f$ v\'erifie la {\it partie d'existence du crit\`ere valuatif}
si, pour tout diagramme (\ref{equation-diagram}) et tout $\gamma$
comme dans la D\'efinition \ref{definition-fill-in-diagram},
il existe une extension de corps $K'/K$, un anneau de valuation $A' \subset K'$
dominant $A$, tels que la cat\'egorie des fl\`eches pointill\'ees pour le
rectangle ext\'erieur du diagramme
$$
\xymatrix{
\Spec(K') \ar[r] \ar@/^2em/[rr]_{x'} \ar[d]_{j'} &
\Spec(K) \ar[d]_j \ar[r]_-x &
\mathcal{X} \ar[d]^f \\
\Spec(A') \ar[r] \ar@/_2em/[rr]^{y'} &
\Spec(A) \ar[r]^-y &
\mathcal{Y}
}
$$
et la $2$-fl\`eche induite $\gamma' : y' \circ j' \to f \circ x'$ soit non vide.
\end{definition}

\noindent
Nous avons d\'ej\`a vu, dans le cas des morphismes d'espaces alg\'ebriques,
qu'il est n\'ecessaire d'autoriser des extensions des corps des fractions
pour obtenir la bonne notion de crit\`ere valuatif.
Voir Morphismes d'espaces, Exemple
\ref{spaces-morphisms-example-finite-separable-needed}.
Toutefois, pour les morphismes entre espaces alg\'ebriques s\'epar\'es, une telle
extension n'est pas n\'ecessaire
(Morphismes d'espaces, Lemme \ref{spaces-morphisms-lemma-usual-enough}).
Cependant, pour les morphismes entre champs alg\'ebriques, une extension peut
\^etre n\'ecessaire m\^eme si $\mathcal{X}$ et $\mathcal{Y}$ sont tous deux s\'epar\'es.
Consid\'erons par exemple le morphisme de champs alg\'ebriques
$$
[\Spec(\mathbf{C})/G] \to \Spec(\mathbf{C})
$$
au-dessus du sch\'ema de base $\Spec(\mathbf{C})$,
o\`u $G$ est un groupe d'ordre $2$. La source et le but sont tous deux des
champs alg\'ebriques s\'epar\'es, et le morphisme est propre. Il
v\'erifie donc les parties d'unicit\'e et d'existence du crit\`ere valuatif
(voir le Lemme \ref{lemma-criterion-proper}).
Mais, d'autre part, il existe un diagramme
$$
\xymatrix{
\Spec(K) \ar[r] \ar[d] & [\Spec(\mathbf{C})/G] \ar[d] \\
\Spec(A) \ar[r] & \Spec(\mathbf{C})
}
$$
pour lequel aucune fl\`eche pointill\'ee n'existe lorsque $A = \mathbf{C}[[t]]$ et
$K = \mathbf{C}((t))$. En effet, la fl\`eche horizontale sup\'erieure est donn\'ee
par un torseur sous $G$ sur le spectre de $K = \mathbf{C}((t))$. Comme tout torseur sous $G$
sur le spectre de l'anneau local strictement hens\'elien $A = \mathbf{C}[[t]]$ est trivial, nous voyons
que, si une fl\`eche pointill\'ee existe toujours, alors tout torseur sous $G$ sur
$K$ est trivial. Ce n'est pas le cas, car $G = \{+1, -1\}$
et, d'apr\`es la th\'eorie de Kummer, les torseurs sous $G$ sur $K$ sont class\'es par
$K^*/(K^*)^2$, qui n'est pas trivial.

\begin{lemma}
\label{lemma-base-change-existence}
Le changement de base, par un morphisme quelconque de champs alg\'ebriques,
d'un morphisme de champs alg\'ebriques qui v\'erifie la
partie d'existence du crit\`ere valuatif est un morphisme de champs alg\'ebriques
qui v\'erifie la partie d'existence du crit\`ere valuatif.
\end{lemma}

\begin{proof}
Cela r\'esulte du Lemme \ref{lemma-cat-dotted-arrows-base-change}
et de la d\'efinition.
\end{proof}

\begin{lemma}
\label{lemma-composition-existence}
Le compos\'e de morphismes de champs alg\'ebriques qui v\'erifient la
partie d'existence du crit\`ere valuatif est encore un
morphisme de champs alg\'ebriques qui v\'erifie la
partie d'existence du crit\`ere valuatif.
\end{lemma}

\begin{proof}
Cela r\'esulte du Lemme \ref{lemma-cat-dotted-arrows-composition}
et de la d\'efinition.
\end{proof}

\begin{lemma}
\label{lemma-existence-representable}
Soit $f : \mathcal{X} \to \mathcal{Y}$ un morphisme de champs alg\'ebriques
qui est repr\'esentable par des espaces alg\'ebriques. Alors les conditions suivantes sont \`equivalentes :
\begin{enumerate}
\item $f$ v\'erifie la partie d'existence du crit\`ere valuatif ;
\item pour tout sch\'ema $T$ et tout morphisme $T \to \mathcal{Y}$,
le morphisme $\mathcal{X} \times_\mathcal{Y} T \to T$ v\'erifie
la partie d'existence du crit\`ere valuatif en tant que morphisme
d'espaces alg\'ebriques.
\end{enumerate}
\end{lemma}

\begin{proof}
Cela r\'esulte du Lemme \ref{lemma-cat-dotted-arrows-base-change}
et de la d\'efinition.
\end{proof}

\begin{lemma}
\label{lemma-closed-immersion-valuative-criteria}
Une immersion ferm\'ee de champs alg\'ebriques v\'erifie \`a la fois
les parties d'existence et d'unicit\'e du crit\`ere valuatif.
\end{lemma}

\begin{proof}
D\'emonstration omise. Indication : se ramener au cas d'une immersion ferm\'ee de
sch\'emas \`a l'aide des Lemmes \ref{lemma-uniqueness-representable} et
\ref{lemma-existence-representable}.
\end{proof}




\section{Critère valuatif pour la seconde diagonale}
\label{section-valuative-second}

\noindent
La réciproque a déjà été démontrée dans le
Lemme \ref{lemma-cat-dotted-arrows}.
Le critère proprement dit est le suivant.

\begin{lemma}
\label{lemma-setoids-and-diagonal}
Soit $f : \mathcal{X} \to \mathcal{Y}$ un morphisme de champs algébriques.
Si le morphisme $\Delta_f$ est quasi-séparé et si, pour tout diagramme
(\ref{equation-diagram}) et tout choix de $\gamma$ comme dans la
Définition \ref{definition-fill-in-diagram},
la catégorie des flèches pointillées
est un setoïde, alors $\Delta_f$ est séparé.
\end{lemma}

\begin{proof}
Nous allons donner une démonstration détaillée, mais nous recommandons vivement au
lecteur de chercher sa propre démonstration en s'inspirant de l'argument
donné dans la preuve du Lemme \ref{lemma-cat-dotted-arrows}.

\medskip\noindent
Supposons que $\Delta_f$ soit quasi-séparé et que, pour tout diagramme
(\ref{equation-diagram}) et tout choix de $\gamma$ comme dans la
Définition \ref{definition-fill-in-diagram},
la catégorie des flèches pointillées soit un setoïde.
D'après le Lemme \ref{lemma-diagonal-diagonal}, il suffit de montrer que
$e : \mathcal{X} \to \mathcal{I}_{\mathcal{X}/\mathcal{Y}}$
est une immersion fermée. D'après le
Lemme \ref{lemma-first-diagonal-separated-second-diagonal-closed},
il suffit même de montrer que $e = \Delta_{f, 2}$ est
universellement fermé.
L'un ou l'autre de ces lemmes montre aussi que $e = \Delta_{f, 2}$ est quasi-compact
en vertu de notre hypothèse de quasi-séparation de $\Delta_f$.

\medskip\noindent
Dans ce paragraphe, nous allons montrer que $e$ vérifie la partie
d'existence du critère valuatif. Considérons le diagramme $2$-commutatif en flèches pleines
$$
\xymatrix{
\Spec(K) \ar[r]_x \ar[d]_j & \mathcal{X} \ar[d]^e \\
\Spec(A) \ar[r]^{(a, \theta)} & \mathcal{I}_{\mathcal{X}/\mathcal{Y}}
}
$$
et soit $\alpha : (a, \theta) \circ j \to e \circ x$ un $2$-morphisme quelconque
témoignant de la $2$-commutativité du diagramme (nous employons $\alpha$ à la place
de la lettre $\gamma$ utilisée dans la Définition \ref{definition-fill-in-diagram}).
Notons que $f \circ \theta = \text{id}$; nous utiliserons ce fait ci-dessous.
On a $e \circ x = (x, \text{id}_x)$ et
$(a, \theta) \circ j = (a \circ j, \theta \star \text{id}_j)$.
Ainsi, $\alpha$ est une $2$-flèche $\alpha : a \circ j \to x$
compatible avec $\theta \star \text{id}_j$ et $\text{id}_x$.
Posons $y = f \circ x$ et $\beta = \text{id}_{f \circ a}$.
En reprenant en sens inverse les arguments donnés dans la preuve du
Lemme \ref{lemma-cat-dotted-arrows},
on voit que $\theta$ est un automorphisme de la
flèche pointillée $(a, \alpha, \beta)$ pour laquelle
$$
\gamma : y \circ j \to f \circ x
\quad\text{égal à}\quad
\text{id}_f \star \alpha : f \circ a \circ j \to f \circ x
$$
D'autre part, $\text{id}_a$ est aussi un automorphisme; on en conclut
$\theta = \text{id}_a$ grâce à l'hypothèse sur $f$.
On peut alors prendre pour flèche pointillée du diagramme affiché ci-dessus
le morphisme $a : \Spec(A) \to \mathcal{X}$ muni des deux $2$-morphismes suivants :
le premier, $(a, \text{id}_a) \circ j \to (x, \text{id}_x)$, est donné par $\alpha$,
et le second, $(a, \theta) \to e \circ a$, par $\text{id}_a$.

\medskip\noindent
D'après le Lemme \ref{lemma-base-change-existence}, tout changement de base de $e$
vérifie la partie d'existence du critère valuatif.
Comme $e$ est représentable par des espaces algébriques, il suffit de
montrer que $e$ est universellement fermé après changement de base
par un morphisme $I \to \mathcal{I}_{\mathcal{X}/\mathcal{Y}}$
surjectif et lisse, où $I$ est un espace algébrique
(voir Propriétés des champs, section
\ref{stacks-properties-section-properties-morphisms}).
Ce changement de base $e' : X' \to I'$ est un
morphisme quasi-compact d'espaces algébriques qui
vérifie la partie d'existence du critère valuatif;
il est donc universellement fermé d'après
Morphismes d'espaces, Lemme
\ref{spaces-morphisms-lemma-quasi-compact-existence-universally-closed}.
\end{proof}






\section{Critère valuatif pour la diagonale}
\label{section-valuative-diagonal}

\noindent
Le résultat est le Lemme \ref{lemma-uniqueness-and-diagonal}.
Nous commençons par énoncer et démontrer un lemme auxiliaire formel.

\begin{lemma}
\label{lemma-helper-diagonal}
Soit $f : \mathcal{X} \to \mathcal{Y}$ un morphisme de champs algébriques.
Considérons le diagramme $2$-commutatif en flèches pleines
$$
\xymatrix{
\Spec(K) \ar[rr]_-x \ar[d]_j & &
\mathcal{X} \ar[d]^{\Delta_f} \\
\Spec(A) \ar[rr]^{(a_1, a_2, \varphi)} \ar@{..>}[rru] & &
\mathcal{X} \times_\mathcal{Y} \mathcal{X}
}
$$
où $A$ est un anneau de valuation de corps des fractions $K$. Soit
$\gamma : (a_1, a_2, \varphi) \circ j \longrightarrow \Delta_f \circ x$
un $2$-morphisme témoignant de la $2$-commutativité du diagramme.
Alors :
\begin{enumerate}
\item En écrivant $\gamma = (\alpha_1, \alpha_2)$ avec
$\alpha_i : a_i \circ j \to x$, on obtient deux flèches pointillées
$(a_1, \alpha_1, \text{id})$ et $(a_2, \alpha_2, \varphi)$ dans
le diagramme
$$
\xymatrix{
\Spec(K) \ar[r]_-x \ar[d]_j & \mathcal{X} \ar[d]^f \\
\Spec(A) \ar[r]^-{f \circ a_1} \ar@{..>}[ru] & \mathcal{Y}
}
$$
\item La catégorie des flèches pointillées associée au diagramme initial
et à $\gamma$ est un setoïde; l'ensemble des classes d'isomorphisme
de ses objets s'identifie à l'ensemble des morphismes
$(a_1, \alpha_1, \text{id}) \to (a_2, \alpha_2, \varphi)$ de
la catégorie des flèches pointillées.
\end{enumerate}
\end{lemma}

\begin{proof}
Comme $\Delta_f$ est représentable par des espaces algébriques (la diagonale
de $\Delta_f$ est donc séparée), le Lemme
\ref{lemma-cat-dotted-arrows} montre que la catégorie des flèches pointillées
du premier diagramme commutatif du lemme est un setoïde.
Toutes les autres assertions du lemme résultent de calculs
de diagrammes dans une $2$-catégorie, que nous omettons.
\end{proof}

\begin{lemma}
\label{lemma-uniqueness-and-diagonal}
Soit $f : \mathcal{X} \to \mathcal{Y}$ un morphisme de champs algébriques.
Supposons que $f$ soit quasi-séparé.
Si $f$ vérifie la partie d'unicité du critère valuatif,
alors $f$ est séparé.
\end{lemma}

\begin{proof}
L'hypothèse sur $f$ signifie que $\Delta_f$ est quasi-compact
et quasi-séparé (Définition \ref{definition-separated}).
Il nous faut montrer que $\Delta_f$ est propre.
Le Lemme \ref{lemma-setoids-and-diagonal} montre que $\Delta_f$
est séparé. D'après le Lemme \ref{lemma-properties-diagonal},
on sait que $\Delta_f$ est localement de type fini.
Pour achever la preuve, il reste à montrer que
$\Delta_f$ est universellement fermé. Un argument formel
(voir le Lemme \ref{lemma-helper-diagonal}) montre que
la partie d'unicité du critère valuatif entraîne
l'existence d'une flèche pointillée dans tout diagramme en flèches pleines de la forme suivante :
$$
\xymatrix{
\Spec(K) \ar[d] \ar[r] & \mathcal{X} \ar[d]^{\Delta_f} \\
\Spec(A) \ar[r] \ar@{..>}[ru] & \mathcal{X} \times_\mathcal{Y} \mathcal{X}
}
$$
Cette propriété étant préservée par tout changement de base,
tout changement de base de $\Delta_f$ par un espace algébrique
muni d'un morphisme vers $\mathcal{X} \times_\mathcal{Y} \mathcal{X}$
vérifie la partie d'existence du critère valuatif et
est donc universellement fermé d'après le critère valuatif
pour les morphismes d'espaces algébriques; voir
Morphismes d'espaces, Lemme
\ref{spaces-morphisms-lemma-quasi-compact-existence-universally-closed}.
\end{proof}

\noindent
Voici une réciproque.

\begin{lemma}
\label{lemma-converse-uniqueness-and-diagonal}
Soit $f : \mathcal{X} \to \mathcal{Y}$ un morphisme de champs algébriques.
Si $f$ est séparé, alors $f$ vérifie la
partie d'unicité du critère valuatif.
\end{lemma}

\begin{proof}
Puisque $f$ est séparé, toutes les catégories de flèches pointillées
sont des setoïdes d'après le Lemme \ref{lemma-cat-dotted-arrows}.
Considérons un diagramme
$$
\xymatrix{
\Spec(K) \ar[r]_-x \ar[d]_j & \mathcal{X} \ar[d]^f \\
\Spec(A) \ar[r]^-y \ar@{..>}[ru] & \mathcal{Y}
}
$$
et un $2$-morphisme $\gamma : y \circ j \to f \circ x$ comme dans la
Définition \ref{definition-fill-in-diagram}. Considérons deux
objets $(a, \alpha, \beta)$ et $(a', \beta', \alpha')$
de la catégorie des flèches pointillées. Pour achever la preuve, il
faut montrer que ces objets sont isomorphes. L'isomorphisme
$$
f \circ a \xrightarrow{\beta^{-1}} y \xrightarrow{\beta'} f \circ a'
$$
signifie que $(a, a', \beta' \circ \beta^{-1})$ définit un morphisme
$\Spec(A) \to \mathcal{X} \times_\mathcal{Y} \mathcal{X}$.
D'autre part, $\alpha$ et $\alpha'$ définissent
une $2$-flèche
$$
(a, a', \beta' \circ \beta^{-1}) \circ j =
(a \circ j, a' \circ j,
(\beta' \star \text{id}_j) \circ (\beta \star \text{id}_j)^{-1})
\xrightarrow{(\alpha, \alpha')} (x, x, \text{id}) = \Delta_f \circ x
$$
On utilise ici que $(a, \alpha, \beta)$ et $(a', \alpha', \beta')$
sont toutes deux des flèches pointillées associées à $\gamma$.
On obtient un diagramme commutatif
$$
\xymatrix{
\Spec(K) \ar[d]_j \ar[rr]_x & & \mathcal{X} \ar[d]^{\Delta_f} \\
\Spec(A) \ar[rr]^{(a, a', \beta' \circ \beta^{-1})} & &
\mathcal{X} \times_\mathcal{Y} \mathcal{X}
}
$$
dont la $2$-commutativité est donnée par $(\alpha, \alpha')$. Or
$\Delta_f$ est représentable par des espaces algébriques
(Lemme \ref{lemma-properties-diagonal})
et propre puisque $f$ est séparé. Le
Lemme \ref{lemma-existence-representable}
et le critère valuatif de propreté pour les espaces algébriques
(Morphismes d'espaces, Lemme \ref{spaces-morphisms-lemma-characterize-proper})
montrent donc qu'il existe une flèche pointillée.
En explicitant la construction, on voit que cela signifie que
$(a, \alpha, \beta)$ et $(a', \alpha', \beta')$
sont isomorphes dans la catégorie des flèches pointillées, comme souhaité.
\end{proof}





 
\section{Critère valuatif pour les morphismes universellement fermés}
\sectionmark{Critère valuatif de fermeture universelle}
\label{section-valutive-criterion}

\noindent
Voici un énoncé.

\begin{lemma}
\label{lemma-quasi-compact-existence-universally-closed}
Soit $f : \mathcal{X} \to \mathcal{Y}$ un morphisme de champs algébriques.
Supposons que
\begin{enumerate}
\item $f$ soit quasi-compact, et que
\item $f$ vérifie la partie d'existence du critère valuatif.
\end{enumerate}
Alors $f$ est universellement fermé.
\end{lemma}

\begin{proof}
D'après les Lemmes \ref{lemma-base-change-quasi-compact}
et \ref{lemma-base-change-existence},
les propriétés (1) et (2) sont préservées par
tout changement de base. D'après le Lemme \ref{lemma-universally-closed-local},
il suffit de montrer que $|T \times_\mathcal{Y} \mathcal{X}| \to |T|$
est fermée chaque fois que $T$ est un schéma affine muni d'un morphisme vers $\mathcal{Y}$.
Il suffit donc de montrer ceci : si $f : \mathcal{X} \to Y$ est un
morphisme quasi-compact d'un champ algébrique vers un schéma affine
qui vérifie la partie d'existence du critère valuatif,
alors $|f|$ est fermée. Soit $T \subset |\mathcal{X}|$ une partie fermée.
Pour conclure la démonstration, il faut montrer que $f(T)$ est fermé.

\medskip\noindent
Soit $\mathcal{Z} \subset \mathcal{X}$ le sous-champ algébrique réduit
induit par $T$ (Propriétés des champs,
Définition \ref{stacks-properties-definition-reduced-induced-stack}).
Alors $i : \mathcal{Z} \to \mathcal{X}$ est une immersion fermée
et il faut montrer que l'image de $|\mathcal{Z}| \to |Y|$
est fermée. Comme une immersion fermée est quasi-compacte
(Lemme \ref{lemma-closed-immersion-quasi-compact})
et vérifie la partie d'existence du critère valuatif
(Lemme \ref{lemma-closed-immersion-valuative-criteria}),
comme le composé de morphismes quasi-compacts est quasi-compact
(Lemme \ref{lemma-composition-quasi-compact})
et comme la composition préserve la propriété de vérifier
la partie d'existence du critère valuatif
(Lemme \ref{lemma-composition-existence}),
nous en concluons qu'il suffit de montrer ceci : si $f : \mathcal{X} \to Y$
est un morphisme quasi-compact d'un champ algébrique vers un schéma affine
qui vérifie la partie d'existence du critère valuatif,
alors $|f|(|\mathcal{X}|)$ est fermé.

\medskip\noindent
Puisque $\mathcal{X}$ est quasi-compact (car quasi-compact au-dessus du
schéma affine $Y$), nous pouvons choisir un schéma affine $U$ et un morphisme
lisse et surjectif $U \to \mathcal{X}$
(Propriétés des champs, Lemme
\ref{stacks-properties-lemma-quasi-compact-stack}).
Supposons que $y \in Y$ appartienne à l'adhérence de l'image de $U \to Y$
(autrement dit, à l'adhérence de l'image de $|f|$).
Alors, d'après
Morphismes, Lemme \ref{morphisms-lemma-reach-points-scheme-theoretic-image},
il existe un anneau de valuation $A$ de corps des fractions $K$
et un diagramme commutatif
$$
\xymatrix{
\Spec(K) \ar[r] \ar[d] & U \ar[d] \\
\Spec(A) \ar[r] & Y
}
$$
tel que le point fermé de $\Spec(A)$ ait pour image $y$. Par hypothèse,
il existe une extension $K'/K$, un anneau de valuation $A' \subset K'$
dominant $A$, ainsi que la flèche pointillée du diagramme suivant :
$$
\xymatrix{
\Spec(K') \ar[r] \ar[d] &
\Spec(K) \ar[r] \ar[d] &
U \ar[d] \ar[r] &
\mathcal{X} \ar[d]^f \\
\Spec(A') \ar[r] \ar@{..>}[rrru] &
\Spec(A) \ar[r] &
Y  \ar@{=}[r] &
Y
}
$$
Ainsi, $y$ appartient à l'image de $|f|$, ce qui conclut.
\end{proof}

\noindent
Voici une réciproque.

\begin{lemma}
\label{lemma-converse-existence-universally-closed}
Soit $f : \mathcal{X} \to \mathcal{Y}$ un morphisme de champs algébriques.
Supposons que
\begin{enumerate}
\item $f$ soit quasi-séparé, et que
\item $f$ soit universellement fermé.
\end{enumerate}
Alors $f$ vérifie la partie d'existence du critère valuatif.
\end{lemma}

\begin{proof}
Considérons un diagramme en flèches pleines
$$
\xymatrix{
\Spec(K) \ar[r]_-x \ar[d]_j & \mathcal{X} \ar[d]^f \\
\Spec(A) \ar[r]^-y \ar@{..>}[ru] & \mathcal{Y}
}
$$
où $A$ est un anneau de valuation de corps des fractions $K$
et où $\gamma : y \circ j \longrightarrow f \circ x$ est comme dans la
Définition \ref{definition-fill-in-diagram}. D'après le
Lemme \ref{lemma-cat-dotted-arrows-base-change},
pour trouver une flèche pointillée (après avoir éventuellement remplacé
$K$ par une extension et $A$ par un anneau de valuation le dominant),
nous pouvons remplacer $\mathcal{Y}$ par $\Spec(A)$ et $\mathcal{X}$
par $\Spec(A) \times_\mathcal{Y} \mathcal{X}$. Bien entendu,
nous utilisons ici le fait que les propriétés d'être
quasi-séparé et universellement fermé sont préservées par changement de base.
Nous sommes ainsi ramenés au cas traité au paragraphe suivant.

\medskip\noindent
Considérons un diagramme en flèches pleines
$$
\xymatrix{
\Spec(K) \ar[r]_-x \ar[d]_j & \mathcal{X} \ar[d]^f \\
\Spec(A) \ar@{=}[r] \ar@{..>}[ru] & \Spec(A)
}
$$
où $A$ est un anneau de valuation de corps des fractions $K$, comme dans la
Définition \ref{definition-fill-in-diagram}.
D'après le Lemme \ref{lemma-quasi-compact-permanence} et puisque
$f$ est quasi-séparé, nous savons que
le morphisme $x$ est quasi-compact.
Comme $f$ est universellement fermé, nous savons en particulier
que $|f|(\overline{\{x\}})$ est fermé dans $\Spec(A)$.
Puisque cette image contient le point générique de $\Spec(A)$,
il existe un point $x' \in |\mathcal{X}|$ dans l'adhérence
de $x$ qui s'envoie sur le point fermé de $\Spec(A)$.
D'après le Lemme \ref{lemma-reach-points-scheme-theoretic-image},
il existe un diagramme commutatif
$$
\xymatrix{
\Spec(K') \ar[r] \ar[d] & \Spec(K) \ar[d] \\
\Spec(A') \ar[r] & \mathcal{X}
}
$$
tel que le point fermé de $\Spec(A')$ s'envoie sur $x' \in |\mathcal{X}|$.
Il s'ensuit que $\Spec(A') \to \Spec(A)$ envoie le point fermé
sur le point fermé, c'est-à-dire que $A'$ domine $A$, ce qui achève la démonstration.
\end{proof}





\section{Critère valuatif de propreté}
\label{section-valutive-criterion-properness}

\noindent
Voici l'énoncé.

\begin{lemma}
\label{lemma-criterion-proper}
Soit $f : \mathcal{X} \to \mathcal{Y}$ un morphisme de champs algébriques.
Supposons que $f$ soit de type fini et quasi-séparé.
Alors les conditions suivantes sont équivalentes :
\begin{enumerate}
\item $f$ est propre, et
\item $f$ vérifie à la fois les parties d'unicité et d'existence
du critère valuatif.
\end{enumerate}
\end{lemma}

\begin{proof}
Un morphisme propre est exactement un morphisme séparé, de type fini et
universellement fermé. Le lemme résulte donc des Lemmes
\ref{lemma-uniqueness-and-diagonal},
\ref{lemma-converse-uniqueness-and-diagonal},
\ref{lemma-quasi-compact-existence-universally-closed} et
\ref{lemma-converse-existence-universally-closed}.
\end{proof}





 
 
\section{Morphismes localement d'intersection compl\`ete}
\label{section-lci}

\noindent
La propri\'et\'e ``\^etre un morphisme localement d'intersection compl\`ete''
des morphismes d'espaces alg\'ebriques est
locale pour la topologie lisse sur la source et sur le but; voir
Descente sur les espaces, Lemme \ref{spaces-descent-lemma-smooth-local-source-target}
et
Compl\'ements sur les morphismes d'espaces, Lemmes
\ref{spaces-more-morphisms-lemma-descending-property-lci} et
\ref{spaces-more-morphisms-lemma-lci-syntomic-local-source}.
D'apr\`es le Lemme \ref{lemma-local-source-target} ci-dessus, nous pouvons d\'efinir
comme suit la notion de morphisme localement d'intersection compl\`ete
pour les champs alg\'ebriques; elle co\"incide avec la
notion d\'ej\`a d\'efinie dans
Compl\'ements sur les morphismes d'espaces,
section \ref{spaces-more-morphisms-section-lci},
lorsque la source et le but sont tous deux des espaces alg\'ebriques.

\begin{definition}
\label{definition-lci}
Soit $f : \mathcal{X} \to \mathcal{Y}$ un morphisme de champs alg\'ebriques.
On dit que $f$ est un {\it morphisme localement d'intersection compl\`ete} ou
{\it de Koszul} si les conditions \`equivalentes du
Lemme \ref{lemma-local-source-target}
sont satisfaites pour $\mathcal{P} = \text{localement d'intersection compl\`ete}$.
\end{definition}

\begin{lemma}
\label{lemma-composition-lci}
Le compos\'e de morphismes localement d'intersection compl\`ete est
un morphisme localement d'intersection compl\`ete.
\end{lemma}

\begin{proof}
Le résultat découle de
la Remarque \ref{remark-composition}
et du
Lemme \ref{spaces-more-morphisms-lemma-composition-lci}
de Compl\'ements sur les morphismes d'espaces.
\end{proof}

\begin{lemma}
\label{lemma-flat-base-change-lci}
Un changement de base plat d'un morphisme localement d'intersection compl\`ete est
un morphisme localement d'intersection compl\`ete.
\end{lemma}

\begin{proof}
D\'emonstration omise. Indication : raisonner exactement comme dans la Remarque \ref{remark-base-change}
(mais seulement lorsque $\mathcal{Y}' \to \mathcal{Y}$ est plat) en utilisant
Compl\'ements sur les morphismes d'espaces, Lemme
\ref{spaces-more-morphisms-lemma-flat-base-change-lci}.
\end{proof}

\begin{lemma}
\label{lemma-lci-permanence}
Soit
$$
\xymatrix{
\mathcal{X} \ar[rr]_f \ar[rd] & & \mathcal{Y} \ar[ld] \\
& \mathcal{Z}
}
$$
un diagramme commutatif de morphismes de champs alg\'ebriques.
Supposons que $\mathcal{Y} \to \mathcal{Z}$ soit lisse et que
$\mathcal{X} \to \mathcal{Z}$ soit un morphisme localement d'intersection compl\`ete.
Alors $f : \mathcal{X} \to \mathcal{Y}$ est un morphisme
localement d'intersection compl\`ete.
\end{lemma}

\begin{proof}
Choisissons un sch\'ema $W$ et un morphisme lisse surjectif $W \to \mathcal{Z}$.
Choisissons un sch\'ema $V$ et un morphisme lisse surjectif
$V \to W \times_\mathcal{Z} \mathcal{Y}$.
Choisissons un sch\'ema $U$ et un morphisme lisse surjectif
$U \to V \times_\mathcal{Y} \mathcal{X}$.
Alors $U \to W$ est un morphisme localement d'intersection compl\`ete de sch\'emas et
$V \to W$ est un morphisme lisse de sch\'emas. Par le r\'esultat pour les sch\'emas
(Compl\'ements sur les morphismes, Lemme \ref{more-morphisms-lemma-lci-permanence}),
nous concluons que $U \to V$ est un morphisme localement d'intersection compl\`ete.
Par d\'efinition, cela signifie que $f$ est un morphisme localement d'intersection compl\`ete.
\end{proof}





\section{Morphismes préservant les stabilisateurs}
\label{section-stabilizer-preserving}

\noindent
Dans la littérature, on dit qu'un morphisme $f : \mathcal{X} \to \mathcal{Y}$ de
champs algébriques {\it préserve les stabilisateurs} ou
{\it reflète les points fixes} si le morphisme induit
$\mathcal{I}_\mathcal{X} \to
\mathcal{X} \times_\mathcal{Y} \mathcal{I}_\mathcal{Y}$
est un isomorphisme. Un tel morphisme induit un isomorphisme
entre les groupes d'automorphismes (remarque \ref{remark-identify-automorphism-groups})
en tout point de $\mathcal{X}$.
Dans cette section, nous démontrons quelques lemmes élémentaires sur cette notion.

\begin{lemma}
\label{lemma-stabilizer-preserving}
Soit $f : \mathcal{X} \to \mathcal{Y}$ un morphisme de champs algébriques.
Si $\mathcal{I}_\mathcal{X} \to
\mathcal{X} \times_\mathcal{Y} \mathcal{I}_\mathcal{Y}$ est un isomorphisme,
alors $f$ est représentable par des espaces algébriques.
\end{lemma}

\begin{proof}
Cela résulte immédiatement du lemme \ref{lemma-second-diagonal}.
\end{proof}

\begin{remark}
\label{remark-get-property-auts-from-diagonal}
Soit $f : \mathcal{X} \to \mathcal{Y}$ un morphisme de champs algébriques.
Soit $U \to \mathcal{X}$ un morphisme dont la source est un espace algébrique.
Soit $G \to H$ l'image réciproque du morphisme
$\mathcal{I}_\mathcal{X} \to
\mathcal{X} \times_\mathcal{Y} \mathcal{I}_\mathcal{Y}$
sur $U$. Si $\Delta_f$ est non ramifié, \'etale, etc.,
il en va de même de $G \to H$. En effet,
$$
\xymatrix{
U \times_\mathcal{X} U \ar[r] \ar[d] & \mathcal{X} \ar[d]^{\Delta_f} \\
U \times_\mathcal{Y} U \ar[r] & \mathcal{X} \times_\mathcal{Y} \mathcal{X}
}
$$
est cartésien et le morphisme $G \to H$ se déduit de la
flèche verticale de gauche par le changement de base défini par la diagonale $U \to U \times U$.
Comparer avec la démonstration du lemme \ref{lemma-separated-implies-isom}.
\end{remark}

\begin{lemma}
\label{lemma-aut-iso-unramified}
Soit $f : \mathcal{X} \to \mathcal{Y}$ un morphisme non ramifié
de champs algébriques. Les conditions suivantes sont équivalentes :
\begin{enumerate}
\item $\mathcal{I}_\mathcal{X} \to
\mathcal{X} \times_\mathcal{Y} \mathcal{I}_\mathcal{Y}$
est un isomorphisme, et
\item $f$ induit un isomorphisme entre les groupes d'automorphismes en $x$ et en $f(x)$
(remarque \ref{remark-identify-automorphism-groups}) pour tout
$x \in |\mathcal{X}|$.
\end{enumerate}
\end{lemma}

\begin{proof}
Choisissons un schéma $U$ et un morphisme lisse surjectif $U \to \mathcal{X}$.
Notons $G \to H$ l'image réciproque du morphisme
$\mathcal{I}_\mathcal{X} \to
\mathcal{X} \times_\mathcal{Y} \mathcal{I}_\mathcal{Y}$
sur $U$. D'après la remarque \ref{remark-get-property-auts-from-diagonal} et le
lemme \ref{lemma-characterize-unramified}, le morphisme $G \to H$ est \'etale.
La condition (1) équivaut à dire que
$G \to H$ est un isomorphisme (cela résulte par exemple de l'application du
lemme de Propriétés des champs
\ref{stacks-properties-lemma-check-property-covering}).
La condition (2) équivaut à dire que
pour tout $u \in U$, le morphisme de fibres $G_u \to H_u$
est un isomorphisme. Ainsi, l'implication (1) $\Rightarrow$ (2) est triviale.
Si (2) est satisfaite, alors $G \to H$ est un morphisme d'espaces algébriques surjectif, universellement injectif,
et \'etale. Un tel morphisme est un isomorphisme d'après le
lemme de Morphismes d'espaces
\ref{spaces-morphisms-lemma-etale-universally-injective-open}.
\end{proof}

\begin{lemma}
\label{lemma-stabilizer-preserving-unramified}
\begin{reference}
\cite[Proposition 3.5]{rydh_quotients} et
\cite[Proposition 2.5]{alper_quotient}
\end{reference}
Soit $f : \mathcal{X} \to \mathcal{Y}$ un morphisme de champs algébriques.
Supposons que
\begin{enumerate}
\item $f$ soit représentable par des espaces algébriques et non ramifié, et
\item $\mathcal{I}_\mathcal{Y} \to \mathcal{Y}$ soit propre.
\end{enumerate}
Alors l'ensemble des $x \in |\mathcal{X}|$ tels que $f$ induise un
isomorphisme entre les groupes d'automorphismes en $x$ et en $f(x)$
(remarque \ref{remark-identify-automorphism-groups}) est ouvert.
Notons $\mathcal{U} \subset \mathcal{X}$ le sous-champ ouvert correspondant ;
alors le morphisme
$\mathcal{I}_\mathcal{U} \to
\mathcal{U} \times_\mathcal{Y} \mathcal{I}_\mathcal{Y}$
est un isomorphisme.
\end{lemma}

\begin{proof}
Choisissons un schéma $U$ et un morphisme lisse surjectif $U \to \mathcal{X}$.
Notons $G \to H$ l'image réciproque du morphisme
$\mathcal{I}_\mathcal{X} \to
\mathcal{X} \times_\mathcal{Y} \mathcal{I}_\mathcal{Y}$
sur $U$. D'après la remarque \ref{remark-get-property-auts-from-diagonal} et le
lemme \ref{lemma-characterize-unramified}, le morphisme
$G \to H$ est \'etale. Comme $f$ est représentable par des espaces algébriques,
nous voyons que $G \to H$ est un monomorphisme. Ainsi, $G \to H$ est une immersion
ouverte, voir
le lemme de Morphismes d'espaces
\ref{spaces-morphisms-lemma-etale-universally-injective-open}.
Par hypothèse, $H \to U$ est propre.

\medskip\noindent
Ces préparatifs étant faits, démontrons le lemme.
L'image réciproque du sous-ensemble de $|\mathcal{X}|$ considéré dans le lemme est
manifestement l'ensemble des $u \in U$ tels que $G_u \to H_u$ soit un isomorphisme
(puisque, après tout, $G_u$ est un sous-groupe ouvert de l'espace algébrique en groupes $H_u$).
Cet ensemble est ouvert, car son complément est
l'image du fermé $|H| \setminus |G|$, et l'application $|H| \to |U|$
est fermée. D'après le
lemme de Propriétés des champs \ref{stacks-properties-lemma-open-substacks},
nous pouvons considérer le sous-champ ouvert correspondant $\mathcal{U}$ de $\mathcal{X}$.
La dernière assertion du lemme découle de l'application du
lemme \ref{lemma-aut-iso-unramified}
à $\mathcal{U} \to \mathcal{Y}$.
\end{proof}

\begin{lemma}
\label{lemma-base-change-stabilizer-preserving}
Soit
$$
\xymatrix{
\mathcal{X}' \ar[r] \ar[d]_{f'} & \mathcal{X} \ar[d]^f \\
\mathcal{Y}' \ar[r] & \mathcal{Y}
}
$$
un diagramme cartésien de champs algébriques.
\begin{enumerate}
\item Soit $x' \in |\mathcal{X}'|$, d'image $x \in |\mathcal{X}|$.
Si $f$ induit un isomorphisme entre les groupes d'automorphismes en
$x$ et en $f(x)$ (remarque \ref{remark-identify-automorphism-groups}), alors
$f'$ induit un isomorphisme entre les groupes d'automorphismes en $x'$ et en $f(x')$.
\item Si $\mathcal{I}_\mathcal{X} \to
\mathcal{X} \times_\mathcal{Y} \mathcal{I}_\mathcal{Y}$ est un isomorphisme,
alors $\mathcal{I}_{\mathcal{X}'} \to
\mathcal{X}' \times_{\mathcal{Y}'} \mathcal{I}_{\mathcal{Y}'}$
est un isomorphisme.
\end{enumerate}
\end{lemma}

\begin{proof}
Démonstration omise.
\end{proof}

\begin{lemma}
\label{lemma-stabilizer-preserving-points-cartesian}
Soit
$$
\xymatrix{
\mathcal{X}' \ar[r] \ar[d]_{f'} & \mathcal{X} \ar[d]^f \\
\mathcal{Y}' \ar[r]^g & \mathcal{Y}
}
$$
un diagramme cartésien de champs algébriques. Si $f$ induit en tout point un isomorphisme
entre les groupes d'automorphismes
(remarque \ref{remark-identify-automorphism-groups}),
alors l'application
$$
\Mor(\Spec(k), \mathcal{X}')
\longrightarrow
\Mor(\Spec(k), \mathcal{Y}') \times \Mor(\Spec(k), \mathcal{X})
$$
est injective sur les classes d'isomorphisme pour tout corps $k$.
\end{lemma}

\begin{proof}
Il s'agit de montrer qu'étant donné $(y', x)$, il existe au plus un $x'$
qui ait cette image.
D'après notre construction des $2$-produits fibrés, un morphisme
$x'$ est donné par un triplet $(x, y', \alpha)$,
où $\alpha : g \circ y' \to f \circ x$ est un $2$-morphisme.
Supposons maintenant donné un second triplet $(x, y', \beta)$ de cette sorte.
Alors $\alpha$ et $\beta$ diffèrent par un point à valeurs dans $k$, noté
$\epsilon$, de l'espace algébrique en groupes d'automorphismes $G_{f(x)}$.
Puisque, par hypothèse, $f$ induit un isomorphisme $G_x \to G_{f(x)}$,
nous pouvons relever $\epsilon$ en un point à valeurs dans $k$, noté
$\gamma$, de $G_x$. Alors $(\gamma, \text{id}) : (x, y', \alpha) \to
(x, y', \beta)$ est l'isomorphisme recherché.
\end{proof}

\begin{lemma}
\label{lemma-etale-iso}
Soit $f : \mathcal{X} \to \mathcal{Y}$ un morphisme de champs algébriques.
Supposons que $f$ soit \'etale, que $f$ induise un isomorphisme
entre les groupes d'automorphismes en tout point
(remarque \ref{remark-identify-automorphism-groups}),
et que, pour tout corps algébriquement clos $k$, le foncteur
$$
f : \Mor(\Spec(k), \mathcal{X}) \longrightarrow \Mor(\Spec(k), \mathcal{Y})
$$
soit une équivalence. Alors $f$ est un isomorphisme.
\end{lemma}

\begin{proof}
D'après le lemme \ref{lemma-universally-injective}, $f$ est
universellement injectif. En combinant les
lemmes \ref{lemma-stabilizer-preserving} et
\ref{lemma-aut-iso-unramified},
nous voyons que $f$ est représentable par des espaces algébriques.
Par conséquent, $f$ est une immersion ouverte d'après le lemme de Morphismes d'espaces
\ref{spaces-morphisms-lemma-etale-universally-injective-open}.
Pour conclure, remarquons que la condition de l'énoncé garantit également
que $f$ est surjectif.
\end{proof}






 
\section{Normalisation}
\label{section-normalization}

\noindent
Cette section est l'analogue de
Morphismes d'espaces, section \ref{spaces-morphisms-section-normalization}.

\begin{lemma}
\label{lemma-prepare-normalization}
Soit $\mathcal{X}$ un champ algébrique. Les conditions suivantes sont équivalentes :
\begin{enumerate}
\item il existe un morphisme lisse surjectif $U \to \mathcal{X}$, où $U$
est un schéma, tel que tout ouvert quasi-compact de $U$ ait
un nombre fini de composantes irréductibles,
\item pour tout schéma $U$ et tout morphisme lisse
$U \to \mathcal{X}$, tout ouvert quasi-compact de $U$ a un nombre fini de
composantes irréductibles,
\item pour tout espace algébrique $Y$ et tout morphisme lisse $Y \to \mathcal{X}$,
l'espace $Y$ satisfait aux conditions équivalentes de
Morphismes d'espaces, Lemme \ref{spaces-morphisms-lemma-prepare-normalization},
et
\item pour tout champ algébrique quasi-compact $\mathcal{Y}$ lisse sur
$\mathcal{X}$, l'espace $|\mathcal{Y}|$ a un nombre fini de composantes
irréductibles.
\end{enumerate}
\end{lemma}

\begin{proof}
L'équivalence de (1), (2) et (3) résulte de
Descente, Lemme \ref{descent-lemma-locally-finite-nr-irred-local-fppf},
Propriétés des champs, Lemme \ref{stacks-properties-lemma-type-property}, et de
Morphismes d'espaces, Lemme \ref{spaces-morphisms-lemma-prepare-normalization}.
Ces références montrent aussi clairement que la condition (4)
entraîne la condition (1). Réciproquement, supposons que les conditions équivalentes
(1), (2) et (3) soient satisfaites, et soit $\mathcal{Y} \to \mathcal{X}$ un
morphisme lisse de champs algébriques, avec $\mathcal{Y}$
quasi-compact. Nous pouvons alors choisir un schéma affine $V$ et un morphisme lisse
surjectif $V \to \mathcal{Y}$ d'après Propriétés des champs, Lemme
\ref{stacks-properties-lemma-quasi-compact-stack}.
Comme $V$ a un nombre fini de composantes irréductibles d'après (2), et que
l'application $|V| \to |\mathcal{Y}|$ est surjective et continue, on en conclut que
$|\mathcal{Y}|$ a un nombre fini de composantes irréductibles d'après Topologie, Lemme
\ref{topology-lemma-surjective-continuous-irreducible-components}.
\end{proof}

\begin{lemma}
\label{lemma-normalization}
Soit $\mathcal{X}$ un champ algébrique satisfaisant aux
conditions équivalentes du Lemme \ref{lemma-prepare-normalization}.
Alors il existe un morphisme intégral de champs algébriques
$$
\mathcal{X}^\nu \longrightarrow \mathcal{X}
$$
tel que, pour tout schéma $U$ et tout morphisme lisse $U \to \mathcal{X}$,
le produit fibré $\mathcal{X}^\nu \times_\mathcal{X} U$
soit le normalisé de $U$.
\end{lemma}

\begin{proof}
Soit $U \to \mathcal{X}$ un morphisme lisse surjectif, où $U$ est un schéma.
Posons $R = U \times_\mathcal{X} U$. Rappelons que l'on obtient un groupoïde lisse
$(U, R, s, t, c)$ en espaces algébriques et une présentation
$\mathcal{X} = [U/R]$ de $\mathcal{X}$ ; voir Champs algébriques,
Lemmes \ref{algebraic-lemma-map-space-into-stack} et
\ref{algebraic-lemma-stack-presentation}, ainsi que la
Définition \ref{algebraic-definition-presentation}.
L'hypothèse sur $\mathcal{X}$ signifie que le normalisé
$U^\nu$ de $U$ est défini ; voir Morphismes, Définition
\ref{morphisms-definition-normalization}.
D'après Morphismes d'espaces, Lemme
\ref{spaces-morphisms-lemma-normalization},
la formation du normalisé commute au changement de base par un morphisme lisse
d'espaces algébriques. Ainsi, le normalisé $R^\nu$ de $R$
est isomorphe tant à $R \times_{s, U} U^\nu$ qu'à $U^\nu \times_{U, t} R$.
Nous obtenons donc deux morphismes lisses
$s^\nu : R^\nu \to U^\nu$ et $t^\nu : R^\nu \to U^\nu$
d'espaces algébriques. Un calcul formel sur les produits fibrés montre que
$R^\nu \times_{s^\nu, U^\nu, t^\nu} R^\nu$ est le normalisé
de $R \times_{s, U, t} R$. Par conséquent, le morphisme lisse
$c : R \times_{s, U, t} R \to R$ se prolonge lui aussi en $c^\nu$.
De même, l'inverse $i : R \to R$ (qui est un isomorphisme) induit
un isomorphisme $i^\nu : R^\nu \to R^\nu$. Enfin, l'identité
$e : U \to R$ admet un relèvement $e^\nu : U^\nu \to R^\nu$, par exemple
parce que $e$ est une section de $s$ et que $R^\nu = R \times_{U, s} U^\nu$.
Nous affirmons que $(U^\nu, R^\nu, s^\nu, t^\nu, c^\nu)$ est
un groupoïde lisse en espaces algébriques. Pour le voir, il faut vérifier
les axiomes (1), (2)(a), (2)(b), (3)(a) et (3)(b) de
Groupoïdes, section \ref{groupoids-section-groupoids}
pour $(U^\nu, R^\nu, s^\nu, t^\nu, c^\nu, e^\nu, i^\nu)$.
Pour (1), par exemple, il faut vérifier que les deux morphismes
$a, b : R^\nu \times_{s^\nu, U^\nu, t^\nu} R^\nu
\times_{s^\nu, U^\nu, t^\nu} R^\nu \to R^\nu$ que l'on obtient coïncident.
Cela résulte du fait que les deux morphismes correspondants
$R \times_{s, U, t} R \times_{s, U, t} R \to R$
coïncident, et que les morphismes $a$ et $b$ sont les prolongements uniques
de ce morphisme aux normalisés correspondants. Les autres
axiomes se traitent de même.

\medskip\noindent
Considérons le champ algébrique $\mathcal{X}^\nu = [U^\nu/R^\nu]$
(Champs algébriques, Théorème
\ref{algebraic-theorem-smooth-groupoid-gives-algebraic-stack}).
Puisque nous disposons d'un morphisme
$(U^\nu, R^\nu, s^\nu, t^\nu, c^\nu) \to (U, R, s, t, c)$
de groupoïdes en espaces algébriques, nous obtenons un morphisme
$\nu : \mathcal{X}^\nu \to \mathcal{X}$ de champs algébriques.
Comme $R^\nu = R \times_{s, U} U^\nu$, on voit que
$U^\nu = \mathcal{X}^\nu \times_\mathcal{X} U$ d'après
Groupoïdes en espaces, Lemme
\ref{spaces-groupoids-lemma-criterion-fibre-product}.
En particulier, puisque $U^\nu \to U$ est intégral, $\nu$ est intégral.
Nous omettons de vérifier que la propriété de changement de base
énoncée dans le lemme est satisfaite pour tout morphisme lisse d'un schéma vers
$\mathcal{X}$.
\end{proof}

\noindent
Ceci conduit à la définition suivante.

\begin{definition}
\label{definition-normalization}
Soit $\mathcal{X}$ un champ algébrique satisfaisant aux
conditions équivalentes du Lemme \ref{lemma-prepare-normalization}.
Nous définissons la {\it normalisation} de $\mathcal{X}$ comme le morphisme
$$
\nu : \mathcal{X}^\nu \longrightarrow \mathcal{X}
$$
construit dans le Lemme \ref{lemma-normalization}.
\end{definition}
 
  
   
    
     
      




\section{Points et spécialisations}
\label{section-specializations}

\noindent
Cette section est l'analogue de la section
\ref{decent-spaces-section-specializations} d'Espaces décents.

\begin{lemma}
\label{lemma-specialization}
Soit $\mathcal{X}$ un champ algébrique. Soit $f : U \to \mathcal{X}$
un morphisme lisse, où $U$ est un espace algébrique. Soit
$x' \leadsto x$ une spécialisation de points de $|\mathcal{X}|$.
Soit $u \in |U|$ tel que $f(u) = x$. Si $(\mathcal{X}, x')$ vérifie les
conditions équivalentes de Propriétés des champs,
Lemme \ref{stacks-properties-lemma-UR-quasi-compact-above-x},
alors il existe une spécialisation
$u' \leadsto u$ dans $|U|$ telle que $f(u') = x'$.
\end{lemma}

\begin{proof}
Choisissons un morphisme \'etale $(U_1, u_1) \to (U, u)$, où $U_1$
est un schéma affine. Remplaçons alors $U$ par $U_1$.
Nous pouvons donc supposer que $U$ est un schéma affine.
Considérons l'espace algébrique $R = U \times_\mathcal{X} U$
muni des projections lisses $t, s : R \to U$.
Choisissons un point $w \in U$ d'image $x'$ ; cela est possible puisque
l'application $f : |U| \to |\mathcal{X}|$ est ouverte.
D'après notre hypothèse sur $x'$, la fibre
$F' = t^{-1}(w) = R \times_{t, U} w$
de $t : R \to U$ au-dessus de $w$ est un espace algébrique quasi-compact.
Choisissons un schéma affine $T$ et un morphisme \'etale surjectif
$T \to F'$. Le fait que $x' \leadsto x$ signifie que $u$ appartient à
l'adhérence de l'image du morphisme
$$
T \to F' \to R \xrightarrow{s} U
$$
En effet, cette image est la fibre de $|U| \to |\mathcal{X}|$
au-dessus de $x'$ ; si $u \in V \subset |U|$, pour un ouvert disjoint
de cette fibre, alors $f(V)$ est un voisinage ouvert de $x$
qui ne contient pas $x'$ ; contradiction.
Ainsi, d'après Morphismes, Lemme
\ref{morphisms-lemma-reach-points-scheme-theoretic-image},
il existe $u' \in |U|$, dans la fibre
de $|U| \to |\mathcal{X}|$ au-dessus de $x'$, qui se spécialise en $u$.
\end{proof}


\section{Champs algébriques décents}
\label{section-reasonable-decent}

\noindent
Cette section est l'analogue de
Espaces décents, section \ref{decent-spaces-section-reasonable-decent}.
En particulier, la définition suivante est compatible avec
la notion d'espace algébrique décent qui y est définie.

\begin{definition}
\label{definition-decent}
Soit $\mathcal{X}$ un champ algébrique. Nous disons que $\mathcal{X}$ est
{\it décent} si, pour tout $x \in |\mathcal{X}|$, les conditions équivalentes
de Propriétés des champs, Lemme
\ref{stacks-properties-lemma-UR-quasi-compact-above-x} sont satisfaites.
\end{definition}

\noindent
Certains reformuleraient cette définition en disant que
tout point de $\mathcal{X}$ est quasi-compact.
Une condition un peu plus forte consisterait à demander que tout morphisme
$\Spec(k) \to \mathcal{X}$ appartenant à la classe d'équivalence de $x$ soit
quasi-séparé aussi bien que quasi-compact.

\begin{lemma}
\label{lemma-quasi-separated-decent}
Un champ algébrique quasi-séparé $\mathcal{X}$ est décent.
Plus généralement, si $\Delta : \mathcal{X} \to \mathcal{X} \times \mathcal{X}$
est quasi-compact, alors $\mathcal{X}$ est décent.
\end{lemma}

\begin{proof}
En effet, si $\mathcal{X}$ est quasi-séparé, alors
tout morphisme $f : T \to \mathcal{X}$ dont la source est un schéma
quasi-compact $T$ est quasi-compact, voir
Lemme \ref{lemma-quasi-compact-permanence}.
Si l'on sait que $\Delta$ est quasi-compact, on utilise la description
$$
T \times_{f, \mathcal{X}, f'} T' =
(T \times T')
\times_{(f, f'), \mathcal{X} \times \mathcal{X}, \Delta}
\mathcal{X}
$$
pour le démontrer. Nous omettons les détails.
\end{proof}

\begin{lemma}
\label{lemma-representable-decent}
Soit $f : \mathcal{X} \to \mathcal{Y}$ un morphisme
de champs algébriques. Supposons que $\mathcal{Y}$ soit décent
et que $f$ soit représentable (par des schémas), ou que $f$ soit représentable
par des espaces algébriques et quasi-séparé.
Alors $\mathcal{X}$ est décent.
\end{lemma}

\begin{proof}
Soit $x \in |\mathcal{X}|$, d'image $y \in |\mathcal{Y}|$.
Choisissons un morphisme $y : \Spec(k) \to \mathcal{Y}$
dans la classe d'équivalence définissant $y$.
Posons $\mathcal{X}_y = \Spec(k) \times_{y, \mathcal{Y}} \mathcal{X}$.
Choisissons un point $x' \in |\mathcal{X}_y|$ d'image $x$, voir
Propriétés des champs, Lemme \ref{stacks-properties-lemma-points-cartesian}.
Choisissons un morphisme $x' : \Spec(k') \to \mathcal{X}_y$ dans la classe
d'équivalence de $x'$. Considérons le diagramme
$$
\xymatrix{
\Spec(k') \ar[r]_{x'} &
\mathcal{X}_y \ar[r] \ar[d] &
\mathcal{X} \ar[d] \\
&
\Spec(k) \ar[r]^y &
\mathcal{Y}
}
$$
Le morphisme $y$ est quasi-compact puisque $\mathcal{Y}$ est décent.
Par conséquent, $\mathcal{X}_y \to \mathcal{X}$ est quasi-compact, car il
s'obtient par changement de base
(Lemme \ref{lemma-base-change-quasi-compact}).
Pour conclure, il suffit donc de montrer que $x'$ est quasi-compact
(Lemme \ref{lemma-composition-quasi-compact}).
Si $f$ est représentable, alors $\mathcal{X}_y$ est un schéma
et $x'$ est quasi-compact.
Si $f$ est représentable par des espaces algébriques et quasi-séparé,
alors $\mathcal{X}_y$ est un espace algébrique quasi-séparé
et donc décent (Espaces décents, Lemme
\ref{decent-spaces-lemma-properties-trivial-implications}).
\end{proof}

\begin{lemma}
\label{lemma-qc-decent}
Soit $f : \mathcal{X} \to \mathcal{Y}$ un morphisme de champs algébriques.
Si $f$ est quasi-compact et surjectif et si $\mathcal{X}$ est décent,
alors $\mathcal{Y}$ est décent.
\end{lemma}

\begin{proof}
Soit $x : \Spec(k) \to \mathcal{X}$ un morphisme, où $k$ est un corps,
et notons $y = f \circ x$. Puisque $f$ est surjectif, tout point
de $|\mathcal{Y}|$ s'obtient de cette manière, voir Propriétés des champs,
Lemme \ref{stacks-properties-lemma-characterize-surjective}.
Soient $T$ un schéma affine et $T \to \mathcal{Y}$ un morphisme.
Il suffit de montrer que
$T \times_{\mathcal{Y}, y} \Spec(k)$ est quasi-compact, voir
Lemme \ref{lemma-check-quasi-compact-covering}.
Nous avons
$$
(T \times_{\mathcal{Y}} \mathcal{X}) \times_{\mathcal{X}, x} \Spec(k) =
T \times_{\mathcal{Y}, y} \Spec(k)
$$
Le morphisme $T \times_{\mathcal{Y}} \mathcal{X} \to T$ est quasi-compact ;
ainsi, $T \times_\mathcal{Y} \mathcal{X}$ est quasi-compact. Comme $x$ est
un morphisme quasi-compact puisque $\mathcal{X}$ est décent,
le produit fibré affiché est quasi-compact.
\end{proof}

\begin{lemma}
\label{lemma-gerbe-decent}
Soit $f : \mathcal{X} \to \mathcal{Y}$ un morphisme de champs algébriques.
Si $\mathcal{X}$ est une gerbe sur $\mathcal{Y}$ et si $\mathcal{X}$ est décent,
alors $\mathcal{Y}$ est décent.
\end{lemma}

\begin{proof}
Supposons que $\mathcal{X}$ soit une gerbe sur $\mathcal{Y}$ et que
$\mathcal{X}$ soit décent. Remarquons que $f$ est un homéomorphisme universel
d'après le Lemme \ref{lemma-gerbe-bijection-points}.
Le lemme résulte donc du Lemme \ref{lemma-qc-decent}.
\end{proof}


\section{Points des champs alg\'ebriques d\'ecents}
\label{section-points-decent}

\noindent
Cette section est l'analogue de la section
\ref{decent-spaces-section-points} d'Espaces d\'ecents.
Nous ne savons pas si l'espace topologique associ\'e
\`a un champ alg\'ebrique d\'ecent est toujours sobre ; voir
la Proposition \ref{proposition-decent-sober} pour un r\'esultat
un peu plus faible.

\begin{lemma}
\label{lemma-kolmogorov}
Soit $\mathcal{X}$ un champ alg\'ebrique d\'ecent.
Alors $|\mathcal{X}|$ est un espace de Kolmogoroff (voir
Topologie, D\'efinition \ref{topology-definition-generic-point}).
\end{lemma}

\begin{proof}
Soient $x_1, x_2 \in |\mathcal{X}|$ tels que $x_1 \leadsto x_2$ et
$x_2 \leadsto x_1$. Nous devons montrer que $x_1 = x_2$.
Soit $\mathcal{Z} \subset \mathcal{X}$ le sous-champ ferm\'e
r\'eduit tel que $|\mathcal{Z}|$ soit \`egal \`a
$\overline{\{x_1\}} = \overline{\{x_2\}}$.
D'apr\`es le Lemme \ref{lemma-representable-decent},
le champ $\mathcal{Z}$ est d\'ecent.
En rempla\c{c}ant $\mathcal{X}$ par $\mathcal{Z}$, nous nous ramenons
au cas examin\'e dans le paragraphe suivant.

\medskip\noindent
Supposons $|\mathcal{X}|$ irr\'eductible, ayant $x_1$ et $x_2$
pour points g\'en\'eriques. Choisissons un sch\'ema affine $U$, des points $u_1, u_2 \in U$
et un morphisme lisse $f : U \to \mathcal{X}$ tels que
$f(u_i) = x_i$. Il existe alors un troisi\`eme point $u_3 \in U$
qui est le point g\'en\'erique d'une composante irr\'eductible de $U$
et dont l'image $x_3 \in |\mathcal{X}|$ est {\it elle aussi} un
point g\'en\'erique de $|\mathcal{X}|$. En effet, il suffit de
choisir pour $u_3$ un point g\'en\'erique quelconque d'une composante irr\'eductible
passant par $u_1$ (ou par $u_2$, si l'on pr\'ef\`ere).
Dans le paragraphe suivant, nous montrerons que $x_1 = x_3$ et $x_2 = x_3$,
ce qui donnera le r\'esultat voulu.

\medskip\noindent
Par sym\'etrie, il suffit de prouver que $x_1 = x_3$. Puisque $x_1$
est un point g\'en\'erique de $|\mathcal{X}|$, on a une sp\'ecialisation
$x_1 \leadsto x_3$.
D'apr\`es le Lemme \ref{lemma-specialization}, il existe une sp\'ecialisation
$u'_1 \leadsto u_3$ dans $U$ (!) dont l'image est $x_1 \leadsto x_3$.
Or $u_3$ est le point g\'en\'erique d'une composante irr\'eductible ;
on a donc $u'_1 = u_3$, comme voulu.
\end{proof}

\begin{lemma}
\label{lemma-decent-locally-noetherian-sober}
Soit $\mathcal{X}$ un champ alg\'ebrique d\'ecent et localement noeth\'erien.
Alors $|\mathcal{X}|$ est un espace topologique sobre et localement noeth\'erien.
\end{lemma}

\begin{proof}
D'apr\`es le Lemme \ref{lemma-Noetherian-topology}, l'espace topologique $|\mathcal{X}|$
est localement noeth\'erien. D'apr\`es le Lemme \ref{lemma-kolmogorov}, l'espace topologique
$|\mathcal{X}|$ est un espace de Kolmogoroff. D'apr\`es le
Lemme \ref{lemma-locally-Noetherian-quasi-sober}, l'espace topologique
$|\mathcal{X}|$ est quasi-sobre. Cela ach\`eve la d\'emonstration ; voir
Topologie, D\'efinition \ref{topology-definition-generic-point}.
\end{proof}

\begin{proposition}
\label{proposition-decent-sober}
Soit $\mathcal{X}$ un champ alg\'ebrique d\'ecent tel que
$\mathcal{I}_\mathcal{X} \to \mathcal{X}$ soit quasi-compact.
Alors $|\mathcal{X}|$ est sobre.
\end{proposition}

\begin{proof}
D'apr\`es le Lemme \ref{lemma-kolmogorov}, $|\mathcal{X}|$ est un espace de Kolmogoroff
(en fait, nous allons le red\'emontrer).
Soit $T \subset |\mathcal{X}|$ une partie ferm\'ee irr\'eductible.
Nous devons montrer que $T$ poss\`ede un point g\'en\'erique.
Soit $\mathcal{Z} \subset \mathcal{X}$ le sous-champ ferm\'e induit
correspondant \`a $T$, muni de sa structure r\'eduite ; voir Propri\'et\'es des champs, D\'efinition
\ref{stacks-properties-definition-reduced-induced-stack}.
Puisque $\mathcal{Z} \to \mathcal{X}$ est une immersion ferm\'ee,
$\mathcal{Z}$ est un champ alg\'ebrique d\'ecent ; voir le
Lemme \ref{lemma-representable-decent}.
De plus, le morphisme $\mathcal{I}_\mathcal{Z} \to \mathcal{Z}$
est le changement de base de $\mathcal{I}_\mathcal{X} \to \mathcal{X}$
(Lemme \ref{lemma-monomorphism-cartesian-square-inertia}).
Par cons\'equent, $\mathcal{I}_\mathcal{Z} \to \mathcal{Z}$ est quasi-compact
(Lemme \ref{lemma-base-change-quasi-compact}).
Nous nous ramenons ainsi au cas examin\'e dans le paragraphe suivant.

\medskip\noindent
Supposons que $\mathcal{X}$ soit d\'ecent, que
$\mathcal{I}_\mathcal{X} \to \mathcal{X}$ soit quasi-compact,
que $\mathcal{X}$ soit r\'eduit et que $|\mathcal{X}|$ soit irr\'eductible.
Nous devons montrer que $|\mathcal{X}|$ poss\`ede un point g\'en\'erique.
D'apr\`es la Proposition \ref{proposition-open-stratum},
il existe un sous-champ ouvert dense $\mathcal{U} \subset \mathcal{X}$
qui est une gerbe. Autrement dit, $|\mathcal{U}| \subset |\mathcal{X}|$
est ouvert et dense.
Nous pouvons donc supposer en outre que $\mathcal{X}$ est une gerbe.
Supposons que $\mathcal{X} \to X$ munisse $\mathcal{X}$ d'une structure de gerbe sur
l'espace alg\'ebrique $X$. Alors $|\mathcal{X}| \cong |X|$ d'apr\`es le
Lemme \ref{lemma-gerbe-bijection-points}.
En particulier, $X$ est quasi-compact et $|X|$ est irr\'eductible.
De plus, le Lemme \ref{lemma-gerbe-decent} montre
que $X$ est un espace alg\'ebrique d\'ecent.
Alors $|\mathcal{X}| = |X|$ est sobre d'apr\`es
Espaces d\'ecents, Proposition \ref{decent-spaces-proposition-reasonable-sober},
et poss\`ede donc un point g\'en\'erique (unique).
\end{proof}




%
\section{Champs alg\'ebriques int\`egres}
\label{section-integral-stacks}

\noindent
Cette section est l'analogue de Espaces sur des corps, section
\ref{spaces-over-fields-section-integral-spaces}.
Les consid\'erations de cette section et le r\'esultat de la
Proposition \ref{proposition-decent-sober} nous am\`enent \`a d\'efinir
comme suit un champ alg\'ebrique int\`egre (et cette d\'efinition est compatible
avec les d\'efinitions d\'ej\`a existantes des sch\'emas int\`egres et
des espaces alg\'ebriques int\`egres).

\begin{definition}
\label{definition-integral-algebraic-stack}
On dit qu'un champ alg\'ebrique $\mathcal{X}$ est
{\it int\`egre} s'il est r\'eduit et d\'ecent, si
$\mathcal{I}_\mathcal{X} \to \mathcal{X}$ est quasi-compact,
et si $|\mathcal{X}|$ est irr\'eductible.
\end{definition}

\noindent
Remarquons que si $\mathcal{X}$ est quasi-s\'epar\'e, il suffit,
pour qu'il soit int\`egre, que $\mathcal{X}$ soit r\'eduit et que
$|\mathcal{X}|$ soit irr\'eductible; voir le Lemme \ref{lemma-totaro-case}.

\begin{lemma}
\label{lemma-structure-integral-stack}
Soit $\mathcal{X}$ un champ alg\'ebrique int\`egre. Alors
\begin{enumerate}
\item $|\mathcal{X}|$ est sobre et irr\'eductible, et poss\`ede un unique point g\'en\'erique,
\item il existe un sous-champ ouvert $\mathcal{U} \subset \mathcal{X}$
qui est une gerbe sur un sch\'ema int\`egre $U$.
\end{enumerate}
\end{lemma}

\begin{proof}
La Proposition \ref{proposition-decent-sober} montre que $|\mathcal{X}|$
est sobre. Il est bien entendu aussi irr\'eductible et poss\`ede donc un unique
point g\'en\'erique $x$ (par d\'efinition de la sobri\'et\'e).
La Proposition \ref{proposition-open-stratum} montre
qu'il existe un sous-champ ouvert dense $\mathcal{U} \subset \mathcal{X}$
qui est une gerbe sur un espace alg\'ebrique $U$.
Alors $U$ est un espace alg\'ebrique d\'ecent d'apr\`es le
Lemme \ref{lemma-gerbe-decent} (et puisque $\mathcal{U}$
est d\'ecent d'apr\`es le Lemme \ref{lemma-representable-decent}).
Comme $|U| = |\mathcal{U}|$, on voit que $|U|$ est irr\'eductible.
Enfin, puisque $\mathcal{U}$ est r\'eduit, le morphisme
$\mathcal{U} \to U$ se factorise par $U_{red}$; voir
Propri\'et\'es des champs, Lemme \ref{stacks-properties-lemma-map-into-reduction}.
Or $\mathcal{U} \to U$ est plat, localement de pr\'esentation finie
et surjectif (Lemme \ref{lemma-gerbe-fppf}); il s'ensuit que
$U = U_{red}$, c'est-\`a-dire que $U$ est r\'eduit (un petit d\'etail est omis).
Ainsi, $U$ est un espace alg\'ebrique int\`egre; voir
Espaces sur des corps, D\'efinition
\ref{spaces-over-fields-definition-integral-algebraic-space}.
Enfin, nous pouvons remplacer $U$ (et, corr\'elativement, $\mathcal{U}$)
par un sous-espace ouvert et supposer que $U$ est un sch\'ema int\`egre; voir
la discussion de Espaces sur des corps, section
\ref{spaces-over-fields-section-integral-spaces}.
\end{proof}

\begin{lemma}
\label{lemma-totaro-case}
Soit $\mathcal{X}$ un champ alg\'ebrique r\'eduit et quasi-s\'epar\'e
dont l'espace topologique sous-jacent $|\mathcal{X}|$ est irr\'eductible.
Alors $\mathcal{X}$ est int\`egre.
\end{lemma}

\begin{proof}
Si $\mathcal{X}$ est quasi-s\'epar\'e, alors $\mathcal{X}$ est d\'ecent
d'apr\`es le Lemme \ref{lemma-quasi-separated-decent}.
Si $\mathcal{X}$ est quasi-s\'epar\'e, alors
$\Delta : \mathcal{X} \to \mathcal{X} \times \mathcal{X}$ est quasi-compact;
par cons\'equent, $\mathcal{I}_\mathcal{X} \to \mathcal{X}$ est quasi-compact
comme changement de base de $\Delta$ par $\Delta$; voir le
Lemme \ref{lemma-base-change-quasi-compact}.
On voit ainsi que toutes les hypoth\`eses de la
D\'efinition \ref{definition-integral-algebraic-stack} sont satisfaites
(et aussi que l'on peut remplacer ``quasi-s\'epar\'e'' par
``$\Delta_\mathcal{X}$ est quasi-compact'').
\end{proof}

\begin{lemma}
\label{lemma-decent-irreducible-closed}
Soit $\mathcal{X}$ un champ alg\'ebrique d\'ecent tel que
$\mathcal{I}_\mathcal{X} \to \mathcal{X}$ soit quasi-compact.
Il existe des bijections canoniques entre les ensembles suivants :
\begin{enumerate}
\item l'ensemble des points de $\mathcal{X}$, c'est-\`a-dire $|\mathcal{X}|$,
\item l'ensemble des parties ferm\'ees irr\'eductibles de $|\mathcal{X}|$,
\item l'ensemble des sous-champs ferm\'es int\`egres de $\mathcal{X}$.
\end{enumerate}
La bijection de (1) vers (2) envoie $x$ sur $\overline{\{x\}}$.
La bijection de (3) vers (2) envoie $\mathcal{Z}$ sur $|\mathcal{Z}|$.
\end{lemma}

\begin{proof}
La premi\`ere application d\'efinit une bijection entre (1) et (2), puisque
$|\mathcal{X}|$ est sobre d'apr\`es la Proposition \ref{proposition-decent-sober}.
\'Etant donn\'ee une partie ferm\'ee irr\'eductible $T \subset |\mathcal{X}|$,
il existe un unique sous-champ ferm\'e r\'eduit $\mathcal{Z} \subset \mathcal{X}$
tel que $|\mathcal{Z}| = T$ : $\mathcal{Z}$ est pr\'ecis\'ement
la structure r\'eduite de sous-champ induite sur $T$; voir
Propri\'et\'es des champs, D\'efinition
\ref{stacks-properties-definition-reduced-induced-stack}.
Alors $\mathcal{Z}$ est un champ alg\'ebrique int\`egre : il est
d\'ecent (Lemme \ref{lemma-representable-decent}), le
morphisme $\mathcal{I}_\mathcal{Z} \to \mathcal{Z}$ est quasi-compact
(comme changement de base de $\mathcal{I}_\mathcal{X} \to \mathcal{X}$; voir
le Lemme \ref{lemma-monomorphism-cartesian-square-inertia}),
$\mathcal{Z}$ est r\'eduit, et $|\mathcal{Z}|$ est irr\'eductible.
\end{proof}






\section{Gerbes r\'esiduelles}
\label{section-residual-gerbe}

\noindent
Cette section prolonge Propri\'et\'es des champs, section
\ref{stacks-properties-section-residual-gerbe}.

\begin{lemma}
\label{lemma-gerbe-residual-gerbe}
Soit $\pi : \mathcal{X} \to \mathcal{Y}$ un morphisme de champs alg\'ebriques.
Soit $x \in |\mathcal{X}|$, d'image $y \in |\mathcal{Y}|$.
Supposons que la gerbe r\'esiduelle $\mathcal{Z}_y \subset \mathcal{Y}$
de $\mathcal{Y}$ en $y$ existe et que $\mathcal{X}$ soit une gerbe
sur $\mathcal{Y}$. Alors
$\mathcal{Z}_x = \mathcal{Z}_y \times_\mathcal{Y} \mathcal{X}$ est
la gerbe r\'esiduelle de $\mathcal{X}$ en $x$.
\end{lemma}

\begin{proof}
Le morphisme $\mathcal{Z}_x \to \mathcal{X}$ est un monomorphisme, car
il se d\'eduit par changement de base du monomorphisme $\mathcal{Z}_y \to \mathcal{Y}$.
Le morphisme $\pi$ est un hom\'eomorphisme universel d'apr\`es le
Lemme \ref{lemma-gerbe-bijection-points}; par cons\'equent, $|\mathcal{Z}_x| = \{x\}$.
Enfin, le morphisme $\mathcal{Z}_x \to \mathcal{Z}_y$ est lisse,
car il se d\'eduit par changement de base du morphisme lisse $\pi$, voir le
Lemme \ref{lemma-gerbe-smooth}.
Puisque $\mathcal{Z}_y$ est r\'eduit et localement noeth\'erien, il en
est donc de m\^eme de $\mathcal{Z}_x$ (d\'etails omis). Ainsi, $\mathcal{Z}_x$ est
la gerbe r\'esiduelle de $\mathcal{X}$ en $x$ d'apr\`es
Propri\'et\'es des champs, D\'efinition
\ref{stacks-properties-definition-residual-gerbe}.
\end{proof}

\begin{lemma}
\label{lemma-factor-through-residual-space}
Soit $f : \mathcal{Y} \to \mathcal{X}$ un morphisme de champs alg\'ebriques.
Soit $x \in |\mathcal{X}|$ un point. Supposons que
\begin{enumerate}
\item $\mathcal{X}$ soit d\'ecent ou localement noeth\'erien (ou les deux),
\item $\mathcal{I}_\mathcal{X} \to \mathcal{X}$ soit quasi-compact,
\item $|f|(|\mathcal{Y}|)$ soit contenu dans $\{x\} \subset |\mathcal{X}|$, et
\item $\mathcal{Y}$ soit r\'eduit.
\end{enumerate}
Alors $f$ se factorise par la gerbe r\'esiduelle $\mathcal{Z}_x$
de $\mathcal{X}$ en $x$ (dont l'existence est assur\'ee par le
Lemme \ref{lemma-every-point-residual-gerbe} ou le
Lemme \ref{lemma-every-point-residual-gerbe-locally-Noetherian}).
\end{lemma}

\begin{proof}
Soit $T = \overline{\{x\}} \subset |\mathcal{X}|$ l'adh\'erence de $x$.
D'apr\`es Propri\'et\'es des champs, Lemme
\ref{stacks-properties-lemma-reduced-closed-substack},
il existe un sous-champ ferm\'e r\'eduit $\mathcal{X}' \subset \mathcal{X}$
tel que $T = |\mathcal{X}'|$. D'apr\`es Propri\'et\'es des champs, Lemme
\ref{stacks-properties-lemma-map-into-reduction}, le
morphisme $f$ se factorise par $\mathcal{X}'$. Si $\mathcal{X}$ est d\'ecent,
alors, d'apr\`es le Lemme \ref{lemma-representable-decent}, le champ
$\mathcal{X}'$ est d\'ecent. Si $\mathcal{X}$ est localement noeth\'erien,
alors $\mathcal{X}'$ est localement noeth\'erien (d\'etails omis).
Le morphisme $\mathcal{I}_{\mathcal{X}'} \to \mathcal{X}'$
se d\'eduit par changement de base de $\mathcal{I}_\mathcal{X} \to \mathcal{X}$ d'apr\`es le
Lemme \ref{lemma-monomorphism-cartesian-square-inertia};
le Lemme \ref{lemma-base-change-quasi-compact} montre donc que
$\mathcal{I}_{\mathcal{X}'} \to \mathcal{X}'$ est quasi-compact.
On est ainsi ramen\'e au cas examin\'e dans le paragraphe suivant.

\medskip\noindent
Supposons que $\mathcal{X}$ soit r\'eduit et que $x \in |\mathcal{X}|$
soit un point g\'en\'erique. La Proposition \ref{proposition-open-stratum}
entra\^ine l'existence d'un sous-champ ouvert dense
$\mathcal{U} \subset \mathcal{X}$ qui est une gerbe.
Notons que $x \in |\mathcal{U}|$.
En reprenant les arguments ci-dessus, on se ram\`ene au cas examin\'e
dans le paragraphe suivant.

\medskip\noindent
Supposons que $\mathcal{X} \to X$ soit une gerbe sur l'espace alg\'ebrique $X$.
Si $\mathcal{X}$ est d\'ecent, alors les
Lemmes \ref{lemma-gerbe-bijection-points} et \ref{lemma-qc-decent}
montrent que l'espace $X$ est d\'ecent. Si $\mathcal{X}$ est localement noeth\'erien,
alors $X$ est localement noeth\'erien par descente fppf (d\'etails omis).
Le r\'esultat correspondant vaut donc pour $X$, voir Espaces d\'ecents, Lemme
\ref{decent-spaces-lemma-factor-through-residual-space-decent} ou
\ref{decent-spaces-lemma-factor-through-residual-space-Noetherian}
(un petit d\'etail est omis).
En appliquant le Lemme \ref{lemma-gerbe-residual-gerbe}, on conclut que le r\'esultat
vaut aussi pour $\mathcal{X}$.
\end{proof}

\begin{remark}
\label{lemma-factor-through-residual-space-Noetherian}
Nous ignorons si le Lemme \ref{lemma-factor-through-residual-space}
reste vrai lorsqu'on suppose seulement que $\mathcal{X}$ est localement noeth\'erien,
c'est-\`a-dire lorsqu'on supprime l'hypoth\`ese de quasi-compacit\'e de l'inertie. Dans ce cas,
si $x$ est un point ferm\'e, le r\'esultat est certainement vrai, comme le montre le
lemme beaucoup plus simple qui suit.
\end{remark}

\begin{lemma}
\label{lemma-residual-space-closed}
Soit $\mathcal{X}$ un champ alg\'ebrique localement noeth\'erien.
Soit $x \in |\mathcal{X}|$, dont la gerbe r\'esiduelle est
$\mathcal{Z}_x \subset \mathcal{X}$
(Lemme \ref{lemma-every-point-residual-gerbe-locally-Noetherian}).
Alors $x$ est un point ferm\'e de $|\mathcal{X}|$ si et seulement si
le morphisme $\mathcal{Z}_x \to \mathcal{X}$ est une immersion ferm\'ee.
\end{lemma}

\begin{proof}
Si $\mathcal{Z}_x \to \mathcal{X}$ est une immersion ferm\'ee, alors
$x$ est un point ferm\'e de $|\mathcal{X}|$, voir par exemple le
Lemme \ref{lemma-closed-immersion-proper}.
R\'eciproquement, supposons que $x$ soit un point ferm\'e de $|\mathcal{X}|$.
Soit $\mathcal{Z} \subset \mathcal{X}$
le sous-champ ferm\'e r\'eduit tel que $|\mathcal{Z}| = \{x\}$
(Propri\'et\'es des champs,
Lemme \ref{stacks-properties-lemma-reduced-closed-substack}).
Alors $\mathcal{Z}$ est un champ alg\'ebrique localement noeth\'erien
d'apr\`es les Lemmes \ref{lemma-immersion-locally-finite-type} et
\ref{lemma-locally-finite-type-locally-noetherian}.
Comme $\mathcal{Z}$ est aussi r\'eduit et que $|\mathcal{Z}| = \{x\}$,
il s'ensuit que $\mathcal{Z} = \mathcal{Z}_x$ est, par d\'efinition, la
gerbe r\'esiduelle.
\end{proof}











\begin{multicols}{2}[\section{Autres chapitres}]
\noindent
Préliminaires
\begin{enumerate}
\item \hyperref[introduction-section-phantom]{Introduction}
\item \hyperref[conventions-section-phantom]{Conventions}
\item \hyperref[sets-section-phantom]{Théorie des ensembles}
\item \hyperref[categories-section-phantom]{Catégories}
\item \hyperref[topology-section-phantom]{Topologie}
\item \hyperref[sheaves-section-phantom]{Faisceaux sur les espaces}
\item \hyperref[sites-section-phantom]{Sites et faisceaux}
\item \hyperref[stacks-section-phantom]{Champs}
\item \hyperref[fields-section-phantom]{Corps}
\item \hyperref[algebra-section-phantom]{Algèbre commutative}
\item \hyperref[brauer-section-phantom]{Groupes de Brauer}
\item \hyperref[homology-section-phantom]{Algèbre homologique}
\item \hyperref[derived-section-phantom]{Catégories dérivées}
\item \hyperref[simplicial-section-phantom]{Méthodes simpliciales}
\item \hyperref[more-algebra-section-phantom]{Compléments d'algèbre}
\item \hyperref[smoothing-section-phantom]{Lissage des morphismes d'anneaux}
\item \hyperref[modules-section-phantom]{Faisceaux de modules}
\item \hyperref[sites-modules-section-phantom]{Modules sur les sites}
\item \hyperref[injectives-section-phantom]{Objets injectifs}
\item \hyperref[cohomology-section-phantom]{Cohomologie des faisceaux}
\item \hyperref[sites-cohomology-section-phantom]{Cohomologie sur les sites}
\item \hyperref[dga-section-phantom]{Algèbre différentielle graduée}
\item \hyperref[dpa-section-phantom]{Algèbre à puissances divisées}
\item \hyperref[sdga-section-phantom]{Faisceaux différentiels gradués}
\item \hyperref[hypercovering-section-phantom]{Hyperrecouvrements}
\end{enumerate}
Schémas
\begin{enumerate}
\setcounter{enumi}{25}
\item \hyperref[schemes-section-phantom]{Schémas}
\item \hyperref[constructions-section-phantom]{Constructions de schémas}
\item \hyperref[properties-section-phantom]{Propriétés des schémas}
\item \hyperref[morphisms-section-phantom]{Morphismes de schémas}
\item \hyperref[coherent-section-phantom]{Cohomologie des schémas}
\item \hyperref[divisors-section-phantom]{Diviseurs}
\item \hyperref[limits-section-phantom]{Limites de schémas}
\item \hyperref[varieties-section-phantom]{Variétés}
\item \hyperref[topologies-section-phantom]{Topologies sur les schémas}
\item \hyperref[descent-section-phantom]{Descente}
\item \hyperref[perfect-section-phantom]{Catégories dérivées de schémas}
\item \hyperref[more-morphisms-section-phantom]{Compléments sur les morphismes}
\item \hyperref[flat-section-phantom]{Compléments sur la platitude}
\item \hyperref[groupoids-section-phantom]{Schémas en groupoïdes}
\item \hyperref[more-groupoids-section-phantom]{Compléments sur les schémas en groupoïdes}
\item \hyperref[etale-section-phantom]{Morphismes étales de schémas}
\end{enumerate}
Thèmes de théorie des schémas
\begin{enumerate}
\setcounter{enumi}{41}
\item \hyperref[chow-section-phantom]{Homologie de Chow}
\item \hyperref[intersection-section-phantom]{Théorie de l'intersection}
\item \hyperref[pic-section-phantom]{Schémas de Picard des courbes}
\item \hyperref[weil-section-phantom]{Théories cohomologiques de Weil}
\item \hyperref[adequate-section-phantom]{Modules adéquats}
\item \hyperref[dualizing-section-phantom]{Complexes dualisants}
\item \hyperref[duality-section-phantom]{Dualité pour les schémas}
\item \hyperref[discriminant-section-phantom]{Discriminants et différentes}
\item \hyperref[derham-section-phantom]{Cohomologie de de Rham}
\item \hyperref[local-cohomology-section-phantom]{Cohomologie locale}
\item \hyperref[algebraization-section-phantom]{Géométrie algébrique et formelle}
\item \hyperref[curves-section-phantom]{Courbes algébriques}
\item \hyperref[resolve-section-phantom]{Résolution des surfaces}
\item \hyperref[models-section-phantom]{Réduction semi-stable}
\item \hyperref[functors-section-phantom]{Foncteurs et morphismes}
\item \hyperref[equiv-section-phantom]{Catégories dérivées de variétés}
\item \hyperref[pione-section-phantom]{Groupes fondamentaux de schémas}
\item \hyperref[etale-cohomology-section-phantom]{Cohomologie étale}
\item \hyperref[crystalline-section-phantom]{Cohomologie cristalline}
\item \hyperref[proetale-section-phantom]{Cohomologie pro-étale}
\item \hyperref[relative-cycles-section-phantom]{Cycles relatifs}
\item \hyperref[more-etale-section-phantom]{Compléments de cohomologie étale}
\item \hyperref[trace-section-phantom]{La formule des traces}
\end{enumerate}
Espaces algébriques
\begin{enumerate}
\setcounter{enumi}{64}
\item \hyperref[spaces-section-phantom]{Espaces algébriques}
\item \hyperref[spaces-properties-section-phantom]{Propriétés des espaces algébriques}
\item \hyperref[spaces-morphisms-section-phantom]{Morphismes d'espaces algébriques}
\item \hyperref[decent-spaces-section-phantom]{Espaces algébriques décents}
\item \hyperref[spaces-cohomology-section-phantom]{Cohomologie des espaces algébriques}
\item \hyperref[spaces-limits-section-phantom]{Limites d'espaces algébriques}
\item \hyperref[spaces-divisors-section-phantom]{Diviseurs sur les espaces algébriques}
\item \hyperref[spaces-over-fields-section-phantom]{Espaces algébriques sur des corps}
\item \hyperref[spaces-topologies-section-phantom]{Topologies sur les espaces algébriques}
\item \hyperref[spaces-descent-section-phantom]{Descente et espaces algébriques}
\item \hyperref[spaces-perfect-section-phantom]{Catégories dérivées d'espaces}
\item \hyperref[spaces-more-morphisms-section-phantom]{Compléments sur les morphismes d'espaces}
\item \hyperref[spaces-flat-section-phantom]{Platitude sur les espaces algébriques}
\item \hyperref[spaces-groupoids-section-phantom]{Groupoïdes dans les espaces algébriques}
\item \hyperref[spaces-more-groupoids-section-phantom]{Compléments sur les groupoïdes dans les espaces}
\item \hyperref[bootstrap-section-phantom]{Amorçage}
\item \hyperref[spaces-pushouts-section-phantom]{Sommes amalgamées d'espaces algébriques}
\end{enumerate}
Thèmes de géométrie
\begin{enumerate}
\setcounter{enumi}{81}
\item \hyperref[spaces-chow-section-phantom]{Groupes de Chow des espaces}
\item \hyperref[groupoids-quotients-section-phantom]{Quotients de groupoïdes}
\item \hyperref[spaces-more-cohomology-section-phantom]{Compléments sur la cohomologie des espaces}
\item \hyperref[spaces-simplicial-section-phantom]{Espaces simpliciaux}
\item \hyperref[spaces-duality-section-phantom]{Dualité pour les espaces}
\item \hyperref[formal-spaces-section-phantom]{Espaces algébriques formels}
\item \hyperref[restricted-section-phantom]{Algébrisation des espaces formels}
\item \hyperref[spaces-resolve-section-phantom]{Résolution des surfaces revisitée}
\end{enumerate}
Théorie des déformations
\begin{enumerate}
\setcounter{enumi}{89}
\item \hyperref[formal-defos-section-phantom]{Théorie formelle des déformations}
\item \hyperref[defos-section-phantom]{Théorie des déformations}
\item \hyperref[cotangent-section-phantom]{Le complexe cotangent}
\item \hyperref[examples-defos-section-phantom]{Problèmes de déformation}
\end{enumerate}
Champs algébriques
\begin{enumerate}
\setcounter{enumi}{93}
\item \hyperref[algebraic-section-phantom]{Champs algébriques}
\item \hyperref[examples-stacks-section-phantom]{Exemples de champs}
\item \hyperref[stacks-sheaves-section-phantom]{Faisceaux sur les champs algébriques}
\item \hyperref[criteria-section-phantom]{Critères de représentabilité}
\item \hyperref[artin-section-phantom]{Axiomes d'Artin}
\item \hyperref[quot-section-phantom]{Espaces de Quot et de Hilbert}
\item \hyperref[stacks-properties-section-phantom]{Propriétés des champs algébriques}
\item \hyperref[stacks-morphisms-section-phantom]{Morphismes de champs algébriques}
\item \hyperref[stacks-limits-section-phantom]{Limites de champs algébriques}
\item \hyperref[stacks-cohomology-section-phantom]{Cohomologie des champs algébriques}
\item \hyperref[stacks-perfect-section-phantom]{Catégories dérivées de champs}
\item \hyperref[stacks-introduction-section-phantom]{Introduction aux champs algébriques}
\item \hyperref[stacks-more-morphisms-section-phantom]{Compléments sur les morphismes de champs}
\item \hyperref[stacks-geometry-section-phantom]{La géométrie des champs}
\end{enumerate}
Thèmes de théorie des espaces de modules
\begin{enumerate}
\setcounter{enumi}{107}
\item \hyperref[moduli-section-phantom]{Champs de modules}
\item \hyperref[moduli-curves-section-phantom]{Espaces de modules de courbes}
\end{enumerate}
Divers
\begin{enumerate}
\setcounter{enumi}{109}
\item \hyperref[examples-section-phantom]{Exemples}
\item \hyperref[exercises-section-phantom]{Exercices}
\item \hyperref[guide-section-phantom]{Guide bibliographique}
\item \hyperref[desirables-section-phantom]{Desiderata}
\item \hyperref[coding-section-phantom]{Style de programmation}
\item \hyperref[obsolete-section-phantom]{Obsolète}
\item \hyperref[fdl-section-phantom]{Licence de documentation libre GNU}
\item \hyperref[index-section-phantom]{Index généré automatiquement}
\end{enumerate}
\end{multicols}


\bibliography{my}
\bibliographystyle{amsalpha}

\end{document}
